% === START OF CHUNK preamble+ch1-2 ===
%%%%%%%%%%%%%%%%%%%%%%%%%%%%%%%%%%%%%%%%%%%%%%%%%%%%%%%%%%%%%
%% ALGEBRAIC TOPOLOGY — A Graduate Course
%% Preamble + Chapter 1 + Chapter 2
%%%%%%%%%%%%%%%%%%%%%%%%%%%%%%%%%%%%%%%%%%%%%%%%%%%%%%%%%%%%%

\documentclass[12pt,a4paper,openany]{book}

%% ─── Encoding and fonts ─────────────────────────────────
\usepackage[utf8]{inputenc}
\usepackage[T1]{fontenc}
\usepackage{lmodern}

%% ─── Page geometry ──────────────────────────────────────
\usepackage[margin=2.5cm]{geometry}

%% ─── Mathematics ────────────────────────────────────────
\usepackage{amsmath,amssymb,amsthm,mathrsfs,mathtools}

%% ─── Graphics and diagrams ──────────────────────────────
\usepackage{tikz}
\usetikzlibrary{arrows.meta,calc,positioning,patterns,decorations.markings,decorations.pathmorphing,cd}
\usepackage{pgfplots}
\pgfplotsset{compat=1.18}
\usepackage{tikz-cd}

%% ─── Colored boxes ──────────────────────────────────────
\usepackage[most]{tcolorbox}
\tcbuselibrary{theorems,skins,breakable}

%% ─── Hyperlinks ─────────────────────────────────────────
\usepackage[colorlinks=true,linkcolor=blue!60!black,citecolor=green!50!black]{hyperref}

%% ─── Index ──────────────────────────────────────────────
\usepackage{makeidx}
\makeindex

%% ─── Lists, headers, captions, tables ───────────────────
\usepackage{enumitem}
\usepackage{fancyhdr}
\usepackage{caption}
\usepackage{booktabs}
\usepackage{array}
\usepackage{longtable}
\usepackage{microtype}
\usepackage{xcolor}
\usepackage{minitoc}

%% ─── Page style ─────────────────────────────────────────
\pagestyle{fancy}
\fancyhf{}
\fancyhead[LE]{\leftmark}
\fancyhead[RO]{\rightmark}
\fancyfoot[C]{\thepage}
\renewcommand{\headrulewidth}{0.4pt}
\renewcommand{\footrulewidth}{0pt}
\setlength{\headheight}{15pt}

%% ─── Theorem environments (amsthm) ─────────────────────
\theoremstyle{definition}
\newtheorem{definition}{Definition}[chapter]
\newtheorem{example}[definition]{Example}
\newtheorem{exercise}{Exercise}[chapter]

\theoremstyle{plain}
\newtheorem{theorem}[definition]{Theorem}
\newtheorem{proposition}[definition]{Proposition}
\newtheorem{lemma}[definition]{Lemma}
\newtheorem{corollary}[definition]{Corollary}

\theoremstyle{remark}
\newtheorem{remark}[definition]{Remark}
\newtheorem{notation}[definition]{Notation}

%% ─── tcolorbox wrapping ────────────────────────────────
\tcolorboxenvironment{definition}{enhanced,breakable,colback=blue!5,colframe=blue!70!black,fonttitle=\bfseries,before skip=10pt,after skip=10pt,left=8pt,right=8pt}
\tcolorboxenvironment{theorem}{enhanced,breakable,colback=red!5,colframe=red!70!black,fonttitle=\bfseries,before skip=10pt,after skip=10pt,left=8pt,right=8pt}
\tcolorboxenvironment{proposition}{enhanced,breakable,colback=orange!5,colframe=orange!70!black,fonttitle=\bfseries,before skip=10pt,after skip=10pt,left=8pt,right=8pt}
\tcolorboxenvironment{lemma}{enhanced,breakable,colback=yellow!5,colframe=yellow!50!black,fonttitle=\bfseries,before skip=10pt,after skip=10pt,left=8pt,right=8pt}
\tcolorboxenvironment{corollary}{enhanced,breakable,colback=green!5,colframe=green!50!black,fonttitle=\bfseries,before skip=10pt,after skip=10pt,left=8pt,right=8pt}
\tcolorboxenvironment{remark}{enhanced,breakable,colback=gray!5,colframe=gray!50!black,before skip=8pt,after skip=8pt,left=6pt,right=6pt}
\tcolorboxenvironment{example}{enhanced,breakable,colback=purple!5,colframe=purple!50!black,before skip=8pt,after skip=8pt,left=6pt,right=6pt}

%% ─── Custom commands ───────────────────────────────────
\newcommand{\N}{\mathbb{N}}
\newcommand{\Z}{\mathbb{Z}}
\newcommand{\Q}{\mathbb{Q}}
\newcommand{\R}{\mathbb{R}}
\newcommand{\C}{\mathbb{C}}
\newcommand{\RP}{\mathbb{R}\mathrm{P}}
\newcommand{\CP}{\mathbb{C}\mathrm{P}}
\DeclareMathOperator{\Hom}{Hom}
\DeclareMathOperator{\im}{im}
\DeclareMathOperator{\coker}{coker}
\DeclareMathOperator{\Aut}{Aut}
\DeclareMathOperator{\Id}{Id}
\newcommand{\abs}[1]{\left|#1\right|}
\newcommand{\set}[1]{\left\{#1\right\}}
\newcommand{\rel}{\operatorname{rel}}
\newcommand{\F}{\mathbb{F}}

%% ─── Title information ─────────────────────────────────
%%%%%%%%%%%%%%%%%%%%%%%%%%%%%%%%%%%%%%%%%%%%%%%%%%%%%%%%%%%%%
\begin{document}
\dominitoc
%%%%%%%%%%%%%%%%%%%%%%%%%%%%%%%%%%%%%%%%%%%%%%%%%%%%%%%%%%%%%

\begin{titlepage}
\begin{tikzpicture}[remember picture, overlay]
  \fill[left color=blue!15!white, right color=blue!5!white]
    (current page.south west) rectangle (current page.north east);
  \fill[color=orange!70!yellow]
    ([yshift=-2cm]current page.north west) rectangle ([yshift=-2.3cm]current page.north east);
  \fill[color=blue!80!black, rounded corners=8pt]
    ([xshift=3cm, yshift=-4cm]current page.north west) rectangle
    ([xshift=-3cm, yshift=-10cm]current page.north east);
  \node[anchor=center, text=white, font=\Huge\bfseries]
    at ([yshift=-6cm]current page.north) {Algebraic Topology};
  \node[anchor=center, text=orange!80!yellow, font=\Large]
    at ([yshift=-7.5cm]current page.north) {Lecture Notes};
  \node[anchor=center, text=white!80, font=\large]
    at ([yshift=-8.8cm]current page.north) {Master M1 --- 2025--2026};
  \node[anchor=center, text=blue!80!black, font=\Large\itshape]
    at ([yshift=-12cm]current page.north) {Ya\'e Ulrich Gaba};
  \draw[orange!70!yellow, line width=2pt]
    ([xshift=5cm, yshift=-13.5cm]current page.north west) --
    ([xshift=-5cm, yshift=-13.5cm]current page.north east);
  \node[anchor=center, text width=12cm, align=center, text=blue!60!black, font=\itshape]
    at ([yshift=-16cm]current page.north) {``Algebraic topology is the art of turning geometry problems into algebra problems.''\\[4pt]--- Henri Poincar\'e};
  \node[anchor=south, text=gray, font=\small]
    at ([yshift=1.5cm]current page.south) {\today};
\end{tikzpicture}
\end{titlepage}
\tableofcontents

%%%%%%%%%%%%%%%%%%%%%%%%%%%%%%%%%%%%%%%%%%%%%%%%%%%%%%%%%%%%%
%% NOTATION TABLE
%%%%%%%%%%%%%%%%%%%%%%%%%%%%%%%%%%%%%%%%%%%%%%%%%%%%%%%%%%%%%
\chapter*{Notation}
\addcontentsline{toc}{chapter}{Notation}

\noindent
The following notation is used throughout this text.

\bigskip

\begin{longtable}{>{\raggedright\arraybackslash}p{3.8cm} p{9.5cm}}
\toprule
\textbf{Symbol} & \textbf{Meaning} \\
\midrule
\endhead
$\N, \Z, \Q, \R, \C$ & Natural numbers, integers, rationals, reals, complex numbers \\[4pt]
$I = [0,1]$ & The unit interval \\[4pt]
$S^n$ & The $n$-sphere $\set{x \in \R^{n+1} : \abs{x} = 1}$ \\[4pt]
$D^n$ & The closed $n$-disk $\set{x \in \R^n : \abs{x} \le 1}$ \\[4pt]
$\RP^n$, $\CP^n$ & Real and complex projective $n$-space \\[4pt]
$T^n$ & The $n$-torus $\underbrace{S^1 \times \cdots \times S^1}_{n}$ \\[4pt]
$X \times Y$ & Product of topological spaces \\[4pt]
$X / {\sim}$ & Quotient space by equivalence relation $\sim$ \\[4pt]
$X \vee Y$ & Wedge sum (one-point union) \\[4pt]
$X \simeq Y$ & $X$ is homotopy equivalent to $Y$ \\[4pt]
$X \cong Y$ & $X$ is homeomorphic to $Y$ \\[4pt]
$f \simeq g$ & $f$ is homotopic to $g$ \\[4pt]
$[f]$ & Homotopy class of the map $f$ \\[4pt]
$\pi_1(X, x_0)$ & Fundamental group of $X$ at basepoint $x_0$ \\[4pt]
$\pi_n(X, x_0)$ & $n$-th homotopy group of $X$ at basepoint $x_0$ \\[4pt]
$H_n(X)$, $H^n(X)$ & $n$-th singular homology and cohomology groups \\[4pt]
$f_*$, $f^*$ & Induced map on homology, cohomology \\[4pt]
$\Hom(A, B)$ & Set (or group) of homomorphisms from $A$ to $B$ \\[4pt]
$\im f$, $\ker f$ & Image and kernel of a homomorphism $f$ \\[4pt]
$\coker f$ & Cokernel of $f$: quotient of codomain by $\im f$ \\[4pt]
$\Aut(G)$ & Automorphism group of $G$ \\[4pt]
$\Id_X$ & Identity map on $X$ \\[4pt]
$\mathbf{Top}$ & Category of topological spaces and continuous maps \\[4pt]
$\mathbf{hTop}$ & Homotopy category of topological spaces \\[4pt]
$\mathbf{Grp}$ & Category of groups and group homomorphisms \\[4pt]
$\mathbf{Ab}$ & Category of abelian groups \\[4pt]
$\mathbf{Set}$ & Category of sets and functions \\[4pt]
$\star$ & Concatenation of paths \\[4pt]
$\bar{\gamma}$ & Reverse path: $\bar{\gamma}(t) = \gamma(1-t)$ \\[4pt]
$c_{x_0}$ & Constant path at $x_0$ \\[4pt]
\bottomrule
\end{longtable}

%%%%%%%%%%%%%%%%%%%%%%%%%%%%%%%%%%%%%%%%%%%%%%%%%%%%%%%%%%%%%
%% PREFACE
%%%%%%%%%%%%%%%%%%%%%%%%%%%%%%%%%%%%%%%%%%%%%%%%%%%%%%%%%%%%%
\chapter*{Preface}
\addcontentsline{toc}{chapter}{Preface}

Algebraic topology is the study of topological spaces by means of algebraic
invariants. Its central theme is the translation of topological problems---often
geometric and intuitive, yet notoriously hard to resolve by purely topological
methods---into algebraic ones, where powerful computational machinery is
available.

This course is intended for beginning graduate students who have a solid
background in point-set topology and abstract algebra. We assume familiarity
with topological spaces, continuous maps, compactness, connectedness, quotient
spaces, groups, rings, modules, and exact sequences. Some exposure to category
theory is helpful but not strictly required; we introduce the necessary
categorical language as it arises.

The first chapter provides a rapid review of point-set topology, with an
emphasis on the concepts and constructions that are most relevant to the
algebraic theory: homotopy, retracts, CW complexes, and the homotopy category.
The second chapter introduces the fundamental group $\pi_1$, our first and most
classical algebraic invariant. We prove its basic properties, compute
$\pi_1(S^1) \cong \Z$ from first principles, and derive the two-dimensional
Brouwer fixed-point theorem as a striking application.

Subsequent chapters will develop covering space theory, the Seifert--van Kampen
theorem, singular homology, cohomology, and further topics. Throughout, we
emphasize rigorous proofs, concrete computations, and the interplay between
algebra and geometry.

\medskip

\noindent\textit{Prerequisites.} Point-set topology (at the level of Munkres,
Chapters~1--4) and abstract algebra (groups, rings, homomorphisms, quotient
groups).

\medskip

\noindent\textit{Recommended references.}
\begin{itemize}[leftmargin=2em]
\item A.~Hatcher, \textit{Algebraic Topology}, Cambridge University Press, 2002.
\item J.\,P.~May, \textit{A Concise Course in Algebraic Topology}, University
  of Chicago Press, 1999.
\item W.\,S.~Massey, \textit{A Basic Course in Algebraic Topology}, Springer, 1991.
\item G.\,E.~Bredon, \textit{Topology and Geometry}, Springer, 1993.
\item E.\,H.~Spanier, \textit{Algebraic Topology}, Springer, 1966.
\end{itemize}

%%%%%%%%%%%%%%%%%%%%%%%%%%%%%%%%%%%%%%%%%%%%%%%%%%%%%%%%%%%%%
%%
%%  CHAPTER 1 : TOPOLOGY REVIEW AND MOTIVATIONS
%%
%%%%%%%%%%%%%%%%%%%%%%%%%%%%%%%%%%%%%%%%%%%%%%%%%%%%%%%%%%%%%
\chapter{Topology Review and Motivations}\label{ch:review}

The purpose of this chapter is twofold. First, we recall the basic notions of
point-set topology that form the foundation for everything that follows.
Second, we introduce the language and the guiding philosophy of algebraic
topology, emphasizing the passage from topological spaces to algebraic
invariants.

%% ──────────────────────────────────────────────────────
\section{Topological Spaces and Continuous Maps}\label{sec:top-spaces}

We begin by recalling the fundamental definitions.

\begin{definition}[Topological space]\label{def:topological-space}
A \emph{topological space}\index{topological space} is a pair $(X, \tau)$ where
$X$ is a set and $\tau \subseteq \mathcal{P}(X)$ is a collection of subsets,
called \emph{open sets}, satisfying:
\begin{enumerate}[label=(\roman*)]
\item $\varnothing \in \tau$ and $X \in \tau$.
\item If $\{U_\alpha\}_{\alpha \in A} \subseteq \tau$, then
  $\bigcup_{\alpha \in A} U_\alpha \in \tau$.
\item If $U_1, \ldots, U_n \in \tau$, then $U_1 \cap \cdots \cap U_n \in \tau$.
\end{enumerate}
\end{definition}

\begin{definition}[Continuous map]\label{def:continuous-map}
A map $f \colon X \to Y$ between topological spaces is \emph{continuous}\index{continuous map}
if for every open set $V \subseteq Y$, the preimage $f^{-1}(V)$ is open in $X$.
\end{definition}

\begin{definition}[Homeomorphism]\label{def:homeomorphism}
A \emph{homeomorphism}\index{homeomorphism} is a continuous bijection
$f \colon X \to Y$ whose inverse $f^{-1} \colon Y \to X$ is also continuous.
If such a map exists, we write $X \cong Y$ and say that $X$ and $Y$ are
\emph{homeomorphic}.
\end{definition}

\begin{remark}\label{rem:homeo-equiv}
Homeomorphism is an equivalence relation on topological spaces. A topological
invariant is any property or quantity that is preserved under homeomorphism. One
of the central goals of topology is to find invariants that are powerful enough
to distinguish non-homeomorphic spaces.
\end{remark}

\begin{example}[Metric spaces]\label{ex:metric-spaces}
Every metric space $(X, d)$ carries a natural topology whose open sets are
unions of open balls $B(x, r) = \set{y \in X : d(x,y) < r}$. In particular,
$\R^n$ with the Euclidean metric, the spheres $S^n$, and the disks $D^n$ are
all topological spaces in this way.
\end{example}

\begin{example}[Quotient topology]\label{ex:quotient-topology}
Given a topological space $X$ and an equivalence relation $\sim$ on $X$, the
\emph{quotient space}\index{quotient space} $X/{\sim}$ carries the finest
topology making the projection $q \colon X \to X/{\sim}$ continuous. That is, a
set $U \subseteq X/{\sim}$ is open if and only if $q^{-1}(U)$ is open in $X$.
Many important spaces arise as quotients: for instance, the torus $T^2 \cong
\R^2 / \Z^2$, and the real projective plane $\RP^2 \cong S^2 / (x \sim -x)$.
\end{example}


%% ──────────────────────────────────────────────────────
\section{Connectedness and Path-Connectedness}\label{sec:connectedness}

\begin{definition}[Connected space]\label{def:connected}
A topological space $X$ is \emph{connected}\index{connected space} if it cannot
be written as a disjoint union $X = U \sqcup V$ of two nonempty open sets.
Equivalently, the only subsets of $X$ that are both open and closed are
$\varnothing$ and $X$ itself.
\end{definition}

\begin{definition}[Path-connected space]\label{def:path-connected}
A space $X$ is \emph{path-connected}\index{path-connected space} if for every
pair of points $x, y \in X$, there exists a continuous map $\gamma \colon I \to X$
(a \emph{path}) with $\gamma(0) = x$ and $\gamma(1) = y$.
\end{definition}

\begin{proposition}[Path-connected implies connected]\label{prop:path-connected-implies-connected}
Every path-connected space is connected.
\end{proposition}

\begin{proof}
Suppose $X$ is path-connected and $X = U \sqcup V$ with $U, V$ open and
nonempty. Pick $x \in U$ and $y \in V$, and let $\gamma \colon I \to X$ be a
path from $x$ to $y$. Then $I = \gamma^{-1}(U) \sqcup \gamma^{-1}(V)$ is a
decomposition of $I$ into two disjoint nonempty open sets, contradicting the
connectedness of $I$. Thus no such decomposition exists.
\end{proof}

\begin{example}[Connected but not path-connected]\label{ex:topologist-sine}
The \emph{topologist's sine curve}\index{topologist's sine curve}
\[
  X = \set{(x, \sin(1/x)) : x \in (0,1]} \cup \set{(0,y) : y \in [-1,1]}
\]
is connected but not path-connected. The closure of the graph of $\sin(1/x)$ is
connected (being the closure of a connected set), but there is no path in $X$
from $(0,0)$ to $(1, \sin 1)$.
\end{example}

\begin{center}
\begin{tikzpicture}[scale=3]
  \draw[->] (-0.1,0) -- (1.2,0) node[right] {$x$};
  \draw[->] (0,-1.3) -- (0,1.3) node[above] {$y$};
  \draw[thick,blue!70!black,domain=0.032:1,samples=800] plot (\x, {sin(1/\x r)});
  \draw[thick,red!70!black] (0,-1) -- (0,1);
  \node[below right] at (0.5,-1.2) {\small Topologist's sine curve};
\end{tikzpicture}
\captionof{figure}{The topologist's sine curve: connected but not path-connected.}\label{fig:topologist-sine}
\end{center}


%% ──────────────────────────────────────────────────────
\section{Compactness and the Hausdorff Property}\label{sec:compact-hausdorff}

\begin{definition}[Compact space]\label{def:compact}
A topological space $X$ is \emph{compact}\index{compact space} if every open
cover of $X$ admits a finite subcover: whenever $X = \bigcup_{\alpha \in A}
U_\alpha$ with each $U_\alpha$ open, there exist $\alpha_1, \ldots, \alpha_n
\in A$ such that $X = U_{\alpha_1} \cup \cdots \cup U_{\alpha_n}$.
\end{definition}

\begin{definition}[Hausdorff space]\label{def:hausdorff}
A space $X$ is \emph{Hausdorff}\index{Hausdorff space} (or $T_2$) if for any
two distinct points $x, y \in X$, there exist disjoint open sets $U, V$ with
$x \in U$ and $y \in V$.
\end{definition}

\begin{theorem}[Compact subsets of Hausdorff spaces are closed]\label{thm:compact-hausdorff-closed}
Let $X$ be a Hausdorff space and let $K \subseteq X$ be compact. Then $K$ is
closed in $X$.
\end{theorem}

\begin{proof}
We show that $X \setminus K$ is open. Let $x \in X \setminus K$. For each $y
\in K$, choose disjoint open sets $U_y \ni x$ and $V_y \ni y$. The collection
$\{V_y\}_{y \in K}$ covers $K$, so by compactness there exist $y_1, \ldots,
y_n$ with $K \subseteq V_{y_1} \cup \cdots \cup V_{y_n}$. Then $U = U_{y_1}
\cap \cdots \cap U_{y_n}$ is an open neighborhood of $x$ disjoint from $K$.
Hence $X \setminus K$ is open.
\end{proof}

\begin{proposition}[{Continuous image of a compact set is compact}]\label{prop:continuous-compact}
If $f \colon X \to Y$ is continuous and $X$ is compact, then $f(X)$ is compact.
\end{proposition}

\begin{proof}
Let $\{V_\alpha\}$ be an open cover of $f(X)$. Then $\{f^{-1}(V_\alpha)\}$ is
an open cover of $X$. By compactness of $X$, finitely many
$f^{-1}(V_{\alpha_1}), \ldots, f^{-1}(V_{\alpha_n})$ cover $X$, whence
$V_{\alpha_1}, \ldots, V_{\alpha_n}$ cover $f(X)$.
\end{proof}

\begin{corollary}[{Bijective continuous map from compact to Hausdorff is a homeomorphism}]\label{cor:compact-to-hausdorff}
If $f \colon X \to Y$ is a continuous bijection, $X$ is compact, and $Y$ is
Hausdorff, then $f$ is a homeomorphism.
\end{corollary}

\begin{proof}
It suffices to show that $f$ is a closed map. Let $C \subseteq X$ be closed.
Since $X$ is compact, $C$ is compact, so $f(C)$ is compact by
Proposition~\ref{prop:continuous-compact}, hence closed in $Y$ by
Theorem~\ref{thm:compact-hausdorff-closed}.
\end{proof}


%% ──────────────────────────────────────────────────────
\section{Homotopy and Homotopy Equivalence}\label{sec:homotopy}

The key insight of algebraic topology is that homeomorphism is too rigid a
notion of equivalence for many purposes. \emph{Homotopy equivalence} is a
coarser but more flexible relation that captures the essential ``shape'' of a
space.

\begin{definition}[Homotopy]\label{def:homotopy}
Let $f, g \colon X \to Y$ be continuous maps. A \emph{homotopy}\index{homotopy}
from $f$ to $g$ is a continuous map $H \colon X \times I \to Y$ such that
$H(x, 0) = f(x)$ and $H(x, 1) = g(x)$ for all $x \in X$. We write $f
\simeq g$ and say $f$ is \emph{homotopic} to $g$.
\end{definition}

\begin{center}
\begin{tikzpicture}[scale=1.2]
  \fill[blue!8] (0,0) rectangle (4,2.5);
  \draw[thick] (0,0) rectangle (4,2.5);
  \node[left] at (0,0) {$0$};
  \node[left] at (0,2.5) {$1$};
  \node[below] at (0,0) {$X$};
  \draw[-{Stealth},thick] (4.5,1.25) -- node[above] {$H$} (5.5,1.25);
  \draw[thick,blue!70!black] (0.2,0) .. controls (1.5,0.5) and (3,0.2) .. (3.8,0);
  \node[below,blue!70!black] at (2,-0.2) {$f = H(-,0)$};
  \draw[thick,red!70!black] (0.2,2.5) .. controls (1,2.0) and (3,2.8) .. (3.8,2.5);
  \node[above,red!70!black] at (2,2.7) {$g = H(-,1)$};
  \foreach \t in {0.3,0.6,0.9,1.2,1.5,1.8,2.1} {
    \draw[gray!50,thin] (0.2,\t) .. controls (1.2,\t+0.2) and (3,\t-0.1) .. (3.8,\t);
  }
  \begin{scope}[xshift=6cm,yshift=0.25cm]
    \draw[thick,fill=green!10] (0,0) ellipse (1.2 and 1);
    \node at (0,0) {$Y$};
  \end{scope}
\end{tikzpicture}
\captionof{figure}{A homotopy $H \colon X \times I \to Y$ deforming $f$ into $g$.}\label{fig:homotopy}
\end{center}

\begin{proposition}[Homotopy is an equivalence relation]\label{prop:homotopy-equiv-rel}
For any spaces $X$ and $Y$, homotopy $\simeq$ is an equivalence relation on the
set of continuous maps from $X$ to $Y$.
\end{proposition}

\begin{proof}
\emph{Reflexivity:} $H(x,t) = f(x)$ is a homotopy from $f$ to itself.

\emph{Symmetry:} If $H$ is a homotopy from $f$ to $g$, then $\bar{H}(x,t) =
H(x, 1-t)$ is a homotopy from $g$ to $f$.

\emph{Transitivity:} If $H$ is a homotopy from $f$ to $g$ and $K$ is a
homotopy from $g$ to $h$, define
\[
  L(x, t) = \begin{cases} H(x, 2t) & \text{if } 0 \le t \le \tfrac{1}{2}, \\
  K(x, 2t - 1) & \text{if } \tfrac{1}{2} \le t \le 1. \end{cases}
\]
Since $H(x,1) = g(x) = K(x,0)$, the map $L$ is well-defined and continuous by
the pasting lemma. It is a homotopy from $f$ to $h$.
\end{proof}

\begin{definition}[Homotopy equivalence]\label{def:homotopy-equivalence}
A continuous map $f \colon X \to Y$ is a \emph{homotopy
equivalence}\index{homotopy equivalence} if there exists a continuous map $g
\colon Y \to X$ such that $g \circ f \simeq \Id_X$ and $f \circ g \simeq
\Id_Y$. In this case, $g$ is called a \emph{homotopy inverse} of $f$, and we
write $X \simeq Y$.
\end{definition}

\begin{definition}[Contractible space]\label{def:contractible}
A space $X$ is \emph{contractible}\index{contractible space} if it is homotopy
equivalent to a point, i.e., $X \simeq \{\star\}$. Equivalently, the identity
map $\Id_X$ is homotopic to a constant map.
\end{definition}

\begin{example}[Contractible spaces]\label{ex:contractible}
The Euclidean space $\R^n$ is contractible: the homotopy $H(x,t) = (1-t)x$
contracts $\R^n$ to the origin. More generally, any convex subset of $\R^n$ is
contractible. The open disk, the closed disk $D^n$, and any star-shaped region
are all contractible.
\end{example}


%% ──────────────────────────────────────────────────────
\section{Retracts and Deformation Retracts}\label{sec:retracts}

\begin{definition}[Retraction and retract]\label{def:retraction}
Let $A \subseteq X$ be a subspace. A \emph{retraction}\index{retraction} of $X$
onto $A$ is a continuous map $r \colon X \to A$ such that $r(a) = a$ for all $a
\in A$ (equivalently, $r \circ \iota = \Id_A$ where $\iota \colon A
\hookrightarrow X$ is the inclusion). If such a map exists, $A$ is called a
\emph{retract}\index{retract} of $X$.
\end{definition}

\begin{definition}[Deformation retract]\label{def:deformation-retract}
A subspace $A \subseteq X$ is a \emph{deformation retract}\index{deformation
retract} of $X$ if there exists a retraction $r \colon X \to A$ such that
$\iota \circ r \simeq \Id_X$, i.e., the composition $X \xrightarrow{r} A
\xhookrightarrow{\iota} X$ is homotopic to the identity on $X$.

The retract is a \emph{strong deformation retract}\index{strong deformation
retract} if the homotopy $H \colon X \times I \to X$ from $\Id_X$ to $\iota
\circ r$ can be chosen so that $H(a, t) = a$ for all $a \in A$ and all $t \in
I$.
\end{definition}

\begin{example}[$S^{n-1}$ as a deformation retract of $\R^n \setminus \{0\}$]\label{ex:sphere-retract}
The map $r \colon \R^n \setminus \{0\} \to S^{n-1}$ defined by $r(x) = x /
\abs{x}$ is a retraction. The homotopy
\[
  H(x, t) = (1-t)x + t \frac{x}{\abs{x}} = \bigl(1 - t + t/\abs{x}\bigr) x
\]
shows that $S^{n-1}$ is a strong deformation retract of $\R^n \setminus \{0\}$.
In particular, $S^{n-1} \simeq \R^n \setminus \{0\}$.
\end{example}

\begin{center}
\begin{tikzpicture}[scale=1.5]
  \draw[fill=blue!5] (0,0) circle (1.8);
  \draw[thick,blue!70!black] (0,0) circle (1);
  \fill[white] (0,0) circle (0.06);
  \draw (0,0) circle (0.06);
  \node[below] at (0,-0.12) {\small $0$};
  \foreach \angle in {20,70,130,200,260,320} {
    \coordinate (P) at (\angle:1.5);
    \coordinate (Q) at (\angle:1);
    \draw[-{Stealth},red!60!black,thick] (P) -- (Q);
  }
  \node at (0,-2.2) {$r(x) = x/\abs{x}$};
  \node[blue!70!black] at (1.15,0.8) {$S^{n-1}$};
\end{tikzpicture}
\captionof{figure}{Radial retraction of $\R^n \setminus \{0\}$ onto $S^{n-1}$.}\label{fig:radial-retract}
\end{center}

\begin{proposition}[Deformation retract implies homotopy equivalence]\label{prop:def-retract-htpy-equiv}
If $A$ is a deformation retract of $X$, then $A \simeq X$.
\end{proposition}

\begin{proof}
The inclusion $\iota \colon A \hookrightarrow X$ and the retraction $r \colon X
\to A$ satisfy $r \circ \iota = \Id_A$ and $\iota \circ r \simeq \Id_X$, which
is precisely the definition of a homotopy equivalence.
\end{proof}


%% ──────────────────────────────────────────────────────
\section{CW Complexes}\label{sec:CW-complexes}

CW complexes form the most important class of spaces in algebraic topology. They
are built inductively by attaching cells of increasing dimension, providing a
combinatorial handle on their topology.

\begin{definition}[CW complex]\label{def:CW-complex}
A \emph{CW complex}\index{CW complex} is a topological space $X$ together with
a filtration
\[
  X^0 \subseteq X^1 \subseteq X^2 \subseteq \cdots \subseteq X = \bigcup_{n \ge 0} X^n,
\]
where:
\begin{enumerate}[label=(\roman*)]
\item $X^0$ is a discrete set of points (the \emph{$0$-cells}).
\item For each $n \ge 1$, the \emph{$n$-skeleton} $X^n$ is obtained from
  $X^{n-1}$ by attaching $n$-cells: for each cell $e^n_\alpha$, there is a
  continuous \emph{attaching map} $\varphi_\alpha \colon S^{n-1} \to X^{n-1}$
  and
  \[
    X^n = X^{n-1} \cup_{\{\varphi_\alpha\}} \bigsqcup_\alpha D^n_\alpha,
  \]
  where the $n$-disk $D^n_\alpha$ is glued to $X^{n-1}$ along its boundary via
  $\varphi_\alpha$.
\item $X$ has the \emph{weak topology}: a subset $A \subseteq X$ is closed if
  and only if $A \cap X^n$ is closed in $X^n$ for each $n$.
\end{enumerate}
The letters ``CW'' stand for \emph{Closure-finite} (each cell meets only
finitely many other cells) and \emph{Weak topology}.
\end{definition}

\begin{example}[CW structure on $S^n$]\label{ex:CW-sphere}
The sphere $S^n$ has a CW structure with exactly two cells: one $0$-cell $e^0$
(a point) and one $n$-cell $e^n$, attached via the map $\varphi \colon S^{n-1}
\to e^0$ that collapses the entire boundary to the point. This gives $S^n \cong
D^n / \partial D^n$.
\end{example}

\begin{example}[CW structure on the torus $T^2$]\label{ex:CW-torus}
The torus $T^2$ admits a CW structure with one $0$-cell $v$, two $1$-cells $a$
and $b$ (both loops based at $v$), and one $2$-cell $e^2$ attached via the word
$aba^{-1}b^{-1}$.
\end{example}

\begin{center}
\begin{tikzpicture}[scale=1.8]
  % Square representing the 2-cell
  \fill[blue!5] (0,0) rectangle (2,2);
  \draw[thick,red!70!black,-{Stealth}] (0,0) -- (1,0) node[below,midway] {$a$};
  \draw[thick,red!70!black,-{Stealth}] (1,0) -- (2,0);
  \draw[thick,red!70!black,-{Stealth}] (0,2) -- (1,2) node[above,midway] {$a$};
  \draw[thick,red!70!black,-{Stealth}] (1,2) -- (2,2);
  \draw[thick,blue!70!black,-{Stealth}] (0,0) -- (0,1) node[left,midway] {$b$};
  \draw[thick,blue!70!black,-{Stealth}] (0,1) -- (0,2);
  \draw[thick,blue!70!black,-{Stealth}] (2,0) -- (2,1) node[right,midway] {$b$};
  \draw[thick,blue!70!black,-{Stealth}] (2,1) -- (2,2);
  \fill (0,0) circle (1.2pt) node[below left] {$v$};
  \fill (2,0) circle (1.2pt) node[below right] {$v$};
  \fill (0,2) circle (1.2pt) node[above left] {$v$};
  \fill (2,2) circle (1.2pt) node[above right] {$v$};
  \node at (1,1) {$e^2$};
\end{tikzpicture}
\captionof{figure}{The standard CW structure on the torus: one vertex, two edges, one face.}\label{fig:CW-torus}
\end{center}

\begin{example}[Real projective space $\RP^n$]\label{ex:CW-RPn}
The real projective space $\RP^n$ has a CW structure with exactly one cell in
each dimension $0, 1, \ldots, n$. The $k$-cell is attached to $\RP^{k-1}$ via
the quotient map $S^{k-1} \to \RP^{k-1}$, which identifies antipodal points.
Thus
\[
  \RP^n = e^0 \cup e^1 \cup e^2 \cup \cdots \cup e^n.
\]
\end{example}

\begin{remark}\label{rem:CW-advantages}
CW complexes enjoy many pleasant properties. They are always Hausdorff and
paracompact. A compact subset of a CW complex meets only finitely many cells.
Moreover, any CW complex is homotopy equivalent to a space built from simplices
(a simplicial complex), but the CW decomposition is typically far more
economical: for instance, $S^n$ requires only $2$ cells as a CW complex but at
least $2(n+1)$ simplices as a simplicial complex.
\end{remark}


%% ──────────────────────────────────────────────────────
\section{The Homotopy Category}\label{sec:homotopy-category}

We briefly introduce the categorical perspective, which clarifies the
functorial nature of algebraic topology.

\begin{definition}[Category $\mathbf{Top}$]\label{def:category-Top}
The category $\mathbf{Top}$\index{category!$\mathbf{Top}$} has topological
spaces as objects and continuous maps as morphisms. Composition is the usual
composition of functions.
\end{definition}

\begin{definition}[Homotopy category $\mathbf{hTop}$]\label{def:category-hTop}
The \emph{homotopy category}\index{homotopy category} $\mathbf{hTop}$ has the
same objects as $\mathbf{Top}$, but the morphisms from $X$ to $Y$ are homotopy
classes $[X, Y]$ of continuous maps. Composition is induced by the composition
of representatives:
\[
  [g] \circ [f] = [g \circ f].
\]
This is well-defined because homotopy is compatible with composition.
\end{definition}

\begin{remark}\label{rem:hTop-isomorphisms}
The isomorphisms in $\mathbf{hTop}$ are precisely the homotopy equivalences.
Thus two spaces are ``the same'' in $\mathbf{hTop}$ if and only if they are
homotopy equivalent.
\end{remark}

A \emph{functor}\index{functor} $F \colon \mathbf{hTop} \to \mathbf{Grp}$
assigns a group $F(X)$ to each space $X$ and a group homomorphism $F([f])
\colon F(X) \to F(Y)$ to each homotopy class $[f]$, preserving identity and
composition. The fundamental group $\pi_1$ is the first example of such a
functor; homology groups $H_n$ provide further examples.

\begin{center}
\begin{tikzcd}[column sep=large,row sep=large]
X \ar[r, "f"] \ar[d, "g"'] & Y \ar[d, "h"] \\
Z \ar[r, "k"'] & W
\end{tikzcd}
\qquad
\begin{tikzcd}[column sep=large,row sep=large]
F(X) \ar[r, "F(f)"] \ar[d, "F(g)"'] & F(Y) \ar[d, "F(h)"] \\
F(Z) \ar[r, "F(k)"'] & F(W)
\end{tikzcd}
\captionof{figure}{A functor $F$ maps a commutative diagram in $\mathbf{hTop}$ to a commutative diagram in $\mathbf{Grp}$.}\label{fig:functor-diagram}
\end{center}


%% ──────────────────────────────────────────────────────
\section{Motivations for Algebraic Topology}\label{sec:motivations}

The fundamental question driving algebraic topology is:

\medskip
\begin{center}
\emph{When are two topological spaces ``essentially the same,'' and how can we
tell them apart?}
\end{center}
\medskip

Purely topological arguments can be remarkably difficult. Consider the
following questions:
\begin{enumerate}[label=(\arabic*)]
\item Is $\R^2$ homeomorphic to $\R^3$? (No---but this is surprisingly hard to
  prove without algebraic tools.)
\item Is $S^1$ a retract of $D^2$? (No---proved below using the fundamental group.)
\item Does every continuous map $f \colon D^2 \to D^2$ have a fixed point?
  (Yes---the Brouwer fixed-point theorem, proved in Chapter~\ref{ch:fundamental-group}.)
\item Can the $2$-sphere $S^2$ be combed? That is, does there exist a nonvanishing
  continuous tangent vector field on $S^2$? (No---the hairy ball theorem.)
\end{enumerate}

The strategy of algebraic topology is to associate to each space $X$ an algebraic
object---a group, a ring, a module---in such a way that continuous maps induce
homomorphisms, and homotopy equivalent spaces receive isomorphic invariants. If
two spaces have non-isomorphic invariants, they cannot be homotopy equivalent
(and a fortiori cannot be homeomorphic).

\begin{center}
\begin{tikzpicture}[
    box/.style={draw, rounded corners, minimum width=3cm, minimum height=1cm, align=center},
    arr/.style={-{Stealth}, thick}
  ]
  \node[box, fill=blue!8] (top) at (0,0) {Topological\\problem};
  \node[box, fill=red!8] (alg) at (6,0) {Algebraic\\problem};
  \node[box, fill=green!8] (sol) at (6,-2.5) {Algebraic\\solution};
  \node[box, fill=orange!8] (tsol) at (0,-2.5) {Topological\\conclusion};
  \draw[arr] (top) -- node[above] {\small functor} (alg);
  \draw[arr] (alg) -- node[right] {\small computation} (sol);
  \draw[arr] (sol) -- node[below] {\small interpretation} (tsol);
  \draw[arr,dashed,gray] (top) -- node[left] {\small \textcolor{gray}{direct?}} (tsol);
\end{tikzpicture}
\captionof{figure}{The algebraic topology paradigm: translate, compute, interpret.}\label{fig:AT-paradigm}
\end{center}

The main invariants we will study are:
\begin{itemize}
\item The \textbf{fundamental group} $\pi_1(X, x_0)$: captures information about
  loops in $X$ up to continuous deformation.
\item \textbf{Higher homotopy groups} $\pi_n(X, x_0)$: higher-dimensional
  analogues of $\pi_1$.
\item \textbf{Singular homology} $H_n(X; G)$: abelian groups measuring
  $n$-dimensional ``holes'' in $X$.
\item \textbf{Singular cohomology} $H^n(X; G)$: the dual theory, with a rich
  multiplicative structure (the cup product).
\end{itemize}


%% ──────────────────────────────────────────────────────
\section{Exercises}\label{sec:ch1-exercises}

\begin{exercise}\label{exo:ch1-1}
Let $X$ be a topological space with the discrete topology (every subset is
open). Show that $X$ is connected if and only if $\abs{X} \le 1$.
\end{exercise}

\begin{exercise}\label{exo:ch1-2}
Let $f \colon X \to Y$ be continuous, and suppose $X$ is connected. Show that
$f(X)$ is connected.
\end{exercise}

\begin{exercise}\label{exo:ch1-3}
Prove that $S^n$ is path-connected for all $n \ge 1$.
\emph{Hint:} For $n = 1$, parameterize the circle. For $n \ge 2$, use the fact
that $S^n \setminus \{p\} \cong \R^n$ (stereographic projection) and that $\R^n$
is path-connected.
\end{exercise}

\begin{exercise}\label{exo:ch1-4}
Show that the product $X \times Y$ is connected (resp.\ path-connected) if and
only if both $X$ and $Y$ are connected (resp.\ path-connected).
\end{exercise}

\begin{exercise}\label{exo:ch1-5}
Let $X$ and $Y$ be homotopy equivalent. Show that if $X$ is path-connected,
then so is $Y$.
\end{exercise}

\begin{exercise}\label{exo:ch1-6}
Show that the M\"{o}bius band is homotopy equivalent to $S^1$.
\emph{Hint:} Exhibit the central circle as a deformation retract.
\end{exercise}

\begin{exercise}\label{exo:ch1-7}
Describe a CW structure on the Klein bottle $K$ with one $0$-cell, two
$1$-cells, and one $2$-cell. What is the attaching word for the $2$-cell?
\end{exercise}

\begin{exercise}\label{exo:ch1-8}
Show that the punctured plane $\R^2 \setminus \{0\}$ is homotopy equivalent to
$S^1$, and that $\R^3 \setminus \{0\}$ is homotopy equivalent to $S^2$.
Generalize to $\R^n \setminus \{0\} \simeq S^{n-1}$.
\end{exercise}

\begin{exercise}\label{exo:ch1-9}
Let $X = S^1 \vee S^1$ be the wedge (one-point union) of two circles. Describe
a CW structure on $X$ and compute the Euler characteristic $\chi(X) = c_0 - c_1
+ c_2 - \cdots$ where $c_k$ is the number of $k$-cells.
\end{exercise}

\begin{exercise}[{Mapping cylinder}]\label{exo:ch1-10}
Given a continuous map $f \colon X \to Y$, the \emph{mapping cylinder} is the
quotient space
\[
  M_f = \bigl( (X \times I) \sqcup Y \bigr) / \bigl( (x,1) \sim f(x) \bigr).
\]
Show that $Y$ is a strong deformation retract of $M_f$. Conclude that any
continuous map is homotopic to an inclusion, up to homotopy equivalence.
\end{exercise}


%%%%%%%%%%%%%%%%%%%%%%%%%%%%%%%%%%%%%%%%%%%%%%%%%%%%%%%%%%%%%
%%
%%  CHAPTER 2 : THE FUNDAMENTAL GROUP
%%
%%%%%%%%%%%%%%%%%%%%%%%%%%%%%%%%%%%%%%%%%%%%%%%%%%%%%%%%%%%%%
\chapter{The Fundamental Group --- Definition and First Properties}\label{ch:fundamental-group}

The fundamental group is the oldest and most intuitive invariant of algebraic
topology, introduced by Henri Poincar\'{e} in 1895. It captures the essential
information about loops in a space up to continuous deformation.

%% ──────────────────────────────────────────────────────
\section{Paths and Loops}\label{sec:paths-loops}

Throughout this chapter, $I = [0,1]$ denotes the unit interval.

\begin{definition}[Path]\label{def:path}
A \emph{path}\index{path} in a topological space $X$ is a continuous map
$\gamma \colon I \to X$. The point $\gamma(0)$ is the \emph{initial point} and
$\gamma(1)$ is the \emph{terminal point}. We say $\gamma$ is a path \emph{from}
$\gamma(0)$ \emph{to} $\gamma(1)$.
\end{definition}

\begin{definition}[Loop]\label{def:loop}
A \emph{loop}\index{loop} in $X$ based at a point $x_0 \in X$ is a path
$\gamma \colon I \to X$ with $\gamma(0) = \gamma(1) = x_0$. The point $x_0$ is
called the \emph{basepoint}.
\end{definition}

\begin{definition}[Constant path]\label{def:constant-path}
For any $x_0 \in X$, the \emph{constant path} at $x_0$ is the map $c_{x_0}
\colon I \to X$ defined by $c_{x_0}(t) = x_0$ for all $t$.
\end{definition}

\begin{definition}[Reverse path]\label{def:reverse-path}
Given a path $\gamma \colon I \to X$, its \emph{reverse}\index{reverse path} is
the path $\bar{\gamma} \colon I \to X$ defined by $\bar{\gamma}(t) =
\gamma(1-t)$.
\end{definition}

\begin{definition}[Concatenation of paths]\label{def:concatenation}
Given paths $\gamma, \delta \colon I \to X$ with $\gamma(1) = \delta(0)$, the
\emph{concatenation}\index{concatenation of paths} $\gamma \star \delta \colon I
\to X$ is defined by
\[
  (\gamma \star \delta)(t) = \begin{cases}
    \gamma(2t) & \text{if } 0 \le t \le \tfrac{1}{2}, \\
    \delta(2t - 1) & \text{if } \tfrac{1}{2} \le t \le 1.
  \end{cases}
\]
This is continuous by the pasting lemma, since $\gamma(1) = \delta(0)$.
\end{definition}

\begin{center}
\begin{tikzpicture}[scale=1.3]
  \draw[thick,blue!70!black,-{Stealth}] (0,0) .. controls (0.5,1.5) and (1.5,1) .. (2.5,0.8)
    node[midway,above left] {$\gamma$};
  \draw[thick,red!70!black,-{Stealth}] (2.5,0.8) .. controls (3.5,0.5) and (4,1.5) .. (5,0)
    node[midway,above right] {$\delta$};
  \fill (0,0) circle (1.5pt) node[below] {$x_0$};
  \fill (2.5,0.8) circle (1.5pt) node[above] {$x_1$};
  \fill (5,0) circle (1.5pt) node[below] {$x_2$};
  \node at (2.5,-0.8) {$\gamma \star \delta$: traverse $\gamma$ then $\delta$, at double speed};
\end{tikzpicture}
\captionof{figure}{Concatenation of paths $\gamma$ (from $x_0$ to $x_1$) and $\delta$ (from $x_1$ to $x_2$).}\label{fig:concatenation}
\end{center}


%% ──────────────────────────────────────────────────────
\section{Path Homotopy}\label{sec:path-homotopy}

We now introduce the key equivalence relation on paths.

\begin{definition}[Path homotopy]\label{def:path-homotopy}
Two paths $\gamma_0, \gamma_1 \colon I \to X$ with the same initial point
$\gamma_0(0) = \gamma_1(0) = x_0$ and the same terminal point $\gamma_0(1) =
\gamma_1(1) = x_1$ are \emph{path-homotopic}\index{path homotopy} if there
exists a continuous map $H \colon I \times I \to X$ such that:
\begin{enumerate}[label=(\roman*)]
\item $H(s, 0) = \gamma_0(s)$ and $H(s, 1) = \gamma_1(s)$ for all $s \in I$,
\item $H(0, t) = x_0$ and $H(1, t) = x_1$ for all $t \in I$.
\end{enumerate}
We write $\gamma_0 \simeq \gamma_1 \rel{\partial I}$ and denote the homotopy
class of $\gamma$ by $[\gamma]$.
\end{definition}

Condition (ii) is what distinguishes path homotopy from free homotopy: the
endpoints are held fixed throughout the deformation.

\begin{center}
\begin{tikzpicture}[scale=1.4]
  \fill[purple!6] (0,0) .. controls (0.5,1.5) and (2,2) .. (4,0.5)
    -- (4,0.5) .. controls (2,-1) and (0.5,-0.5) .. (0,0);
  \draw[thick,blue!70!black] (0,0) .. controls (0.5,1.5) and (2,2) .. (4,0.5)
    node[midway,above] {$\gamma_0$};
  \draw[thick,red!70!black] (0,0) .. controls (0.5,-0.5) and (2,-1) .. (4,0.5)
    node[midway,below] {$\gamma_1$};
  \foreach \t in {0.2,0.4,0.6,0.8} {
    \draw[gray!40,thin] (0,0) .. controls ({0.5},{1.5-2*\t}) and ({2},{2-3*\t}) .. (4,0.5);
  }
  \fill (0,0) circle (1.5pt) node[left] {$x_0$};
  \fill (4,0.5) circle (1.5pt) node[right] {$x_1$};
  \node[purple!70!black] at (2,0.3) {\small $H$};
\end{tikzpicture}
\captionof{figure}{A path homotopy $H$ deforming $\gamma_0$ into $\gamma_1$ with fixed endpoints.}\label{fig:path-homotopy}
\end{center}

\begin{proposition}[Path homotopy is an equivalence relation]\label{prop:path-homotopy-equiv}
Path homotopy is an equivalence relation on the set of paths from $x_0$ to $x_1$.
\end{proposition}

\begin{proof}
The proof is analogous to that of Proposition~\ref{prop:homotopy-equiv-rel}.
Reflexivity uses $H(s,t) = \gamma(s)$. Symmetry uses $\bar{H}(s,t) = H(s,
1-t)$. Transitivity uses the concatenation of homotopies along the $t$-parameter
with the pasting lemma.
\end{proof}


%% ──────────────────────────────────────────────────────
\section{The Fundamental Group}\label{sec:fundamental-group-def}

\begin{theorem}[{The fundamental group $\pi_1(X, x_0)$}]\label{thm:fundamental-group}
Let $X$ be a topological space and $x_0 \in X$. The set of path-homotopy
classes of loops based at $x_0$, equipped with the operation
\[
  [\gamma] \cdot [\delta] = [\gamma \star \delta],
\]
is a group, denoted $\pi_1(X, x_0)$\index{fundamental group} and called the
\emph{fundamental group} of $X$ at the basepoint $x_0$.
\end{theorem}

\begin{proof}
We must verify well-definedness, associativity, the existence of an identity
element, and the existence of inverses.

\textbf{Well-definedness.} We must show that if $\gamma_0 \simeq \gamma_1
\rel{\partial I}$ and $\delta_0 \simeq \delta_1 \rel{\partial I}$, then
$\gamma_0 \star \delta_0 \simeq \gamma_1 \star \delta_1 \rel{\partial I}$. Let
$H$ be a homotopy from $\gamma_0$ to $\gamma_1$ and $K$ a homotopy from
$\delta_0$ to $\delta_1$. Define
\[
  L(s, t) = \begin{cases}
    H(2s, t) & \text{if } 0 \le s \le \tfrac{1}{2}, \\
    K(2s-1, t) & \text{if } \tfrac{1}{2} \le s \le 1.
  \end{cases}
\]
Then $L$ is a well-defined path homotopy from $\gamma_0 \star \delta_0$ to
$\gamma_1 \star \delta_1$.

\textbf{Associativity.} We need $(\gamma \star \delta) \star \varepsilon \simeq
\gamma \star (\delta \star \varepsilon) \rel{\partial I}$. Define
\[
  H(s, t) = \begin{cases}
    \gamma\!\left(\dfrac{4s}{1+t}\right) & \text{if } 0 \le s \le \dfrac{1+t}{4}, \\[8pt]
    \delta(4s - 1 - t) & \text{if } \dfrac{1+t}{4} \le s \le \dfrac{2+t}{4}, \\[8pt]
    \varepsilon\!\left(\dfrac{4s - 2 - t}{2 - t}\right) & \text{if } \dfrac{2+t}{4} \le s \le 1.
  \end{cases}
\]
One checks that $H$ is continuous (by the pasting lemma applied to the three
pieces), that $H(s,0) = ((\gamma \star \delta) \star \varepsilon)(s)$, that
$H(s,1) = (\gamma \star (\delta \star \varepsilon))(s)$, and that $H(0,t) =
x_0$ and $H(1,t) = \varepsilon(1)$ for all $t$. The key idea is that the
reparameterization amounts to changing where the ``time'' is spent along each of
the three paths.

\textbf{Identity.} The constant loop $c_{x_0}$ serves as the identity: we claim
$\gamma \star c_{x_0} \simeq \gamma \simeq c_{x_0} \star \gamma \rel{\partial
I}$. For the first, define
\[
  H(s, t) = \begin{cases}
    \gamma\!\left(\dfrac{2s}{1+t}\right) & \text{if } 0 \le s \le \dfrac{1+t}{2}, \\[6pt]
    x_0 & \text{if } \dfrac{1+t}{2} \le s \le 1.
  \end{cases}
\]
Then $H(s,0) = (\gamma \star c_{x_0})(s)$ and $H(s,1) = \gamma(s)$. The other
case is similar.

\textbf{Inverses.} For any loop $\gamma$ based at $x_0$, we claim $[\gamma
\star \bar{\gamma}] = [c_{x_0}]$. Define
\[
  H(s, t) = \begin{cases}
    \gamma(2s) & \text{if } 0 \le s \le \tfrac{t}{2}, \\
    \gamma(t) & \text{if } \tfrac{t}{2} \le s \le 1 - \tfrac{t}{2}, \\
    \gamma(2 - 2s) & \text{if } 1 - \tfrac{t}{2} \le s \le 1.
  \end{cases}
\]
At $t = 0$, this is the constant path $c_{x_0}$. At $t = 1$, this is $\gamma
\star \bar{\gamma}$. Thus $\gamma \star \bar{\gamma} \simeq c_{x_0}$. A similar
construction shows $\bar{\gamma} \star \gamma \simeq c_{x_0}$.
\end{proof}

\begin{remark}\label{rem:pi1-pointed}
The fundamental group $\pi_1(X, x_0)$ depends on the choice of basepoint $x_0$.
However, as we will see shortly, for path-connected spaces the groups at
different basepoints are isomorphic (though not canonically so).
\end{remark}

\begin{notation}\label{not:pi1}
When $X$ is path-connected and the basepoint is unimportant (e.g., when stating
that the group is trivial or abelian), we may write $\pi_1(X)$ instead of
$\pi_1(X, x_0)$.
\end{notation}

\begin{definition}[Simply connected space]\label{def:simply-connected}
A path-connected space $X$ is \emph{simply connected}\index{simply connected} if
$\pi_1(X, x_0) = \{e\}$ for some (hence every) basepoint $x_0 \in X$. Equivalently,
$X$ is simply connected if it is path-connected and every loop in $X$ is
path-homotopic to a constant loop.
\end{definition}


%% ──────────────────────────────────────────────────────
\section{Functoriality of $\pi_1$}\label{sec:functoriality}

The fundamental group is not merely an assignment of a group to each pointed
space; it is a \emph{functor} from the category of pointed topological spaces to
the category of groups.

\begin{definition}[Pointed space and pointed map]\label{def:pointed-space}
A \emph{pointed topological space}\index{pointed space} is a pair $(X, x_0)$
where $X$ is a topological space and $x_0 \in X$ is a chosen basepoint. A
\emph{pointed map} $f \colon (X, x_0) \to (Y, y_0)$ is a continuous map
$f \colon X \to Y$ with $f(x_0) = y_0$.
\end{definition}

\begin{theorem}[Functoriality of $\pi_1$]\label{thm:pi1-functor}
Let $f \colon (X, x_0) \to (Y, y_0)$ be a pointed map. The assignment
\[
  f_* \colon \pi_1(X, x_0) \to \pi_1(Y, y_0), \qquad f_*([\gamma]) = [f \circ \gamma],
\]
is a well-defined group homomorphism, called the \emph{induced
homomorphism}\index{induced homomorphism}. Moreover:
\begin{enumerate}[label=(\roman*)]
\item $(\Id_X)_* = \Id_{\pi_1(X, x_0)}$.
\item $(g \circ f)_* = g_* \circ f_*$.
\end{enumerate}
In other words, $\pi_1$ is a (covariant) functor from the category of pointed
topological spaces to the category of groups.
\end{theorem}

\begin{proof}
\emph{Well-definedness.} If $\gamma_0 \simeq \gamma_1 \rel{\partial I}$ via a
homotopy $H$, then $f \circ H$ is a path homotopy from $f \circ \gamma_0$ to
$f \circ \gamma_1$.

\emph{Homomorphism.} We have $f \circ (\gamma \star \delta) = (f \circ \gamma)
\star (f \circ \delta)$, since
\[
  (f \circ (\gamma \star \delta))(t) = f((\gamma \star \delta)(t)) =
  \begin{cases}
    f(\gamma(2t)) & \text{if } t \le \tfrac{1}{2}, \\
    f(\delta(2t-1)) & \text{if } t \ge \tfrac{1}{2},
  \end{cases}
  = ((f \circ \gamma) \star (f \circ \delta))(t).
\]
Hence $f_*([\gamma] \cdot [\delta]) = f_*([\gamma]) \cdot f_*([\delta])$.

Properties (i) and (ii) are immediate from the definitions.
\end{proof}

The functoriality is neatly expressed by the following commutative diagram. If
$f \colon (X, x_0) \to (Y, y_0)$ and $g \colon (Y, y_0) \to (Z, z_0)$ are
pointed maps, then:

\begin{center}
\begin{tikzcd}[column sep=large, row sep=large]
(X, x_0) \ar[r, "f"] \ar[dr, "g \circ f"'] & (Y, y_0) \ar[d, "g"] \\
  & (Z, z_0)
\end{tikzcd}
\qquad$\xmapsto{\quad \pi_1 \quad}$\qquad
\begin{tikzcd}[column sep=large, row sep=large]
\pi_1(X, x_0) \ar[r, "f_*"] \ar[dr, "(g \circ f)_*"'] & \pi_1(Y, y_0) \ar[d, "g_*"] \\
  & \pi_1(Z, z_0)
\end{tikzcd}
\captionof{figure}{Functoriality of $\pi_1$: the functor preserves composition.}\label{fig:pi1-functor}
\end{center}


%% ──────────────────────────────────────────────────────
\section{Change of Basepoint}\label{sec:change-basepoint}

\begin{theorem}[Change of basepoint isomorphism]\label{thm:change-basepoint}
Let $X$ be a topological space and let $\alpha \colon I \to X$ be a path from
$x_0$ to $x_1$. The map
\[
  \beta_\alpha \colon \pi_1(X, x_0) \to \pi_1(X, x_1), \qquad \beta_\alpha([\gamma]) = [\bar{\alpha} \star \gamma \star \alpha],
\]
is a group isomorphism.
\end{theorem}

\begin{proof}
First, $\beta_\alpha$ is well-defined: if $\gamma \simeq \gamma' \rel{\partial
I}$, then $\bar{\alpha} \star \gamma \star \alpha \simeq \bar{\alpha} \star
\gamma' \star \alpha \rel{\partial I}$.

Next, $\beta_\alpha$ is a homomorphism:
\begin{align*}
  \beta_\alpha([\gamma] \cdot [\delta])
  &= [\bar{\alpha} \star (\gamma \star \delta) \star \alpha] \\
  &= [\bar{\alpha} \star \gamma \star (c_{x_0}) \star \delta \star \alpha] \\
  &\simeq [\bar{\alpha} \star \gamma \star (\alpha \star \bar{\alpha}) \star \delta \star \alpha] \\
  &= [(\bar{\alpha} \star \gamma \star \alpha) \star (\bar{\alpha} \star \delta \star \alpha)] \\
  &= \beta_\alpha([\gamma]) \cdot \beta_\alpha([\delta]).
\end{align*}

Finally, $\beta_{\bar{\alpha}} \colon \pi_1(X, x_1) \to \pi_1(X, x_0)$ is
the inverse of $\beta_\alpha$:
\[
  \beta_{\bar{\alpha}} \circ \beta_\alpha([\gamma])
  = [\alpha \star (\bar{\alpha} \star \gamma \star \alpha) \star \bar{\alpha}]
  = [(\alpha \star \bar{\alpha}) \star \gamma \star (\alpha \star \bar{\alpha})]
  = [c_{x_0} \star \gamma \star c_{x_0}]
  = [\gamma]. \qedhere
\]
\end{proof}

\begin{center}
\begin{tikzpicture}[scale=1.4]
  \fill (0,0) circle (1.5pt) node[below] {$x_0$};
  \fill (3,0) circle (1.5pt) node[below] {$x_1$};
  \draw[thick,green!50!black,-{Stealth}] (0,0) .. controls (1,0.5) and (2,0.5) .. (3,0)
    node[midway,above] {$\alpha$};
  \draw[thick,blue!70!black,{Stealth}-{Stealth}] (0,0) .. controls (-0.5,1.5) and (0.5,2) .. (0,0)
    node[left,midway] {$\gamma$};
  \draw[thick,red!70!black,{Stealth}-{Stealth}] (3,0) .. controls (2.5,1.5) and (3.5,2) .. (3,0)
    node[right,midway] {$\bar\alpha\star\gamma\star\alpha$};
\end{tikzpicture}
\captionof{figure}{Change of basepoint: a loop $\gamma$ at $x_0$ is conjugated by the path $\alpha$ to yield a loop at $x_1$.}\label{fig:change-basepoint}
\end{center}

\begin{corollary}\label{cor:pi1-path-connected}
If $X$ is path-connected, then $\pi_1(X, x_0) \cong \pi_1(X, x_1)$ for any two
points $x_0, x_1 \in X$.
\end{corollary}

\begin{remark}[Non-canonical isomorphism]\label{rem:non-canonical}
The isomorphism $\beta_\alpha$ depends on the choice of path $\alpha$: different
paths from $x_0$ to $x_1$ may yield different isomorphisms. The isomorphism is
canonical (independent of $\alpha$) if and only if $\pi_1(X, x_0)$ is abelian.
Indeed, if $\alpha$ and $\alpha'$ are two paths from $x_0$ to $x_1$, then
$\beta_{\alpha'} \circ \beta_\alpha^{-1}$ is conjugation by the loop $\alpha
\star \bar{\alpha}' \in \pi_1(X, x_1)$.
\end{remark}


%% ──────────────────────────────────────────────────────
\section{Homotopy Invariance}\label{sec:homotopy-invariance}

The fundamental group is a homotopy invariant: homotopy equivalent spaces have
isomorphic fundamental groups. This is one of the most important properties of
$\pi_1$, and we give a complete proof.

\begin{theorem}[{$\pi_1$ is a homotopy invariant}]\label{thm:pi1-homotopy-invariant}
Let $f \colon (X, x_0) \to (Y, y_0)$ be a homotopy equivalence. Then
\[
  f_* \colon \pi_1(X, x_0) \to \pi_1(Y, y_0)
\]
is a group isomorphism.
\end{theorem}

To prove this, we first establish a key lemma relating the induced maps of
homotopic maps.

\begin{lemma}[Homotopic maps induce conjugate homomorphisms]\label{lem:homotopic-maps-pi1}
Let $f, g \colon (X, x_0) \to Y$ be continuous maps (not necessarily pointed),
and let $H \colon X \times I \to Y$ be a homotopy from $f$ to $g$. Define the
path $\alpha(t) = H(x_0, t)$ from $f(x_0)$ to $g(x_0)$. Then for every loop
$\gamma$ based at $x_0$:
\[
  g_*([\gamma]) = [\bar{\alpha}] \cdot f_*([\gamma]) \cdot [\alpha]
  \quad \text{in } \pi_1(Y, g(x_0)).
\]
That is, $g_* = \beta_{\bar{\alpha}} \circ f_*$ where $\beta_{\bar{\alpha}}$ is
the change-of-basepoint isomorphism.
\end{lemma}

\begin{proof}
Consider the map $\Gamma \colon I \times I \to Y$ defined by $\Gamma(s, t) =
H(\gamma(s), t)$. The four boundary segments of $I \times I$ map as follows:
\begin{itemize}
\item Bottom $(t = 0)$: $\Gamma(s, 0) = f(\gamma(s))$, the loop $f \circ
  \gamma$.
\item Top $(t = 1)$: $\Gamma(s, 1) = g(\gamma(s))$, the loop $g \circ \gamma$.
\item Left $(s = 0)$: $\Gamma(0, t) = H(x_0, t) = \alpha(t)$, the path $\alpha$.
\item Right $(s = 1)$: $\Gamma(1, t) = H(x_0, t) = \alpha(t)$, the path $\alpha$.
\end{itemize}
By reparameterizing the boundary of the square, we obtain a path homotopy
$\alpha \star (f \circ \gamma) \star \bar{\alpha} \simeq g \circ \gamma
\rel{\partial I}$. Hence $[g \circ \gamma] = [\alpha] \cdot [f \circ \gamma]
\cdot [\bar{\alpha}]$, i.e., $g_*([\gamma]) = [\alpha] \cdot f_*([\gamma])
\cdot [\bar{\alpha}]$.

Equivalently, $[\bar{\alpha}] \cdot g_*([\gamma]) \cdot [\alpha] =
f_*([\gamma])$, which gives $g_* = \beta_{\bar{\alpha}} \circ f_*$ with the
convention stated.
\end{proof}

\begin{proof}[Proof of Theorem~\ref{thm:pi1-homotopy-invariant}]
Let $g \colon Y \to X$ be a homotopy inverse to $f$, so $g \circ f \simeq
\Id_X$ and $f \circ g \simeq \Id_Y$.

\emph{Step 1: $g_* \circ f_*$ is an isomorphism.} By
Lemma~\ref{lem:homotopic-maps-pi1} applied to the homotopy $g \circ f \simeq
\Id_X$, there is a path $\alpha$ in $X$ from $g(f(x_0))$ to $x_0$ such that
\[
  (\Id_X)_* = \beta_{\bar{\alpha}} \circ (g \circ f)_* = \beta_{\bar{\alpha}} \circ g_* \circ f_*.
\]
Since $(\Id_X)_* = \Id_{\pi_1(X, x_0)}$ and $\beta_{\bar{\alpha}}$ is an
isomorphism, $g_* \circ f_*$ is an isomorphism from $\pi_1(X, x_0)$ to
$\pi_1(X, g(f(x_0)))$.

\emph{Step 2: $f_* \circ g_*$ is an isomorphism.} Similarly, the homotopy $f
\circ g \simeq \Id_Y$ yields a path $\beta$ in $Y$ such that $f_* \circ g_* =
\beta_\beta$, which is an isomorphism.

\emph{Step 3: $f_*$ is an isomorphism.} Since $g_* \circ f_*$ is injective (as
a composition equal to an isomorphism up to conjugation), $f_*$ is injective.
Since $f_* \circ g_*$ is surjective, $f_*$ is surjective. Hence $f_*$ is a
bijective homomorphism, i.e., an isomorphism.
\end{proof}

\begin{corollary}\label{cor:contractible-pi1-trivial}
If $X$ is contractible, then $\pi_1(X, x_0)$ is trivial for every $x_0 \in X$.
\end{corollary}

\begin{proof}
A contractible space is homotopy equivalent to a point, and $\pi_1(\{\star\}) =
\{e\}$.
\end{proof}

\begin{corollary}[$\pi_1(\R^n) = 0$]\label{cor:pi1-Rn}
For all $n \ge 1$, $\pi_1(\R^n, x_0) = \{e\}$ for any basepoint $x_0$.
\end{corollary}

\begin{proof}
$\R^n$ is contractible (via the straight-line homotopy to the origin), so the
result follows from Corollary~\ref{cor:contractible-pi1-trivial}.
\end{proof}


%% ──────────────────────────────────────────────────────
\section{The Fundamental Group of the Circle}\label{sec:pi1-S1}

The computation of $\pi_1(S^1) \cong \Z$ is one of the cornerstones of
algebraic topology. We present the classical proof using lifting theory for the
covering map $p \colon \R \to S^1$.

\begin{definition}[Covering map $p \colon \R \to S^1$]\label{def:covering-R-S1}
Define $p \colon \R \to S^1$ by $p(t) = (\cos 2\pi t, \sin 2\pi t)$, viewing
$S^1 \subseteq \R^2$. This is a covering map: every point of $S^1$ has a
neighborhood $U$ such that $p^{-1}(U)$ is a disjoint union of open sets, each
mapped homeomorphically onto $U$ by $p$.
\end{definition}

\begin{center}
\begin{tikzpicture}[scale=1]
  % The real line
  \draw[thick,-{Stealth}] (-3.5,2) -- (3.5,2);
  \foreach \k in {-3,-2,-1,0,1,2,3} {
    \fill (\k,2) circle (1.5pt);
    \node[above] at (\k,2) {\small $\k$};
  }
  \node[left] at (-3.5,2) {$\R$};
  % The circle
  \draw[thick] (0,-0.5) circle (1);
  \fill (1,-0.5) circle (1.5pt) node[right] {\small $1$};
  \node[left] at (-1.3,-0.5) {$S^1$};
  % The map p
  \draw[-{Stealth},thick,blue!70!black] (0,1.6) -- node[right] {$p$} (0,0.7);
\end{tikzpicture}
\captionof{figure}{The covering map $p \colon \R \to S^1$, $t \mapsto e^{2\pi i t}$.}\label{fig:covering-S1}
\end{center}

\begin{lemma}[Path lifting]\label{lem:path-lifting}
Let $\gamma \colon I \to S^1$ be a path with $\gamma(0) = x_0 \in S^1$, and
let $\tilde{x}_0 \in p^{-1}(x_0)$. Then there exists a unique path
$\tilde{\gamma} \colon I \to \R$ such that $p \circ \tilde{\gamma} = \gamma$
and $\tilde{\gamma}(0) = \tilde{x}_0$.
\end{lemma}

\begin{proof}
\emph{Existence.} Cover $S^1$ by open sets $\{U_j\}$ that are evenly covered by
$p$. Since $I$ is compact, we can find a partition $0 = t_0 < t_1 < \cdots < t_m
= 1$ such that each $\gamma([t_{k-1}, t_k])$ lies in some $U_{j_k}$. We lift
$\gamma$ inductively on each subinterval $[t_{k-1}, t_k]$: since $U_{j_k}$ is
evenly covered, $p^{-1}(U_{j_k}) = \bigsqcup_\ell \tilde{U}_{j_k,\ell}$ with
$p|_{\tilde{U}_{j_k,\ell}}$ a homeomorphism. The lift of $\gamma|_{[t_{k-1},
t_k]}$ is determined by choosing the sheet containing $\tilde{\gamma}(t_{k-1})$.

\emph{Uniqueness.} If $\tilde{\gamma}$ and $\tilde{\gamma}'$ are two lifts with
the same initial point, then $\{t \in I : \tilde{\gamma}(t) =
\tilde{\gamma}'(t)\}$ is both open (by local homeomorphism) and closed (by
continuity), hence is all of $I$ since it contains $0$ and $I$ is connected.
\end{proof}

\begin{lemma}[Homotopy lifting]\label{lem:homotopy-lifting}
Let $H \colon I \times I \to S^1$ be a path homotopy (i.e., $H(0,t) = x_0$ and
$H(1,t) = x_1$ for all $t$), and let $\tilde{x}_0 \in p^{-1}(x_0)$. Then there
exists a unique lift $\tilde{H} \colon I \times I \to \R$ with $p \circ
\tilde{H} = H$ and $\tilde{H}(0,0) = \tilde{x}_0$. Moreover, $\tilde{H}$ is a
path homotopy in $\R$: $\tilde{H}(0, t) = \tilde{x}_0$ and $\tilde{H}(1, t) =
\tilde{H}(1, 0)$ for all $t$.
\end{lemma}

\begin{proof}
The existence and uniqueness follow by the same Lebesgue number argument applied
to the square $I \times I$. Subdivide $I \times I$ into small squares, each
mapping into an evenly covered neighborhood, and lift inductively.

For the boundary conditions: $t \mapsto \tilde{H}(0, t)$ is a lift of the
constant path $t \mapsto x_0$ starting at $\tilde{x}_0$, hence is constant (by
uniqueness of lifts). Similarly, $t \mapsto \tilde{H}(1, t)$ is a lift of the
constant path at $x_1$, hence is constant.
\end{proof}

\begin{definition}[Degree of a loop]\label{def:degree-loop}
For a loop $\gamma$ in $S^1$ based at $1 = (1,0)$, let $\tilde{\gamma}$ be the
unique lift to $\R$ with $\tilde{\gamma}(0) = 0$. The integer $\tilde{\gamma}(1)
\in \Z$ is called the \emph{degree}\index{degree of a loop} (or \emph{winding
number}) of $\gamma$, denoted $\deg(\gamma)$.
\end{definition}

\begin{remark}\label{rem:degree-well-defined}
Since $p(\tilde{\gamma}(1)) = \gamma(1) = 1 = p(0)$ and $p^{-1}(1) = \Z$, we
have $\tilde{\gamma}(1) \in \Z$. By Lemma~\ref{lem:homotopy-lifting}, if
$\gamma \simeq \gamma' \rel{\partial I}$, then $\deg(\gamma) = \deg(\gamma')$,
so the degree depends only on the homotopy class.
\end{remark}

\begin{theorem}[$\pi_1(S^1) \cong \Z$]\label{thm:pi1-S1}
The map $\Phi \colon \pi_1(S^1, 1) \to \Z$ defined by $\Phi([\gamma]) =
\deg(\gamma)$ is a group isomorphism.
\end{theorem}

\begin{proof}
We verify that $\Phi$ is a well-defined group homomorphism and a bijection.

\textbf{Well-definedness.} If $[\gamma] = [\gamma']$, then $\gamma \simeq
\gamma' \rel{\partial I}$. By Lemma~\ref{lem:homotopy-lifting}, the lifts
$\tilde{\gamma}$ and $\tilde{\gamma}'$ (both starting at $0$) satisfy
$\tilde{\gamma}(1) = \tilde{\gamma}'(1)$. So $\deg(\gamma) = \deg(\gamma')$.

\textbf{Homomorphism.} Let $\gamma$ and $\delta$ be loops at $1$, with lifts
$\tilde{\gamma}$ and $\tilde{\delta}$ starting at $0$. Write $n =
\tilde{\gamma}(1)$ and $m = \tilde{\delta}(1)$. Define $\tilde{\delta}_n(t) =
\tilde{\delta}(t) + n$; this is a lift of $\delta$ starting at $n$. The
concatenation $\tilde{\gamma} \star \tilde{\delta}_n$ is a path in $\R$ from $0$
to $m + n$, lifting $\gamma \star \delta$. By uniqueness, $\widetilde{\gamma
\star \delta} = \tilde{\gamma} \star \tilde{\delta}_n$, and so
\[
  \Phi([\gamma \star \delta]) = \widetilde{\gamma \star \delta}(1) = m + n = \Phi([\gamma]) + \Phi([\delta]).
\]

\textbf{Surjectivity.} For each $n \in \Z$, let $\omega_n(t) = (\cos 2\pi nt,
\sin 2\pi nt)$. This is a loop at $1$, and its lift starting at $0$ is
$\tilde{\omega}_n(t) = nt$, which ends at $n$. Hence $\Phi([\omega_n]) = n$.

\textbf{Injectivity.} Suppose $\Phi([\gamma]) = 0$, i.e., $\tilde{\gamma}(1) =
0$. Then $\tilde{\gamma}$ is a loop in $\R$ based at $0$. Since $\R$ is simply
connected (being contractible), $\tilde{\gamma} \simeq c_0 \rel{\partial I}$ in
$\R$. Applying $p$ to this homotopy, we get $\gamma = p \circ \tilde{\gamma}
\simeq p \circ c_0 = c_1 \rel{\partial I}$ in $S^1$. Hence $[\gamma] =
[c_1]$, the identity element.
\end{proof}

\begin{remark}[Geometric interpretation]\label{rem:winding-number}
The integer $n = \Phi([\gamma])$ is the \emph{winding number}: it counts how
many times $\gamma$ wraps around $S^1$, with sign indicating direction. The loop
$\omega_n$ wraps $n$ times counterclockwise (or $|n|$ times clockwise if $n < 0$).
\end{remark}

\begin{center}
\begin{tikzpicture}[scale=0.9]
  % n = 0
  \begin{scope}[xshift=-4cm]
    \draw[thick] (0,0) circle (1);
    \fill (1,0) circle (2pt);
    \node[below] at (0,-1.5) {$n = 0$};
    \draw[blue!70!black,thick,{Stealth}-{Stealth}] (1,0) -- (1,0);
    \node[right] at (1.1,0) {\small $\bullet$};
  \end{scope}
  % n = 1
  \begin{scope}
    \draw[thick] (0,0) circle (1);
    \fill (1,0) circle (2pt);
    \node[below] at (0,-1.5) {$n = 1$};
    \draw[blue!70!black,thick,-{Stealth},
      decoration={markings, mark=at position 0.5 with {\arrow{Stealth}}},
      postaction={decorate}]
      (1,0) arc[start angle=0,end angle=350,radius=1];
  \end{scope}
  % n = 2
  \begin{scope}[xshift=4cm]
    \draw[thick] (0,0) circle (1);
    \fill (1,0) circle (2pt);
    \node[below] at (0,-1.5) {$n = 2$};
    \draw[blue!70!black,thick,
      decoration={markings, mark=at position 0.25 with {\arrow{Stealth}},
                            mark=at position 0.75 with {\arrow{Stealth}}},
      postaction={decorate}]
      (1,0) arc[start angle=0,end angle=710,radius=1];
  \end{scope}
  % n = -1
  \begin{scope}[xshift=8cm]
    \draw[thick] (0,0) circle (1);
    \fill (1,0) circle (2pt);
    \node[below] at (0,-1.5) {$n = -1$};
    \draw[red!70!black,thick,
      decoration={markings, mark=at position 0.5 with {\arrow{Stealth}}},
      postaction={decorate}]
      (1,0) arc[start angle=0,end angle=-350,radius=1];
  \end{scope}
\end{tikzpicture}
\captionof{figure}{Loops on $S^1$ with winding numbers $0$, $1$, $2$, and $-1$.}\label{fig:winding-numbers}
\end{center}


%% ──────────────────────────────────────────────────────
\section{Applications: Brouwer Fixed-Point Theorem in Dimension 2}\label{sec:brouwer}

As a striking application of the computation $\pi_1(S^1) \cong \Z$, we prove
the Brouwer fixed-point theorem for the disk $D^2$.

\begin{lemma}[No retraction $D^2 \to S^1$]\label{lem:no-retraction}
There is no retraction $r \colon D^2 \to S^1$, i.e., no continuous map $r$
with $r|_{S^1} = \Id_{S^1}$.
\end{lemma}

\begin{proof}
Suppose for contradiction that $r \colon D^2 \to S^1$ is a retraction. Let
$\iota \colon S^1 \hookrightarrow D^2$ be the inclusion. Then $r \circ \iota =
\Id_{S^1}$, so by functoriality of $\pi_1$:
\[
  r_* \circ \iota_* = (\Id_{S^1})_* = \Id_{\pi_1(S^1)}.
\]
Now consider the diagram of induced homomorphisms:
\begin{center}
\begin{tikzcd}
\pi_1(S^1, x_0) \ar[r, "\iota_*"] \ar[rr, bend right=20, "\Id"'] &
\pi_1(D^2, x_0) \ar[r, "r_*"] &
\pi_1(S^1, x_0)
\end{tikzcd}
\end{center}
Since $D^2$ is contractible, $\pi_1(D^2, x_0) = \{e\}$. Hence $\iota_*$ maps
$\pi_1(S^1) \cong \Z$ to the trivial group, and $r_* \circ \iota_*$ is the
zero map. But we need $r_* \circ \iota_* = \Id$, which is impossible since
$\Z \ne \{0\}$. Contradiction.
\end{proof}

\begin{theorem}[Brouwer fixed-point theorem, dimension 2]\label{thm:brouwer-2d}
Every continuous map $f \colon D^2 \to D^2$ has a fixed point.
\end{theorem}

\begin{proof}
Suppose for contradiction that $f$ has no fixed point, i.e., $f(x) \ne x$ for
all $x \in D^2$. For each $x \in D^2$, consider the ray starting at $f(x)$,
passing through $x$, and continuing until it hits $S^1$. Define $r(x)$ to be
this intersection point.

Explicitly: $r(x) = x + t(x) \cdot (x - f(x))$ where $t(x) \ge 0$ is the
unique non-negative real number such that $\abs{r(x)} = 1$. To find $t(x)$, we
solve
\[
  \abs{x + t(x - f(x))}^2 = 1,
\]
which is a quadratic in $t$ with a unique non-negative solution (given by the
quadratic formula, since $\abs{x - f(x)} > 0$ by hypothesis). The function $t(x)$
is continuous (by continuity of the quadratic formula), hence $r$ is continuous.

Moreover, if $x \in S^1$, then $\abs{x} = 1$, so $t(x) = 0$ and $r(x) = x$.
Thus $r \colon D^2 \to S^1$ is a retraction, contradicting
Lemma~\ref{lem:no-retraction}.
\end{proof}

\begin{center}
\begin{tikzpicture}[scale=1.8]
  \draw[thick,fill=blue!3] (0,0) circle (1);
  \node at (-0.3,-0.3) {\small $D^2$};
  \coordinate (X) at (0.3,0.2);
  \coordinate (FX) at (-0.2,0.5);
  \fill[blue!70!black] (X) circle (1.2pt) node[below right] {\small $x$};
  \fill[red!70!black] (FX) circle (1.2pt) node[above left] {\small $f(x)$};
  % Ray from f(x) through x to S^1
  \pgfmathsetmacro{\dx}{0.3-(-0.2)}
  \pgfmathsetmacro{\dy}{0.2-0.5}
  % direction: (0.5, -0.3), normalized approx (0.857, -0.514)
  \coordinate (RX) at ({0.3+0.7*0.857},{0.2+0.7*(-0.514)});
  \draw[thick,orange!70!black,-{Stealth}] (FX) -- (X);
  \draw[thick,orange!70!black,-{Stealth}] (X) -- (0.88,-0.15);
  \fill[green!50!black] (0.88,-0.15) circle (1.2pt) node[right] {\small $r(x)$};
  \node[thick] at (0,-1.4) {\small Ray from $f(x)$ through $x$ to $S^1$};
\end{tikzpicture}
\captionof{figure}{Construction of the retraction $r$ in the proof of the Brouwer fixed-point theorem.}\label{fig:brouwer-proof}
\end{center}

\begin{remark}\label{rem:brouwer-higher}
The Brouwer fixed-point theorem holds in all dimensions: every continuous map
$f \colon D^n \to D^n$ has a fixed point. For $n \ge 3$, the proof requires
homology theory (since $\pi_1(S^{n-1}) = 0$ for $n \ge 3$ and the argument
above breaks down). We will return to this in the chapter on singular homology.
\end{remark}

\begin{corollary}[No-retraction theorem]\label{cor:no-retract-general}
$S^1$ is not a retract of $D^2$.
\end{corollary}

\begin{proof}
This is precisely Lemma~\ref{lem:no-retraction}.
\end{proof}


%% ──────────────────────────────────────────────────────
\section{Further First Computations}\label{sec:first-computations}

Using the tools developed so far, we can compute the fundamental groups of
several basic spaces.

\begin{proposition}[$\pi_1(S^n) = 0$ for $n \ge 2$]\label{prop:pi1-Sn}
For $n \ge 2$, the sphere $S^n$ is simply connected: $\pi_1(S^n, x_0) = \{e\}$.
\end{proposition}

\begin{proof}
Let $\gamma \colon I \to S^n$ be a loop based at $x_0$. We claim $\gamma$ is
homotopic to a loop that misses some point $q \in S^n$. If $\gamma$ is not
surjective, this is immediate. If $\gamma$ is surjective, we use the
compactness of $I$ and the fact that $\gamma$ cannot fill $S^n$ in a
sufficiently strong sense: by a Lebesgue number argument and a small
perturbation (for $n \ge 2$, the complement of the image of a curve in $S^n$
is dense), $\gamma$ can be homotoped to a non-surjective loop.

Once $\gamma$ misses some point $q$, it lies in $S^n \setminus \{q\} \cong
\R^n$ (via stereographic projection). Since $\R^n$ is contractible, $\gamma$ is
null-homotopic in $S^n \setminus \{q\}$, hence in $S^n$.
\end{proof}

\begin{proposition}[Fundamental group of a product]\label{prop:pi1-product}
For any pointed spaces $(X, x_0)$ and $(Y, y_0)$,
\[
  \pi_1(X \times Y, (x_0, y_0)) \cong \pi_1(X, x_0) \times \pi_1(Y, y_0).
\]
\end{proposition}

\begin{proof}
The projections $p_1 \colon X \times Y \to X$ and $p_2 \colon X \times Y \to Y$
induce a homomorphism
\[
  \Psi \colon \pi_1(X \times Y) \to \pi_1(X) \times \pi_1(Y), \qquad
  [\gamma] \mapsto ((p_1)_*[\gamma],\, (p_2)_*[\gamma]).
\]
The inverse is given by $\Psi^{-1}([\alpha], [\beta]) = [\alpha \times \beta]$,
where $(\alpha \times \beta)(t) = (\alpha(t), \beta(t))$. One checks directly
that $\Psi$ and $\Psi^{-1}$ are inverse homomorphisms.
\end{proof}

\begin{corollary}[Fundamental group of the torus]\label{cor:pi1-torus}
$\pi_1(T^2) = \pi_1(S^1 \times S^1) \cong \Z \times \Z \cong \Z^2$.
More generally, $\pi_1(T^n) \cong \Z^n$.
\end{corollary}

\begin{example}[Fundamental group of the punctured plane]\label{ex:pi1-punctured-plane}
Since $\R^2 \setminus \{0\}$ is homotopy equivalent to $S^1$ (via radial
retraction, Example~\ref{ex:sphere-retract}), we have $\pi_1(\R^2 \setminus
\{0\}) \cong \pi_1(S^1) \cong \Z$ by homotopy invariance. The generator is the
loop that winds once around the origin.
\end{example}

\begin{example}[Cylinder and M\"{o}bius band]\label{ex:cylinder-mobius}
The cylinder $S^1 \times I$ deformation retracts onto $S^1 \times \{0\}$, so
$\pi_1(S^1 \times I) \cong \Z$. The M\"{o}bius band also deformation retracts
onto its central circle (which is a copy of $S^1$), so it too has fundamental
group $\Z$.
\end{example}


%% ──────────────────────────────────────────────────────
\section{Topological Consequences of Functoriality}\label{sec:consequences-functoriality}

The functoriality of $\pi_1$ has many immediate and powerful consequences. We
collect several here.

\begin{proposition}[Retract inherits fundamental group]\label{prop:retract-pi1}
Let $A$ be a retract of $X$ with retraction $r \colon X \to A$, and let $\iota
\colon A \hookrightarrow X$ be the inclusion. Then for any $a_0 \in A$:
\begin{enumerate}[label=(\roman*)]
\item $\iota_* \colon \pi_1(A, a_0) \to \pi_1(X, a_0)$ is injective.
\item $r_* \colon \pi_1(X, a_0) \to \pi_1(A, a_0)$ is surjective.
\item $\pi_1(X, a_0) \cong \pi_1(A, a_0) \oplus \ker(r_*)$ when $\pi_1(X, a_0)$
  is abelian (in general, $\pi_1(A, a_0)$ is a retract of $\pi_1(X, a_0)$ in the
  category of groups).
\end{enumerate}
\end{proposition}

\begin{proof}
Since $r \circ \iota = \Id_A$, functoriality gives $r_* \circ \iota_* =
\Id_{\pi_1(A, a_0)}$. This immediately implies that $\iota_*$ is injective and
$r_*$ is surjective. For (iii), when $\pi_1(X, a_0)$ is abelian, the map
$\iota_* \circ r_*$ is an idempotent endomorphism of $\pi_1(X, a_0)$, giving
the splitting.
\end{proof}

\begin{proposition}[Homeomorphic spaces have isomorphic $\pi_1$]\label{prop:homeo-pi1}
If $f \colon X \to Y$ is a homeomorphism with $f(x_0) = y_0$, then $f_* \colon
\pi_1(X, x_0) \to \pi_1(Y, y_0)$ is an isomorphism.
\end{proposition}

\begin{proof}
Since $f$ is a homeomorphism, it is in particular a homotopy equivalence (with
homotopy inverse $f^{-1}$). By Theorem~\ref{thm:pi1-homotopy-invariant}, $f_*$
is an isomorphism. Alternatively, one can check directly: $(f^{-1})_*$ is a
two-sided inverse of $f_*$ by functoriality.
\end{proof}

\begin{proposition}[$\R^2 \not\cong \R^n$ for $n \ne 2$]\label{prop:R2-not-Rn}
The Euclidean spaces $\R^2$ and $\R^n$ are not homeomorphic for $n \ne 2$.
\end{proposition}

\begin{proof}
Suppose $f \colon \R^2 \to \R^n$ were a homeomorphism. Restricting to
complements of a point, $f$ would give a homeomorphism $\R^2 \setminus \{0\}
\cong \R^n \setminus \{f(0)\}$. But $\pi_1(\R^2 \setminus \{0\}) \cong \Z$,
whereas for $n = 1$, $\R \setminus \{p\}$ is disconnected (so not even
path-connected), and for $n \ge 3$, $\pi_1(\R^n \setminus \{p\}) \cong
\pi_1(S^{n-1}) = \{e\}$ by Proposition~\ref{prop:pi1-Sn}. In each case the
fundamental groups are non-isomorphic, so no such homeomorphism can exist.
\end{proof}

\begin{remark}\label{rem:invariance-of-domain}
The general result that $\R^m \cong \R^n$ implies $m = n$ is called
\emph{invariance of domain}\index{invariance of domain} (or invariance of
dimension). A full proof for all $m, n$ requires homology theory. The argument
above settles the case $m = 2$.
\end{remark}

\begin{theorem}[Fundamental theorem of algebra, topological proof sketch]\label{thm:FTA-sketch}
Every non-constant polynomial $p(z) \in \C[z]$ has a root in $\C$.
\end{theorem}

\begin{proof}[Proof sketch]
Write $p(z) = z^n + a_{n-1}z^{n-1} + \cdots + a_0$ with $n \ge 1$. Suppose for
contradiction that $p(z) \ne 0$ for all $z \in \C$. Define $f_R \colon S^1 \to
S^1$ by
\[
  f_R(e^{i\theta}) = \frac{p(Re^{i\theta})}{\abs{p(Re^{i\theta})}}.
\]
For $R$ large, $p(Re^{i\theta}) \approx R^n e^{in\theta}$, so $f_R$ is
homotopic to the map $e^{i\theta} \mapsto e^{in\theta}$, which has degree $n$.
For $R = 0$, $f_0$ is the constant map $p(0)/\abs{p(0)}$, which has degree $0$.
Since the degree is a homotopy invariant and $R \mapsto f_R$ gives a homotopy
from $f_0$ to $f_R$ (using the assumption that $p$ has no zeros), we obtain $n
= 0$, a contradiction.
\end{proof}

\begin{center}
\begin{tikzpicture}[scale=1.1]
  % z-plane
  \begin{scope}
    \draw[->] (-1.8,0) -- (1.8,0) node[right] {\small $\Re$};
    \draw[->] (0,-1.8) -- (0,1.8) node[above] {\small $\Im$};
    \draw[thick,blue!70!black,
      decoration={markings, mark=at position 0.25 with {\arrow{Stealth}},
                            mark=at position 0.75 with {\arrow{Stealth}}},
      postaction={decorate}]
      (0,0) circle (1.3);
    \node[blue!70!black] at (1.7,1.0) {\small $|z|=R$};
    \node[below] at (0,-2.1) {$z$-plane};
  \end{scope}
  % Arrow
  \draw[-{Stealth},thick] (2.5,0) -- node[above] {\small $p$} (4,0);
  % w-plane
  \begin{scope}[xshift=6.5cm]
    \draw[->] (-1.8,0) -- (1.8,0) node[right] {\small $\Re$};
    \draw[->] (0,-1.8) -- (0,1.8) node[above] {\small $\Im$};
    \draw[thick,red!70!black,
      decoration={markings, mark=at position 0.15 with {\arrow{Stealth}},
                            mark=at position 0.55 with {\arrow{Stealth}}},
      postaction={decorate}]
      (0:1.2) .. controls (40:1.5) and (70:0.8) .. (100:1.3)
      .. controls (140:1.0) and (170:1.5) .. (200:1.1)
      .. controls (240:0.9) and (280:1.4) .. (320:1.2)
      .. controls (345:1.3) and (355:1.2) .. (360:1.2);
    \node[below] at (0,-2.1) {$w$-plane};
    \node[red!70!black] at (1.6,1.1) {\small image};
  \end{scope}
\end{tikzpicture}
\captionof{figure}{The image of a large circle under $p(z)$ winds $n$ times around the origin.}\label{fig:FTA}
\end{center}

\begin{proposition}[Fundamental group of a topological group is abelian]\label{prop:pi1-topological-group}
If $G$ is a topological group with identity element $e$, then $\pi_1(G, e)$ is
abelian.
\end{proposition}

\begin{proof}
Let $\gamma$ and $\delta$ be loops in $G$ based at $e$. Define
\[
  H(s, t) = \gamma(s) \cdot \delta(t),
\]
where $\cdot$ denotes the group operation in $G$. This is continuous since $G$ is
a topological group. We compute:
\begin{align*}
  H(s, 0) &= \gamma(s) \cdot e = \gamma(s), & H(s, 1) &= \gamma(s) \cdot e = \gamma(s), \\
  H(0, t) &= e \cdot \delta(t) = \delta(t), & H(1, t) &= e \cdot \delta(t) = \delta(t).
\end{align*}
Now consider the two paths along the boundary of $I \times I$ from $(0,0)$ to
$(1,1)$: the path along the bottom then right gives $\gamma \star \delta$, and
the path along the left then top gives $\delta \star \gamma$. The map $H$
provides a homotopy between these two paths (relative to endpoints), showing
$[\gamma] \cdot [\delta] = [\delta] \cdot [\gamma]$.

More precisely, define $\alpha(t) = H(t, t) = \gamma(t) \cdot \delta(t)$ and note
that the paths $s \mapsto H(s, 0)$ then $t \mapsto H(1, t)$ and $t \mapsto H(0,
t)$ then $s \mapsto H(s, 1)$ give, respectively, $\gamma \star \delta$ and
$\delta \star \gamma$. Since $H$ is a homotopy of the square with the right
boundary values, one obtains $[\gamma \star \delta] = [\delta \star \gamma]$.
\end{proof}

\begin{remark}\label{rem:eckmann-hilton}
The proof above is an instance of the \emph{Eckmann--Hilton
argument}\index{Eckmann--Hilton argument}: when a set carries two binary
operations, both with the same unit, and each distributing over the other, then
the two operations coincide and are commutative. Here the two operations are
concatenation of loops and pointwise multiplication in $G$.
\end{remark}


%% ──────────────────────────────────────────────────────
\section{Summary and Outlook}\label{sec:ch2-summary}

Let us summarize the main results of this chapter.

\begin{center}
\renewcommand{\arraystretch}{1.4}
\begin{tabular}{@{}lll@{}}
\toprule
\textbf{Space $X$} & $\pi_1(X)$ & \textbf{Justification} \\
\midrule
$\R^n$ & $\{e\}$ & Contractible \\
$D^n$ & $\{e\}$ & Contractible \\
$S^n$, $n \ge 2$ & $\{e\}$ & Proposition~\ref{prop:pi1-Sn} \\
$S^1$ & $\Z$ & Theorem~\ref{thm:pi1-S1} \\
$T^n = (S^1)^n$ & $\Z^n$ & Corollary~\ref{cor:pi1-torus} \\
$\R^2 \setminus \{0\}$ & $\Z$ & Homotopy equivalent to $S^1$ \\
$S^2 \times S^1$ & $\Z$ & Product formula \\
M\"{o}bius band & $\Z$ & Deformation retract onto $S^1$ \\
Cylinder $S^1 \times I$ & $\Z$ & Deformation retract onto $S^1$ \\
\bottomrule
\end{tabular}
\captionof{table}{Fundamental groups computed in Chapter~\ref{ch:fundamental-group}.}\label{tab:pi1-summary}
\end{center}

In the next chapter, we will develop the theory of covering spaces, which
provides a powerful geometric tool for computing and understanding the
fundamental group. We will then prove the Seifert--van Kampen theorem, which
allows us to compute the fundamental group of a space built from simpler pieces
whose fundamental groups are known. These tools will enable us to compute
$\pi_1$ of surfaces, wedge sums, and many other spaces.


%% ──────────────────────────────────────────────────────
\section{Exercises}\label{sec:ch2-exercises}

\begin{exercise}\label{exo:ch2-1}
Let $\gamma, \delta \colon I \to X$ be two paths from $x_0$ to $x_1$. Show that
$\gamma \simeq \delta \rel{\partial I}$ if and only if $\gamma \star
\bar{\delta}$ is null-homotopic (i.e., $[\gamma \star \bar{\delta}] = e$ in
$\pi_1(X, x_0)$).
\end{exercise}

\begin{exercise}\label{exo:ch2-2}
Let $f \colon (X, x_0) \to (Y, y_0)$ be a pointed map and let $A \subseteq X$
be a retract of $X$ with $x_0 \in A$. Show that $f_* \colon \pi_1(X, x_0) \to
\pi_1(Y, y_0)$ maps $\iota_*(\pi_1(A, x_0))$ injectively into $\pi_1(Y, y_0)$,
where $\iota \colon A \hookrightarrow X$ is the inclusion. In particular,
$\pi_1(A, x_0)$ is (isomorphic to) a subgroup of $\pi_1(X, x_0)$.
\end{exercise}

\begin{exercise}\label{exo:ch2-3}
Show that $S^1$ is not simply connected by proving directly (without using
$\pi_1(S^1) \cong \Z$) that the loop $\omega_1(t) = (\cos 2\pi t, \sin 2\pi t)$
is not null-homotopic.
\emph{Hint:} Use the degree/lifting approach.
\end{exercise}

\begin{exercise}\label{exo:ch2-4}
Let $f \colon S^1 \to S^1$ be a continuous map. The \emph{degree} of $f$,
denoted $\deg(f)$, is the integer $n$ such that $f_* \colon \pi_1(S^1) \to
\pi_1(S^1)$ is multiplication by $n$ (i.e., $f_*(k) = nk$ under the
identification $\pi_1(S^1) \cong \Z$). Show that:
\begin{enumerate}[label=(\alph*)]
\item If $f$ has no fixed point, then $\deg(f) = -1$.
  \emph{Hint:} show $f \simeq -\Id$ (the antipodal map).
\item The antipodal map $a \colon S^1 \to S^1$, $a(z) = -z$, has degree $-1$.
\item If $f$ is not surjective, then $\deg(f) = 0$.
\end{enumerate}
\end{exercise}

\begin{exercise}\label{exo:ch2-5}
\emph{(Borsuk--Ulam in dimension 1.)} Show that for every continuous map
$f \colon S^1 \to \R$, there exists a point $x \in S^1$ with $f(x) = f(-x)$.
\emph{Hint:} Consider $g(x) = f(x) - f(-x)$ and use the intermediate value
theorem.
\end{exercise}

\begin{exercise}\label{exo:ch2-6}
Let $X$ be a path-connected space. Show that $\pi_1(X, x_0)$ is abelian if and
only if for all paths $\alpha, \beta$ from $x_0$ to $x_1$, the change of
basepoint isomorphisms $\beta_\alpha$ and $\beta_\beta$ agree (i.e., $\beta_\alpha
= \beta_\beta \colon \pi_1(X, x_0) \to \pi_1(X, x_1)$).
\end{exercise}

\begin{exercise}\label{exo:ch2-7}
Compute $\pi_1(X)$ for the following spaces (justify your answers using results
from this chapter):
\begin{enumerate}[label=(\alph*)]
\item $X = \R^3 \setminus \{0\}$.
\item $X = S^1 \times \R^2$.
\item $X = S^2 \times S^1$.
\item $X = \R^2 \setminus \{p_1, \ldots, p_k\}$ for $k$ distinct points (just
  state the answer; a full computation requires the Seifert--van Kampen
  theorem).
\end{enumerate}
\end{exercise}

\begin{exercise}\label{exo:ch2-8}
Use the Brouwer fixed-point theorem (dimension 2) to show that every $2 \times
2$ real matrix $A$ with all entries positive has a positive real eigenvalue.
\emph{Hint:} Consider the action of $A$ on the simplex $\Delta = \{(x,y) \in
\R^2 : x + y = 1,\, x \ge 0,\, y \ge 0\}$.
\end{exercise}

\begin{exercise}\label{exo:ch2-9}
Let $f, g \colon X \to S^1$ be continuous maps. Define $f \cdot g \colon X \to
S^1$ by $(f \cdot g)(x) = f(x) \cdot g(x)$, using the group multiplication in
$S^1 \subseteq \C$. Show that $(f \cdot g)_* = f_* + g_*$ as homomorphisms
$\pi_1(X, x_0) \to \pi_1(S^1) \cong \Z$.
\end{exercise}

\begin{exercise}\label{exo:ch2-10}
Let $p \colon \R \to S^1$ be the standard covering map. Show that for any simply
connected, locally path-connected space $Y$ and any continuous map $f \colon Y
\to S^1$, there exists a lift $\tilde{f} \colon Y \to \R$ with $p \circ
\tilde{f} = f$. \emph{Hint:} Fix a basepoint $y_0 \in Y$, choose $\tilde{f}(y_0)
\in p^{-1}(f(y_0))$, and for any $y \in Y$ define $\tilde{f}(y)$ using the lift
of $f \circ \gamma$ where $\gamma$ is a path from $y_0$ to $y$. Show this is
well-defined using simple connectivity.
\end{exercise}

% === END OF CHUNK preamble+ch1-2 ===
% === START OF CHUNK ch3-4 ===

%%%%%%%%%%%%%%%%%%%%%%%%%%%%%%%%%%%%%%%%%%%%%%%%%%%%%%%%%%%%
%%  Algebraic Topology — Chapters 3 & 4
%%  Graduate Level (Master / PhD)
%%  Chapter 3: The Seifert--Van Kampen Theorem
%%  Chapter 4: Covering Spaces and the Fundamental Group
%%%%%%%%%%%%%%%%%%%%%%%%%%%%%%%%%%%%%%%%%%%%%%%%%%%%%%%%%%%%

%% ══════════════════════════════════════════════════════════
%%  CHAPTER 3: THE SEIFERT–VAN KAMPEN THEOREM
%% ══════════════════════════════════════════════════════════

\chapter{The Seifert--Van Kampen Theorem}\label{ch:svk}

The computation of the fundamental group of a circle, while beautiful, relied
on the particular geometry of $S^1$.  For more complicated spaces---wedge sums,
surfaces, CW~complexes---one needs a systematic tool that expresses
$\pi_1(X)$ in terms of simpler pieces.  The \emph{Seifert--Van Kampen
theorem}\index{Seifert--Van Kampen theorem} is precisely such a tool: it
describes the fundamental group of a union $X = U_1 \cup U_2$ in terms of
$\pi_1(U_1)$, $\pi_1(U_2)$, and $\pi_1(U_1 \cap U_2)$.

The algebraic prerequisite is the notion of \emph{free product} and, more
generally, \emph{amalgamated free product} of groups.  We begin with these
constructions.

%% ──────────────────────────────────────────────────────────
\section{Free groups and free products}\label{sec:free-products}

\begin{definition}[Free group]\label{def:free-group}
  \index{free group}
  Let $S$ be a set.  The \textbf{free group} on $S$, denoted $F(S)$, is a group
  $F(S)$ together with a map $\iota\colon S \to F(S)$ satisfying the following
  universal property: for every group $G$ and every map $f\colon S \to G$,
  there exists a unique group homomorphism $\varphi\colon F(S) \to G$ such that
  $\varphi \circ \iota = f$.
  \[
    \begin{tikzcd}
      S \ar[r,"\iota"] \ar[dr,"f"'] & F(S) \ar[d,dashed,"\exists!\,\varphi"] \\
      & G
    \end{tikzcd}
  \]
\end{definition}

\begin{remark}[Existence and explicit construction]\label{rem:free-group-explicit}
  Concretely, $F(S)$ consists of all reduced words in the alphabet
  $S \cup S^{-1}$: finite sequences $s_1^{\varepsilon_1} s_2^{\varepsilon_2}
  \cdots s_n^{\varepsilon_n}$ with $s_i \in S$, $\varepsilon_i \in \{+1,-1\}$,
  and no adjacent pair $s\, s^{-1}$ or $s^{-1}\, s$.  The group operation is
  concatenation followed by free reduction.  The empty word serves as the
  identity.  The universal property is then easily verified.
\end{remark}

\begin{definition}[Free product of groups]\label{def:free-product}
  \index{free product!of groups}
  Let $\{G_\alpha\}_{\alpha \in A}$ be a family of groups.  The
  \textbf{free product} $\displaystyle\mathop{\scalebox{1.3}{$\ast$}}_{\alpha
  \in A} G_\alpha$ is a group $G$ together with homomorphisms
  $\iota_\alpha \colon G_\alpha \to G$ satisfying the following universal
  property: for every group $H$ and every family of homomorphisms
  $\varphi_\alpha \colon G_\alpha \to H$, there exists a unique homomorphism
  $\Phi\colon G \to H$ such that $\Phi \circ \iota_\alpha = \varphi_\alpha$
  for all $\alpha$.
  \[
    \begin{tikzcd}
      G_\alpha \ar[r,"\iota_\alpha"] \ar[dr,"\varphi_\alpha"'] &
      \displaystyle\mathop{\scalebox{1.1}{$\ast$}}_\alpha G_\alpha
        \ar[d,dashed,"\exists!\,\Phi"] \\
      & H
    \end{tikzcd}
  \]
\end{definition}

\begin{remark}[Explicit construction]\label{rem:free-product-explicit}
  Elements of $G_1 * G_2$ are represented by reduced words
  $g_1 g_2 g_3 \cdots g_n$ where each $g_i$ belongs to one of the factors
  $G_1$ or $G_2$, successive letters come from \emph{different} factors, and no
  letter is the identity.  Multiplication is concatenation followed by
  amalgamation of adjacent letters from the same factor and deletion of
  resulting identity elements.
\end{remark}

\begin{example}[Free product $\Z * \Z$]\label{ex:Z-star-Z}
  The free product $\Z * \Z$ is the free group on two generators.  Writing
  $a$ for the generator of the first copy and $b$ for the second, a typical
  reduced word is $a^2 b^{-1} a^3 b^2$.  Note that $\Z * \Z$ is
  \emph{not} abelian: $ab \neq ba$.
\end{example}

\begin{proposition}[Basic properties of free products]\label{prop:free-product-props}
  Let $G$ and $H$ be groups.
  \begin{enumerate}[label=(\roman*)]
    \item The inclusions $\iota_G\colon G \hookrightarrow G * H$ and
      $\iota_H\colon H \hookrightarrow G * H$ are injective.
    \item $G * H \cong H * G$.
    \item $G * \{e\} \cong G$.
    \item The abelianization of $G*H$ is $G^{\mathrm{ab}} \oplus H^{\mathrm{ab}}$.
  \end{enumerate}
\end{proposition}

\begin{proof}
  Parts (i)--(iii) are straightforward from the reduced-word description.
  For (iv), the abelianization is the largest abelian quotient.  The universal
  property of $G*H$, combined with the fact that every pair of homomorphisms
  into an abelian group $A$ factors uniquely, yields
  $(G*H)^{\mathrm{ab}} \cong G^{\mathrm{ab}} \oplus H^{\mathrm{ab}}$.
\end{proof}

%% ──────────────────────────────────────────────────────────
\section{Amalgamated free products}\label{sec:amalgamated}

\begin{definition}[Amalgamated free product]\label{def:amalgamated-free-product}
  \index{amalgamated free product}\index{free product!amalgamated}
  Let $G_1, G_2$ be groups and let $K$ be a group with homomorphisms
  $j_1\colon K \to G_1$ and $j_2\colon K \to G_2$.  The \textbf{amalgamated
  free product} $G_1 *_K G_2$ is a group $G$ together with homomorphisms
  $\iota_1\colon G_1 \to G$ and $\iota_2\colon G_2 \to G$ satisfying
  $\iota_1 \circ j_1 = \iota_2 \circ j_2$ and the following universal
  property: for every group $H$ and homomorphisms $\varphi_1\colon G_1 \to H$,
  $\varphi_2\colon G_2 \to H$ with $\varphi_1 \circ j_1 = \varphi_2 \circ
  j_2$, there exists a unique homomorphism $\Phi\colon G \to H$ making the
  diagram commute.
  \[
    \begin{tikzcd}[column sep=2.5em, row sep=2.5em]
      & K \ar[dl,"j_1"'] \ar[dr,"j_2"] & \\
      G_1 \ar[dr,"\iota_1"'] \ar[ddr,bend right,"\varphi_1"'] & &
      G_2 \ar[dl,"\iota_2"] \ar[ddl,bend left,"\varphi_2"] \\
      & G_1 *_K G_2 \ar[d,dashed,"\exists!\,\Phi"] & \\
      & H &
    \end{tikzcd}
  \]
  Equivalently, $G_1 *_K G_2$ is the \textbf{pushout}\index{pushout!of groups}
  in the category of groups:
  \[
    \begin{tikzcd}
      K \ar[r,"j_2"] \ar[d,"j_1"'] & G_2 \ar[d,"\iota_2"] \\
      G_1 \ar[r,"\iota_1"'] & G_1 *_K G_2 \ar[ul, phantom, "\ulcorner" very near start]
    \end{tikzcd}
  \]
\end{definition}

\begin{remark}[Explicit construction]\label{rem:amalgamated-explicit}
  The amalgamated free product can be constructed as
  \[
    G_1 *_K G_2 \;=\; (G_1 * G_2) \,\big/\, N,
  \]
  where $N$ is the normal closure of $\set{j_1(k)\, j_2(k)^{-1} : k \in K}$.
  In other words, we take the free product and impose the relations
  $j_1(k) = j_2(k)$ for all $k \in K$.
\end{remark}

\begin{example}[Trivial amalgamation]\label{ex:trivial-amalgamation}
  When $K = \{e\}$, the amalgamated free product $G_1 *_{\{e\}} G_2$ is just the
  ordinary free product $G_1 * G_2$.
\end{example}

\begin{example}[Amalgamation with $K = \Z$]\label{ex:amalgam-Z}
  Let $G_1 = G_2 = \Z$ with $j_1(1) = 2$ and $j_2(1) = 3$.  Then
  $G_1 *_K G_2 = \langle a, b \mid a^2 = b^3 \rangle$, which is the trefoil
  knot group.\index{trefoil knot group}
\end{example}

%% ──────────────────────────────────────────────────────────
\section{Statement and proof of the Seifert--Van Kampen theorem}\label{sec:svk-statement}

We now state the central theorem of this chapter.

\begin{theorem}[Seifert--Van Kampen]\label{thm:svk}
  \index{Seifert--Van Kampen theorem!statement}
  Let $X$ be a topological space and let $U_1, U_2 \subseteq X$ be open
  subsets such that $X = U_1 \cup U_2$ and $U_0 := U_1 \cap U_2$ is
  path-connected and nonempty.  Choose a basepoint $x_0 \in U_0$.  Let
  $j_k\colon \pi_1(U_0, x_0) \to \pi_1(U_k, x_0)$ and
  $\iota_k\colon \pi_1(U_k, x_0) \to \pi_1(X, x_0)$ be the homomorphisms
  induced by the inclusions $U_0 \hookrightarrow U_k$ and
  $U_k \hookrightarrow X$.  Then the diagram
  \[
    \begin{tikzcd}[column sep=2em]
      \pi_1(U_0, x_0) \ar[r,"j_2"] \ar[d,"j_1"'] &
        \pi_1(U_2, x_0) \ar[d,"\iota_2"] \\
      \pi_1(U_1, x_0) \ar[r,"\iota_1"'] &
        \pi_1(X, x_0) \ar[ul, phantom, "\ulcorner" very near start]
    \end{tikzcd}
  \]
  is a pushout in the category of groups.  Equivalently,
  \[
    \pi_1(X, x_0) \;\cong\; \pi_1(U_1, x_0) \,*_{\pi_1(U_0, x_0)}\, \pi_1(U_2, x_0).
  \]
\end{theorem}

\begin{remark}[Hypotheses]\label{rem:svk-hypotheses}
  The hypothesis that $U_0 = U_1 \cap U_2$ is path-connected is essential.
  If $U_0$ has multiple path components, a more general version of the theorem
  holds involving groupoids (see Exercise~\ref{exo:svk-groupoid}).  Some
  authors also require each of $U_1, U_2, U_0$ to be path-connected; this is
  sufficient for most applications.
\end{remark}

\begin{proof}[Proof of Theorem~\ref{thm:svk}]
  \index{Seifert--Van Kampen theorem!proof}
  We must show that $\pi_1(X, x_0)$ satisfies the universal property of the
  amalgamated free product $\pi_1(U_1) *_{\pi_1(U_0)} \pi_1(U_2)$.

  \medskip
  \noindent\textbf{Step 1: Surjectivity.}
  We show that the natural map
  $\Phi\colon \pi_1(U_1) *_{\pi_1(U_0)} \pi_1(U_2) \to \pi_1(X, x_0)$ is
  surjective.  Let $[\gamma] \in \pi_1(X, x_0)$.  By the Lebesgue number
  lemma, we can subdivide $[0,1]$ into subintervals $[t_{i-1}, t_i]$,
  $i = 1, \ldots, n$, such that each $\gamma\!\big|_{[t_{i-1}, t_i]}$ has
  image contained in either $U_1$ or $U_2$.

  For each division point $t_i$ with $0 < i < n$, if $\gamma(t_i) \notin U_0$
  then the preceding and succeeding subpaths must lie in the same $U_k$, and we
  can merge them.  So we may assume $\gamma(t_i) \in U_1 \cup U_2$ for all $i$,
  and whenever the path crosses from $U_1$ to $U_2$ (or vice versa), the
  crossing point lies in $U_0$.

  Choose a path $\alpha_i$ in $U_0$ from $x_0$ to $\gamma(t_i)$ for each $i$
  with $\gamma(t_i) \in U_0$ (this is possible since $U_0$ is
  path-connected).  Then $\gamma$ is homotopic to a product of loops
  $\alpha_{i-1} \cdot \gamma_i \cdot \overline{\alpha}_i$, each of which lies
  entirely in $U_1$ or $U_2$.  Hence $[\gamma]$ is in the image of $\Phi$.

  \medskip
  \noindent\textbf{Step 2: Injectivity.}
  Suppose $w = g_1 g_2 \cdots g_m$ is a reduced word in
  $\pi_1(U_1) * \pi_1(U_2)$ that maps to the identity in $\pi_1(X, x_0)$.
  We must show that $w$ can be reduced to the empty word using only the
  amalgamation relations $j_1(k) = j_2(k)$, $k \in \pi_1(U_0)$.

  Represent each $g_i$ by a loop $\gamma_i$ in $U_1$ or $U_2$, so that
  $\gamma_1 \cdot \gamma_2 \cdots \gamma_m$ is null-homotopic in $X$ via some
  homotopy $H\colon [0,1] \times [0,1] \to X$.

  By the Lebesgue number lemma applied to the open cover $\{H^{-1}(U_1),
  H^{-1}(U_2)\}$ of $[0,1]^2$, we can subdivide the square into small
  rectangles, each mapping into $U_1$ or $U_2$.  The restriction of $H$ to
  each horizontal strip provides a homotopy between consecutive ``layers'' of
  loops.

  Tracking these subdivisions carefully, one shows that the word $w$ can be
  transformed to the empty word by a sequence of moves, each of which either:
  \begin{itemize}
    \item inserts or deletes a pair $g\, g^{-1}$ within the same factor, or
    \item replaces $j_1(k)$ by $j_2(k)$ (or vice versa) for some
      $k \in \pi_1(U_0)$.
  \end{itemize}
  These are precisely the relations defining the amalgamated free product.
  Hence $w$ represents the identity in $\pi_1(U_1) *_{\pi_1(U_0)} \pi_1(U_2)$,
  proving injectivity.

  \medskip
  \noindent\textbf{Step 3: Universal property.}
  Steps~1 and~2 together show that $\Phi$ is an isomorphism.  The universal
  property of $\pi_1(X, x_0)$ as a pushout now follows from the universal
  property of the amalgamated free product.
\end{proof}

%% ──────────────────────────────────────────────────────────
\section{Applications}\label{sec:svk-applications}

\begin{remark}[Presentations from the Seifert--Van Kampen theorem]\label{rem:svk-presentations}
  In practice, the Seifert--Van Kampen theorem is often used in the following
  form.  Suppose $\pi_1(U_1) = \langle S_1 \mid R_1 \rangle$,
  $\pi_1(U_2) = \langle S_2 \mid R_2 \rangle$, and
  $\pi_1(U_0) = \langle S_0 \mid R_0 \rangle$ with $j_1$ and $j_2$
  expressed on generators.  Then
  \[
    \pi_1(X) \;\cong\; \langle S_1 \cup S_2 \mid R_1 \cup R_2 \cup
      \set{j_1(s) = j_2(s) : s \in S_0} \rangle.
  \]
  That is, we combine all generators and all relations, and add new relations
  equating the images of the generators of $\pi_1(U_0)$ in the two factors.
\end{remark}

\begin{corollary}[Collapsing a simply connected subspace]\label{cor:svk-collapse}
  Let $A \subseteq X$ be a simply connected subspace such that $A$ is
  closed, $X$ is a nice space, and $A$ has a path-connected open neighborhood
  that deformation-retracts onto $A$.  Then the quotient map
  $q\colon X \to X/A$ induces an isomorphism $\pi_1(X) \cong \pi_1(X/A)$.
\end{corollary}

\subsection{Fundamental group of the wedge sum}\label{subsec:wedge-sum}

\begin{definition}[Wedge sum]\label{def:wedge-sum}
  \index{wedge sum}
  Let $(X, x_0)$ and $(Y, y_0)$ be pointed spaces.  The \textbf{wedge sum}
  (or one-point union) is
  $X \vee Y := (X \sqcup Y) / (x_0 \sim y_0)$.
\end{definition}

\begin{center}
  \begin{tikzpicture}[scale=0.9]
    % First circle
    \draw[thick, blue!70!black] (0,0) circle (1);
    \node[below left, blue!70!black] at (-0.7,-0.7) {$S^1$};
    % Second circle
    \draw[thick, red!70!black] (2,0) circle (1);
    \node[below right, red!70!black] at (2.7,-0.7) {$S^1$};
    % Wedge point
    \fill (1,0) circle (2.5pt);
    \node[above] at (1,0.15) {$x_0$};
    % Label
    \node[below] at (1,-1.5) {$S^1 \vee S^1$};
  \end{tikzpicture}
\end{center}

\begin{theorem}[Fundamental group of a wedge sum]\label{thm:pi1-wedge}
  \index{fundamental group!of wedge sum}
  Let $X$ and $Y$ be path-connected spaces with non-degenerate basepoints
  (meaning each basepoint has a contractible open neighborhood).  Then
  \[
    \pi_1(X \vee Y) \;\cong\; \pi_1(X) * \pi_1(Y).
  \]
\end{theorem}

\begin{proof}
  Let $p$ denote the wedge point.  By the non-degeneracy hypothesis, there
  exist contractible open neighborhoods $V_X \ni x_0$ in $X$ and
  $V_Y \ni y_0$ in $Y$.  Set
  \[
    U_1 = X \cup V_Y, \qquad U_2 = V_X \cup Y, \qquad U_0 = U_1 \cap U_2 = V_X \cup V_Y,
  \]
  where the unions are taken inside $X \vee Y$.  Then $U_0$ is contractible
  (it is the wedge of two contractible sets along a point), so
  $\pi_1(U_0) = \{e\}$.  Moreover, $U_1$ deformation-retracts onto $X$ and
  $U_2$ deformation-retracts onto $Y$.  By the Seifert--Van Kampen theorem,
  \[
    \pi_1(X \vee Y) \;\cong\; \pi_1(U_1) *_{\{e\}} \pi_1(U_2) \;\cong\;
    \pi_1(X) * \pi_1(Y). \qedhere
  \]
\end{proof}

\begin{corollary}[$\pi_1(S^1 \vee S^1)$]\label{cor:pi1-bouquet}
  \index{fundamental group!of figure eight}
  $\pi_1(S^1 \vee S^1) \cong \Z * \Z$, the free group on two generators.
  More generally, $\pi_1\!\left(\bigvee_{n} S^1\right) \cong F_n$, the free
  group of rank $n$.
\end{corollary}

\subsection{The real projective plane}\label{subsec:RP2}

\begin{theorem}[$\pi_1(\RP^2)$]\label{thm:pi1-RP2}
  \index{fundamental group!of projective plane}\index{projective plane}
  $\pi_1(\RP^2) \cong \Z/2\Z$.
\end{theorem}

\begin{proof}
  Recall that $\RP^2$ is obtained from $D^2$ by identifying antipodal points
  on the boundary $S^1$.  Equivalently, $\RP^2$ has a CW~structure with one
  $0$-cell, one $1$-cell, and one $2$-cell attached by a map of degree~$2$.
  Concretely, let $a$ denote the loop generating $\pi_1$ of the $1$-skeleton
  (which is $S^1$).  The $2$-cell is attached along $a^2$.

  By the Seifert--Van Kampen theorem (cell attachment version,
  Proposition~\ref{prop:cell-attachment}),
  \[
    \pi_1(\RP^2) \;\cong\; \langle a \mid a^2 = 1 \rangle \;\cong\; \Z/2\Z. \qedhere
  \]
\end{proof}

\begin{center}
  \begin{tikzpicture}[scale=1.2]
    % Disk with identification
    \draw[thick] (0,0) circle (1.5);
    % Arrows showing identification
    \draw[thick, -Stealth, blue!70!black] (1.5,0) arc[start angle=0, end angle=85, radius=1.5];
    \draw[thick, -Stealth, blue!70!black] (-1.5,0) arc[start angle=180, end angle=265, radius=1.5];
    \draw[thick, -Stealth, red!70!black] (0,1.5) arc[start angle=90, end angle=175, radius=1.5];
    \draw[thick, -Stealth, red!70!black] (0,-1.5) arc[start angle=270, end angle=355, radius=1.5];
    % Labels
    \node[right, blue!70!black] at (1.55,0.7) {$a$};
    \node[left, blue!70!black] at (-1.55,-0.7) {$a$};
    \node[above left, red!70!black] at (-0.7,1.55) {$a$};
    \node[below right, red!70!black] at (0.7,-1.55) {$a$};
    \node at (0,0) {$\RP^2$};
  \end{tikzpicture}
\end{center}

\subsection{The Klein bottle}\label{subsec:klein}

\begin{theorem}[$\pi_1$ of the Klein bottle]\label{thm:pi1-klein}
  \index{fundamental group!of Klein bottle}\index{Klein bottle}
  The Klein bottle $K$ satisfies
  \[
    \pi_1(K) \;\cong\; \langle a, b \mid abab^{-1} = 1 \rangle.
  \]
\end{theorem}

\begin{proof}
  The Klein bottle is obtained from a square by the identifications shown below:

  \begin{center}
    \begin{tikzpicture}[scale=1.8]
      \draw[thick, blue!70!black, -Stealth] (0,0) -- (1,0);
      \draw[thick, blue!70!black, -Stealth] (2,0) -- (1,0); % reversed!
      \draw[thick, red!70!black, -Stealth] (0,0) -- (0,1);
      \draw[thick, red!70!black, -Stealth] (2,0) -- (2,1);
      \draw[thick, blue!70!black, -Stealth] (0,2) -- (1,2);
      \draw[thick, blue!70!black, -Stealth] (2,2) -- (1,2); % reversed!
      \draw[thick, red!70!black, -Stealth] (0,1) -- (0,2);
      \draw[thick, red!70!black, -Stealth] (2,1) -- (2,2);
      % Labels
      \node[below, blue!70!black] at (1,0) {$a$};
      \node[above, blue!70!black] at (1,2) {$a$};
      \node[left, red!70!black] at (0,1) {$b$};
      \node[right, red!70!black] at (2,1) {$b$};
    \end{tikzpicture}
  \end{center}

  The Klein bottle has a CW~structure with one $0$-cell, two $1$-cells $a$
  and $b$, and one $2$-cell attached along the word $abab^{-1}$.  The
  $1$-skeleton is $S^1 \vee S^1$, so $\pi_1$ of the $1$-skeleton is $F(a,b)$.
  Attaching the $2$-cell kills the element $abab^{-1}$, giving
  \[
    \pi_1(K) \cong \langle a, b \mid abab^{-1} \rangle. \qedhere
  \]
\end{proof}

\begin{remark}[Klein bottle group]\label{rem:klein-nonabelian}
  The group $\pi_1(K) = \langle a, b \mid aba = b \rangle$ is a non-abelian
  infinite group.  Its abelianization is $\Z \oplus \Z/2\Z$.  In particular,
  $K$ is not orientable (the torsion detects non-orientability).
\end{remark}

\subsection{Orientable surfaces}\label{subsec:surfaces}

\begin{theorem}[Fundamental group of orientable surfaces]\label{thm:pi1-surface}
  \index{fundamental group!of surface}\index{surface group}
  Let $\Sigma_g$ denote the closed orientable surface of genus $g \geq 1$.
  Then
  \[
    \pi_1(\Sigma_g) \;\cong\;
    \left\langle a_1, b_1, \ldots, a_g, b_g \;\middle|\;
    \prod_{i=1}^{g} [a_i, b_i] = 1 \right\rangle,
  \]
  where $[a_i, b_i] = a_i b_i a_i^{-1} b_i^{-1}$.
\end{theorem}

\begin{proof}
  The surface $\Sigma_g$ is obtained from a $4g$-gon by identifying edges in
  pairs according to the word
  $a_1 b_1 a_1^{-1} b_1^{-1} a_2 b_2 a_2^{-1} b_2^{-1} \cdots a_g b_g
  a_g^{-1} b_g^{-1}$.  All vertices are identified to a single point.

  \begin{center}
    \begin{tikzpicture}[scale=1.6]
      % Octagon for genus 2
      \def\n{8}
      \foreach \i in {1,...,\n} {
        \coordinate (P\i) at ({90+360/\n*(\i-1)}:1.5);
      }
      % Draw edges with labels (genus 2 identification)
      \draw[thick, blue!70!black, -Stealth] (P1) -- (P2)
        node[midway, above right] {$a_1$};
      \draw[thick, red!70!black, -Stealth] (P2) -- (P3)
        node[midway, above left] {$b_1$};
      \draw[thick, blue!70!black, Stealth-] (P3) -- (P4)
        node[midway, left] {$a_1$};
      \draw[thick, red!70!black, Stealth-] (P4) -- (P5)
        node[midway, below left] {$b_1$};
      \draw[thick, green!50!black, -Stealth] (P5) -- (P6)
        node[midway, below left] {$a_2$};
      \draw[thick, orange!80!black, -Stealth] (P6) -- (P7)
        node[midway, below right] {$b_2$};
      \draw[thick, green!50!black, Stealth-] (P7) -- (P8)
        node[midway, right] {$a_2$};
      \draw[thick, orange!80!black, Stealth-] (P8) -- (P1)
        node[midway, above right] {$b_2$};
      % Center label
      \node at (0,0) {$\Sigma_2$};
    \end{tikzpicture}
  \end{center}

  The CW~structure has one $0$-cell, $2g$ oriented $1$-cells
  $a_1, b_1, \ldots, a_g, b_g$, and one $2$-cell.  The $1$-skeleton is a wedge
  of $2g$ circles, so $\pi_1$ of the $1$-skeleton is $F(a_1, b_1, \ldots,
  a_g, b_g)$.  Attaching the $2$-cell along the boundary word
  $\prod_{i=1}^g [a_i, b_i]$ gives the result by the cell attachment
  proposition.
\end{proof}

\begin{corollary}[$\pi_1$ of the torus]\label{cor:pi1-torus}
  \index{fundamental group!of torus}
  $\pi_1(T^2) = \pi_1(\Sigma_1) \cong \langle a, b \mid [a,b] = 1 \rangle
  \cong \Z \times \Z$.
\end{corollary}

\begin{remark}[Euler characteristic check]\label{rem:euler-char-surfaces}
  \index{Euler characteristic!of surfaces}
  The surface $\Sigma_g$ has a CW~structure with $1$ vertex, $2g$ edges, and
  $1$ face, so $\chi(\Sigma_g) = 1 - 2g + 1 = 2 - 2g$.  The abelianization
  of $\pi_1(\Sigma_g)$ is $\Z^{2g}$ (see Exercise~\ref{exo:svk-abelianize}),
  which is consistent with $H_1(\Sigma_g; \Z) \cong \Z^{2g}$ and the
  relation $\chi = \sum (-1)^k \mathrm{rank}(H_k) = 1 - 2g + 1 = 2 - 2g$.
\end{remark}

\begin{proposition}[Fundamental group of connected sums]\label{prop:pi1-connected-sum}
  \index{connected sum!fundamental group}
  Let $M$ and $N$ be closed connected $n$-manifolds with $n \geq 3$.  Then
  \[
    \pi_1(M \mathbin{\#} N) \;\cong\; \pi_1(M) * \pi_1(N).
  \]
  For surfaces ($n = 2$), the same formula holds at the level of the
  $1$-skeleton, but the extra $2$-cell relation must be accounted for.
\end{proposition}

\begin{proof}
  The connected sum $M \# N$ is obtained by removing open balls $B_M
  \subset M$ and $B_N \subset N$ and gluing along the resulting boundary
  spheres $S^{n-1}$.  Take $U_1 = M \setminus \{p\}$ and
  $U_2 = N \setminus \{q\}$ (where $p, q$ are interior points of the removed
  balls), thickened slightly so that $U_0 = U_1 \cap U_2$ is homotopy
  equivalent to $S^{n-1}$.  For $n \geq 3$, $\pi_1(S^{n-1}) = \{e\}$,
  so the Seifert--Van Kampen theorem gives a free product.
\end{proof}

\subsection{Cell attachments}\label{subsec:cell-attachment}

\begin{proposition}[Effect of attaching a $2$-cell]\label{prop:cell-attachment}
  \index{cell attachment}\index{CW complex!cell attachment}
  Let $X$ be a path-connected space and let $Y = X \cup_\varphi e^2$ be
  obtained by attaching a $2$-cell along a map $\varphi\colon S^1 \to X$.
  Let $[\varphi] \in \pi_1(X, x_0)$ be the class of the attaching map
  (after choosing a path from $\varphi(1)$ to $x_0$ if needed).  Then
  \[
    \pi_1(Y, x_0) \;\cong\; \pi_1(X, x_0) \,\big/\, \langle\!\langle [\varphi] \rangle\!\rangle,
  \]
  where $\langle\!\langle [\varphi] \rangle\!\rangle$ denotes the normal
  closure.
\end{proposition}

\begin{proof}
  We apply the Seifert--Van Kampen theorem.  Let $U_1 = X \cup
  (\text{open collar of } e^2)$ and let $U_2$ be the open disk $e^2$.  Then
  $U_1$ deformation-retracts onto $X$, $U_2$ is contractible, and
  $U_0 = U_1 \cap U_2$ deformation-retracts onto $S^1$.  The map
  $\pi_1(U_0) \to \pi_1(U_1)$ sends the generator of $\pi_1(S^1) \cong \Z$ to
  $[\varphi]$, while $\pi_1(U_0) \to \pi_1(U_2) = \{e\}$ is trivial.
  Therefore,
  \[
    \pi_1(Y) \cong \pi_1(X) *_\Z \{e\} \cong \pi_1(X) / \langle\!\langle [\varphi] \rangle\!\rangle. \qedhere
  \]
\end{proof}

\begin{corollary}[Fundamental group of CW complexes]\label{cor:pi1-CW}
  \index{CW complex!fundamental group}
  Let $X$ be a CW~complex.  Then $\pi_1(X)$ depends only on the $2$-skeleton
  $X^{(2)}$ and is given by
  \[
    \pi_1(X) \;\cong\; \left\langle
      \text{$1$-cells} \;\middle|\; \text{attaching words of $2$-cells}
    \right\rangle.
  \]
  Moreover, attaching cells of dimension $\geq 3$ does not change $\pi_1$.
\end{corollary}

\begin{proof}
  The $1$-skeleton $X^{(1)}$ is a graph, so $\pi_1(X^{(1)})$ is a free group
  (on generators corresponding to $1$-cells not in a maximal tree, by
  Proposition~\ref{prop:free-product-props}).  Attaching $2$-cells imposes
  the stated relations by iterated application of
  Proposition~\ref{prop:cell-attachment}.

  For cells of dimension $n \geq 3$: the attaching map $S^{n-1} \to X^{(n-1)}$
  satisfies $\pi_1(S^{n-1}) = 0$ for $n-1 \geq 2$, so the amalgamation is
  trivial and $\pi_1$ does not change.
\end{proof}

\begin{proposition}[Every finitely presented group is realizable]\label{prop:every-fp-group}
  \index{finitely presented group}
  For every finitely presented group $G = \langle g_1, \ldots, g_n \mid
  r_1, \ldots, r_m \rangle$, there exists a compact CW~complex $X$ with
  $\pi_1(X) \cong G$.
\end{proposition}

\begin{proof}
  Take $X^{(1)} = \bigvee_{i=1}^n S^1$ and attach $m$ copies of $2$-cells
  along the relators $r_1, \ldots, r_m$.  By
  Corollary~\ref{cor:pi1-CW}, $\pi_1(X) \cong G$.
\end{proof}

%% ──────────────────────────────────────────────────────────
\section{Exercises}\label{sec:svk-exercises}

\begin{exercise}\label{exo:svk-Sn}
  Use the Seifert--Van Kampen theorem to show that $\pi_1(S^n) = \{e\}$ for
  $n \geq 2$.

  \textit{Hint:} Write $S^n$ as the union of two open sets, each
  homeomorphic to $\R^n$, whose intersection is homotopy equivalent to
  $S^{n-1}$ (which is connected for $n \geq 2$).
\end{exercise}

\begin{exercise}\label{exo:svk-torus-direct}
  Give a direct proof using the Seifert--Van Kampen theorem that
  $\pi_1(T^2) \cong \Z \times \Z$, by decomposing the torus into two
  suitable open sets.  Compare with the CW~complex approach.
\end{exercise}

\begin{exercise}\label{exo:svk-RP2-wedge}
  Compute $\pi_1(\RP^2 \vee S^1)$.
\end{exercise}

\begin{exercise}\label{exo:svk-non-orientable}
  \index{non-orientable surface}
  Let $N_k$ denote the connected sum of $k$ copies of $\RP^2$ (the
  non-orientable surface of genus $k$).  Show that
  \[
    \pi_1(N_k) \cong \langle c_1, \ldots, c_k \mid c_1^2 c_2^2 \cdots c_k^2 = 1 \rangle.
  \]
  Verify that $N_1 = \RP^2$ and $N_2$ is the Klein bottle.
\end{exercise}

\begin{exercise}\label{exo:svk-theta}
  Let $\Theta$ be the \emph{theta curve}: the space formed by two vertices
  $p, q$ joined by three distinct arcs.  Compute $\pi_1(\Theta)$.
\end{exercise}

\begin{exercise}\label{exo:svk-dunce}
  \index{dunce hat}
  The \textbf{dunce hat} is obtained from a triangle by identifying all three
  edges via $a \to a \to a$ (all edges are identified with the same
  orientation).  Show that the dunce hat is simply connected.
\end{exercise}

\begin{exercise}\label{exo:svk-abelianize}
  Let $G = \pi_1(\Sigma_g)$.  Show that $G^{\mathrm{ab}} \cong \Z^{2g}$.
  Deduce that surfaces of different genus are not homotopy equivalent.
\end{exercise}

\begin{exercise}\label{exo:svk-genus2}
  Let $\Sigma_2$ be the closed orientable surface of genus $2$.  Show that
  $\Sigma_2$ can be decomposed as the union of two tori-with-a-hole.  Use
  the Seifert--Van Kampen theorem directly (not via the polygon model) to
  recover the presentation
  $\langle a_1, b_1, a_2, b_2 \mid [a_1, b_1][a_2, b_2] = 1 \rangle$.
\end{exercise}

\begin{exercise}\label{exo:svk-mapping-torus}
  \index{mapping torus}
  Let $f\colon S^1 \to S^1$ be the map $z \mapsto z^n$.  The
  \textbf{mapping torus} $M_f$ is obtained from $S^1 \times [0,1]$ by
  identifying $(z,0) \sim (f(z), 1)$.  Show that
  $\pi_1(M_f) \cong \langle a, t \mid t a t^{-1} = a^n \rangle$.
\end{exercise}

\begin{exercise}\label{exo:svk-groupoid}
  \index{Seifert--Van Kampen theorem!groupoid version}
  \textbf{(Groupoid version.)}
  Formulate and prove a version of the Seifert--Van Kampen theorem when
  $U_1 \cap U_2$ has two path components.  Apply it to compute
  $\pi_1(S^1)$ by decomposing $S^1$ into two open arcs.
\end{exercise}

%% ══════════════════════════════════════════════════════════
%%  CHAPTER 4: COVERING SPACES AND THE FUNDAMENTAL GROUP
%% ══════════════════════════════════════════════════════════

\chapter{Covering Spaces and the Fundamental Group}\label{ch:coverings}

Covering space theory provides one of the most elegant interactions between
algebra and topology.  The key result---the \emph{Galois
correspondence}\index{Galois correspondence}---establishes a bijection between
connected coverings of a ``nice'' space $X$ and subgroups of $\pi_1(X)$,
paralleling the correspondence between field extensions and subgroups of a
Galois group.  Along the way, we develop the crucial lifting properties that
make the theory work.

%% ──────────────────────────────────────────────────────────
\section{Definitions and first examples}\label{sec:covering-def}

\begin{definition}[Covering space]\label{def:covering}
  \index{covering space}
  A \textbf{covering space} (or \textbf{cover}) of a topological space $X$
  is a topological space $\widetilde{X}$ together with a surjective continuous
  map $p\colon \widetilde{X} \to X$ such that for every $x \in X$, there
  exists an open neighborhood $U$ of $x$ (called an
  \textbf{evenly covered neighborhood}\index{evenly covered neighborhood})
  for which
  \[
    p^{-1}(U) = \bigsqcup_{\alpha \in A} V_\alpha,
  \]
  where each $V_\alpha$ is open in $\widetilde{X}$ and
  $p|_{V_\alpha}\colon V_\alpha \xrightarrow{\;\sim\;} U$ is a homeomorphism.
  The sets $V_\alpha$ are called \textbf{sheets}\index{sheets (covering)} over
  $U$.
\end{definition}

\begin{definition}[Degree of a covering]\label{def:covering-degree}
  \index{covering space!degree}
  If $X$ is connected, the cardinality $\abs{p^{-1}(x)}$ is the same for all
  $x \in X$.  This common value is called the \textbf{degree} (or
  \textbf{number of sheets}) of the covering.
\end{definition}

\begin{example}[The exponential map]\label{ex:exp-covering}
  \index{covering space!of the circle}
  The map $p\colon \R \to S^1$ defined by $p(t) = e^{2\pi i t}$ is a
  covering of infinite degree.  For any open arc $U \subset S^1$ of length
  less than $2\pi$, the preimage $p^{-1}(U)$ is a disjoint union of open
  intervals in $\R$, each mapped homeomorphically onto $U$.

  \begin{center}
    \begin{tikzpicture}[scale=0.85]
      % The helix
      \draw[thick, blue!70!black, domain=0:720, samples=200, smooth]
        plot ({1.2*cos(\x)}, {0.03*\x - 0.5});
      % The circle
      \draw[thick, red!70!black] (0,-3.5) circle (1.2);
      % Projection arrow
      \draw[-Stealth, thick, dashed] (0,-0.5) -- (0,-2.2)
        node[midway, right] {$p$};
      % Labels
      \node[right, blue!70!black] at (1.4, 5) {$\R$};
      \node[right, red!70!black] at (1.4, -3.5) {$S^1$};
    \end{tikzpicture}
  \end{center}
\end{example}

\begin{example}[Finite cyclic coverings of $S^1$]\label{ex:cyclic-covering-S1}
  \index{covering space!cyclic}
  The map $p_n\colon S^1 \to S^1$ given by $z \mapsto z^n$ is an $n$-sheeted
  covering.
\end{example}

\begin{example}[Covering of $\RP^2$]\label{ex:covering-RP2}
  \index{covering space!of projective plane}
  The quotient map $p\colon S^2 \to \RP^2$ sending each point to its
  equivalence class under the antipodal identification is a $2$-sheeted
  covering.  Here $S^2$ is simply connected, and the two sheets over any
  evenly covered set correspond to a point and its antipode.
\end{example}

\begin{example}[Covering of the torus]\label{ex:covering-torus}
  \index{covering space!of torus}
  The map $p\colon \R^2 \to T^2 = \R^2/\Z^2$ is the universal covering.
  For each $n \geq 1$, the torus $T^2$ also has an $n^2$-sheeted covering
  $T^2 \to T^2$ induced by the linear map $(x,y) \mapsto (nx, ny)$.

  \begin{center}
    \begin{tikzpicture}[scale=0.6]
      % Grid for R^2
      \draw[step=1, gray!40, thin] (-0.5,-0.5) grid (4.5,4.5);
      \draw[thick, ->] (-0.5,0) -- (4.8,0);
      \draw[thick, ->] (0,-0.5) -- (0,4.8);
      % Fundamental domain
      \fill[blue!10] (0,0) rectangle (1,1);
      \draw[thick, blue!70!black] (0,0) rectangle (1,1);
      % Arrow
      \draw[-Stealth, thick] (5.5,2) -- (7,2) node[midway, above] {$p$};
      % Torus (schematic)
      \draw[thick, red!70!black] (9,2) ellipse (1.5 and 1);
      \draw[thick, red!70!black] (8.2,2) arc[start angle=180, end angle=360,
        x radius=0.8, y radius=0.35];
      \draw[thick, red!70!black, dashed] (8.2,2) arc[start angle=180, end angle=0,
        x radius=0.8, y radius=0.35];
      % Labels
      \node[above right] at (4.5, 4.5) {$\R^2$};
      \node[below, red!70!black] at (9,0.7) {$T^2$};
    \end{tikzpicture}
  \end{center}
\end{example}

\begin{definition}[Morphism of coverings]\label{def:covering-morphism}
  \index{covering space!morphism}
  Let $p_1\colon \widetilde{X}_1 \to X$ and $p_2\colon \widetilde{X}_2 \to X$
  be coverings of $X$.  A \textbf{morphism of coverings} is a continuous map
  $f\colon \widetilde{X}_1 \to \widetilde{X}_2$ such that $p_2 \circ f = p_1$:
  \[
    \begin{tikzcd}
      \widetilde{X}_1 \ar[rr,"f"] \ar[dr,"p_1"'] & & \widetilde{X}_2 \ar[dl,"p_2"] \\
      & X &
    \end{tikzcd}
  \]
  An isomorphism of coverings is a morphism that is also a homeomorphism.
\end{definition}

%% ──────────────────────────────────────────────────────────
\section{Lifting properties}\label{sec:lifting}

The fundamental property of covering spaces is the ability to lift paths
and homotopies.

\begin{lemma}[Path lifting]\label{lem:path-lifting}
  \index{path lifting}\index{lifting!of paths}
  Let $p\colon \widetilde{X} \to X$ be a covering, let
  $\gamma\colon [0,1] \to X$ be a path, and let
  $\tilde{x}_0 \in p^{-1}(\gamma(0))$.  Then there exists a unique path
  $\tilde{\gamma}\colon [0,1] \to \widetilde{X}$ such that
  $p \circ \tilde{\gamma} = \gamma$ and $\tilde{\gamma}(0) = \tilde{x}_0$.
  \[
    \begin{tikzcd}
      & \widetilde{X} \ar[d,"p"] \\
      {[0,1]} \ar[r,"\gamma"'] \ar[ur,dashed,"\exists!\,\tilde\gamma"] & X
    \end{tikzcd}
  \]
\end{lemma}

\begin{proof}
  \textbf{Existence.}
  Cover $\gamma([0,1])$ by evenly covered neighborhoods.  By compactness of
  $[0,1]$, find a partition $0 = t_0 < t_1 < \cdots < t_n = 1$ such that each
  $\gamma([t_{i-1}, t_i])$ is contained in an evenly covered set $U_i$.

  We construct $\tilde{\gamma}$ inductively.  On $[0, t_1]$: the preimage
  $p^{-1}(U_1)$ decomposes into sheets; exactly one sheet $V_1$ contains
  $\tilde{x}_0$.  Define $\tilde{\gamma} = (p|_{V_1})^{-1} \circ \gamma$ on
  $[0, t_1]$.  Proceeding step by step, each lift is determined by the
  requirement that $\tilde{\gamma}$ be continuous.

  \textbf{Uniqueness.}
  If $\tilde{\gamma}_1$ and $\tilde{\gamma}_2$ are two lifts with the same
  initial point, the set $\{t : \tilde{\gamma}_1(t) = \tilde{\gamma}_2(t)\}$
  is both open and closed in $[0,1]$ (by the homeomorphism property of
  sheets), nonempty (it contains $0$), hence equal to $[0,1]$.
\end{proof}

\begin{lemma}[Homotopy lifting]\label{lem:homotopy-lifting}
  \index{homotopy lifting}\index{lifting!of homotopies}
  Let $p\colon \widetilde{X} \to X$ be a covering, let
  $H\colon [0,1] \times [0,1] \to X$ be a homotopy with
  $H(0,s) = x_0$ for all $s$, and let $\tilde{x}_0 \in p^{-1}(x_0)$.
  Then there exists a unique lift
  $\widetilde{H}\colon [0,1] \times [0,1] \to \widetilde{X}$ with
  $p \circ \widetilde{H} = H$ and $\widetilde{H}(0,s) = \tilde{x}_0$ for all
  $s$.
  \[
    \begin{tikzcd}
      {[0,1] \times \{0\}} \ar[r,"\tilde{\gamma}_0"] \ar[d,hook] &
        \widetilde{X} \ar[d,"p"] \\
      {[0,1] \times [0,1]} \ar[r,"H"'] \ar[ur,dashed,"\widetilde{H}"] & X
    \end{tikzcd}
  \]
\end{lemma}

\begin{proof}
  The argument is analogous to path lifting.  Subdivide $[0,1]^2$ into small
  squares, each mapping into an evenly covered set, and lift inductively, one
  square at a time, working row by row or column by column.  Uniqueness follows
  from the same open-and-closed argument.
\end{proof}

\begin{corollary}[Lifting loops and $p_*$]\label{cor:lifting-loops}
  Let $p\colon \widetilde{X} \to X$ be a covering and
  $\tilde{x}_0 \in p^{-1}(x_0)$.  A loop $\gamma$ at $x_0$ lifts to a
  loop at $\tilde{x}_0$ if and only if
  $[\gamma] \in p_*(\pi_1(\widetilde{X}, \tilde{x}_0))$.
\end{corollary}

\begin{proof}
  The lift $\tilde{\gamma}$ starts at $\tilde{x}_0$ and ends at some point
  $\tilde{x}_1 \in p^{-1}(x_0)$.  By definition, $\tilde{\gamma}$ is a loop
  if and only if $\tilde{x}_1 = \tilde{x}_0$, which happens if and only if
  $[\gamma]$ is in the image of $p_*$.
\end{proof}

\begin{example}[Lifting loops to $\R$]\label{ex:lifting-loops-R}
  Consider the covering $p\colon \R \to S^1$, $t \mapsto e^{2\pi i t}$, with
  basepoint lift $\tilde{x}_0 = 0$.  The loop $\gamma_n(t) = e^{2\pi i n t}$
  (winding $n$ times) lifts to the path $\tilde{\gamma}_n(t) = nt$ in $\R$.
  This lift is a loop (i.e., $\tilde{\gamma}_n(1) = \tilde{\gamma}_n(0) = 0$)
  only when $n = 0$, consistent with $p_*(\pi_1(\R)) = \{0\} \leq \Z$.

  For the $k$-sheeted covering $p_k\colon S^1 \to S^1$, $z \mapsto z^k$,
  the loop $\gamma_n$ lifts to a loop if and only if $k \mid n$, consistent
  with $p_{k*}(\pi_1(S^1)) = k\Z \leq \Z$.
\end{example}

\begin{theorem}[Monodromy theorem]\label{thm:monodromy-thm}
  \index{monodromy theorem}
  Let $p\colon \widetilde{X} \to X$ be a covering and let $\gamma_0, \gamma_1$
  be paths in $X$ from $x_0$ to $x_1$ that are homotopic relative to
  endpoints.  Let $\tilde{\gamma}_0, \tilde{\gamma}_1$ be their lifts starting
  at the same point $\tilde{x}_0 \in p^{-1}(x_0)$.  Then
  $\tilde{\gamma}_0(1) = \tilde{\gamma}_1(1)$, and the lifts are homotopic
  relative to endpoints in $\widetilde{X}$.
\end{theorem}

\begin{proof}
  Let $H$ be a homotopy from $\gamma_0$ to $\gamma_1$ relative to endpoints.
  By the homotopy lifting lemma (Lemma~\ref{lem:homotopy-lifting}), $H$ lifts
  to $\widetilde{H}$ with $\widetilde{H}(\cdot, 0) = \tilde{\gamma}_0$.
  Since $H(1,s) = x_1$ for all $s$, the path $s \mapsto \widetilde{H}(1,s)$
  lies in the discrete fiber $p^{-1}(x_1)$; by continuity, it is constant.
  Thus $\tilde{\gamma}_0(1) = \widetilde{H}(1,0) = \widetilde{H}(1,1) =
  \tilde{\gamma}_1(1)$.  Similarly, $\widetilde{H}(\cdot, 1)$ is the unique
  lift of $\gamma_1$ starting at $\tilde{x}_0$, so
  $\widetilde{H}(\cdot, 1) = \tilde{\gamma}_1$.  Hence $\widetilde{H}$ is a
  homotopy from $\tilde{\gamma}_0$ to $\tilde{\gamma}_1$ relative to
  endpoints.
\end{proof}

%% ──────────────────────────────────────────────────────────
\section{The monodromy action}\label{sec:monodromy-action}

\begin{definition}[Monodromy action]\label{def:monodromy-action}
  \index{monodromy action}
  Let $p\colon \widetilde{X} \to X$ be a covering and let $x_0 \in X$.  The
  \textbf{monodromy action} of $\pi_1(X, x_0)$ on the fiber
  $F = p^{-1}(x_0)$ is defined by
  \[
    [\gamma] \cdot \tilde{x}_0 \;:=\; \tilde{\gamma}(1),
  \]
  where $\tilde{\gamma}$ is the unique lift of $\gamma$ starting at
  $\tilde{x}_0$.  This is well-defined by the monodromy theorem
  (Theorem~\ref{thm:monodromy-thm}).
\end{definition}

\begin{proposition}[Properties of the monodromy action]\label{prop:monodromy-props}
  The monodromy action is a right action of $\pi_1(X, x_0)$ on $F$.
  Moreover:
  \begin{enumerate}[label=(\roman*)]
    \item If $\widetilde{X}$ is path-connected, the action is
      \textbf{transitive}.\index{transitive action}
    \item The stabilizer of $\tilde{x}_0 \in F$ is the subgroup
      $p_*(\pi_1(\widetilde{X}, \tilde{x}_0)) \leq \pi_1(X, x_0)$.
  \end{enumerate}
\end{proposition}

\begin{proof}
  The constant loop lifts to the constant path, so the identity acts
  trivially.  If $\gamma$ and $\delta$ are loops at $x_0$, then the lift of
  $\gamma \cdot \delta$ starting at $\tilde{x}_0$ is the concatenation of the
  lift of $\gamma$ starting at $\tilde{x}_0$ with the lift of $\delta$
  starting at $\tilde{\gamma}(1)$.  This verifies the action axiom.

  (i) If $\widetilde{X}$ is path-connected and $\tilde{x}_0, \tilde{x}_1 \in
  F$, choose a path $\tilde{\alpha}$ in $\widetilde{X}$ from $\tilde{x}_0$ to
  $\tilde{x}_1$.  Then $\gamma := p \circ \tilde{\alpha}$ is a loop at $x_0$
  and $[\gamma] \cdot \tilde{x}_0 = \tilde{x}_1$.

  (ii) $[\gamma] \cdot \tilde{x}_0 = \tilde{x}_0$ if and only if the lift
  $\tilde{\gamma}$ is a loop at $\tilde{x}_0$, i.e.,
  $[\gamma] = p_*([\tilde{\gamma}]) \in p_*(\pi_1(\widetilde{X},
  \tilde{x}_0))$.
\end{proof}

\begin{corollary}[Index formula]\label{cor:index-formula}
  \index{covering space!degree}
  If $p\colon \widetilde{X} \to X$ is a connected covering, then
  \[
    \abs{p^{-1}(x_0)} \;=\; [\pi_1(X, x_0) : p_*(\pi_1(\widetilde{X}, \tilde{x}_0))].
  \]
  In particular, the number of sheets equals the index of the image of
  $\pi_1(\widetilde{X})$ in $\pi_1(X)$.
\end{corollary}

\begin{proof}
  By the orbit-stabilizer theorem applied to the transitive monodromy action,
  $\abs{F} = [\pi_1(X, x_0) : \mathrm{Stab}(\tilde{x}_0)] = [\pi_1(X, x_0) :
  p_*(\pi_1(\widetilde{X}, \tilde{x}_0))]$.
\end{proof}

%% ──────────────────────────────────────────────────────────
\section{The induced homomorphism and the lifting criterion}\label{sec:induced-homo}

\begin{proposition}[Injectivity of $p_*$]\label{prop:p-star-injective}
  \index{covering space!induced homomorphism}
  Let $p\colon \widetilde{X} \to X$ be a covering and $\tilde{x}_0 \in
  p^{-1}(x_0)$.  The induced homomorphism
  $p_*\colon \pi_1(\widetilde{X}, \tilde{x}_0) \to \pi_1(X, x_0)$ is
  injective.
\end{proposition}

\begin{proof}
  Suppose $p_*([\tilde{\gamma}]) = [p \circ \tilde{\gamma}] = [c_{x_0}]$ in
  $\pi_1(X, x_0)$.  Then $p \circ \tilde{\gamma}$ is null-homotopic in $X$
  via some homotopy $H$.  Lift $H$ to $\widetilde{H}$ in $\widetilde{X}$.
  Since $H(1, s) = x_0$ for all $s$, the path $s \mapsto \widetilde{H}(1,s)$
  lies in the discrete fiber $p^{-1}(x_0)$ and is therefore constant.
  Hence $\widetilde{H}$ is a null-homotopy of $\tilde{\gamma}$ in
  $\widetilde{X}$.
\end{proof}

\begin{theorem}[General lifting criterion]\label{thm:lifting-criterion}
  \index{lifting!criterion}
  Let $p\colon \widetilde{X} \to X$ be a covering with $\widetilde{X}$
  path-connected.  Let $Y$ be a connected and locally path-connected space,
  $y_0 \in Y$, and $f\colon (Y, y_0) \to (X, x_0)$ a continuous map.
  Choose $\tilde{x}_0 \in p^{-1}(x_0)$.  Then there exists a lift
  $\tilde{f}\colon (Y, y_0) \to (\widetilde{X}, \tilde{x}_0)$ with
  $p \circ \tilde{f} = f$ if and only if
  \[
    f_*\!\big(\pi_1(Y, y_0)\big) \;\subseteq\; p_*\!\big(\pi_1(\widetilde{X}, \tilde{x}_0)\big).
  \]
  Moreover, when it exists, such a lift is unique.
  \[
    \begin{tikzcd}
      & (\widetilde{X}, \tilde{x}_0) \ar[d,"p"] \\
      (Y, y_0) \ar[r,"f"'] \ar[ur,dashed,"\tilde{f}"] & (X, x_0)
    \end{tikzcd}
  \]
\end{theorem}

\begin{proof}
  \textbf{Necessity.}  If $\tilde{f}$ exists, then $f_* = p_* \circ
  \tilde{f}_*$, so $f_*(\pi_1(Y)) \subseteq p_*(\pi_1(\widetilde{X}))$.

  \textbf{Sufficiency.}  For $y \in Y$, choose a path $\alpha$ in $Y$ from
  $y_0$ to $y$.  Then $f \circ \alpha$ is a path in $X$ starting at $x_0$.
  Lift it to $\widetilde{f \circ \alpha}$ starting at $\tilde{x}_0$, and
  define $\tilde{f}(y) := \widetilde{f \circ \alpha}(1)$.

  We must show this is well-defined.  If $\beta$ is another path from $y_0$
  to $y$, then $\alpha \cdot \overline{\beta}$ is a loop at $y_0$.  Its image
  $f_*([\alpha \cdot \overline{\beta}])$ lies in
  $p_*(\pi_1(\widetilde{X}, \tilde{x}_0))$ by hypothesis.  By the monodromy
  theorem, the lifts of $f \circ \alpha$ and $f \circ \beta$ have the same
  endpoint.

  \textbf{Continuity.}  Local path-connectedness of $Y$ ensures that for any
  $y \in Y$ and any evenly covered neighborhood $U$ of $f(y)$, we can find a
  path-connected neighborhood $W$ of $y$ mapping into $U$, on which
  $\tilde{f}$ is continuous (it equals $(p|_V)^{-1} \circ f$ on $W$ for an
  appropriate sheet $V$).

  \textbf{Uniqueness.}  Two lifts $\tilde{f}_1, \tilde{f}_2$ agreeing at $y_0$
  agree everywhere by the same open-and-closed argument as in
  Lemma~\ref{lem:path-lifting}.
\end{proof}

%% ──────────────────────────────────────────────────────────
\section{The universal covering space}\label{sec:universal-cover}

\begin{definition}[Simply connected covering]\label{def:universal-cover}
  \index{universal covering space}\index{covering space!universal}
  A covering $p\colon \widetilde{X} \to X$ is called a \textbf{universal
  covering} (or \textbf{universal cover}) if $\widetilde{X}$ is simply
  connected.
\end{definition}

\begin{theorem}[Existence of the universal cover]\label{thm:universal-cover-existence}
  \index{universal covering space!existence}
  Let $X$ be a connected, locally path-connected, and semi-locally simply
  connected space.  Then $X$ admits a universal covering
  $p\colon \widetilde{X} \to X$, unique up to isomorphism of coverings.
\end{theorem}

\begin{proof}[Proof sketch]
  Fix $x_0 \in X$.  Define $\widetilde{X}$ as the set of homotopy classes
  (relative to endpoints) of paths in $X$ starting at $x_0$:
  \[
    \widetilde{X} = \set{[\gamma] : \gamma \text{ is a path in } X
      \text{ with } \gamma(0) = x_0},
  \]
  with $p([\gamma]) = \gamma(1)$.  The topology on $\widetilde{X}$ is
  generated by sets of the form
  \[
    U_{[\gamma], V} = \set{[\gamma \cdot \eta] : \eta \text{ is a path in } V
      \text{ with } \eta(0) = \gamma(1)},
  \]
  where $V$ is a path-connected open subset of $X$ such that
  $\pi_1(V) \to \pi_1(X)$ is trivial (guaranteed by the semi-local simple
  connectivity hypothesis).

  One checks that $p$ is a covering map, $\widetilde{X}$ is simply connected,
  and any two universal covers are isomorphic via the lifting criterion.
\end{proof}

\begin{remark}[Semi-local simple connectivity]\label{rem:semi-local}
  \index{semi-locally simply connected}
  A space $X$ is \textbf{semi-locally simply connected} if every point has a
  neighborhood $U$ such that the inclusion-induced map
  $\pi_1(U) \to \pi_1(X)$ is trivial.  This does \emph{not} require $U$ to
  be simply connected.  The Hawaiian earring is the standard example of a
  space that fails this condition and has no universal cover.
\end{remark}

\begin{example}[Universal covers]\label{ex:universal-covers}
  \index{universal covering space!examples}
  \begin{enumerate}[label=(\roman*)]
    \item $\R \to S^1$, $t \mapsto e^{2\pi i t}$, is the universal cover
      of $S^1$.
    \item $S^2 \to \RP^2$ is the universal cover of $\RP^2$.
    \item $S^n \to \RP^n$ for all $n \geq 1$.
    \item $\R^2 \to T^2 = S^1 \times S^1$ is the universal cover of the torus.
    \item The universal cover of $S^1 \vee S^1$ is an infinite $4$-regular
      tree (the Cayley graph of $F_2$).\index{Cayley graph}
  \end{enumerate}
\end{example}

\begin{proposition}[Universal property of the universal cover]\label{prop:universal-cover-universal}
  Let $p\colon \widetilde{X} \to X$ be a universal covering and let
  $q\colon Y \to X$ be any connected covering.  Then there exists a covering
  map $r\colon \widetilde{X} \to Y$ with $q \circ r = p$:
  \[
    \begin{tikzcd}
      \widetilde{X} \ar[dr,"p"'] \ar[rr,dashed,"r"] & & Y \ar[dl,"q"] \\
      & X &
    \end{tikzcd}
  \]
\end{proposition}

\begin{proof}
  Since $\widetilde{X}$ is simply connected,
  $p_*(\pi_1(\widetilde{X})) = \{e\} \subseteq q_*(\pi_1(Y))$.  By
  the lifting criterion (Theorem~\ref{thm:lifting-criterion}), $p$ lifts to
  $r\colon \widetilde{X} \to Y$.  One checks that $r$ is itself a covering map.
\end{proof}

%% ──────────────────────────────────────────────────────────
\section{The Galois correspondence}\label{sec:galois-correspondence}

We now state the main classification theorem.

\begin{definition}[Locally path-connected and semi-locally simply connected]\label{def:nice-space}
  We say that a space $X$ is \textbf{nice} if it is connected, locally
  path-connected, and semi-locally simply connected.  All CW~complexes and
  manifolds are nice.
\end{definition}

\begin{theorem}[Classification of coverings --- Galois correspondence]\label{thm:galois-correspondence}
  \index{Galois correspondence!for coverings}
  \index{covering space!classification}
  Let $X$ be a nice space with basepoint $x_0$.  There is a bijection
  \[
    \left\{
    \begin{array}{c}
      \text{Connected coverings}\\
      p\colon (\widetilde{X}, \tilde{x}_0) \to (X, x_0)\\
      \text{up to isomorphism}
    \end{array}
    \right\}
    \;\xleftrightarrow{\;1:1\;}\;
    \left\{
    \begin{array}{c}
      \text{Subgroups of } \pi_1(X, x_0)\\
      \text{up to conjugacy}
    \end{array}
    \right\}
  \]
  given by $p \mapsto p_*\big(\pi_1(\widetilde{X}, \tilde{x}_0)\big)$.
  If we fix a basepoint lift $\tilde{x}_0$, then the correspondence is with
  actual subgroups (not conjugacy classes).

  Under this correspondence:
  \begin{enumerate}[label=(\roman*)]
    \item The universal cover corresponds to the trivial subgroup $\{e\}$.
    \item The trivial covering $\Id\colon X \to X$ corresponds to the whole
      group $\pi_1(X, x_0)$.
    \item The degree of the covering equals the index of the subgroup:
      $\abs{p^{-1}(x_0)} = [\pi_1(X) : H]$ where $H = p_*(\pi_1(\widetilde{X}))$.
    \item Larger subgroups correspond to ``smaller'' (fewer sheets) coverings.
  \end{enumerate}
\end{theorem}

\begin{proof}
  \textbf{Injectivity.}  If $p_1$ and $p_2$ are connected coverings with
  $p_{1*}(\pi_1(\widetilde{X}_1, \tilde{x}_{0,1})) =
  p_{2*}(\pi_1(\widetilde{X}_2, \tilde{x}_{0,2}))$, then by the lifting
  criterion applied in both directions, there exist lifts
  $f\colon \widetilde{X}_1 \to \widetilde{X}_2$ and
  $g\colon \widetilde{X}_2 \to \widetilde{X}_1$ with $g \circ f = \Id$ and
  $f \circ g = \Id$ (by uniqueness of lifts).  Hence $p_1 \cong p_2$.

  \textbf{Surjectivity.}  Given a subgroup $H \leq \pi_1(X, x_0)$, we
  construct a covering with $p_*(\pi_1(\widetilde{X})) = H$.  Start with the
  universal cover $\tilde{p}\colon \widetilde{X} \to X$ and define
  $\widetilde{X}_H = \widetilde{X} / {\sim_H}$, where
  $[\gamma_1] \sim_H [\gamma_2]$ if and only if $\gamma_1(1) = \gamma_2(1)$
  and $[\gamma_1 \cdot \overline{\gamma_2}] \in H$.  The projection
  $p_H\colon \widetilde{X}_H \to X$ is a covering with
  $p_{H*}(\pi_1(\widetilde{X}_H)) = H$.
\end{proof}

\begin{example}[Coverings of $S^1$]\label{ex:coverings-S1-classification}
  \index{covering space!of the circle!classification}
  Since $\pi_1(S^1) \cong \Z$, the subgroups are $n\Z$ for $n \geq 0$.
  \begin{itemize}
    \item $n = 0$: subgroup $\{0\}$, giving the universal cover
      $\R \to S^1$.
    \item $n \geq 1$: subgroup $n\Z$, giving the $n$-sheeted cover
      $S^1 \xrightarrow{z \mapsto z^n} S^1$.
  \end{itemize}
  Since $\Z$ is abelian, every subgroup is normal, and there are no conjugacy
  ambiguities.
\end{example}

\begin{example}[Coverings of $S^1 \vee S^1$]\label{ex:coverings-bouquet}
  \index{covering space!of wedge of circles}
  Since $\pi_1(S^1 \vee S^1) \cong F_2$, which has a vast collection of
  subgroups (including free groups of every countable rank, by the
  Nielsen--Schreier theorem\index{Nielsen--Schreier theorem}), the figure
  eight has a very rich collection of coverings.  For instance:
  \begin{itemize}
    \item The commutator subgroup $[F_2, F_2]$ gives the covering
      corresponding to the maximal abelian quotient $\Z^2$; this is an
      infinite-sheeted cover whose group of deck transformations is $\Z^2$.
    \item Each normal subgroup $N \trianglelefteq F_2$ with $F_2/N$ finite
      gives a finite regular covering.
  \end{itemize}
\end{example}

\begin{example}[Lattice of coverings of $\RP^2 \vee S^1$]\label{ex:coverings-RP2-wedge}
  Since $\pi_1(\RP^2 \vee S^1) \cong \Z/2\Z * \Z$, which is an infinite
  non-abelian group, the space $\RP^2 \vee S^1$ has a rich family of
  coverings.  The subgroup lattice of $\Z/2\Z * \Z$ includes:
  \begin{itemize}
    \item The trivial subgroup $\{e\}$: universal cover (infinite tree-like
      structure with $S^2$-shaped ``bubbles'').
    \item The commutator subgroup $[\Z/2\Z * \Z,\, \Z/2\Z * \Z]$: an
      infinite-index normal subgroup.
    \item Finite-index subgroups giving finite-sheeted coverings.
  \end{itemize}
\end{example}

\begin{remark}[Functoriality of the correspondence]\label{rem:galois-functor}
  \index{Galois correspondence!functoriality}
  The Galois correspondence is order-reversing with respect to inclusion:
  if $H_1 \leq H_2 \leq \pi_1(X)$ with corresponding coverings $p_1$ and
  $p_2$, then there exists a covering map $\widetilde{X}_1 \to \widetilde{X}_2$
  making the diagram commute:
  \[
    \begin{tikzcd}
      \widetilde{X}_1 \ar[rr] \ar[dr,"p_1"'] & & \widetilde{X}_2 \ar[dl,"p_2"] \\
      & X &
    \end{tikzcd}
  \]
  The number of sheets of $\widetilde{X}_1 \to \widetilde{X}_2$ is $[H_2 : H_1]$.
\end{remark}

%% ──────────────────────────────────────────────────────────
\section{Deck transformations}\label{sec:deck-transformations}

\begin{definition}[Deck transformation]\label{def:deck-transformation}
  \index{deck transformation}\index{covering transformation|see {deck transformation}}
  Let $p\colon \widetilde{X} \to X$ be a covering.  A
  \textbf{deck transformation} (or \textbf{covering transformation}) is a
  homeomorphism $\varphi\colon \widetilde{X} \to \widetilde{X}$ such that
  $p \circ \varphi = p$.  The set of all deck transformations forms a group
  under composition, denoted $\Aut(p)$ or $\mathrm{Deck}(\widetilde{X}/X)$.
  \[
    \begin{tikzcd}
      \widetilde{X} \ar[rr,"\varphi"] \ar[dr,"p"'] & & \widetilde{X} \ar[dl,"p"] \\
      & X &
    \end{tikzcd}
  \]
\end{definition}

\begin{proposition}[Fixed-point-free action]\label{prop:deck-free}
  \index{deck transformation!free action}
  If $\widetilde{X}$ is connected, then every non-identity deck transformation
  acts without fixed points: if $\varphi(\tilde{x}) = \tilde{x}$ for some
  $\tilde{x} \in \widetilde{X}$, then $\varphi = \Id$.
\end{proposition}

\begin{proof}
  The set $\set{\tilde{x} \in \widetilde{X} : \varphi(\tilde{x}) = \tilde{x}}$
  is closed (as $\widetilde{X}$ is Hausdorff) and open (locally, $\varphi$
  permutes the sheets, so if it fixes a point it fixes the entire sheet).
  By connectedness, this set is either empty or all of $\widetilde{X}$.
\end{proof}

\begin{proposition}[Structure of the deck group]\label{prop:deck-structure}
  \index{deck transformation!structure}
  Let $p\colon \widetilde{X} \to X$ be a connected covering with
  $\tilde{x}_0 \in p^{-1}(x_0)$, and let $H = p_*(\pi_1(\widetilde{X},
  \tilde{x}_0))$.  Then:
  \begin{enumerate}[label=(\roman*)]
    \item $\Aut(p)$ is isomorphic to $N(H)/H$, where
      $N(H) = \set{g \in \pi_1(X, x_0) : gHg^{-1} = H}$ is the normalizer
      of $H$.
    \item $\abs{\Aut(p)} \leq \abs{p^{-1}(x_0)}$, with equality if and only if
      $H$ is normal in $\pi_1(X, x_0)$.
  \end{enumerate}
\end{proposition}

\begin{proof}
  (i) A deck transformation $\varphi$ is determined by $\varphi(\tilde{x}_0)$,
  which must lie in $p^{-1}(x_0)$.  Given $[\gamma] \in \pi_1(X, x_0)$,
  define $\varphi_{[\gamma]}(\tilde{x}_0) = \tilde{\gamma}(1)$ (the endpoint
  of the lift of $\gamma$).  For $\varphi_{[\gamma]}$ to extend to a deck
  transformation, one needs $\varphi_{[\gamma]}$ to be well-defined on all of
  $\widetilde{X}$, which happens precisely when
  $[\gamma] H [\gamma]^{-1} = H$, i.e., $[\gamma] \in N(H)$.  The map
  $[\gamma] \mapsto \varphi_{[\gamma]}$ has kernel $H$, giving
  $\Aut(p) \cong N(H)/H$.

  (ii) By the orbit-stabilizer theorem, $\abs{\Aut(p)} =
  \abs{N(H)/H} \leq [\pi_1(X) : H] = \abs{p^{-1}(x_0)}$.  Equality holds
  if and only if $N(H) = \pi_1(X, x_0)$, i.e., $H \trianglelefteq
  \pi_1(X, x_0)$.
\end{proof}

%% ──────────────────────────────────────────────────────────
\section{Normal (regular) coverings}\label{sec:normal-coverings}

\begin{definition}[Normal covering]\label{def:normal-covering}
  \index{covering space!normal (regular)}\index{regular covering|see {normal covering}}
  A connected covering $p\colon \widetilde{X} \to X$ is called
  \textbf{normal} (or \textbf{regular} or \textbf{Galois}) if
  $p_*(\pi_1(\widetilde{X}, \tilde{x}_0))$ is a normal subgroup of
  $\pi_1(X, x_0)$.  Equivalently, $\Aut(p)$ acts transitively on each fiber.
\end{definition}

\begin{theorem}[Characterization of normal coverings]\label{thm:normal-covering-char}
  \index{covering space!normal!characterization}
  For a connected covering $p\colon \widetilde{X} \to X$ with $X$ nice, the
  following are equivalent:
  \begin{enumerate}[label=(\roman*)]
    \item $p$ is normal.
    \item $H = p_*(\pi_1(\widetilde{X}))$ is normal in $\pi_1(X)$.
    \item $\Aut(p)$ acts transitively on every fiber $p^{-1}(x)$.
    \item $\widetilde{X} / \Aut(p) \cong X$.
    \item For every loop $\gamma$ in $X$, either every lift of $\gamma$ is a
      loop, or no lift is a loop.
  \end{enumerate}
  When these hold, $\Aut(p) \cong \pi_1(X) / H$.
\end{theorem}

\begin{proof}
  The equivalence (i)$\Leftrightarrow$(ii) is by definition.

  (ii)$\Leftrightarrow$(iii):  By Proposition~\ref{prop:deck-structure},
  $\abs{\Aut(p)} = [N(H):H]$, and $\Aut(p)$ acts transitively on fibers if
  and only if $\abs{\Aut(p)} = [G:H]$ (where $G = \pi_1(X)$), i.e.,
  $N(H) = G$, i.e., $H \trianglelefteq G$.

  (iii)$\Rightarrow$(iv):  If $\Aut(p)$ acts transitively on fibers, the
  natural map $\widetilde{X}/\Aut(p) \to X$ is a homeomorphism.

  (iv)$\Rightarrow$(iii) is clear.

  (ii)$\Leftrightarrow$(v):  A lift of $\gamma$ starting at $\tilde{x}$ is a
  loop if and only if $[\gamma] \in p_*(\pi_1(\widetilde{X}, \tilde{x}))$.
  As $\tilde{x}$ varies over the fiber, the subgroup
  $p_*(\pi_1(\widetilde{X}, \tilde{x}))$ runs over the conjugates of $H$.
  The condition that all lifts are simultaneously loops (or non-loops)
  is equivalent to $H$ being invariant under conjugation, i.e., $H$ is normal.
\end{proof}

\begin{example}[Normal coverings]\label{ex:normal-coverings}
  \begin{enumerate}[label=(\roman*)]
    \item Every universal covering is normal (the trivial subgroup is always
      normal).  For the universal cover, $\Aut(p) \cong \pi_1(X)$.
    \item $S^1 \xrightarrow{z^n} S^1$ is normal with
      $\Aut(p) \cong \Z/n\Z$ (rotation by $2\pi/n$).
    \item $S^2 \to \RP^2$ is normal with $\Aut(p) \cong \Z/2\Z$ (the
      antipodal map).
    \item If $\pi_1(X)$ is abelian, every connected covering of $X$ is
      normal.
  \end{enumerate}
\end{example}

\begin{proposition}[Recovering $\pi_1$ from the universal cover]\label{prop:pi1-from-universal}
  \index{fundamental group!from universal cover}
  If $p\colon \widetilde{X} \to X$ is the universal cover of a nice space $X$,
  then
  \[
    \pi_1(X, x_0) \;\cong\; \Aut(p).
  \]
\end{proposition}

\begin{proof}
  Since $\widetilde{X}$ is simply connected,
  $H = p_*(\pi_1(\widetilde{X})) = \{e\}$, which is normal.  By
  Theorem~\ref{thm:normal-covering-char},
  $\Aut(p) \cong \pi_1(X)/\{e\} = \pi_1(X)$.
\end{proof}

\begin{example}[Deck transformations of $\R \to S^1$]\label{ex:deck-R-S1}
  \index{deck transformation!of $\R \to S^1$}
  The deck transformations of $p\colon \R \to S^1$, $t \mapsto e^{2\pi i t}$,
  are the translations $t \mapsto t + n$ for $n \in \Z$.  Hence
  $\Aut(p) \cong \Z \cong \pi_1(S^1)$.
\end{example}

%% ──────────────────────────────────────────────────────────
\section{Applications and constructions}\label{sec:covering-applications}

\begin{proposition}[Coverings of graphs]\label{prop:covering-graphs}
  \index{covering space!of graphs}
  Every covering space of a graph is a graph.  Every subgroup of a free group
  is free.
\end{proposition}

\begin{proof}
  A graph is a $1$-dimensional CW~complex.  The preimage of a CW~structure
  under a covering map inherits a CW~structure (sheets of cells are cells).
  Hence coverings of graphs are graphs.

  Now let $F = \pi_1(\Gamma)$ for a graph $\Gamma$, and let $H \leq F$.  By
  the Galois correspondence, $H = p_*(\pi_1(\widetilde{\Gamma}))$ for some
  covering $\widetilde{\Gamma}$ of $\Gamma$.  Since $\widetilde{\Gamma}$ is
  also a graph, $\pi_1(\widetilde{\Gamma})$ is free.  Since $p_*$ is
  injective, $H \cong \pi_1(\widetilde{\Gamma})$ is free.
\end{proof}

\begin{corollary}[Nielsen--Schreier theorem]\label{cor:nielsen-schreier}
  \index{Nielsen--Schreier theorem}
  Every subgroup of a free group is free.  Moreover, if $F$ has finite rank $n$
  and $H$ has finite index $d$, then $H$ has rank $d(n-1) + 1$.
\end{corollary}

\begin{proof}
  The first statement was proved above.  For the rank formula: if $F =
  \pi_1(\Gamma)$ where $\Gamma$ has one vertex and $n$ edges (so
  $\chi(\Gamma) = 1 - n$), and $H$ corresponds to a $d$-sheeted covering
  $\widetilde{\Gamma}$, then $\chi(\widetilde{\Gamma}) = d \cdot
  \chi(\Gamma) = d(1-n)$.  Since $\widetilde{\Gamma}$ is a graph with
  Euler characteristic $1 - \mathrm{rank}(H)$, we get
  $\mathrm{rank}(H) = d(n-1) + 1$.
\end{proof}

\begin{proposition}[Coverings and group actions]\label{prop:covering-group-action}
  \index{covering space!from group action}
  Let $G$ be a group acting on a topological space $\widetilde{X}$ by
  homeomorphisms.  Suppose the action is \textbf{properly discontinuous}:
  every $\tilde{x} \in \widetilde{X}$ has a neighborhood $U$ such that
  $g \cdot U \cap U = \varnothing$ for all $g \neq e$.  Then the quotient map
  $p\colon \widetilde{X} \to \widetilde{X}/G$ is a normal covering with
  $\Aut(p) \cong G$.
\end{proposition}

\begin{proof}
  The properly discontinuous condition ensures that $p$ is a covering: each
  $U$ as above is evenly covered, with sheets $\{g \cdot U\}_{g \in G}$.
  Clearly each $g \in G$ defines a deck transformation, and the action is
  transitive on fibers by construction.  Hence $p$ is normal with
  $\Aut(p) = G$.
\end{proof}

\begin{example}[Lens spaces]\label{ex:lens-spaces}
  \index{lens space}
  Let $p, q$ be coprime positive integers with $1 \leq q < p$.  Consider
  $S^3 \subset \C^2$ and the $\Z/p\Z$-action generated by
  $(z_1, z_2) \mapsto (e^{2\pi i/p} z_1, e^{2\pi i q/p} z_2)$.  This action
  is properly discontinuous and free, so $S^3 \to L(p,q) := S^3/(\Z/p\Z)$ is
  a covering.  Since $S^3$ is simply connected,
  $\pi_1(L(p,q)) \cong \Z/p\Z$.
\end{example}

\begin{example}[Coverings of the Klein bottle]\label{ex:coverings-klein}
  \index{Klein bottle!coverings}
  Recall $\pi_1(K) = \langle a, b \mid abab^{-1} \rangle$.  The abelianization
  is $\Z \oplus \Z/2\Z$.  The subgroup $\langle a, b^2, bab^{-1} \rangle$
  has index $2$ and corresponds to the torus $T^2$ as a $2$-sheeted covering
  of $K$.

  \begin{center}
    \begin{tikzpicture}[scale=0.8]
      % Klein bottle (schematic square)
      \draw[thick] (0,0) rectangle (3,3);
      \draw[thick, blue!70!black, -Stealth] (0,0) -- (1.5,0);
      \draw[thick, blue!70!black, Stealth-] (3,0) -- (1.5,0);
      \draw[thick, red!70!black, -Stealth] (0,0) -- (0,1.5);
      \draw[thick, red!70!black, -Stealth] (3,0) -- (3,1.5);
      \draw[thick, blue!70!black, -Stealth] (0,3) -- (1.5,3);
      \draw[thick, blue!70!black, Stealth-] (3,3) -- (1.5,3);
      \draw[thick, red!70!black, -Stealth] (0,1.5) -- (0,3);
      \draw[thick, red!70!black, -Stealth] (3,1.5) -- (3,3);
      \node at (1.5,1.5) {$K$};
      % Arrow
      \draw[-Stealth, thick] (4,1.5) -- (5.5,1.5) node[midway, above] {$2:1$};
      % Torus covering (double square)
      \draw[thick] (6.5,0) rectangle (9.5,3);
      \draw[thick, dashed, gray] (6.5,1.5) -- (9.5,1.5);
      \draw[thick, green!50!black, -Stealth] (6.5,0) -- (8,0);
      \draw[thick, green!50!black, -Stealth] (6.5,3) -- (8,3);
      \draw[thick, violet, -Stealth] (6.5,0) -- (6.5,1.5);
      \draw[thick, violet, -Stealth] (9.5,0) -- (9.5,1.5);
      \draw[thick, violet, -Stealth] (6.5,1.5) -- (6.5,3);
      \draw[thick, violet, -Stealth] (9.5,1.5) -- (9.5,3);
      \node at (8,0.75) {sheet 1};
      \node at (8,2.25) {sheet 2};
      \node[below] at (8,-0.3) {$T^2$};
    \end{tikzpicture}
  \end{center}
\end{example}

\begin{proposition}[Euler characteristic and coverings]\label{prop:euler-char-covering}
  \index{Euler characteristic!and coverings}
  If $p\colon \widetilde{X} \to X$ is an $n$-sheeted covering of a finite
  CW~complex $X$, then $\chi(\widetilde{X}) = n \cdot \chi(X)$.
\end{proposition}

\begin{proof}
  Each $k$-cell of $X$ lifts to exactly $n$ copies of $k$-cells in
  $\widetilde{X}$.  If $c_k$ denotes the number of $k$-cells of $X$, then
  $\widetilde{X}$ has $n \cdot c_k$ cells of dimension $k$.  Hence
  \[
    \chi(\widetilde{X}) = \sum_k (-1)^k (n \cdot c_k) = n \sum_k (-1)^k c_k
    = n \cdot \chi(X). \qedhere
  \]
\end{proof}

%% ──────────────────────────────────────────────────────────
\section{Summary: the covering space dictionary}\label{sec:dictionary}

The following table summarizes the Galois correspondence.

\begin{center}
\renewcommand{\arraystretch}{1.4}
\begin{tabular}{@{}ll@{}}
  \toprule
  \textbf{Topology} & \textbf{Algebra} \\
  \midrule
  Connected covering $p\colon \widetilde{X} \to X$ &
    Subgroup $H \leq \pi_1(X)$ \\
  Universal cover &
    Trivial subgroup $\{e\}$ \\
  Identity covering $\Id_X$ &
    Whole group $\pi_1(X)$ \\
  Number of sheets $\abs{p^{-1}(x_0)}$ &
    Index $[\pi_1(X) : H]$ \\
  Normal (regular) covering &
    Normal subgroup $H \trianglelefteq \pi_1(X)$ \\
  Deck transformations $\Aut(p)$ &
    $N_{\pi_1(X)}(H)/H$; \; $\pi_1(X)/H$ if normal \\
  Intermediate covering $\widetilde{X}_1 \to \widetilde{X}_2 \to X$ &
    $H_1 \leq H_2 \leq \pi_1(X)$ \\
  \bottomrule
\end{tabular}
\end{center}

%% ──────────────────────────────────────────────────────────
\section{Exercises}\label{sec:covering-exercises}

\begin{exercise}\label{exo:cov-R2-torus}
  Let $p\colon \R^2 \to T^2$ be the universal covering of the torus.
  \begin{enumerate}[label=(\alph*)]
    \item Describe explicitly the deck transformations of $p$.
    \item For each subgroup $H \leq \Z^2 = \pi_1(T^2)$, describe the
      corresponding covering space.
    \item Which coverings of $T^2$ are connected? What are their fundamental
      groups?
  \end{enumerate}
\end{exercise}

\begin{exercise}\label{exo:cov-RP2-lifts}
  Let $p\colon S^2 \to \RP^2$ be the standard $2$-sheeted covering.
  \begin{enumerate}[label=(\alph*)]
    \item Describe all coverings of $\RP^2$.
    \item Let $f\colon T^2 \to \RP^2$ be a continuous map.  Show that $f$
      lifts to a map $\tilde{f}\colon T^2 \to S^2$.
  \end{enumerate}
\end{exercise}

\begin{exercise}\label{exo:cov-3sheeted}
  Classify all connected $3$-sheeted coverings of $S^1 \vee S^1$ (up to
  isomorphism).  Draw the covering graphs and determine which are normal.

  \textit{Hint:} A $3$-sheeted covering corresponds to a transitive action of
  $F_2$ on a set of $3$ elements, i.e., a transitive homomorphism
  $F_2 \to S_3$.
\end{exercise}

\begin{exercise}\label{exo:cov-free-action}
  Let $G$ be a finite group acting freely on a Hausdorff space $X$.
  Show that the action is properly discontinuous.  Conclude that
  $X \to X/G$ is a covering.
\end{exercise}

\begin{exercise}\label{exo:cov-pi1-product}
  Show that $\pi_1(X \times Y, (x_0, y_0)) \cong \pi_1(X, x_0) \times
  \pi_1(Y, y_0)$ by constructing the universal cover of $X \times Y$ from
  the universal covers of $X$ and $Y$.
\end{exercise}

\begin{exercise}\label{exo:cov-klein-coverings}
  \index{Klein bottle!coverings}
  Let $K$ be the Klein bottle.
  \begin{enumerate}[label=(\alph*)]
    \item Show that $T^2$ is a $2$-sheeted covering of $K$.
    \item Find a normal covering of $K$ with deck group $\Z$.
    \item Classify all connected coverings of $K$ with at most $4$ sheets.
  \end{enumerate}
\end{exercise}

\begin{exercise}\label{exo:cov-rank-formula}
  Let $F_n$ be the free group of rank $n$ and $H \leq F_n$ a subgroup of
  index $d < \infty$.  Use the covering space proof of the Nielsen--Schreier
  theorem to show that $\mathrm{rank}(H) = d(n-1) + 1$.  Apply this to find
  the rank of the commutator subgroup $[F_2, F_2]$ (is it finitely generated?).
\end{exercise}

\begin{exercise}\label{exo:cov-cayley-complex}
  \index{Cayley complex}
  Let $G = \langle S \mid R \rangle$ be a finitely presented group.  The
  \textbf{Cayley complex} $\widetilde{X}_G$ is the universal cover of the
  standard $2$-complex $X_G$ (with one vertex, one edge per generator, one
  $2$-cell per relator).  Show that $\widetilde{X}_G$ is simply connected and
  $\Aut(\widetilde{X}_G / X_G) \cong G$.  Describe $\widetilde{X}_G$
  explicitly for $G = \Z/n\Z$ and $G = \Z \times \Z$.
\end{exercise}

\begin{exercise}\label{exo:cov-borsuk-ulam}
  \index{Borsuk--Ulam theorem}
  Use covering space theory to prove the following:
  there is no continuous map $f\colon S^2 \to S^1$ satisfying
  $f(-x) = -f(x)$ for all $x \in S^2$.

  \textit{Hint:} Such an $f$ would induce a map $\bar{f}\colon \RP^2 \to
  \RP^1 = S^1$.  Analyze $\bar{f}_*$ on fundamental groups.
\end{exercise}

\begin{exercise}\label{exo:cov-surface-cover}
  Let $\Sigma_g$ be the closed orientable surface of genus $g \geq 2$.
  \begin{enumerate}[label=(\alph*)]
    \item Show that the universal cover of $\Sigma_g$ is contractible.
      (In fact, it is homeomorphic to $\R^2$---this is much harder to
      prove.)
    \item Construct a covering $\Sigma_h \to \Sigma_g$ for suitable $h > g$.
      What is the relationship between $h$, $g$, and the number of sheets?

      \textit{Hint:} Use Proposition~\ref{prop:euler-char-covering}: if the
      covering has $d$ sheets, then
      $2 - 2h = d(2 - 2g)$, giving $h = d(g-1) + 1$.
  \end{enumerate}
\end{exercise}

\begin{exercise}\label{exo:cov-no-retraction}
  \index{retraction}
  Let $p\colon \widetilde{X} \to X$ be a connected covering with
  $\widetilde{X} \neq X$ (i.e., more than one sheet).  Show that there
  is no continuous map $s\colon X \to \widetilde{X}$ with $p \circ s = \Id_X$
  (i.e., $p$ has no section) unless $X$ is simply connected.

  \textit{Hint:} If $s$ exists, what does $s_* \circ p_*$ tell you about
  $\pi_1(\widetilde{X})$ and $\pi_1(X)$?
\end{exercise}

\begin{exercise}\label{exo:cov-commutator}
  \index{commutator subgroup}
  Let $X = S^1 \vee S^1$ and let $p\colon \widetilde{X} \to X$ be the
  covering corresponding to the commutator subgroup $[F_2, F_2]$ of
  $\pi_1(X) \cong F_2 = \langle a, b \rangle$.
  \begin{enumerate}[label=(\alph*)]
    \item What is $\Aut(p)$?
    \item Show that $\widetilde{X}$ is an infinite graph.  Describe it as a
      subset of $\R^2$: it is the $1$-skeleton of the standard tiling of
      $\R^2$ by unit squares, with the $\Z^2$-action by translation.
    \item What is $\pi_1(\widetilde{X})$?  Is it finitely generated?

      \textit{Hint:} The commutator subgroup $[F_2, F_2]$ has infinite index
      in $F_2$ (since $F_2 / [F_2, F_2] \cong \Z^2$), so the covering is
      infinite-sheeted and $\pi_1(\widetilde{X}) \cong [F_2, F_2]$ is free
      of infinite rank.
  \end{enumerate}
\end{exercise}

\begin{exercise}\label{exo:cov-composition}
  Let $p\colon Y \to X$ and $q\colon Z \to Y$ be coverings.
  \begin{enumerate}[label=(\alph*)]
    \item Show that $p \circ q\colon Z \to X$ is not necessarily a covering
      in general.  Give a counterexample.

      \textit{Hint:} Consider $\R \xrightarrow{p_1} S^1
      \xrightarrow{p_2} S^1$ where $p_2(z) = z^2$.  Is $p_2 \circ p_1$
      a covering?  What goes wrong?
    \item Show that if $p$ has finitely many sheets, then $p \circ q$ is
      indeed a covering.
  \end{enumerate}
\end{exercise}

\begin{exercise}\label{exo:cov-fixed-point}
  \index{Brouwer fixed point theorem}
  Use covering space theory to give another proof that every continuous map
  $f\colon D^2 \to D^2$ has a fixed point (Brouwer's theorem for $n = 2$).

  \textit{Hint:} If $f$ has no fixed point, construct a retraction
  $D^2 \to S^1$ and derive a contradiction using $\pi_1$.
\end{exercise}

\begin{exercise}\label{exo:cov-properly-disc}
  \index{properly discontinuous action}
  Let $X = \R^2 \setminus \{0\}$ and let $\Z$ act on $X$ by
  $n \cdot (x,y) = (2^n x, 2^n y)$.
  \begin{enumerate}[label=(\alph*)]
    \item Verify that this action is properly discontinuous.
    \item Identify the quotient space $X / \Z$.

      \textit{Hint:} It is a familiar compact surface.
    \item Compute $\pi_1(X/\Z)$ using covering space theory.
  \end{enumerate}
\end{exercise}

% === END OF CHUNK ch3-4 ===
% === START OF CHUNK ch5-6 ===
%%%%%%%%%%%%%%%%%%%%%%%%%%%%%%%%%%%%%%%%%%%%%%%%%%%%%%%%%%%%%%%%%%%%
%  Algebraic Topology — Chapters 5–6
%  Singular Homology, Long Exact Sequences, and Excision
%%%%%%%%%%%%%%%%%%%%%%%%%%%%%%%%%%%%%%%%%%%%%%%%%%%%%%%%%%%%%%%%%%%%

%% ================================================================
%%  CHAPTER 5 — Singular Homology: Definition and Properties
%% ================================================================
\chapter{Singular Homology --- Definition and Properties}\label{ch:singular-homology}
\index{singular homology|(}

The passage from topology to algebra is achieved by associating to each
topological space~$X$ a sequence of abelian groups $H_n(X)$, called its
\emph{homology groups}\index{homology groups}.  The construction is entirely
functorial and, remarkably, depends only on the homotopy type of~$X$.
In this chapter we develop the foundations of singular homology theory.

%% ----------------------------------------------------------------
%%  5.1  The standard simplex and singular simplices
%% ----------------------------------------------------------------
\section{The standard simplex and singular simplices}\label{sec:standard-simplex}

\begin{definition}[Standard $n$-simplex]\label{def:standard-simplex}
\index{standard simplex}\index{simplex!standard}
The \textbf{standard $n$-simplex} is the subspace of~$\R^{n+1}$ defined by
\[
  \Delta^n \;=\;
  \Bigl\{\,(t_0,t_1,\dots,t_n)\in\R^{n+1}
  \;\Big|\;
  \sum_{i=0}^{n}t_i=1,\; t_i\ge 0 \text{ for all } i
  \Bigr\}.
\]
We write $e_0,e_1,\dots,e_n$ for the vertices of~$\Delta^n$, where
$e_i$ is the point whose $i$-th coordinate is~$1$ and all other
coordinates are~$0$.
\end{definition}

\begin{definition}[Face maps]\label{def:face-maps}
\index{face map}
For $0\le i\le n$, the \textbf{$i$-th face map} is the affine embedding
\[
  \delta_i^n \colon \Delta^{n-1}\longrightarrow \Delta^n,
  \qquad
  (t_0,\dots,t_{n-1})\longmapsto (t_0,\dots,t_{i-1},0,t_i,\dots,t_{n-1}).
\]
Thus $\delta_i^n$ maps $\Delta^{n-1}$ onto the face of~$\Delta^n$ opposite
the vertex~$e_i$.
\end{definition}

\begin{lemma}[Cosimplicial identity]\label{lem:cosimplicial}
\index{cosimplicial identity}
For $i<j$ we have
\[
  \delta_j^{n+1}\circ\delta_i^n
  \;=\;
  \delta_i^{n+1}\circ\delta_{j-1}^n.
\]
\end{lemma}

\begin{proof}
Both sides send $(t_0,\dots,t_{n-1})$ to the point
$(t_0,\dots,t_{i-1},0,t_i,\dots,t_{j-2},0,t_{j-1},\dots,t_{n-1})$.
A direct coordinate check confirms the identity.
\end{proof}

\begin{definition}[Singular $n$-simplex]\label{def:singular-simplex}
\index{singular simplex}\index{simplex!singular}
Let $X$ be a topological space.  A \textbf{singular $n$-simplex} in~$X$ is
a continuous map
\[
  \sigma \colon \Delta^n \longrightarrow X.
\]
We denote the set of all singular $n$-simplices in~$X$ by $\operatorname{Sing}_n(X)$.
\end{definition}

\begin{example}[Singular simplices in low dimensions]\label{ex:low-dim-simplices}
\begin{itemize}[leftmargin=2em]
  \item A singular $0$-simplex is simply a point of~$X$.
  \item A singular $1$-simplex is a continuous path
    $\sigma\colon [0,1]\to X$.
  \item A singular $2$-simplex is a continuous map from a triangle into~$X$.
\end{itemize}
\end{example}

%% ----------------------------------------------------------------
%%  5.2  Chain groups and the boundary operator
%% ----------------------------------------------------------------
\section{Chain groups and the boundary operator}\label{sec:chain-groups}

\begin{definition}[Singular chain group]\label{def:chain-group}
\index{chain group}\index{singular chain group}
The \textbf{singular chain group} $C_n(X)$ is the free abelian group
generated by $\operatorname{Sing}_n(X)$:
\[
  C_n(X) \;=\; \bigoplus_{\sigma\in\operatorname{Sing}_n(X)} \Z\cdot\sigma.
\]
An element of~$C_n(X)$ is called a \textbf{singular $n$-chain}\index{chain}
and is a finite formal sum $\sum_i n_i\sigma_i$ with $n_i\in\Z$.
We set $C_n(X)=0$ for $n<0$.
\end{definition}

\begin{definition}[Boundary operator]\label{def:boundary-operator}
\index{boundary operator}\index{boundary!operator $\partial_n$}
The \textbf{boundary operator} $\partial_n\colon C_n(X)\to C_{n-1}(X)$
is the homomorphism defined on generators by
\[
  \partial_n(\sigma)
  \;=\;
  \sum_{i=0}^{n}(-1)^i\,\sigma\circ\delta_i^n,
\]
and extended linearly to all of~$C_n(X)$.
\end{definition}

\begin{example}[Boundary of a $1$-simplex]\label{ex:boundary-1-simplex}
If $\sigma\colon\Delta^1\to X$ is a singular $1$-simplex (a path from
$\sigma(e_0)$ to $\sigma(e_1)$), then
\[
  \partial_1(\sigma)
  = \sigma\circ\delta_0^1 - \sigma\circ\delta_1^1
  = \sigma(e_1) - \sigma(e_0).
\]
\end{example}

\begin{example}[Boundary of a $2$-simplex]\label{ex:boundary-2-simplex}
For a singular $2$-simplex $\sigma\colon \Delta^2\to X$,
\[
  \partial_2(\sigma)
  = \sigma\circ\delta_0^2 - \sigma\circ\delta_1^2 + \sigma\circ\delta_2^2,
\]
which is the alternating sum of the three edges of the triangle.
\end{example}

The following theorem is the cornerstone of homological algebra.

\begin{theorem}[$\partial^2=0$]\label{thm:boundary-squared-zero}
\index{boundary!$\partial^2=0$}
For every~$n$, the composition $\partial_{n-1}\circ\partial_n=0$.
Equivalently, $\im\partial_n\subseteq\ker\partial_{n-1}$.
\end{theorem}

\begin{proof}
It suffices to check on a generator $\sigma\in\operatorname{Sing}_n(X)$.  We compute
\begin{align*}
  \partial_{n-1}(\partial_n(\sigma))
  &= \partial_{n-1}\!\Bigl(\sum_{i=0}^{n}(-1)^i\,\sigma\circ\delta_i^n\Bigr)
   = \sum_{i=0}^{n}(-1)^i\sum_{j=0}^{n-1}(-1)^j\,
     \sigma\circ\delta_i^n\circ\delta_j^{n-1} \\
  &= \sum_{0\le j < i\le n}(-1)^{i+j}\,
     \sigma\circ\delta_i^n\circ\delta_j^{n-1}
   + \sum_{0\le i\le j\le n-1}(-1)^{i+j}\,
     \sigma\circ\delta_i^n\circ\delta_j^{n-1}.
\end{align*}
In the first sum, the cosimplicial identity
(Lemma~\ref{lem:cosimplicial}) gives
$\delta_i^n\circ\delta_j^{n-1}=\delta_j^n\circ\delta_{i-1}^{n-1}$
for $j<i$.  Substituting $i'=j$, $j'=i-1$ (so $i'+j'=i+j-1$ and
$j'\ge i'$), this first sum becomes
\[
  \sum_{0\le i'\le j'\le n-1}(-1)^{i'+j'+1}\,
  \sigma\circ\delta_{i'}^n\circ\delta_{j'}^{n-1},
\]
which is exactly the negative of the second sum.  Hence the total is~$0$.
\end{proof}

%% ----------------------------------------------------------------
%%  5.3  The chain complex and homology groups
%% ----------------------------------------------------------------
\section{The chain complex and homology groups}\label{sec:chain-complex}

\begin{definition}[Singular chain complex]\label{def:chain-complex}
\index{chain complex}\index{singular chain complex}
The \textbf{singular chain complex} of~$X$ is the sequence
\[
\begin{tikzcd}
  \cdots \ar[r,"\partial_{n+1}"]
  & C_n(X) \ar[r,"\partial_n"]
  & C_{n-1}(X) \ar[r,"\partial_{n-1}"]
  & \cdots \ar[r,"\partial_2"]
  & C_1(X) \ar[r,"\partial_1"]
  & C_0(X) \ar[r]
  & 0.
\end{tikzcd}
\]
We write $C_\bullet(X)=(C_n(X),\partial_n)_{n\ge 0}$ for this chain complex.
\end{definition}

\begin{definition}[Cycles, boundaries, homology]\label{def:homology-groups}
\index{cycle}\index{boundary}\index{homology groups}
Let $n\ge 0$.  Define:
\begin{itemize}[leftmargin=2em]
  \item The group of \textbf{$n$-cycles}: $Z_n(X)=\ker\partial_n\subseteq C_n(X)$.
  \item The group of \textbf{$n$-boundaries}: $B_n(X)=\im\partial_{n+1}\subseteq C_n(X)$.
\end{itemize}
Since $\partial^2=0$ we have $B_n(X)\subseteq Z_n(X)$, and the
\textbf{$n$-th singular homology group}\index{singular homology!group} of~$X$ is the quotient
\[
  H_n(X) \;=\; \frac{Z_n(X)}{B_n(X)}
  \;=\; \frac{\ker\partial_n}{\im\partial_{n+1}}.
\]
\end{definition}

\begin{remark}[Intuitive meaning]\label{rem:homology-intuition}
An $n$-cycle is a chain ``without boundary,'' while an $n$-boundary
is a chain that bounds an $(n{+}1)$-chain.  Two cycles that differ by a
boundary represent the same homology class.  Thus $H_n(X)$ detects
$n$-dimensional ``holes'' in~$X$ that are not filled in.
\end{remark}

%% ----------------------------------------------------------------
%%  5.4  Functoriality
%% ----------------------------------------------------------------
\section{Functoriality}\label{sec:functoriality}

\begin{definition}[Induced chain map]\label{def:induced-chain-map}
\index{functoriality}\index{chain map!induced}
Let $f\colon X\to Y$ be a continuous map.  For each~$n$, define
$f_\sharp\colon C_n(X)\to C_n(Y)$ on generators by
$f_\sharp(\sigma)=f\circ\sigma$, extended linearly.
\end{definition}

\begin{proposition}[Chain map property]\label{prop:chain-map}
\index{chain map}
The map $f_\sharp$ is a chain map: $\partial_n\circ f_\sharp = f_\sharp\circ\partial_n$
for every~$n$.  Consequently, $f_\sharp$ descends to a well-defined homomorphism
\[
  f_*\colon H_n(X)\longrightarrow H_n(Y).
\]
\end{proposition}

\begin{proof}
On a generator $\sigma\in\operatorname{Sing}_n(X)$:
\[
  \partial_n(f_\sharp(\sigma))
  = \partial_n(f\circ\sigma)
  = \sum_{i=0}^{n}(-1)^i\,(f\circ\sigma)\circ\delta_i^n
  = \sum_{i=0}^{n}(-1)^i\,f\circ(\sigma\circ\delta_i^n)
  = f_\sharp(\partial_n(\sigma)).
\]
Since $f_\sharp$ maps cycles to cycles and boundaries to boundaries,
it induces $f_*$ on quotients.
\end{proof}

\begin{proposition}[Functoriality of homology]\label{prop:functoriality}
\index{functor!homology}
Singular homology defines a functor from the category~$\mathbf{Top}$
of topological spaces to the category~$\mathbf{Ab}$ of abelian groups:
\begin{enumerate}[label=\textup{(\roman*)},leftmargin=2.5em]
  \item $(\Id_X)_* = \Id_{H_n(X)}$.
  \item $(g\circ f)_* = g_*\circ f_*$.
\end{enumerate}
\end{proposition}

\begin{proof}
Both identities follow immediately from the corresponding identities at the
chain level: $(\Id_X)_\sharp=\Id_{C_n(X)}$ and
$(g\circ f)_\sharp = g_\sharp\circ f_\sharp$.
\end{proof}

\begin{corollary}[Homeomorphism invariance]\label{cor:homeo-invariance}
\index{homeomorphism invariance}
If $f\colon X\to Y$ is a homeomorphism, then
$f_*\colon H_n(X)\xrightarrow{\;\sim\;} H_n(Y)$ is an isomorphism
for every~$n$.
\end{corollary}

%% ----------------------------------------------------------------
%%  5.5  Homotopy invariance
%% ----------------------------------------------------------------
\section{Homotopy invariance}\label{sec:homotopy-invariance}

The deepest property of singular homology is that it depends only on
the homotopy type of the space.

\begin{definition}[Chain homotopy]\label{def:chain-homotopy}
\index{chain homotopy}
Let $f_\sharp,g_\sharp\colon C_\bullet(X)\to C_\bullet(Y)$ be chain maps.
A \textbf{chain homotopy} from $f_\sharp$ to $g_\sharp$ is a sequence of
homomorphisms $P_n\colon C_n(X)\to C_{n+1}(Y)$ satisfying
\[
  \partial_{n+1}\circ P_n + P_{n-1}\circ\partial_n
  = g_\sharp - f_\sharp
  \colon C_n(X)\to C_n(Y).
\]
\end{definition}

The key construction is the \emph{prism operator}.

\begin{definition}[Prism operator]\label{def:prism-operator}
\index{prism operator}
Let $F\colon X\times[0,1]\to Y$ be a homotopy from~$f$ to~$g$.
For a singular $n$-simplex $\sigma\colon\Delta^n\to X$, consider the prism
$\Delta^n\times[0,1]$.  We triangulate this prism into $(n{+}1)$-simplices
as follows.  For $0\le i\le n$, define
\[
  \lambda_i\colon\Delta^{n+1}\to \Delta^n\times[0,1]
\]
by
\[
  \lambda_i(e_j) =
  \begin{cases}
    (e_j,0) & \text{if } j\le i,\\
    (e_{j-1},1) & \text{if } j > i.
  \end{cases}
\]
The \textbf{prism operator} $P_n\colon C_n(X)\to C_{n+1}(Y)$ is defined on
generators by
\[
  P_n(\sigma) \;=\; \sum_{i=0}^{n}(-1)^i\,F\circ(\sigma\times\Id)\circ\lambda_i.
\]
\end{definition}

\begin{theorem}[Homotopy invariance]\label{thm:homotopy-invariance}
\index{homotopy invariance}
If $f,g\colon X\to Y$ are homotopic, then $f_*=g_*\colon H_n(X)\to H_n(Y)$
for all~$n$.
\end{theorem}

\begin{proof}
Let $F\colon X\times[0,1]\to Y$ be a homotopy with $F(-,0)=f$ and $F(-,1)=g$.
We claim the prism operator satisfies the chain homotopy identity
\[
  \partial_{n+1}\circ P_n + P_{n-1}\circ\partial_n = g_\sharp - f_\sharp.
\]
By linearity it suffices to verify on a generator $\sigma\in\operatorname{Sing}_n(X)$.
For each prism piece $F\circ(\sigma\times\Id)\circ\lambda_i$, one computes
$\partial_{n+1}$ of this $(n{+}1)$-simplex.  The faces split into three types:
\begin{enumerate}[label=(\alph*),leftmargin=2.5em]
  \item The \emph{top} face (coming from $t=1$), which yields a term in $g_\sharp(\sigma)$.
  \item The \emph{bottom} face (coming from $t=0$), which yields a term in $f_\sharp(\sigma)$.
  \item \emph{Lateral} faces, which cancel in pairs between consecutive prism
    pieces $\lambda_i$ and $\lambda_{i+1}$, thanks to the alternating signs.
\end{enumerate}
More precisely, one verifies the telescoping identity:
\[
  \partial_{n+1}\Bigl(\sum_{i=0}^n (-1)^i\,F\circ(\sigma\times\Id)\circ\lambda_i\Bigr)
  = g_\sharp(\sigma)-f_\sharp(\sigma)
    -\sum_{i=0}^n(-1)^i\,P_{n-1}(\sigma\circ\delta_i^n).
\]
Since $P_{n-1}(\partial_n(\sigma))=\sum_{i=0}^n(-1)^i P_{n-1}(\sigma\circ\delta_i^n)$,
we obtain $\partial_{n+1}\circ P_n + P_{n-1}\circ\partial_n = g_\sharp-f_\sharp$.

Therefore, for any cycle $\alpha\in Z_n(X)$,
\[
  g_\sharp(\alpha)-f_\sharp(\alpha)
  = \partial_{n+1}(P_n(\alpha))+P_{n-1}(\underbrace{\partial_n(\alpha)}_{=\,0})
  = \partial_{n+1}(P_n(\alpha))
  \in B_n(Y),
\]
so $[g_\sharp(\alpha)]=[f_\sharp(\alpha)]$ in $H_n(Y)$, i.e., $f_*=g_*$.
\end{proof}

\begin{corollary}[Homotopy equivalence invariance]\label{cor:htpy-equiv-invariance}
\index{homotopy equivalence}
If $f\colon X\to Y$ is a homotopy equivalence, then
$f_*\colon H_n(X)\xrightarrow{\;\sim\;} H_n(Y)$ is an isomorphism for all~$n$.
In particular, contractible spaces have the homology of a point.
\end{corollary}

\begin{proof}
Let $g\colon Y\to X$ be a homotopy inverse.  Then $g\circ f\simeq\Id_X$
implies $(g\circ f)_*=g_*\circ f_*=\Id$, and similarly
$f_*\circ g_*=\Id$.  So $f_*$ is an isomorphism with inverse~$g_*$.
\end{proof}

%% ----------------------------------------------------------------
%%  5.6  Reduced homology
%% ----------------------------------------------------------------
\section{Reduced homology}\label{sec:reduced-homology}

\begin{definition}[Augmented chain complex]\label{def:augmented-complex}
\index{augmented chain complex}\index{augmentation map}
The \textbf{augmentation map} $\varepsilon\colon C_0(X)\to\Z$ is defined by
$\varepsilon(\sigma)=1$ for every singular $0$-simplex~$\sigma$, extended
linearly.  One checks $\varepsilon\circ\partial_1=0$, so we obtain the
\textbf{augmented chain complex}
\[
\begin{tikzcd}
  \cdots \ar[r,"\partial_2"]
  & C_1(X) \ar[r,"\partial_1"]
  & C_0(X) \ar[r,"\varepsilon"]
  & \Z \ar[r]
  & 0.
\end{tikzcd}
\]
\end{definition}

\begin{definition}[Reduced homology]\label{def:reduced-homology}
\index{reduced homology}\index{homology!reduced}
The \textbf{reduced homology groups} of~$X$ are
\[
  \widetilde{H}_n(X) =
  \begin{cases}
    H_n(X) & \text{if } n\ge 1,\\
    \ker\varepsilon\,/\,\im\partial_1 & \text{if } n=0.
  \end{cases}
\]
Equivalently, $\widetilde{H}_n(X)$ is the homology of the augmented complex.
For non-empty~$X$, we have $H_0(X)\cong\widetilde{H}_0(X)\oplus\Z$.
\end{definition}

\begin{remark}[When to use reduced homology]\label{rem:reduced-vs-unreduced}
Reduced homology eliminates the ``trivial'' copy of~$\Z$ in degree~$0$
that every non-empty path-connected space carries.  This simplifies
many formulas, especially for suspensions and exact sequences.
\end{remark}

%% ----------------------------------------------------------------
%%  5.7  First computations
%% ----------------------------------------------------------------
\section{First computations}\label{sec:first-computations}

\begin{proposition}[Homology of a point]\label{prop:homology-point}
\index{homology!of a point}
Let $\operatorname{pt}=\{*\}$ be a one-point space.  Then
\[
  H_n(\operatorname{pt}) =
  \begin{cases}
    \Z & \text{if } n=0,\\
    0  & \text{if } n\ge 1.
  \end{cases}
\]
Equivalently, $\widetilde{H}_n(\operatorname{pt})=0$ for all~$n$.
\end{proposition}

\begin{proof}
There is exactly one singular $n$-simplex $\sigma_n\colon\Delta^n\to\{*\}$
for every~$n$.  Hence $C_n(\operatorname{pt})\cong\Z$ for all $n\ge 0$,
generated by~$\sigma_n$.
The boundary is $\partial_n(\sigma_n)=\sum_{i=0}^n(-1)^i\sigma_{n-1}$.
The sum $\sum_{i=0}^n(-1)^i$ equals~$1$ when $n$ is even and~$0$ when
$n$ is odd.  Therefore the chain complex takes the form
\[
\begin{tikzcd}
  \cdots \ar[r,"0"] & \Z \ar[r,"\Id"] & \Z \ar[r,"0"] & \Z \ar[r,"\Id"] & \Z \ar[r,"0"] & \Z \ar[r] & 0
\end{tikzcd}
\]
where the maps alternate between $\Id$ and~$0$.  Computing:
\begin{itemize}[leftmargin=2em]
  \item $H_0(\operatorname{pt})=\ker(\partial_0)/\im(\partial_1)=\Z/0=\Z$.
  \item For $n$ odd ($n\ge 1$): $\partial_n=0$ so $\ker\partial_n=\Z$,
    and $\partial_{n+1}=\Id$ so $\im\partial_{n+1}=\Z$, giving $H_n=0$.
  \item For $n$ even ($n\ge 2$): $\partial_n=\Id$ so $\ker\partial_n=0$,
    giving $H_n=0$.
\end{itemize}
Hence $H_n(\operatorname{pt})=0$ for all $n\ge 1$.
\end{proof}

\begin{theorem}[Homology of spheres]\label{thm:homology-spheres}
\index{homology!of spheres}\index{sphere!homology of}
For $n\ge 0$,
\[
  H_k(S^n) \cong
  \begin{cases}
    \Z & \text{if } k=0 \text{ or } k=n,\\
    0  & \text{otherwise},
  \end{cases}
  \qquad
  \widetilde{H}_k(S^n) \cong
  \begin{cases}
    \Z & \text{if } k=n,\\
    0  & \text{otherwise}.
  \end{cases}
\]
\end{theorem}

\begin{proof}
The case $n=0$: $S^0=\{-1,+1\}$ consists of two points, so
$H_0(S^0)\cong\Z^2$ and $\widetilde{H}_0(S^0)\cong\Z$.

For $n\ge 1$, we use the Mayer--Vietoris sequence or the long exact
sequence of the pair (to be established in Chapter~\ref{ch:les-excision}).
We defer the complete inductive proof to
Corollary~\ref{cor:homology-spheres-computed}.
\end{proof}

\begin{proposition}[Homology and path components]\label{prop:h0-components}
\index{homology!$H_0$ and path components}
For any space~$X$, $H_0(X)$ is the free abelian group on the set
$\pi_0(X)$ of path components\index{path components} of~$X$.
In particular, $H_0(X)\cong\Z^{\abs{\pi_0(X)}}$.
\end{proposition}

\begin{proof}
Every $0$-chain is a finite formal sum $\sum n_i x_i$ with $x_i\in X$.
The augmentation $\varepsilon\colon C_0(X)\to\Z$ sends this to $\sum n_i$.
A $1$-chain $\sigma$ has $\partial_1\sigma=\sigma(e_1)-\sigma(e_0)$, so the
boundary of any $1$-chain is a sum of differences of points connected by
paths.

Two points $x,y\in X$ define $0$-chains.  They represent the same class
in~$H_0$ if and only if their difference $y-x$ lies in $B_0(X)=\im\partial_1$,
i.e., there is a finite sequence of paths connecting $x$ to~$y$.
This means $x$ and $y$ lie in the same path component.
Hence $H_0(X)\cong\bigoplus_{\alpha\in\pi_0(X)}\Z$.
\end{proof}

\begin{example}[Homology of $\R^n\setminus\{0\}$]\label{ex:homology-Rn-minus-pt}
\index{homology!of $\R^n\setminus\{0\}$}
Since the inclusion $S^{n-1}\hookrightarrow\R^n\setminus\{0\}$ is a
homotopy equivalence (with homotopy inverse $x\mapsto x/\abs{x}$), we have
\[
  H_k(\R^n\setminus\{0\})\cong H_k(S^{n-1})\cong
  \begin{cases}
    \Z & \text{if } k=0 \text{ or } k=n-1,\\
    0  & \text{otherwise}.
  \end{cases}
\]
This will be crucial for proving invariance of dimension.
\end{example}

\begin{example}[Homology of disjoint unions]\label{ex:disjoint-union}
\index{homology!of disjoint union}
For a disjoint union $X=\bigsqcup_{\alpha}X_\alpha$, every singular simplex
$\sigma\colon\Delta^n\to X$ has its image contained in a single component
$X_\alpha$ (since $\Delta^n$ is connected).  Therefore
\[
  C_n(X)=\bigoplus_\alpha C_n(X_\alpha),\qquad
  H_n(X)\cong\bigoplus_\alpha H_n(X_\alpha).
\]
\end{example}

%% ----------------------------------------------------------------
%%  5.8  Relative homology
%% ----------------------------------------------------------------
\section{Relative homology}\label{sec:relative-homology}

\begin{definition}[Relative chain groups]\label{def:relative-chains}
\index{relative homology}\index{relative chain group}
Let $A\subseteq X$ be a subspace.  The inclusion $A\hookrightarrow X$
induces an injection $C_n(A)\hookrightarrow C_n(X)$.  The \textbf{relative
chain group} is
\[
  C_n(X,A) \;=\; C_n(X)\,/\,C_n(A).
\]
Since $\partial_n$ maps $C_n(A)$ to $C_{n-1}(A)$, it descends to a
boundary operator
$\bar\partial_n\colon C_n(X,A)\to C_{n-1}(X,A)$
with $\bar\partial_{n-1}\circ\bar\partial_n=0$.
\end{definition}

\begin{definition}[Relative homology groups]\label{def:relative-homology}
\index{relative homology!groups}
The \textbf{relative homology groups} are
\[
  H_n(X,A) \;=\; H_n\bigl(C_\bullet(X,A)\bigr)
  \;=\; \frac{\ker\bar\partial_n}{\im\bar\partial_{n+1}}.
\]
\end{definition}

\begin{remark}[Interpretation]\label{rem:relative-interpretation}
A class in $H_n(X,A)$ is represented by a chain in~$X$ whose boundary
lies in~$A$.  We think of relative homology as measuring the
topology of~$X$ ``modulo''~$A$.  The relative group $H_n(X,A)$ detects
``holes in~$X$ relative to~$A$'': cycles that bound in~$A$ are
considered trivial.
\end{remark}

\begin{example}[Relative homology $H_n(X,\varnothing)$]\label{ex:relative-empty}
\index{relative homology!trivial cases}
When $A=\varnothing$, we have $C_n(X,\varnothing)=C_n(X)$, so
$H_n(X,\varnothing)=H_n(X)$.  When $A=X$, the quotient
$C_n(X,X)=0$, so $H_n(X,X)=0$ for all~$n$.
\end{example}

\begin{proposition}[Functoriality of relative homology]\label{prop:relative-functoriality}
\index{relative homology!functoriality}
A continuous map of pairs $f\colon(X,A)\to(Y,B)$ (meaning $f(A)\subseteq B$)
induces homomorphisms $f_*\colon H_n(X,A)\to H_n(Y,B)$ for all~$n$.
These satisfy functoriality: $(\Id)_*=\Id$ and $(g\circ f)_*=g_*\circ f_*$.
Moreover, the connecting homomorphisms $\partial_*$ in the long exact
sequences of the pairs are natural with respect to maps of pairs.
\end{proposition}

\begin{proof}
The map $f_\sharp\colon C_n(X)\to C_n(Y)$ sends $C_n(A)$ to $C_n(B)$,
hence descends to $\bar{f}_\sharp\colon C_n(X,A)\to C_n(Y,B)$.  This
is a chain map, and functoriality follows from the chain level.
Naturality of~$\partial_*$ is a consequence of Theorem~\ref{thm:les-chain}.
\end{proof}

\begin{proposition}[Short exact sequence of chain complexes]\label{prop:ses-chains}
\index{short exact sequence!of chain complexes}
For a pair $(X,A)$, the sequence
\[
\begin{tikzcd}
  0 \ar[r]
  & C_n(A) \ar[r,"i_\sharp"]
  & C_n(X) \ar[r,"j_\sharp"]
  & C_n(X,A) \ar[r]
  & 0
\end{tikzcd}
\]
is exact for every~$n$, where $i\colon A\hookrightarrow X$ is the inclusion
and $j_\sharp$ is the quotient map.
\end{proposition}

\begin{proof}
Exactness at $C_n(A)$: the map $i_\sharp$ is injective since $C_n(A)$
is a free abelian group on $\operatorname{Sing}_n(A)\subseteq\operatorname{Sing}_n(X)$.
Exactness at $C_n(X,A)$: by definition of quotient.
Exactness at $C_n(X)$: $\ker j_\sharp=\im i_\sharp=C_n(A)$.
\end{proof}

%% ----------------------------------------------------------------
%%  5.9  Exercises
%% ----------------------------------------------------------------
\section{Exercises}\label{sec:ex-ch5}

\begin{exercise}\label{exo:simplicial-boundary-2simplex}
\index{boundary operator!computation}
Let $\sigma\colon\Delta^2\to X$ be a singular $2$-simplex.  Verify by
direct computation that $\partial_1(\partial_2(\sigma))=0$.
\end{exercise}

\begin{exercise}\label{exo:h0-discrete}
\index{homology!$H_0$ of discrete space}
Let $X$ be a set with the discrete topology.  Compute $H_n(X)$ for all~$n$.
\end{exercise}

\begin{exercise}\label{exo:homology-contractible}
Show that if $X$ is contractible then $\widetilde{H}_n(X)=0$ for all~$n$.
\end{exercise}

\begin{exercise}\label{exo:chain-homotopy-equivalence}
\index{chain homotopy!equivalence}
Show that a chain homotopy equivalence $f_\sharp\colon C_\bullet(X)\to
C_\bullet(Y)$ induces isomorphisms on homology.  Deduce that
$H_n(X)\cong H_n(Y)$ if $X\simeq Y$.
\end{exercise}

\begin{exercise}\label{exo:cone-acyclic}
\index{cone}\index{acyclic}
Let $CX=(X\times[0,1])/(X\times\{1\})$ be the cone on~$X$.  Show that
$CX$ is contractible, and hence $\widetilde{H}_n(CX)=0$ for all~$n$.
\end{exercise}

\begin{exercise}\label{exo:relative-pair-point}
\index{relative homology!of a pair}
Let $x_0\in X$.  Show that $H_n(X,\{x_0\})\cong\widetilde{H}_n(X)$
for all~$n\ge 0$.
\end{exercise}

\begin{exercise}\label{exo:prism-explicit}
\index{prism operator!explicit}
Let $F\colon [0,1]\times[0,1]\to Y$ be a homotopy of paths.  Write out the
prism operator $P_1$ explicitly and verify the chain homotopy identity
$\partial_2\circ P_1+P_0\circ\partial_1=g_\sharp-f_\sharp$
for a singular $1$-simplex.
\end{exercise}

\begin{exercise}\label{exo:retract-homology}
\index{retract}
Let $A\subseteq X$ be a retract (i.e., there exists $r\colon X\to A$
with $r\circ i=\Id_A$).  Show that $i_*\colon H_n(A)\to H_n(X)$ is
injective and $r_*\colon H_n(X)\to H_n(A)$ is surjective for all~$n$.
\end{exercise}

\begin{exercise}\label{exo:homology-disjoint-union}
\index{homology!of disjoint union}
Prove that $H_n(X\sqcup Y)\cong H_n(X)\oplus H_n(Y)$ for all~$n$.
\end{exercise}

\begin{exercise}\label{exo:euler-char-chain}
\index{Euler characteristic}
Let $C_\bullet$ be a chain complex of finite-dimensional vector spaces
over a field~$k$, with $C_n=0$ for all but finitely many~$n$.  Prove that
\[
  \sum_{n\ge 0}(-1)^n\dim_k C_n
  = \sum_{n\ge 0}(-1)^n\dim_k H_n(C_\bullet).
\]
\textit{Hint}: use the rank--nullity theorem for the short exact sequences
$0\to Z_n\to C_n\to B_{n-1}\to 0$ and $0\to B_n\to Z_n\to H_n\to 0$.
\end{exercise}

\index{singular homology|)}

%% ================================================================
%%  CHAPTER 6 — Long Exact Sequences and Excision
%% ================================================================
\chapter{Long Exact Sequences and Excision}\label{ch:les-excision}
\index{long exact sequence|(}\index{excision|(}

In this chapter we develop the main computational tools of singular
homology: the long exact sequence of a pair and the excision theorem.
Together, they allow the inductive computation of homology groups for a
wide class of spaces.

%% ----------------------------------------------------------------
%%  6.1  Short exact sequences of chain complexes
%% ----------------------------------------------------------------
\section{Short exact sequences of chain complexes}\label{sec:ses-chain}

\begin{definition}[Short exact sequence of chain complexes]\label{def:ses-chain-complexes}
\index{short exact sequence!of chain complexes}
A \textbf{short exact sequence of chain complexes} is a diagram
\[
\begin{tikzcd}
  0 \ar[r]
  & A_\bullet \ar[r,"i"]
  & B_\bullet \ar[r,"p"]
  & C_\bullet \ar[r]
  & 0
\end{tikzcd}
\]
where $i$ and $p$ are chain maps such that for every~$n$,
\[
\begin{tikzcd}
  0 \ar[r]
  & A_n \ar[r,"i_n"]
  & B_n \ar[r,"p_n"]
  & C_n \ar[r]
  & 0
\end{tikzcd}
\]
is a short exact sequence of abelian groups.
\end{definition}

%% ----------------------------------------------------------------
%%  6.2  The connecting homomorphism
%% ----------------------------------------------------------------
\section{The connecting homomorphism}\label{sec:connecting-hom}

\begin{definition}[Connecting homomorphism]\label{def:connecting-hom}
\index{connecting homomorphism}
Given a short exact sequence of chain complexes
$0\to A_\bullet\xrightarrow{i} B_\bullet\xrightarrow{p} C_\bullet\to 0$,
the \textbf{connecting homomorphism}
$\partial_*\colon H_n(C_\bullet)\to H_{n-1}(A_\bullet)$
is defined as follows.

Let $[c]\in H_n(C_\bullet)$ with $\partial^C_n c=0$.
\begin{enumerate}[label=\textup{Step \arabic*.},leftmargin=3.5em]
  \item Since $p_n$ is surjective, choose $b\in B_n$ with $p_n(b)=c$.
  \item Then $p_{n-1}(\partial^B_n b)=\partial^C_n(p_n(b))=\partial^C_n c=0$.
  \item By exactness at $B_{n-1}$, there exists a unique $a\in A_{n-1}$
    with $i_{n-1}(a)=\partial^B_n b$.
  \item One checks $\partial^A_{n-1}a=0$ (since $i_{n-2}(\partial^A_{n-1}a)=\partial^B_{n-1}(\partial^B_n b)=0$ and $i_{n-2}$ is injective).
  \item Set $\partial_*[c]=[a]\in H_{n-1}(A_\bullet)$.
\end{enumerate}

The following diagram summarizes the construction:
\[
\begin{tikzcd}[column sep=large]
  & b \ar[d,mapsto,"\partial^B_n"] \ar[r,mapsto,"p_n"] & c \ar[d,mapsto,"\partial^C_n"] \\
  a \ar[r,mapsto,"i_{n-1}"] & \partial^B_n b \ar[r,mapsto,"p_{n-1}"] & 0
\end{tikzcd}
\]
\end{definition}

\begin{lemma}[Well-definedness of $\partial_*$]\label{lem:connecting-well-defined}
\index{connecting homomorphism!well-definedness}
The class $[a]$ does not depend on the choice of~$b$.
\end{lemma}

\begin{proof}
Suppose $b'$ is another lift with $p_n(b')=c$.  Then $p_n(b-b')=0$, so
$b-b'=i_n(a'')$ for some $a''\in A_n$.  Then
\[
  i_{n-1}(a-a')
  = \partial^B_n b - \partial^B_n b'
  = \partial^B_n(b-b')
  = \partial^B_n(i_n(a''))
  = i_{n-1}(\partial^A_n a'').
\]
Since $i_{n-1}$ is injective, $a-a'=\partial^A_n a''$, so
$[a]=[a']\in H_{n-1}(A_\bullet)$.
\end{proof}

%% ----------------------------------------------------------------
%%  6.3  The Snake Lemma
%% ----------------------------------------------------------------
\section{The Snake Lemma}\label{sec:snake-lemma}

The connecting homomorphism is a special case of the following
fundamental result in homological algebra.

\begin{lemma}[Snake Lemma]\label{lem:snake}
\index{snake lemma}
Consider a commutative diagram of abelian groups with exact rows:
\[
\begin{tikzcd}
  & A \ar[r,"f"] \ar[d,"\alpha"] & B \ar[r,"g"] \ar[d,"\beta"] & C \ar[r] \ar[d,"\gamma"] & 0 \\
  0 \ar[r] & A' \ar[r,"f'"] & B' \ar[r,"g'"] & C'
\end{tikzcd}
\]
Then there is an exact sequence
\[
\begin{tikzcd}[column sep=1.8em]
  \ker\alpha \ar[r] & \ker\beta \ar[r] & \ker\gamma
    \ar[r,"\delta",out=0,in=180,looseness=1.5,overlay]
  & \coker\alpha \ar[r] & \coker\beta \ar[r] & \coker\gamma,
\end{tikzcd}
\]
where $\delta\colon\ker\gamma\to\coker\alpha$ is the \textbf{connecting
map} (or ``snake map'') defined by the diagram chase
$c\mapsto [(\alpha\circ(f)^{-1}\circ(\beta)^{-1}\circ g)(c)]$, made
precise as follows: given $c\in\ker\gamma$, lift to $b\in B$ with
$g(b)=c$; then $g'(\beta(b))=\gamma(g(b))=\gamma(c)=0$, so
$\beta(b)=f'(a')$ for unique $a'\in A'$; set $\delta(c)=[a']\in\coker\alpha$.
\end{lemma}

\begin{proof}
The proof is a diagram chase.  We verify exactness at each term.

\textit{Exactness at $\ker\beta$}:
The map $\ker\alpha\to\ker\beta$ is induced by~$f$.  Its image consists
of elements $f(a)$ with $\alpha(a)=0$.  The map $\ker\beta\to\ker\gamma$
is induced by~$g$.  If $b\in\ker\beta$ maps to $0$ in $\ker\gamma$, then
$g(b)=0$, so $b=f(a)$ for some~$a$ by exactness of the top row.
Then $f'(\alpha(a))=\beta(f(a))=\beta(b)=0$, and $f'$ is injective, so
$\alpha(a)=0$, i.e., $a\in\ker\alpha$.

\textit{Exactness at $\ker\gamma$}:
If $c=g(b)$ with $b\in\ker\beta$, then $\beta(b)=0=f'(0)$, so
$\delta(c)=[0]=0$.  Conversely, if $\delta(c)=0$, then the $a'$
from the construction satisfies $a'=\alpha(a)$ for some~$a$.  Replace
$b$ by $b-f(a)$; then $g(b-f(a))=c$ and $\beta(b-f(a))=\beta(b)-f'(\alpha(a))=f'(a')-f'(a')=0$, so $b-f(a)\in\ker\beta$ maps to~$c$.

\textit{Exactness at $\coker\alpha$}:
First, $\im\delta\subseteq\ker(\coker\alpha\to\coker\beta)$: if $c\in\ker\gamma$
with $\delta(c)=[a']$, then by construction $f'(a')=\beta(b)$ for some~$b$,
so $[f'(a')]=[\beta(b)]=0$ in $\coker\beta$.

Conversely, suppose $[a']\in\coker\alpha$ maps to~$0$ in $\coker\beta$,
i.e., $f'(a')=\beta(b)$ for some $b\in B$.  Set $c=g(b)$.  Then
$\gamma(c)=\gamma(g(b))=g'(\beta(b))=g'(f'(a'))=0$ (since $g'\circ f'=0$
by exactness of the bottom row at~$B'$).  So $c\in\ker\gamma$.  The
lift of~$c$ is~$b$, and $\beta(b)=f'(a')$, so $\delta(c)=[a']$.

\textit{Exactness at $\coker\beta$}:
The map $\coker\alpha\to\coker\beta$ sends $[a']$ to $[f'(a')]$.  The
map $\coker\beta\to\coker\gamma$ sends $[b']$ to $[g'(b')]$.
If $[b']=[f'(a')]$, then $g'(b')=g'(f'(a'))=0$, so $[g'(b')]=0$.
Conversely, if $[g'(b')]=0$, then $g'(b')=\gamma(c_0)$ for some~$c_0$.
Pick $b_0$ with $g(b_0)=c_0$; then $g'(b'-\beta(b_0))=g'(b')-\gamma(g(b_0))=0$.
By exactness at~$B'$, $b'-\beta(b_0)=f'(a')$ for some~$a'$.
Hence $[b']=[f'(a')+\beta(b_0)]=[f'(a')]$ in $\coker\beta$.

Exactness at $\coker\gamma$ is similar and left to the reader.
\end{proof}

%% ----------------------------------------------------------------
%%  6.4  The long exact sequence in homology
%% ----------------------------------------------------------------
\section{The long exact sequence in homology}\label{sec:les-homology}

\begin{theorem}[Long exact sequence of a short exact sequence of chain complexes]\label{thm:les-chain}
\index{long exact sequence!of chain complexes}
Given a short exact sequence of chain complexes
\[
\begin{tikzcd}
  0 \ar[r] & A_\bullet \ar[r,"i"] & B_\bullet \ar[r,"p"] & C_\bullet \ar[r] & 0,
\end{tikzcd}
\]
there is a long exact sequence
\[
\begin{tikzcd}[column sep=1.4em]
  \cdots \ar[r]
  & H_n(A) \ar[r,"i_*"]
  & H_n(B) \ar[r,"p_*"]
  & H_n(C) \ar[r,"\partial_*"]
  & H_{n-1}(A) \ar[r,"i_*"]
  & H_{n-1}(B) \ar[r]
  & \cdots
\end{tikzcd}
\]
extending to the right as
$\cdots\to H_1(C)\xrightarrow{\partial_*} H_0(A)\xrightarrow{i_*} H_0(B)
\xrightarrow{p_*} H_0(C)\to 0$.
Moreover, this construction is natural: a morphism of short exact sequences
induces a morphism of long exact sequences.
\end{theorem}

\begin{proof}
We must prove exactness at three places: at $H_n(B)$, at $H_n(C)$, and
at $H_{n-1}(A)$.

\medskip\noindent
\textbf{Exactness at $H_n(B)$}: $\im i_*\subseteq\ker p_*$.

We have $p_*\circ i_*=(p\circ i)_*=0$ since $p\circ i=0$.

\medskip
Conversely, let $[b]\in\ker p_*$, so $p_n(b)=\partial^C_{n+1}(c')$ for some
$c'\in C_{n+1}$.  Since $p_{n+1}$ is surjective, pick $b'\in B_{n+1}$
with $p_{n+1}(b')=c'$.  Then $p_n(b-\partial^B_{n+1}b')=p_n(b)-\partial^C_{n+1}c'=0$.
By exactness, $b-\partial^B_{n+1}b'=i_n(a)$ for some $a\in A_n$.
Now $i_{n-1}(\partial^A_n a)=\partial^B_n(i_n(a))=\partial^B_n(b)-0=0$
(since $b$ is a cycle and $\partial^B\circ\partial^B=0$), so $\partial^A_n a=0$
by injectivity of~$i_{n-1}$.  Thus $a\in Z_n(A)$ and
$[b]=[i_n(a)]=i_*[a]$.

\medskip\noindent
\textbf{Exactness at $H_n(C)$}: $\im p_*\subseteq\ker\partial_*$.

If $[c]=p_*[b]$, then $c=p_n(b)$ with $\partial^B_n b=0$.
In the construction of~$\partial_*$, we may choose this~$b$ as our
lift, giving $\partial^B_n b=0=i_{n-1}(0)$, so $\partial_*[c]=[0]=0$.

\medskip
Conversely, suppose $\partial_*[c]=0$.  Using the construction, pick
$b$ with $p_n(b)=c$ and let $a\in A_{n-1}$ satisfy
$i_{n-1}(a)=\partial^B_n b$.  If $[a]=0$ in $H_{n-1}(A)$, then
$a=\partial^A_n(a')$ for some $a'\in A_n$.  Set $\tilde{b}=b-i_n(a')$.
Then $p_n(\tilde b)=c$ and
$\partial^B_n\tilde{b}=\partial^B_n b-i_{n-1}(\partial^A_n a')=i_{n-1}(a)-i_{n-1}(a)=0$, so $\tilde{b}\in Z_n(B)$ and
$p_*[\tilde{b}]=[c]$.

\medskip\noindent
\textbf{Exactness at $H_{n-1}(A)$}: $\im\partial_*\subseteq\ker i_*$.

By construction, $i_{n-1}(a)=\partial^B_n b\in B_{n-1}(B)$, so
$i_*[a]=[\partial^B_n b]=0$ in $H_{n-1}(B)$.

\medskip
Conversely, suppose $[a]\in H_{n-1}(A)$ with $i_*[a]=0$, i.e.,
$i_{n-1}(a)=\partial^B_n(b)$ for some $b\in B_n$.  Set $c=p_n(b)$.
Then $\partial^C_n c=p_{n-1}(\partial^B_n b)=p_{n-1}(i_{n-1}(a))=0$,
so $c\in Z_n(C)$.  By the construction of~$\partial_*$ (using this~$b$
as the lift), $\partial_*[c]=[a]$.

\medskip\noindent
\textbf{Naturality}: Given a morphism of short exact sequences
\[
\begin{tikzcd}
  0 \ar[r] & A_\bullet \ar[r] \ar[d,"\alpha"] & B_\bullet \ar[r] \ar[d,"\beta"]
    & C_\bullet \ar[r] \ar[d,"\gamma"] & 0 \\
  0 \ar[r] & A'_\bullet \ar[r] & B'_\bullet \ar[r] & C'_\bullet \ar[r] & 0
\end{tikzcd}
\]
one checks that $\partial_*\circ\gamma_*=\alpha_*\circ\partial_*$ by tracing
the diagram chase through both compositions.  The remaining squares
commute by functoriality of homology.
\end{proof}

%% ----------------------------------------------------------------
%%  6.5  The long exact sequence of a pair
%% ----------------------------------------------------------------
\section{The long exact sequence of a pair}\label{sec:les-pair}

\begin{theorem}[Long exact sequence of a pair]\label{thm:les-pair}
\index{long exact sequence!of a pair}
For any pair $(X,A)$ with $A\subseteq X$, there is a long exact sequence
\[
\begin{tikzcd}[column sep=1.4em]
  \cdots \ar[r]
  & H_n(A) \ar[r,"i_*"]
  & H_n(X) \ar[r,"j_*"]
  & H_n(X,A) \ar[r,"\partial_*"]
  & H_{n-1}(A) \ar[r,"i_*"]
  & H_{n-1}(X) \ar[r]
  & \cdots
\end{tikzcd}
\]
which is natural in the pair $(X,A)$.
\end{theorem}

\begin{proof}
Apply Theorem~\ref{thm:les-chain} to the short exact sequence of
chain complexes from Proposition~\ref{prop:ses-chains}:
\[
  0\to C_\bullet(A)\xrightarrow{i_\sharp}
  C_\bullet(X)\xrightarrow{j_\sharp}
  C_\bullet(X,A)\to 0. \qedhere
\]
\end{proof}

\begin{example}[The pair $(D^n,S^{n-1})$]\label{ex:pair-disk-sphere}
\index{disk!homology relative to sphere}
Consider the pair $(D^n,S^{n-1})$ where $D^n$ is the closed unit disk.
Since $D^n$ is contractible, $\widetilde{H}_k(D^n)=0$ for all~$k$.
The long exact sequence gives
\[
\begin{tikzcd}[column sep=1.3em]
  \cdots \ar[r]
  & H_k(D^n) \ar[r]
  & H_k(D^n,S^{n-1}) \ar[r,"\partial_*"]
  & H_{k-1}(S^{n-1}) \ar[r]
  & H_{k-1}(D^n) \ar[r]
  & \cdots
\end{tikzcd}
\]
For $k\ge 2$, the outer terms vanish, yielding
$H_k(D^n,S^{n-1})\cong H_{k-1}(S^{n-1})$.
\end{example}

%% ----------------------------------------------------------------
%%  6.6  Excision
%% ----------------------------------------------------------------
\section{The excision theorem}\label{sec:excision}

\begin{theorem}[Excision]\label{thm:excision}
\index{excision theorem}
Let $X$ be a topological space and let $Z\subseteq A\subseteq X$ be
subspaces such that $\overline{Z}\subseteq\operatorname{int}(A)$.
Then the inclusion $(X\setminus Z,A\setminus Z)\hookrightarrow(X,A)$
induces isomorphisms
\[
  H_n(X\setminus Z,\, A\setminus Z)
  \;\xrightarrow{\;\cong\;}
  H_n(X,A)
  \qquad\text{for all } n.
\]
Equivalently, if $A,B\subseteq X$ are subspaces whose interiors cover~$X$,
i.e., $X=\operatorname{int}(A)\cup\operatorname{int}(B)$, then
$H_n(B,A\cap B)\xrightarrow{\cong} H_n(X,A)$.
\end{theorem}

The proof uses \emph{barycentric subdivision}\index{barycentric subdivision}
to make chains ``small'' enough that they lie entirely in~$A$ or in
$X\setminus Z$.

\begin{definition}[Barycentric subdivision]\label{def:barycentric-subdivision}
\index{barycentric subdivision!operator}
The \textbf{barycentric subdivision operator}
$\operatorname{sd}_n\colon C_n(X)\to C_n(X)$
is defined inductively.  For a singular $0$-simplex~$\sigma$, set
$\operatorname{sd}_0(\sigma)=\sigma$.  For $n\ge 1$ and
$\sigma\colon\Delta^n\to X$, let $b_\sigma=\sigma(\text{barycenter of }\Delta^n)$
and define
\[
  \operatorname{sd}_n(\sigma)
  = b_\sigma \cdot \operatorname{sd}_{n-1}(\partial_n\sigma),
\]
where $b\cdot c$ denotes the ``cone from~$b$'' operation on chains,
extended linearly.
\end{definition}

\begin{lemma}[Properties of barycentric subdivision]\label{lem:sd-properties}
\index{barycentric subdivision!properties}
\begin{enumerate}[label=\textup{(\roman*)},leftmargin=2.5em]
  \item $\operatorname{sd}$ is a chain map: $\partial_n\circ\operatorname{sd}_n = \operatorname{sd}_{n-1}\circ\partial_n$.
  \item $\operatorname{sd}$ is chain homotopic to the identity:
    there exist homomorphisms $T_n\colon C_n(X)\to C_{n+1}(X)$ with
    $\partial_{n+1}\circ T_n + T_{n-1}\circ\partial_n = \operatorname{sd}_n-\Id$.
  \item After sufficiently many iterations, $\operatorname{sd}^m(\sigma)$ has
    arbitrarily small image for each simplex~$\sigma$: the diameter of each
    simplex in $\operatorname{sd}^m(\Delta^n)$ is at most
    $\bigl(\frac{n}{n+1}\bigr)^m\operatorname{diam}(\Delta^n)$.
\end{enumerate}
\end{lemma}

\begin{proof}[Proof sketch]
Part~(i) follows by induction using the cone construction and the
identity $\partial(b\cdot c)=c-b\cdot\partial c$ (valid when $\partial c$
is defined and the cone makes sense).

Part~(ii) uses the same cone technique: $T_n(\sigma)=b_\sigma\cdot
T_{n-1}(\partial_n\sigma)+b_\sigma\cdot(\operatorname{sd}_n(\sigma)-\sigma)$,
adjusted to make the identity work by induction.

Part~(iii) is a standard estimate.  Each barycentric subdivision replaces
an $n$-simplex with simplices whose diameter is at most
$\frac{n}{n+1}$ times the original.  After $m$ iterations, the factor
is $\bigl(\frac{n}{n+1}\bigr)^m\to 0$.
\end{proof}

\begin{proof}[Proof of the Excision Theorem~\ref{thm:excision}]
We use the formulation: $\set{A,B}$ is an open cover with
$X=\operatorname{int}(A)\cup\operatorname{int}(B)$.  Let
$C_n^{\set{A,B}}(X)\subseteq C_n(X)$ denote the subgroup of chains that
are sums of simplices each mapping into~$A$ or into~$B$.  We need to
show that the inclusion
$C_\bullet^{\set{A,B}}(X)/C_\bullet(A)\hookrightarrow C_\bullet(X)/C_\bullet(A)$
induces isomorphisms on homology.

\medskip\noindent
\textit{Step 1.}
By the Lebesgue number lemma\index{Lebesgue number lemma} (in the compact
setting) or direct argument, for every singular simplex
$\sigma\colon\Delta^n\to X$, there exists $m\ge 0$ such that every simplex
in $\operatorname{sd}^m(\sigma)$ maps into $A$ or into~$B$.  That is,
$\operatorname{sd}^m(\sigma)\in C_n^{\set{A,B}}(X)$.

\medskip\noindent
\textit{Step 2.}
By Lemma~\ref{lem:sd-properties}(ii), $\operatorname{sd}^m$ is chain homotopic
to the identity.  More precisely, there exists $D_n\colon C_n(X)\to C_{n+1}(X)$
with $\operatorname{sd}^m-\Id=\partial_{n+1}\circ D_n+D_{n-1}\circ\partial_n$.

\medskip\noindent
\textit{Step 3 (Surjectivity).}
Given a relative cycle $\alpha\in C_n(X)$ with
$\partial_n\alpha\in C_{n-1}(A)$, we show $[\alpha]$ lies in the image.
Choose~$m$ large enough so that $\operatorname{sd}^m(\alpha)\in
C_n^{\set{A,B}}(X)$.  Then
$\operatorname{sd}^m(\alpha)-\alpha=\partial_{n+1}(D_n\alpha)+D_{n-1}(\partial_n\alpha)$.
Since $\partial_n\alpha\in C_{n-1}(A)$, the term $D_{n-1}(\partial_n\alpha)$
lies in $C_n(A)$.  Hence $\operatorname{sd}^m(\alpha)\equiv\alpha+\partial_{n+1}(D_n\alpha)
\pmod{C_n(A)}$, showing $[\alpha]=[\operatorname{sd}^m(\alpha)]$
in $H_n(X,A)$.  Since $\operatorname{sd}^m(\alpha)\in C_n^{\set{A,B}}(X)$,
surjectivity follows.

\medskip\noindent
\textit{Step 4 (Injectivity).}
Suppose $\alpha\in C_n^{\set{A,B}}(X)$ represents~$0$ in $H_n(X,A)$,
i.e., $\alpha=\partial_{n+1}\beta+\gamma$ with $\beta\in C_{n+1}(X)$
and $\gamma\in C_n(A)$.  Choose~$m$ so that $\operatorname{sd}^m(\beta)\in
C_{n+1}^{\set{A,B}}(X)$.  Then
$\operatorname{sd}^m(\alpha)=\partial_{n+1}(\operatorname{sd}^m\beta)+\operatorname{sd}^m(\gamma)$
and $\alpha-\operatorname{sd}^m(\alpha)=\partial D\alpha+D\partial\alpha$.
Since $\partial\alpha\in C_{n-1}(A)$ (as $\alpha-\gamma=\partial\beta$
implies $\partial\alpha=\partial\gamma\in C_{n-1}(A)$), the chain homotopy
terms stay controlled, and one concludes $[\alpha]=0$
in $H_n(C_\bullet^{\set{A,B}}(X)/C_\bullet(A))$.

\medskip
Finally, the natural identification
$C_\bullet^{\set{A,B}}(X)/C_\bullet(A)\cong C_\bullet(B)/C_\bullet(A\cap B)$
(since a simplex in $C^{\set{A,B}}$ not in $C(A)$ must map into~$B$)
yields the desired isomorphism $H_n(B,A\cap B)\cong H_n(X,A)$.
\end{proof}

%% ----------------------------------------------------------------
%%  6.7  Relative homology and quotient spaces
%% ----------------------------------------------------------------
\section{Relative homology and quotient spaces}\label{sec:relative-quotient}

\begin{theorem}[Collapsing a subspace]\label{thm:relative-quotient}
\index{relative homology!and quotient spaces}
Let $(X,A)$ be a pair such that $A$ is a closed subspace and is a
deformation retract of some open neighborhood in~$X$ (i.e., $A$ is a
``good pair''\index{good pair}).  Then the quotient map
$q\colon(X,A)\to(X/A,A/A)$ induces isomorphisms
\[
  q_*\colon H_n(X,A)\;\xrightarrow{\;\cong\;}\;
  H_n(X/A,\,A/A)\;\cong\;\widetilde{H}_n(X/A)
  \qquad\text{for all } n.
\]
\end{theorem}

\begin{proof}
Let $U$ be an open neighborhood of~$A$ that deformation retracts onto~$A$.
By homotopy invariance, $H_n(U,A)\cong H_n(A,A)=0$.  The long exact
sequence of the triple $(X,U,A)$ then gives $H_n(X,A)\cong H_n(X,U)$.
By excision (removing $A$ from both):
\[
  H_n(X\setminus A,\,U\setminus A)\;\cong\; H_n(X,U).
\]
The quotient map~$q$ restricts to a homeomorphism
$X\setminus A\cong X/A\setminus\{*\}$, and the analogous excision argument
for the quotient gives
\[
  H_n(X/A,\,U/A)\;\cong\; H_n(X/A\setminus\{*\},\,U/A\setminus\{*\}).
\]
Since $U/A$ is contractible (it deformation retracts to the point $A/A$),
$H_n(X/A,\,U/A)\cong\widetilde{H}_n(X/A)$.  Combining all isomorphisms
yields the result.
\end{proof}

\begin{corollary}[Homology of spheres, complete computation]\label{cor:homology-spheres-computed}
\index{sphere!homology of}
For $n\ge 1$,
\[
  \widetilde{H}_k(S^n)
  \cong
  \begin{cases}
    \Z & \text{if } k=n,\\
    0 & \text{otherwise}.
  \end{cases}
\]
\end{corollary}

\begin{proof}
We use the long exact sequence of the pair $(D^n,S^{n-1})$.  Since
$D^n$ is contractible, $\widetilde{H}_k(D^n)=0$ for all~$k$, and
the long exact sequence gives
\[
  \widetilde{H}_k(D^n)\to H_k(D^n,S^{n-1})\xrightarrow{\partial_*}
  \widetilde{H}_{k-1}(S^{n-1})\to\widetilde{H}_{k-1}(D^n),
\]
so $H_k(D^n,S^{n-1})\cong\widetilde{H}_{k-1}(S^{n-1})$.
By Theorem~\ref{thm:relative-quotient}, $H_k(D^n,S^{n-1})\cong
\widetilde{H}_k(D^n/S^{n-1})\cong\widetilde{H}_k(S^n)$.
Therefore $\widetilde{H}_k(S^n)\cong\widetilde{H}_{k-1}(S^{n-1})$.
Induction on~$n$ with base case $\widetilde{H}_0(S^0)\cong\Z$
completes the proof.
\end{proof}

%% ----------------------------------------------------------------
%%  6.8  Applications
%% ----------------------------------------------------------------
\section{Applications}\label{sec:applications-homology}

\subsection{Degree of maps between spheres}\label{subsec:degree}

\begin{definition}[Degree]\label{def:degree}
\index{degree!of a map}
Let $f\colon S^n\to S^n$ be a continuous map with $n\ge 1$.  Since
$H_n(S^n)\cong\Z$, the induced map
$f_*\colon H_n(S^n)\to H_n(S^n)$
is multiplication by an integer $d\in\Z$.  This integer is the
\textbf{degree} of~$f$, denoted $\deg f$.
\end{definition}

\begin{proposition}[Properties of degree]\label{prop:degree-properties}
\index{degree!properties}
\begin{enumerate}[label=\textup{(\roman*)},leftmargin=2.5em]
  \item $\deg(\Id_{S^n})=1$.
  \item $\deg(g\circ f)=\deg(g)\cdot\deg(f)$.
  \item If $f\simeq g$, then $\deg f=\deg g$.
  \item The antipodal map $a\colon S^n\to S^n$, $x\mapsto -x$, has
    $\deg a=(-1)^{n+1}$.
  \item A reflection of~$S^n$ (negating one coordinate) has degree~$-1$.
\end{enumerate}
\end{proposition}

\begin{proof}
(i) Functoriality: $(\Id_{S^n})_*=\Id_{H_n(S^n)}$, which is multiplication
by~$1$.

(ii) $(g\circ f)_*=g_*\circ f_*$ is multiplication by
$\deg g\cdot\deg f$.

(iii) Homotopy invariance (Theorem~\ref{thm:homotopy-invariance}).

(iv) Let $r_i\colon S^n\to S^n$ be the reflection
$(x_0,\dots,x_n)\mapsto(x_0,\dots,-x_i,\dots,x_n)$.  The antipodal map
is $a=r_0\circ r_1\circ\cdots\circ r_n$, a composition of $(n{+}1)$
reflections.  By~(ii), $\deg a=(\deg r_i)^{n+1}$.  So it suffices to
show each $r_i$ has degree~$-1$.

(v) We verify $\deg r_i=-1$ for reflections.  Consider the
generator $[\mu_n]\in H_n(S^n)\cong\Z$ arising from the long exact
sequence of $(D^n_+,S^{n-1})$.  The reflection~$r_i$ reverses the
orientation of the fundamental class: the induced map on the relative
homology $H_n(D^n,S^{n-1})$ is multiplication by~$-1$ (since $r_i$
reverses orientation of~$\R^{n+1}$ restricted to the relevant hemisphere).
Following the isomorphisms through the connecting homomorphism yields
$\deg r_i=-1$.

Combining (iv) and (v): $\deg a=(-1)^{n+1}$.
\end{proof}

\begin{example}[Degree and surjectivity]\label{ex:degree-surjective}
\index{degree!and surjectivity}
If $f\colon S^n\to S^n$ has $\deg f\ne 0$, then $f$ is surjective.
Indeed, if $y\notin f(S^n)$, then $f$ factors through
$S^n\setminus\{y\}\cong\R^n$, which is contractible, so $f_*=0$
and $\deg f=0$.
\end{example}

\subsection{Brouwer fixed point theorem}\label{subsec:brouwer}

\begin{theorem}[Brouwer Fixed Point Theorem]\label{thm:brouwer}
\index{Brouwer fixed point theorem}
Every continuous map $f\colon D^n\to D^n$ has a fixed point, for all
$n\ge 1$.
\end{theorem}

\begin{proof}
Suppose $f$ has no fixed point.  Define $r\colon D^n\to S^{n-1}$ by
letting $r(x)$ be the point where the ray from $f(x)$ through~$x$
meets~$S^{n-1}$.  Then $r$ is continuous and $r|_{S^{n-1}}=\Id_{S^{n-1}}$,
so $r$ is a retraction.

Let $i\colon S^{n-1}\hookrightarrow D^n$ be the inclusion.  Then
$r\circ i=\Id_{S^{n-1}}$, so $(r\circ i)_*=r_*\circ i_*=\Id$ on
$H_{n-1}(S^{n-1})\cong\Z$.  In particular, $i_*$ is injective.

But $D^n$ is contractible, so $H_{n-1}(D^n)=0$ for $n\ge 2$, meaning
$i_*\colon\Z\to 0$ is the zero map---a contradiction.

For $n=1$: $H_0(S^0)\cong\Z^2$ but the retraction forces
$r_*\circ i_*=\Id$ on $\widetilde{H}_0(S^0)\cong\Z$, while
$\widetilde{H}_0(D^1)=0$, giving the same contradiction.
\end{proof}

\begin{corollary}[No-retraction theorem]\label{cor:no-retraction}
\index{no-retraction theorem}
There is no retraction $r\colon D^n\to S^{n-1}$ for $n\ge 1$.
\end{corollary}

\begin{proof}
This was established in the proof above: such a retraction leads to
a contradiction on homology groups.
\end{proof}

\subsection{Invariance of domain}\label{subsec:invariance-domain}

\begin{theorem}[Invariance of Domain]\label{thm:invariance-domain}
\index{invariance of domain}
Let $U\subseteq\R^n$ be open and let $f\colon U\to\R^n$ be a continuous
injection.  Then $f(U)$ is open in~$\R^n$.
\end{theorem}

\begin{proof}
It suffices to show that $f(U)$ is open at each point.  Let $x\in U$;
we show $f(x)$ is an interior point of~$f(U)$.  Choose a closed ball
$\overline{B}\subseteq U$ centered at~$x$.  By the following key lemma,
$f(\overline{B})$ contains a neighborhood of~$f(x)$, so $f(x)\in
\operatorname{int}(f(U))$.
\end{proof}

\begin{lemma}[Key lemma for invariance of domain]\label{lem:invariance-key}
Let $f\colon D^n\to\R^n$ be a continuous injection.  Then $f(x)\notin
f(S^{n-1})$ implies $f(x)\in\operatorname{int}(f(D^n))$ for any
$x\in\operatorname{int}(D^n)$.  More precisely,
$\R^n\setminus f(S^{n-1})$ has exactly two connected components (for
$n\ge 2$), and the one containing $f(x)$ lies in~$f(D^n)$.
\end{lemma}

\begin{proof}[Proof sketch]
Suppose $y\notin f(D^n)$.  Define
$g\colon S^{n-1}\to S^{n-1}$ by
$g(z)=\frac{f(z)-y}{\abs{f(z)-y}}$.
The map~$g$ extends to $D^n\to S^{n-1}$ via
$\hat{g}(z)=\frac{f(z)-y}{\abs{f(z)-y}}$ (well-defined since $y\notin f(D^n)$).
Thus $g=\hat{g}\circ i$ factors through the contractible space~$D^n$,
so $\deg g=0$.  On the other hand, for $y=f(x_0)$ with
$x_0\in\operatorname{int}(D^n)$, one can show $\deg g\ne 0$ by a
deformation argument (deforming~$f$ to the identity), contradicting the
above.  Hence $f(x_0)\notin\R^n\setminus f(D^n)$, and since $f$ is
injective, $f(x_0)\in f(\operatorname{int}(D^n))\subseteq\operatorname{int}(f(D^n))$.
\end{proof}

\begin{corollary}[Topological invariance of dimension]\label{cor:invariance-dimension}
\index{invariance of dimension}
If $m\ne n$, then $\R^m$ is not homeomorphic to~$\R^n$.
\end{corollary}

\begin{proof}
If $m<n$, embed $\R^m\hookrightarrow\R^n$ as $\R^m\times\{0\}$.  This image
is not open in~$\R^n$, contradicting invariance of domain if there
were a homeomorphism $\R^n\xrightarrow{\sim}\R^m\hookrightarrow\R^n$.

Alternatively, one can argue directly via homology: removing a point,
$\R^n\setminus\{0\}\simeq S^{n-1}$, so $H_{n-1}(\R^n\setminus\{0\})\cong\Z$
while $H_{m-1}(\R^m\setminus\{0\})\cong\Z$ for different values of
$m-1\ne n-1$.  A homeomorphism $\R^m\cong\R^n$ would restrict to a
homeomorphism $\R^m\setminus\{0\}\cong\R^n\setminus\{0\}$, yielding
$\Z\cong H_{n-1}(\R^n\setminus\{0\})\cong H_{n-1}(\R^m\setminus\{0\})=0$
for $n-1>m-1$, a contradiction.
\end{proof}

\begin{proposition}[Hairy ball theorem, homological version]\label{prop:hairy-ball}
\index{hairy ball theorem}
The sphere $S^{2n}$ admits no continuous nowhere-vanishing tangent
vector field.
\end{proposition}

\begin{proof}
A nowhere-vanishing tangent vector field on~$S^{2n}$ provides a homotopy
from the identity to the antipodal map (via $x\mapsto(\cos t)\,x+(\sin t)\,v(x)$
for the normalized field~$v$).
But $\deg(\Id)=1$ and $\deg(a)=(-1)^{2n+1}=-1$, contradicting homotopy
invariance of degree.
\end{proof}

%% ----------------------------------------------------------------
%%  6.9  The Mayer--Vietoris sequence
%% ----------------------------------------------------------------
\section{The Mayer--Vietoris sequence}\label{sec:mayer-vietoris}

As an important consequence of excision, we obtain a tool for computing
homology from decompositions.

\begin{theorem}[Mayer--Vietoris sequence]\label{thm:mayer-vietoris}
\index{Mayer--Vietoris sequence}
Let $X=A\cup B$ where $\operatorname{int}(A)\cup\operatorname{int}(B)=X$.
Let $i_A\colon A\cap B\hookrightarrow A$, $i_B\colon A\cap B\hookrightarrow B$,
$j_A\colon A\hookrightarrow X$, $j_B\colon B\hookrightarrow X$ be the
inclusions.  There is a long exact sequence
\[
\begin{tikzcd}[column sep=1.2em]
  \cdots \ar[r]
  & H_n(A\cap B) \ar[r,"\Phi"]
  & H_n(A)\oplus H_n(B) \ar[r,"\Psi"]
  & H_n(X) \ar[r,"\partial_*"]
  & H_{n-1}(A\cap B) \ar[r]
  & \cdots
\end{tikzcd}
\]
where $\Phi(\alpha)=\bigl((i_A)_*(\alpha),\,-(i_B)_*(\alpha)\bigr)$
and $\Psi(\alpha,\beta)=(j_A)_*(\alpha)+(j_B)_*(\beta)$.
\end{theorem}

\begin{proof}
Consider the short exact sequence of chain complexes
\[
\begin{tikzcd}
  0 \ar[r]
  & C_\bullet(A\cap B) \ar[r,"\Phi_\sharp"]
  & C_\bullet(A)\oplus C_\bullet(B) \ar[r,"\Psi_\sharp"]
  & C_\bullet^{\set{A,B}}(X) \ar[r]
  & 0,
\end{tikzcd}
\]
where $C_\bullet^{\set{A,B}}(X)=C_\bullet(A)+C_\bullet(B)\subseteq C_\bullet(X)$.
Exactness at each position is straightforward:
$\Phi_\sharp$ is injective,
$\ker\Psi_\sharp=\im\Phi_\sharp$ (a chain in both $A$ and $B$ lives in $A\cap B$),
and $\Psi_\sharp$ is surjective by definition.

By the proof of the excision theorem, the inclusion
$C_\bullet^{\set{A,B}}(X)\hookrightarrow C_\bullet(X)$ is a chain
homotopy equivalence, hence induces isomorphisms on homology.
Applying Theorem~\ref{thm:les-chain} to the short exact sequence above
and replacing $H_n(C_\bullet^{\set{A,B}}(X))$ with $H_n(X)$ yields the
Mayer--Vietoris sequence.
\end{proof}

\begin{example}[Homology of $S^n$ via Mayer--Vietoris]\label{ex:sphere-mv}
\index{sphere!Mayer--Vietoris computation}
Write $S^n=U_+\cup U_-$ where $U_+$ and $U_-$ are the upper and lower
open hemispheres slightly enlarged to overlap.  Then $U_+\simeq\operatorname{pt}\simeq U_-$
and $U_+\cap U_-\simeq S^{n-1}$.  The Mayer--Vietoris sequence in
reduced homology gives
\[
\begin{tikzcd}[column sep=1.1em]
  \cdots \ar[r]
  & \widetilde{H}_k(U_+)\oplus\widetilde{H}_k(U_-) \ar[r]
  & \widetilde{H}_k(S^n) \ar[r,"\partial_*"]
  & \widetilde{H}_{k-1}(S^{n-1}) \ar[r]
  & \widetilde{H}_{k-1}(U_+)\oplus\widetilde{H}_{k-1}(U_-) \ar[r]
  & \cdots
\end{tikzcd}
\]
Since $\widetilde{H}_*(\operatorname{pt})=0$, we get isomorphisms
$\widetilde{H}_k(S^n)\cong\widetilde{H}_{k-1}(S^{n-1})$ for all~$k$.
Induction gives $\widetilde{H}_n(S^n)\cong\Z$ and
$\widetilde{H}_k(S^n)=0$ for $k\ne n$, confirming
Theorem~\ref{thm:homology-spheres}.
\end{example}

\begin{example}[Homology of the torus]\label{ex:homology-torus}
\index{torus!homology of}
The torus $T^2=S^1\times S^1$ can be decomposed as a union of two
overlapping cylinders.  The Mayer--Vietoris sequence yields
\[
  H_k(T^2)\cong
  \begin{cases}
    \Z & \text{if } k=0 \text{ or } k=2,\\
    \Z^2 & \text{if } k=1,\\
    0 & \text{if } k\ge 3.
  \end{cases}
\]
We leave the details as Exercise~\ref{exo:mv-torus}.
\end{example}

\begin{example}[Homology of the Klein bottle]\label{ex:homology-klein}
\index{Klein bottle!homology}
The Klein bottle~$K$ can be decomposed as a union of two M\"obius bands
$A$ and~$B$, with $A\cap B\simeq S^1\sqcup S^1$.  A careful analysis
of the Mayer--Vietoris sequence (tracking the maps on $H_1$) yields
\[
  H_k(K)\cong
  \begin{cases}
    \Z & \text{if } k=0,\\
    \Z\oplus\Z/2\Z & \text{if } k=1,\\
    0 & \text{if } k\ge 2.
  \end{cases}
\]
The torsion element in $H_1(K)$ reflects the non-orientability of~$K$.
\end{example}

\begin{remark}[Mayer--Vietoris in reduced homology]\label{rem:mv-reduced}
\index{Mayer--Vietoris sequence!reduced}
There is also a Mayer--Vietoris sequence in reduced homology:
\[
\begin{tikzcd}[column sep=1.2em]
  \cdots \ar[r]
  & \widetilde{H}_n(A\cap B) \ar[r]
  & \widetilde{H}_n(A)\oplus\widetilde{H}_n(B) \ar[r]
  & \widetilde{H}_n(X) \ar[r,"\partial_*"]
  & \widetilde{H}_{n-1}(A\cap B) \ar[r]
  & \cdots
\end{tikzcd}
\]
This is obtained by applying the same construction to the augmented chain
complexes.  The reduced version is often more convenient since
$\widetilde{H}_0(\operatorname{pt})=0$.
\end{remark}

%% ----------------------------------------------------------------
%%  6.10  Exercises
%% ----------------------------------------------------------------
\section{Exercises}\label{sec:ex-ch6}

\begin{exercise}\label{exo:connecting-hom-check}
\index{connecting homomorphism!exercise}
Verify directly that the connecting homomorphism
$\partial_*\colon H_n(C_\bullet)\to H_{n-1}(A_\bullet)$
is a group homomorphism (i.e., $\partial_*([c]+[c'])=\partial_*[c]+\partial_*[c']$).
\end{exercise}

\begin{exercise}\label{exo:five-lemma}
\index{five lemma}
Prove the \textbf{Five Lemma}: given a commutative diagram of abelian groups
\[
\begin{tikzcd}
  A_1 \ar[r] \ar[d,"\alpha_1"] & A_2 \ar[r] \ar[d,"\alpha_2"]
  & A_3 \ar[r] \ar[d,"\alpha_3"] & A_4 \ar[r] \ar[d,"\alpha_4"]
  & A_5 \ar[d,"\alpha_5"] \\
  B_1 \ar[r] & B_2 \ar[r] & B_3 \ar[r] & B_4 \ar[r] & B_5
\end{tikzcd}
\]
with exact rows, if $\alpha_1,\alpha_2,\alpha_4,\alpha_5$ are isomorphisms,
then so is~$\alpha_3$.  Show that ``$\alpha_2,\alpha_4$ surjective and
$\alpha_1,\alpha_5$ injective'' suffices if we only want $\alpha_3$
surjective; state and prove the dual for injectivity.
\end{exercise}

\begin{exercise}\label{exo:excision-wedge}
\index{wedge sum!homology}
Let $X=A\vee B$ be the wedge sum (identifying base points).  Assuming
that the base points are ``good'' (deformation retracts of neighborhoods),
use excision to show that
$\widetilde{H}_n(A\vee B)\cong\widetilde{H}_n(A)\oplus\widetilde{H}_n(B)$
for all~$n$.
\end{exercise}

\begin{exercise}\label{exo:suspension-iso}
\index{suspension!homology}
Let $\Sigma X$ denote the (unreduced) suspension of~$X$.  Using the
Mayer--Vietoris sequence, prove that
$\widetilde{H}_n(\Sigma X)\cong\widetilde{H}_{n-1}(X)$
for all~$n$.  Deduce the homology of $S^n$ by induction.
\end{exercise}

\begin{exercise}\label{exo:mv-torus}
\index{torus!Mayer--Vietoris}
Compute the homology of the torus $T^2$ using the Mayer--Vietoris sequence.
Decompose $T^2$ as $A\cup B$ where $A$ and $B$ are each homeomorphic to
$S^1\times(0,1)$ and $A\cap B$ is the disjoint union of two copies
of $S^1\times(0,\tfrac{1}{2})$.  Carefully determine the maps $\Phi$ and $\Psi$
and extract $H_k(T^2)$ for all~$k$.
\end{exercise}

\begin{exercise}\label{exo:brouwer-generalized}
\index{Brouwer fixed point theorem!generalization}
Use the Brouwer fixed point theorem to prove: if $K\subseteq\R^n$ is
a compact convex set and $f\colon K\to K$ is continuous, then $f$ has a
fixed point.
\textit{Hint}: $K$ is homeomorphic to $D^n$ (or a lower-dimensional disk)
if it has non-empty interior.
\end{exercise}

\begin{exercise}\label{exo:degree-computation}
\index{degree!computation}
\begin{enumerate}[label=\textup{(\alph*)},leftmargin=2em]
  \item Show that the map $z\mapsto z^k\colon S^1\to S^1$ (identifying
    $S^1\subseteq\C$) has degree~$k$.
  \item More generally, show that if $f\colon S^n\to S^n$ has no fixed point,
    then $\deg f=(-1)^{n+1}$.
  \item Deduce that a map $f\colon S^{2n}\to S^{2n}$ with
    $\deg f\ne(-1)^{2n+1}=-1$ must have a fixed point.
\end{enumerate}
\end{exercise}

\begin{exercise}\label{exo:borsuk-ulam-dim2}
\index{Borsuk--Ulam theorem}
Prove the \textbf{Borsuk--Ulam theorem in dimension~$2$}: for every
continuous map $f\colon S^2\to\R^2$ there exists $x\in S^2$ with $f(x)=f(-x)$.

\textit{Hint}: suppose not, and consider $g(x)=\frac{f(x)-f(-x)}{\abs{f(x)-f(-x)}}
\colon S^2\to S^1$.  Show $g(-x)=-g(x)$, and use degree theory
to derive a contradiction.
\end{exercise}

\begin{exercise}\label{exo:les-triple}
\index{long exact sequence!of a triple}
Let $A\subseteq B\subseteq X$.  Establish the \textbf{long exact sequence
of the triple}:
\[
\begin{tikzcd}[column sep=1.2em]
  \cdots \ar[r]
  & H_n(B,A) \ar[r]
  & H_n(X,A) \ar[r]
  & H_n(X,B) \ar[r,"\partial_*"]
  & H_{n-1}(B,A) \ar[r]
  & \cdots
\end{tikzcd}
\]
\textit{Hint}: apply the long exact sequence theorem to the short exact
sequence $0\to C_\bullet(B,A)\to C_\bullet(X,A)\to C_\bullet(X,B)\to 0$.
\end{exercise}

\begin{exercise}\label{exo:local-homology}
\index{local homology}
For $x\in X$, the \textbf{local homology groups} are defined as
$H_n(X,X\setminus\{x\})$.  Show that local homology is a topological
invariant: if $f\colon X\to Y$ is a homeomorphism, then
$H_n(X,X\setminus\{x\})\cong H_n(Y,Y\setminus\{f(x)\})$.
Compute $H_n(\R^m,\R^m\setminus\{0\})$ and use it to give another proof
that $\R^m\not\cong\R^n$ for $m\ne n$.
\end{exercise}

\index{long exact sequence|)}
\index{excision|)}

% === END OF CHUNK ch5-6 ===
% === START OF CHUNK ch7-8 ===

%%%%%%%%%%%%%%%%%%%%%%%%%%%%%%%%%%%%%%%%%%%%%%%%%%%%%%%%%%%%%%%%%%%%
%%  CHAPTER 7 — THE MAYER-VIETORIS THEOREM
%%%%%%%%%%%%%%%%%%%%%%%%%%%%%%%%%%%%%%%%%%%%%%%%%%%%%%%%%%%%%%%%%%%%
\chapter{The Mayer--Vietoris Theorem}\label{ch:mayer-vietoris}
\index{Mayer--Vietoris theorem|(}

The Mayer--Vietoris sequence is one of the most powerful computational
tools in algebraic topology. It allows us to compute the homology of
a space $X$ that is written as the union of two subspaces $A$ and $B$,
provided we know the homology of $A$, $B$, and their intersection
$A\cap B$. In spirit it is the homological analogue of the
Seifert--van Kampen theorem for the fundamental group.

%% ──────────────────────────────────────────────────────────────────
\section{Statement of the Mayer--Vietoris sequence}
\label{sec:mv-statement}
\index{Mayer--Vietoris theorem!statement}

Throughout this chapter, all homology groups are taken with
coefficients in an abelian group $G$ (typically $G=\Z$).

\begin{theorem}[Mayer--Vietoris exact sequence]\label{thm:mv-exact}
\index{Mayer--Vietoris theorem!exact sequence}
Let $X$ be a topological space and let $A,B\subseteq X$ be subspaces
whose interiors cover $X$, i.e.\ $X=\mathring{A}\cup\mathring{B}$.
Let $i_A\colon A\cap B\hookrightarrow A$,
$i_B\colon A\cap B\hookrightarrow B$,
$j_A\colon A\hookrightarrow X$, and $j_B\colon B\hookrightarrow X$
denote the inclusion maps. Then there is a long exact sequence
\[
\cdots \xrightarrow{\partial_{n+1}}
H_n(A\cap B)
\xrightarrow{\;\Phi_n\;}
H_n(A)\oplus H_n(B)
\xrightarrow{\;\Psi_n\;}
H_n(X)
\xrightarrow{\;\partial_n\;}
H_{n-1}(A\cap B)
\xrightarrow{\;\Phi_{n-1}\;}
\cdots
\]
where
\begin{itemize}
  \item $\Phi_n(\alpha) = \bigl((i_A)_*(\alpha),\;-(i_B)_*(\alpha)\bigr)$,
  \item $\Psi_n(\alpha,\beta) = (j_A)_*(\alpha) + (j_B)_*(\beta)$,
  \item $\partial_n$ is the connecting homomorphism.
\end{itemize}
The sequence terminates at the right:
\[
\cdots \to H_0(A\cap B) \xrightarrow{\Phi_0}
H_0(A)\oplus H_0(B) \xrightarrow{\Psi_0} H_0(X) \to 0.
\]
\end{theorem}

\begin{remark}[The interior condition]
\index{Mayer--Vietoris theorem!interior condition}
The condition $X = \mathring{A}\cup\mathring{B}$ is essential. It is
automatically satisfied when $A$ and $B$ are open subsets of $X$, or
when $X$ is a CW complex and $A$, $B$ are subcomplexes with $A\cup B=X$.
More generally, it suffices that the inclusion
$C_\bullet(A)+C_\bullet(B)\hookrightarrow C_\bullet(X)$ induces
an isomorphism on homology, which is guaranteed by the excision theorem.
\end{remark}

The Mayer--Vietoris sequence is often displayed as a diagram.
The following commutative diagram illustrates the core short exact
sequence of chain complexes from which it arises:
\[
\begin{tikzcd}[column sep=large]
0 \ar[r]
  & C_\bullet(A\cap B) \ar[r,"\Phi"]
  & C_\bullet(A)\oplus C_\bullet(B) \ar[r,"\Psi"]
  & C_\bullet(A)+C_\bullet(B) \ar[r]
  & 0
\end{tikzcd}
\]

%% ──────────────────────────────────────────────────────────────────
\section{Proof from excision}
\label{sec:mv-proof}
\index{Mayer--Vietoris theorem!proof}
\index{excision theorem}

We now give the full proof of Theorem~\ref{thm:mv-exact}, relying
on the excision theorem and the long exact sequence of a pair.

\begin{proof}[Proof of Theorem~\ref{thm:mv-exact}]
\textbf{Step 1: The short exact sequence of chain complexes.}
Define chain maps
\begin{align*}
  \Phi_\bullet &\colon C_\bullet(A\cap B) \longrightarrow
    C_\bullet(A)\oplus C_\bullet(B),
    &\sigma &\longmapsto \bigl((i_A)_\sharp(\sigma),\;-(i_B)_\sharp(\sigma)\bigr),\\
  \Psi_\bullet &\colon C_\bullet(A)\oplus C_\bullet(B) \longrightarrow
    C_\bullet(A)+C_\bullet(B),
    &(\sigma,\tau) &\longmapsto (j_A)_\sharp(\sigma)+(j_B)_\sharp(\tau).
\end{align*}
Here $C_\bullet(A)+C_\bullet(B)$ denotes the subcomplex of $C_\bullet(X)$
generated by singular simplices whose image lies in $A$ or in $B$.

\emph{$\Phi_\bullet$ is injective.}
If $\Phi_\bullet(\sigma)=(0,0)$, then $(i_A)_\sharp(\sigma)=0$ and
$(i_B)_\sharp(\sigma)=0$. Since $i_A$ and $i_B$ are inclusions,
$\sigma=0$.

\emph{$\Psi_\bullet$ is surjective.}
Every generator of $C_n(A)+C_n(B)$ is a singular simplex with image in
$A$ or in $B$. If $\sigma\colon\Delta^n\to A$, then
$\Psi_n(\sigma,0)=\sigma$; similarly for simplices in $B$.

\emph{$\im\Phi_\bullet=\ker\Psi_\bullet$.}
We have $\Psi_\bullet\circ\Phi_\bullet(\sigma)
= (j_A)_\sharp\circ(i_A)_\sharp(\sigma) - (j_B)_\sharp\circ(i_B)_\sharp(\sigma) = \sigma-\sigma = 0$, so $\im\Phi_\bullet\subseteq\ker\Psi_\bullet$.
Conversely, if $\Psi_n(\sigma,\tau)=0$ in $C_n(A)+C_n(B)$, then
$\sigma=-\tau$ in $C_n(X)$, so every simplex appearing in $\sigma$
has image in both $A$ and $B$, hence in $A\cap B$. Writing
$\gamma$ for this chain viewed in $C_n(A\cap B)$, we get
$\Phi_n(\gamma)=(\sigma,-\tau)=(\sigma,\sigma)$... but the sign gives
$(\sigma,-(-\sigma))$; more carefully, $\sigma = (i_A)_\sharp(\gamma)$
and $-\tau=(i_B)_\sharp(\gamma)$ shows $(\sigma,\tau)=\Phi_n(\gamma)$.

Thus we have a short exact sequence of chain complexes
\[
0 \longrightarrow C_\bullet(A\cap B) \xrightarrow{\;\Phi_\bullet\;}
C_\bullet(A)\oplus C_\bullet(B) \xrightarrow{\;\Psi_\bullet\;}
C_\bullet(A)+C_\bullet(B) \longrightarrow 0.
\]

\textbf{Step 2: Applying the snake lemma.}
By the fundamental theorem on short exact sequences of chain complexes
(the zig-zag lemma), the short exact sequence above induces a long
exact sequence in homology:
\[
\cdots \to H_n(A\cap B) \xrightarrow{\Phi_*}
H_n(A)\oplus H_n(B) \xrightarrow{\Psi_*}
H_n\bigl(C_\bullet(A)+C_\bullet(B)\bigr) \xrightarrow{\partial}
H_{n-1}(A\cap B) \to \cdots
\]

\textbf{Step 3: Identifying the middle term via excision.}
\index{small simplices}
The inclusion $C_\bullet(A)+C_\bullet(B)\hookrightarrow C_\bullet(X)$
induces an isomorphism on homology. This is precisely the content of
the theorem on small simplices: if $\mathcal{U}=\{A,B\}$ is an open
cover of $X$ (or more generally if $X=\mathring{A}\cup\mathring{B}$),
then the inclusion $C_\bullet^{\mathcal{U}}(X)\hookrightarrow C_\bullet(X)$
is a chain homotopy equivalence, where $C_\bullet^{\mathcal{U}}(X)$
consists of chains whose simplices are $\mathcal{U}$-small. Since
$C_\bullet^{\mathcal{U}}(X) = C_\bullet(A)+C_\bullet(B)$, we obtain
\[
  H_n\bigl(C_\bullet(A)+C_\bullet(B)\bigr) \cong H_n(X)
  \quad\text{for all } n\geq 0.
\]

Substituting into the long exact sequence from Step~2 yields the
Mayer--Vietoris exact sequence as stated.

\textbf{Step 4: The connecting homomorphism.}
\index{connecting homomorphism!Mayer--Vietoris}
We describe $\partial_n\colon H_n(X)\to H_{n-1}(A\cap B)$ explicitly.
Let $[z]\in H_n(X)$. Using the isomorphism from Step~3, represent $z$
by a cycle $z=\alpha+\beta$ where $\alpha\in C_n(A)$ and $\beta\in C_n(B)$.
Then $\partial\alpha+\partial\beta=\partial z=0$, so
$\partial\alpha = -\partial\beta$. Since $\partial\alpha\in C_{n-1}(A)$
and $-\partial\beta\in C_{n-1}(B)$, this common chain lies in
$C_{n-1}(A\cap B)$. The connecting homomorphism is
$\partial_n[z] = [\partial\alpha] = [-\partial\beta] \in H_{n-1}(A\cap B)$.
\end{proof}

%% ──────────────────────────────────────────────────────────────────
\section{The reduced Mayer--Vietoris sequence}
\label{sec:mv-reduced}
\index{Mayer--Vietoris theorem!reduced version}

For computations it is often convenient to use reduced homology.

\begin{theorem}[Reduced Mayer--Vietoris sequence]\label{thm:mv-reduced}
Under the same hypotheses as Theorem~\ref{thm:mv-exact}, and assuming
$A\cap B\neq\varnothing$, there is a long exact sequence of reduced
homology groups:
\[
\cdots \to \tilde{H}_n(A\cap B) \xrightarrow{\Phi_n}
\tilde{H}_n(A)\oplus\tilde{H}_n(B) \xrightarrow{\Psi_n}
\tilde{H}_n(X) \xrightarrow{\partial_n}
\tilde{H}_{n-1}(A\cap B) \to \cdots
\to \tilde{H}_0(X) \to 0.
\]
\end{theorem}

\begin{proof}
The augmentation maps $\varepsilon\colon C_0(\,\cdot\,)\to\Z$ are
compatible with the inclusions, yielding a short exact sequence of
augmented chain complexes. The argument proceeds identically to the
unreduced case. The nonemptiness of $A\cap B$ ensures that all
augmentation maps are surjective, so the short exact sequence at the
augmented level remains exact.
\end{proof}

\begin{remark}
When $A\cap B=\varnothing$, the space $X$ is disconnected as
$X=A\sqcup B$ (with $A$ and $B$ open), and the Mayer--Vietoris
sequence degenerates:
$\tilde{H}_n(X)\cong\tilde{H}_n(A)\oplus\tilde{H}_n(B)$ for all $n$.
\end{remark}

%% ──────────────────────────────────────────────────────────────────
\section{Homology of spheres by induction}
\label{sec:mv-spheres}
\index{sphere!homology of}
\index{Mayer--Vietoris theorem!application to spheres}

\begin{theorem}[Homology of spheres]\label{thm:homology-spheres}
For $n\geq 1$,
\[
  \tilde{H}_k(S^n) \cong
  \begin{cases}
    \Z & \text{if } k=n,\\
    0  & \text{if } k\neq n.
  \end{cases}
\]
\end{theorem}

\begin{proof}
We proceed by induction on $n$.

\textbf{Base case $n=1$.}
Decompose $S^1$ as $A=S^1\setminus\{N\}$ and $B=S^1\setminus\{S\}$,
where $N=(0,1)$ and $S=(0,-1)$ are the north and south poles.
Then $A\simeq *\simeq B$ (both are contractible arcs), and
$A\cap B$ is homotopy equivalent to a two-point discrete space
$\{p,q\}$, so $\tilde{H}_0(A\cap B)\cong\Z$.
The reduced Mayer--Vietoris sequence gives
\[
  \cdots\to
  \underbrace{\tilde{H}_1(A)\oplus\tilde{H}_1(B)}_{=\,0}
  \to \tilde{H}_1(S^1) \xrightarrow{\partial}
  \underbrace{\tilde{H}_0(A\cap B)}_{\cong\,\Z}
  \xrightarrow{\Phi_0}
  \underbrace{\tilde{H}_0(A)\oplus\tilde{H}_0(B)}_{=\,0}
  \to \cdots
\]
Exactness forces $\partial\colon\tilde{H}_1(S^1)\xrightarrow{\;\sim\;}\Z$.
For $k\geq 2$ the sequence gives $\tilde{H}_k(S^1)=0$ since
$\tilde{H}_k(A)=\tilde{H}_k(B)=\tilde{H}_{k-1}(A\cap B)=0$.

\textbf{Inductive step.}
Assume the result for $S^{n-1}$, $n\geq 2$. Write $S^n\subseteq\R^{n+1}$
and let
\[
  A = S^n\setminus\{N\},\qquad B = S^n\setminus\{S\},
\]
where $N=(0,\ldots,0,1)$ and $S=(0,\ldots,0,-1)$.
Both $A$ and $B$ are contractible (via stereographic projection onto
$\R^n$), and $A\cap B = S^n\setminus\{N,S\}\simeq S^{n-1}$
(deformation retraction onto the equatorial sphere).

The reduced Mayer--Vietoris sequence reads
\[
  \tilde{H}_k(A)\oplus\tilde{H}_k(B) \to
  \tilde{H}_k(S^n) \xrightarrow{\partial}
  \tilde{H}_{k-1}(A\cap B) \to
  \tilde{H}_{k-1}(A)\oplus\tilde{H}_{k-1}(B).
\]
Since $\tilde{H}_*(A)=\tilde{H}_*(B)=0$, exactness gives isomorphisms
\[
  \tilde{H}_k(S^n) \xrightarrow[\cong]{\;\partial\;}
  \tilde{H}_{k-1}(S^{n-1})
  \quad\text{for all } k.
\]
By the induction hypothesis,
$\tilde{H}_{k-1}(S^{n-1})\cong\Z$ if $k-1=n-1$ (i.e.\ $k=n$) and
vanishes otherwise. The result follows.
\end{proof}

%% ──────────────────────────────────────────────────────────────────
\section{Homology of surfaces}
\label{sec:mv-surfaces}
\index{Mayer--Vietoris theorem!applications to surfaces}

\subsection{The torus}\label{subsec:mv-torus}
\index{torus!homology of}

\begin{example}[Homology of the torus $T^2$]
Recall that $T^2 = S^1\times S^1$. We compute its homology using
Mayer--Vietoris. Think of $T^2$ as a square with opposite sides
identified. Cut the torus along a longitudinal circle to obtain:
\begin{itemize}
  \item $A$ = an open neighborhood of the ``left'' cylinder,
  \item $B$ = an open neighborhood of the ``right'' cylinder,
\end{itemize}
so that $A\simeq S^1\simeq B$ (each piece deformation retracts onto
a circle), and $A\cap B\simeq S^1\sqcup S^1$.

The reduced Mayer--Vietoris sequence in low degrees gives:
\[
\begin{tikzcd}[column sep=small]
0 \ar[r]
  & \tilde{H}_2(T^2) \ar[r,"\partial"]
  & \tilde{H}_1(A\cap B) \ar[r,"\Phi"]
  & \tilde{H}_1(A)\oplus\tilde{H}_1(B) \ar[r,"\Psi"]
  & \tilde{H}_1(T^2) \ar[r,"\partial"]
  & \tilde{H}_0(A\cap B) \ar[r,"\Phi"]
  & \tilde{H}_0(A)\oplus\tilde{H}_0(B)
\end{tikzcd}
\]
Substituting known values:
\[
\begin{tikzcd}[column sep=small]
0 \ar[r]
  & \tilde{H}_2(T^2) \ar[r,"\partial"]
  & \Z\oplus\Z \ar[r,"\Phi"]
  & \Z\oplus\Z \ar[r,"\Psi"]
  & \tilde{H}_1(T^2) \ar[r,"\partial"]
  & \Z \ar[r,"\Phi"]
  & 0
\end{tikzcd}
\]

The map $\Phi_0\colon\tilde{H}_0(A\cap B)\cong\Z\to\tilde{H}_0(A)\oplus
\tilde{H}_0(B)=0$ is the zero map, so exactness gives
$\partial\colon\tilde{H}_1(T^2)\twoheadrightarrow\Z$.

For $\Phi_1\colon\Z^2\to\Z^2$, one computes that it sends the
generators of the two circle components of $A\cap B$ to the
same generator of $\tilde{H}_1(A)$ and $\tilde{H}_1(B)$ respectively.
An analysis of the inclusion maps shows $\Phi_1$ has matrix
$\begin{pmatrix}1&1\\-1&-1\end{pmatrix}$, so $\ker\Phi_1\cong\Z$
and $\im\Phi_1\cong\Z$.

From the exact sequence:
\begin{itemize}
  \item $\tilde{H}_2(T^2)\cong\ker\Phi_1\cong\Z$,
  \item $\tilde{H}_1(T^2)\cong(\Z^2/\im\Phi_1)\oplus\ker(\Phi_0)
         = \Z\oplus\Z$,
  \item $\tilde{H}_k(T^2)=0$ for $k\geq 3$.
\end{itemize}

In summary:
\[
  H_k(T^2) \cong
  \begin{cases}
    \Z   & k=0,2,\\
    \Z^2 & k=1,\\
    0    & k\geq 3.
  \end{cases}
\]
\end{example}

\subsection{Orientable surfaces of genus $g$}
\label{subsec:mv-surfaces-g}
\index{surface!orientable!homology of}
\index{genus}

\begin{theorem}[Homology of $\Sigma_g$]\label{thm:homology-sigma-g}
Let $\Sigma_g$ denote the closed orientable surface of genus $g\geq 0$.
Then
\[
  H_k(\Sigma_g) \cong
  \begin{cases}
    \Z     & k=0,2,\\
    \Z^{2g} & k=1,\\
    0      & k\geq 3.
  \end{cases}
\]
\end{theorem}

\begin{proof}
We use induction on $g$. The case $g=0$ is $\Sigma_0=S^2$, handled
by Theorem~\ref{thm:homology-spheres}. The case $g=1$ is the torus,
computed above. For $g\geq 2$, decompose $\Sigma_g$ along a separating
simple closed curve into a punctured $\Sigma_{g-1}$ (call it $A$) and
a punctured torus (call it $B$), with $A\cap B\simeq S^1$.
Both $A$ and $B$ deformation retract onto wedges of circles:
$A\simeq\bigvee^{2(g-1)} S^1$ and $B\simeq S^1\vee S^1$.
Applying Mayer--Vietoris and chasing the maps yields the result.
Alternatively, one can use the CW structure and cellular homology
to give a direct proof.
\end{proof}

\subsection{The real projective plane}\label{subsec:mv-rp2}
\index{real projective plane!homology of}

\begin{example}[Homology of $\RP^2$]
Decompose $\RP^2$ into a M\"obius band $A$ and a disk $B$, glued
along their boundaries. We have:
\begin{itemize}
  \item $A$ deformation retracts onto $S^1$, so
    $\tilde{H}_1(A)\cong\Z$ and $\tilde{H}_k(A)=0$ for $k\neq 1$.
  \item $B$ is contractible: $\tilde{H}_k(B)=0$ for all $k$.
  \item $A\cap B$ is an annulus, hence $A\cap B\simeq S^1$, so
    $\tilde{H}_1(A\cap B)\cong\Z$ and $\tilde{H}_0(A\cap B)=0$.
\end{itemize}

The relevant piece of the reduced Mayer--Vietoris sequence is:
\[
\begin{tikzcd}[column sep=small]
0 \ar[r]
  & \tilde{H}_2(\RP^2) \ar[r,"\partial"]
  & \Z \ar[r,"\Phi_1"]
  & \Z\oplus 0 \ar[r,"\Psi_1"]
  & \tilde{H}_1(\RP^2) \ar[r]
  & 0
\end{tikzcd}
\]
The inclusion of the boundary circle of the M\"obius band into the
M\"obius band wraps around twice (the boundary of a M\"obius band
goes around the core circle twice). Thus $\Phi_1$ is multiplication
by $2$. Therefore:
\begin{itemize}
  \item $\tilde{H}_2(\RP^2)=\ker(\times 2)=0$,
  \item $\tilde{H}_1(\RP^2)=\Z/2\Z$,
  \item $\tilde{H}_k(\RP^2)=0$ for $k\geq 3$.
\end{itemize}

In summary:
\[
  H_k(\RP^2) \cong
  \begin{cases}
    \Z     & k=0,\\
    \Z/2\Z & k=1,\\
    0      & k\geq 2.
  \end{cases}
\]
\end{example}

\subsection{The Klein bottle}\label{subsec:mv-klein}
\index{Klein bottle!homology of}

\begin{example}[Homology of the Klein bottle $K$]
The Klein bottle can be decomposed as the union of two M\"obius
bands $A$ and $B$ glued along their common boundary circle, so
$A\cap B\simeq S^1$. Each M\"obius band deformation retracts onto
$S^1$, so $\tilde{H}_1(A)\cong\Z\cong\tilde{H}_1(B)$.

The boundary circle of a M\"obius band wraps twice around the core
circle, so the inclusion-induced maps
$\tilde{H}_1(A\cap B)\to\tilde{H}_1(A)$ and
$\tilde{H}_1(A\cap B)\to\tilde{H}_1(B)$ are both multiplication by
$2$. The map $\Phi_1\colon\Z\to\Z\oplus\Z$ is therefore
$\Phi_1(1)=(2,-2)$.

The reduced Mayer--Vietoris sequence gives:
\[
\begin{tikzcd}[column sep=small]
0 \ar[r]
  & \tilde{H}_2(K) \ar[r,"\partial"]
  & \Z \ar[r,"{\Phi_1}"]
  & \Z\oplus\Z \ar[r,"\Psi_1"]
  & \tilde{H}_1(K) \ar[r]
  & 0
\end{tikzcd}
\]

Since $\Phi_1$ is injective ($\ker\Phi_1=0$), we get
$\tilde{H}_2(K)=0$. For $\tilde{H}_1(K)$, we compute
$\coker\Phi_1 = (\Z\oplus\Z)/\langle(2,-2)\rangle \cong \Z\oplus\Z/2\Z$.
(One can verify this by the change of basis $(a,b)\mapsto(a+b,a-b)$.)
Thus:
\[
  H_k(K) \cong
  \begin{cases}
    \Z           & k=0,\\
    \Z\oplus\Z/2\Z & k=1,\\
    0            & k\geq 2.
  \end{cases}
\]
Note that $H_2(K)=0$, reflecting the fact that the Klein bottle is
non-orientable.
\end{example}

%% ──────────────────────────────────────────────────────────────────
\section{Homology of wedge sums}
\label{sec:mv-wedge}
\index{wedge sum!homology of}
\index{Mayer--Vietoris theorem!application to wedge sums}

The Mayer--Vietoris sequence provides a quick proof of the following
useful result.

\begin{proposition}[Homology of a wedge sum]\label{prop:homology-wedge}
Let $X$ and $Y$ be path-connected spaces with basepoints $x_0\in X$
and $y_0\in Y$. Then for all $n\geq 1$,
\[
  \tilde{H}_n(X\vee Y) \cong \tilde{H}_n(X)\oplus\tilde{H}_n(Y).
\]
\end{proposition}

\begin{proof}
Choose open neighborhoods $U$ of $x_0$ in $X$ and $V$ of $y_0$ in $Y$,
both contractible. Let $A = X\cup V\subseteq X\vee Y$ and
$B = U\cup Y\subseteq X\vee Y$ (where we identify the basepoints).
Then $A\simeq X$, $B\simeq Y$, and $A\cap B = U\cup V\simeq *$
(contractible). The reduced Mayer--Vietoris sequence yields
\[
  \cdots \to \underbrace{\tilde{H}_n(A\cap B)}_{=\,0}
  \to \tilde{H}_n(A)\oplus\tilde{H}_n(B)
  \xrightarrow{\Psi} \tilde{H}_n(X\vee Y)
  \xrightarrow{\partial}
  \underbrace{\tilde{H}_{n-1}(A\cap B)}_{=\,0} \to \cdots
\]
so $\Psi$ is an isomorphism, giving
$\tilde{H}_n(X\vee Y)\cong\tilde{H}_n(X)\oplus\tilde{H}_n(Y)$.
\end{proof}

\begin{corollary}\label{cor:wedge-finite}
For a finite wedge sum $\bigvee_{i=1}^k X_i$ of path-connected spaces,
\[
  \tilde{H}_n\Bigl(\bigvee_{i=1}^k X_i\Bigr)
  \cong \bigoplus_{i=1}^k \tilde{H}_n(X_i)
  \quad\text{for all } n\geq 1.
\]
In particular, $\tilde{H}_n\bigl(\bigvee^k S^m\bigr)\cong\Z^k$
for $n=m$ and vanishes for $n\neq m$.
\end{corollary}

%% ──────────────────────────────────────────────────────────────────
\section{The K\"unneth formula (simplified)}
\label{sec:mv-kunneth}
\index{K\"unneth formula}

The Mayer--Vietoris sequence, combined with induction, leads to the
computation of the homology of product spaces. The K\"unneth formula
expresses the homology of a product $X\times Y$ in terms of the
homology of the factors.

\begin{theorem}[K\"unneth formula, free case]\label{thm:kunneth-free}
Let $X$ and $Y$ be topological spaces. If $H_n(Y;\Z)$ is a free abelian
group for all $n$, then there is a natural isomorphism
\[
  H_n(X\times Y;\Z) \cong \bigoplus_{p+q=n}
  H_p(X;\Z)\otimes_\Z H_q(Y;\Z).
\]
\end{theorem}

The isomorphism is realized by the \emph{cross product} (or exterior
product)\index{cross product!homology}:
\[
  \times\colon H_p(X;\Z)\otimes H_q(Y;\Z)\longrightarrow H_{p+q}(X\times Y;\Z),
  \qquad [\alpha]\otimes[\beta]\longmapsto[\alpha\times\beta],
\]
constructed via the Eilenberg--Zilber map\index{Eilenberg--Zilber map}.

\begin{example}[Homology of $T^2$ via K\"unneth]
Using the K\"unneth formula, the homology of $T^2=S^1\times S^1$ is
\begin{align*}
  H_0(T^2) &\cong H_0(S^1)\otimes H_0(S^1) \cong \Z,\\
  H_1(T^2) &\cong \bigl(H_0(S^1)\otimes H_1(S^1)\bigr)
             \oplus\bigl(H_1(S^1)\otimes H_0(S^1)\bigr)
             \cong \Z\oplus\Z,\\
  H_2(T^2) &\cong H_1(S^1)\otimes H_1(S^1) \cong \Z,
\end{align*}
agreeing with our Mayer--Vietoris computation.
\end{example}

\begin{example}[Homology of $S^m\times S^n$]
For $m,n\geq 1$ with $m\neq n$, the K\"unneth formula gives
\[
  H_k(S^m\times S^n;\Z) \cong
  \begin{cases}
    \Z & k=0,\,m,\,n,\,m+n,\\
    0  & \text{otherwise.}
  \end{cases}
\]
When $m=n$, $H_m(S^m\times S^m;\Z)\cong\Z^2$.
This shows, for example, that $S^m\times S^n\not\simeq S^{m+n}$ for
$m,n\geq 1$, since their homology groups differ.
\end{example}

\begin{remark}[The general K\"unneth formula]
When the homology groups have torsion, the full K\"unneth formula
includes a $\operatorname{Tor}$ correction term:
\[
  0 \to \bigoplus_{p+q=n} H_p(X)\otimes H_q(Y) \to H_n(X\times Y)
    \to \bigoplus_{p+q=n-1}\operatorname{Tor}_1^\Z(H_p(X),H_q(Y)) \to 0.
\]
This short exact sequence splits (non-naturally). Key properties of
$\operatorname{Tor}$\index{Tor functor} that are useful in computations:
\begin{enumerate}[label=(\roman*)]
  \item $\operatorname{Tor}_1^\Z(\Z/m\Z,\Z/n\Z)\cong\Z/\gcd(m,n)\Z$.
  \item $\operatorname{Tor}_1^\Z(A,B)=0$ if $A$ or $B$ is free.
  \item $\operatorname{Tor}$ commutes with finite direct sums.
\end{enumerate}
\end{remark}

\begin{example}[Homology of $\RP^2\times\RP^2$]
\index{real projective plane!product}
Since $H_0(\RP^2)=\Z$, $H_1(\RP^2)=\Z/2\Z$, and $H_k(\RP^2)=0$
for $k\geq 2$, the K\"unneth formula gives for $X=\RP^2\times\RP^2$:
\begin{align*}
  H_0(X) &\cong \Z\otimes\Z = \Z,\\
  H_1(X) &\cong (\Z\otimes\Z/2\Z)\oplus(\Z/2\Z\otimes\Z) = (\Z/2\Z)^2,\\
  H_2(X) &\cong (\Z/2\Z\otimes\Z/2\Z)\oplus\operatorname{Tor}_1(\Z/2\Z,\Z)
           \oplus\operatorname{Tor}_1(\Z,\Z/2\Z)\\
         &= \Z/2\Z\oplus 0\oplus 0 = \Z/2\Z,\\
  H_3(X) &\cong \operatorname{Tor}_1(\Z/2\Z,\Z/2\Z) = \Z/2\Z,\\
  H_k(X) &= 0 \quad\text{for } k\geq 4.
\end{align*}
\end{example}

%% ──────────────────────────────────────────────────────────────────
\section{The Euler characteristic}
\label{sec:euler-char}
\index{Euler characteristic}

\begin{definition}[Euler characteristic]
Let $X$ be a topological space with finitely generated homology groups,
all but finitely many of which vanish. The \emph{Euler characteristic}
of $X$ is
\[
  \chi(X) = \sum_{n\geq 0} (-1)^n \operatorname{rank} H_n(X;\Z).
\]
\end{definition}

\begin{proposition}[Additivity of $\chi$ under Mayer--Vietoris]
\label{prop:euler-mv}
\index{Euler characteristic!additivity}
If $X=A\cup B$ with $A\cap B\neq\varnothing$ and all spaces have
well-defined Euler characteristics, then
\[
  \chi(X) = \chi(A) + \chi(B) - \chi(A\cap B).
\]
\end{proposition}

\begin{proof}
This follows from the exactness of the Mayer--Vietoris sequence and
the general fact that in any long exact sequence of finitely generated
abelian groups
$\cdots\to A_n\to B_n\to C_n\to A_{n-1}\to\cdots$,
the alternating sum of ranks vanishes.
\end{proof}

\begin{example}[Euler characteristics of familiar spaces]
\begin{enumerate}[label=(\roman*)]
  \item $\chi(S^n)=1+(-1)^n$, so $\chi(S^{2k})=2$ and $\chi(S^{2k+1})=0$.
  \item $\chi(T^2)=1-2+1=0$.
  \item $\chi(\Sigma_g)=1-2g+1=2-2g$.
  \item $\chi(\RP^2)=1-0+0=1$ (using ranks of the free parts only).
  \item $\chi(K)=1-1+0=0$ for the Klein bottle $K$.
\end{enumerate}
\end{example}

\begin{proposition}[Multiplicativity of $\chi$]\label{prop:euler-product}
\index{Euler characteristic!multiplicativity}
If $X$ and $Y$ have well-defined Euler characteristics, then
\[
  \chi(X\times Y) = \chi(X)\cdot\chi(Y).
\]
\end{proposition}

\begin{proof}
This follows from the K\"unneth formula. If $h_k=\operatorname{rank}H_k(X)$
and $h'_k=\operatorname{rank}H_k(Y)$, then
\[
  \chi(X\times Y)
  = \sum_n(-1)^n\sum_{p+q=n}h_p h'_q
  = \biggl(\sum_p(-1)^p h_p\biggr)\biggl(\sum_q(-1)^q h'_q\biggr)
  = \chi(X)\cdot\chi(Y).\qedhere
\]
\end{proof}

\begin{example}
$\chi(T^n) = \chi(S^1)^n = 0^n = 0$ for all $n\geq 1$.
Also $\chi(S^2\times S^2) = \chi(S^2)^2 = 4$.
\end{example}

%% ──────────────────────────────────────────────────────────────────
\section{Summary of homology computations}
\label{sec:mv-summary}

We collect the results of this chapter in a single table for reference.

\medskip
\begin{center}
\renewcommand{\arraystretch}{1.3}
\begin{tabular}{lccccc}
\toprule
\textbf{Space} & $H_0$ & $H_1$ & $H_2$ & $H_3$ & $\chi$ \\
\midrule
$S^n$ ($n\geq 1$) & $\Z$ & $0$ ($n>1$) & $\cdots$ & $\Z$ in deg.\ $n$ & $1+(-1)^n$ \\
$T^2$ & $\Z$ & $\Z^2$ & $\Z$ & $0$ & $0$ \\
$\Sigma_g$ & $\Z$ & $\Z^{2g}$ & $\Z$ & $0$ & $2-2g$ \\
$\RP^2$ & $\Z$ & $\Z/2\Z$ & $0$ & $0$ & $1$ \\
Klein bottle & $\Z$ & $\Z\oplus\Z/2\Z$ & $0$ & $0$ & $0$ \\
$S^1\vee S^1$ & $\Z$ & $\Z^2$ & $0$ & $0$ & $-1$ \\
\bottomrule
\end{tabular}
\end{center}
\medskip

\begin{remark}
Observe that $H_2(\Sigma_g)\cong\Z$ for all orientable closed surfaces,
reflecting their orientability. The non-orientable surfaces $\RP^2$ and
the Klein bottle have $H_2=0$. More generally, a closed connected
$n$-manifold $M$ satisfies $H_n(M;\Z)\cong\Z$ if and only if $M$ is
orientable.
\index{orientability!and top homology}
\end{remark}

%% ──────────────────────────────────────────────────────────────────
\section{Exercises}
\label{sec:mv-exercises}

\begin{exercise}
Use the Mayer--Vietoris sequence to compute the homology groups of
$S^1\vee S^1$ (the wedge of two circles). \textit{Hint:} Take $A$
and $B$ to be open neighborhoods of each circle, with $A\cap B$
contractible.
\end{exercise}

\begin{exercise}
Compute the homology of $S^1\vee S^2$ using Mayer--Vietoris.
\end{exercise}

\begin{exercise}
Let $X = S^2\vee S^2$. Compute $H_*(X;\Z)$ using Mayer--Vietoris.
Then verify your answer using the general formula for the homology of
a wedge sum.
\end{exercise}

\begin{exercise}
Use the Mayer--Vietoris sequence to show that for any $n\geq 1$,
the sphere $S^n$ is not a retract of $D^{n+1}$. Deduce the
Brouwer fixed point theorem in all dimensions.
\index{Brouwer fixed point theorem}
\end{exercise}

\begin{exercise}
Let $X$ be the space obtained from $S^2$ by identifying the north
and south poles. Compute $H_*(X;\Z)$ by decomposing $X$ appropriately
and applying Mayer--Vietoris.
\end{exercise}

\begin{exercise}
Compute the homology of the genus-$2$ surface $\Sigma_2$ by
decomposing it as the connected sum
$\Sigma_2 = T^2\mathbin{\#} T^2$ and applying Mayer--Vietoris.
Verify that your result agrees with Theorem~\ref{thm:homology-sigma-g}.
\end{exercise}

\begin{exercise}
Let $M$ be the M\"obius band and let $\partial M\cong S^1$ be its
boundary. Compute the relative homology groups $H_*(M,\partial M;\Z)$
using the long exact sequence of the pair and your knowledge of
$H_*(M)$ and $H_*(\partial M)$.
\end{exercise}

\begin{exercise}
Prove that if $X=A\cup B$ with $A$, $B$, and $A\cap B$ all
path-connected, then the Mayer--Vietoris connecting homomorphism
$\partial_1\colon H_1(X)\to H_0(A\cap B)$ is trivial.
\end{exercise}

\begin{exercise}
Show that the $n$-fold torus $T^n = (S^1)^n$ has homology
\[
  H_k(T^n;\Z) \cong \Z^{\binom{n}{k}}.
\]
\textit{Hint:} Use the K\"unneth formula and induction on $n$.
\end{exercise}

\begin{exercise}
Let $X=\RP^2\times\RP^2$. Use the K\"unneth formula (with the
$\operatorname{Tor}$ term) to compute $H_*(X;\Z)$.
\index{real projective plane!product}
\end{exercise}

\index{Mayer--Vietoris theorem|)}

%%%%%%%%%%%%%%%%%%%%%%%%%%%%%%%%%%%%%%%%%%%%%%%%%%%%%%%%%%%%%%%%%%%%
%%  CHAPTER 8 — SINGULAR COHOMOLOGY AND THE CUP PRODUCT
%%%%%%%%%%%%%%%%%%%%%%%%%%%%%%%%%%%%%%%%%%%%%%%%%%%%%%%%%%%%%%%%%%%%
\chapter{Singular Cohomology and the Cup Product}\label{ch:cohomology}
\index{singular cohomology|(}
\index{cup product|(}

Cohomology is the ``dual'' of homology and carries richer algebraic
structure: the cup product endows the cohomology groups with the
structure of a graded ring. This additional structure often distinguishes
spaces that have isomorphic homology groups, making cohomology a strictly
finer invariant.

%% ──────────────────────────────────────────────────────────────────
\section{Cochains and the coboundary operator}
\label{sec:cochains}
\index{cochain}
\index{coboundary operator}

\begin{definition}[Singular cochains]
Let $X$ be a topological space and $G$ an abelian group. The
\emph{group of singular $n$-cochains} is
\[
  C^n(X;G) = \Hom_\Z\bigl(C_n(X;\Z),\,G\bigr),
\]
the group of all homomorphisms from the free abelian group of singular
$n$-chains to $G$. An element $\varphi\in C^n(X;G)$ assigns to each
singular $n$-simplex $\sigma\colon\Delta^n\to X$ a value
$\varphi(\sigma)\in G$.
\end{definition}

\begin{definition}[Coboundary operator]
\index{coboundary operator!definition}
The \emph{coboundary operator}
$\delta^n\colon C^n(X;G)\to C^{n+1}(X;G)$ is the dual of the
boundary map: for $\varphi\in C^n(X;G)$ and
$\sigma\colon\Delta^{n+1}\to X$,
\[
  (\delta^n\varphi)(\sigma) = \varphi(\partial_{n+1}\sigma)
    = \sum_{i=0}^{n+1}(-1)^i \varphi(\sigma\circ d_i),
\]
where $d_i\colon\Delta^n\to\Delta^{n+1}$ is the $i$-th face inclusion.
\end{definition}

\begin{lemma}\label{lem:delta-squared}
$\delta^{n+1}\circ\delta^n = 0$ for all $n\geq 0$.
\end{lemma}

\begin{proof}
For any $\varphi\in C^n(X;G)$ and $(n+2)$-simplex $\sigma$,
\[
  (\delta^{n+1}\delta^n\varphi)(\sigma)
  = (\delta^n\varphi)(\partial_{n+2}\sigma)
  = \varphi(\partial_{n+1}\partial_{n+2}\sigma) = \varphi(0) = 0,
\]
since $\partial^2=0$.
\end{proof}

Thus $(C^\bullet(X;G),\delta)$ is a cochain complex\index{cochain complex}:
\[
\begin{tikzcd}
  C^0(X;G) \ar[r,"\delta^0"]
  & C^1(X;G) \ar[r,"\delta^1"]
  & C^2(X;G) \ar[r,"\delta^2"]
  & \cdots
\end{tikzcd}
\]

%% ──────────────────────────────────────────────────────────────────
\section{Singular cohomology groups}
\label{sec:cohomology-groups}
\index{singular cohomology!definition}

\begin{definition}[Cohomology groups]
The $n$-th \emph{singular cohomology group} of $X$ with coefficients
in $G$ is
\[
  H^n(X;G) = \frac{\ker\bigl(\delta^n\colon C^n(X;G)\to C^{n+1}(X;G)\bigr)}
    {\im\bigl(\delta^{n-1}\colon C^{n-1}(X;G)\to C^n(X;G)\bigr)}
  = \frac{Z^n(X;G)}{B^n(X;G)},
\]
where $Z^n(X;G)=\ker\delta^n$ is the group of \emph{$n$-cocycles}%
\index{cocycle} and $B^n(X;G)=\im\delta^{n-1}$ is the group of
\emph{$n$-coboundaries}\index{coboundary}.
\end{definition}

\begin{remark}
A continuous map $f\colon X\to Y$ induces a chain map
$f_\sharp\colon C_\bullet(X)\to C_\bullet(Y)$, which dualizes to a
cochain map $f^\sharp\colon C^\bullet(Y;G)\to C^\bullet(X;G)$
(note the reversal of direction). This induces
$f^*\colon H^n(Y;G)\to H^n(X;G)$.
Cohomology is thus a \emph{contravariant} functor%
\index{contravariant functor}.
\end{remark}

\begin{example}[Degree zero]
$H^0(X;G)\cong\prod_{\alpha}G$, where the product runs over the
path components of $X$. If $X$ is path-connected, $H^0(X;G)\cong G$.
Indeed, a $0$-cocycle $\varphi\in Z^0(X;G)$ satisfies
$(\delta\varphi)(\sigma)=\varphi(\sigma(1))-\varphi(\sigma(0))=0$
for every path $\sigma$, so $\varphi$ is constant on path components.
\end{example}

\begin{proposition}[Homotopy invariance of cohomology]
\label{prop:coh-homotopy-invariance}
\index{homotopy invariance!cohomology}
If $f,g\colon X\to Y$ are homotopic, then $f^*=g^*\colon H^n(Y;G)\to H^n(X;G)$
for all $n$. In particular, if $X\simeq Y$, then $H^n(X;G)\cong H^n(Y;G)$.
\end{proposition}

\begin{proof}
If $F\colon X\times[0,1]\to Y$ is a homotopy from $f$ to $g$, then
the prism operator $P\colon C_n(X)\to C_{n+1}(Y)$ satisfying
$\partial P + P\partial = g_\sharp - f_\sharp$ dualizes to a cochain
homotopy $P^*\colon C^{n+1}(Y;G)\to C^n(X;G)$ satisfying
$P^*\delta + \delta P^* = g^\sharp - f^\sharp$. Hence
$g^\sharp - f^\sharp$ sends cocycles to coboundaries, so $f^*=g^*$.
\end{proof}

%% ──────────────────────────────────────────────────────────────────
\section{The Universal Coefficient Theorem}
\label{sec:uct}
\index{universal coefficient theorem}

The following theorem relates cohomology to homology.

\begin{theorem}[Universal Coefficient Theorem for cohomology]
\label{thm:uct}
Let $X$ be a topological space and $G$ an abelian group. There is a
natural short exact sequence
\[
  0 \to \operatorname{Ext}^1_\Z(H_{n-1}(X;\Z),\,G)
    \to H^n(X;G)
    \to \Hom_\Z(H_n(X;\Z),\,G) \to 0.
\]
This sequence splits (non-naturally), so
\[
  H^n(X;G) \cong \Hom(H_n(X;\Z),G)
    \oplus\operatorname{Ext}^1_\Z(H_{n-1}(X;\Z),G).
\]
\end{theorem}

\begin{remark}[Key special cases]
\index{Ext functor}
\begin{enumerate}[label=(\roman*)]
  \item If $H_*(X;\Z)$ is free abelian in every degree, then
    $\operatorname{Ext}^1=0$ and
    $H^n(X;G)\cong\Hom(H_n(X;\Z),G)$.
  \item $\operatorname{Ext}^1_\Z(\Z/m\Z,G)\cong G/mG$.
  \item $\operatorname{Ext}^1_\Z(\Z,G)=0$.
  \item $\operatorname{Ext}^1_\Z$ commutes with finite direct sums.
\end{enumerate}
\end{remark}

\begin{proof}[Proof sketch]
Since $C_n(X;\Z)$ is free abelian (it is a free abelian group with
basis the set of singular $n$-simplices), one constructs a free
resolution of each $H_n(X;\Z)$ from the chain complex. The key
algebraic input is the following: given a short exact sequence
of chain complexes $0\to Z_\bullet\to C_\bullet\to B_{\bullet-1}\to 0$
(where $Z_n=\ker\partial_n$ and $B_n=\im\partial_{n+1}$), applying
$\Hom(-,G)$ and taking cohomology yields the desired short exact
sequence. The $\operatorname{Ext}^1$ term arises because $\Hom(-,G)$
is only left exact, not exact, so its first right derived functor
appears.
\end{proof}

\begin{example}[UCT for $\RP^2$]
For $\RP^2$ with $G=\Z$: since $H_0(\RP^2)=\Z$, $H_1(\RP^2)=\Z/2\Z$,
$H_k(\RP^2)=0$ for $k\geq 2$, the UCT gives
\begin{align*}
  H^0(\RP^2;\Z) &\cong \Hom(\Z,\Z)\oplus 0 = \Z,\\
  H^1(\RP^2;\Z) &\cong \Hom(\Z/2\Z,\Z)\oplus \operatorname{Ext}^1(\Z,\Z)
    = 0\oplus 0 = 0,\\
  H^2(\RP^2;\Z) &\cong \Hom(0,\Z)\oplus\operatorname{Ext}^1(\Z/2\Z,\Z)
    = 0\oplus\Z/2\Z = \Z/2\Z.
\end{align*}
Note the ``shift'' of torsion from $H_1$ to $H^2$.
\end{example}

\begin{example}[UCT with $\Z/2\Z$ coefficients]
\index{universal coefficient theorem!example}
For $X=S^n$ with $G=\Z/2\Z$:
\begin{align*}
  H^k(S^n;\Z/2\Z) &\cong \Hom(H_k(S^n),\Z/2\Z)
    \oplus\operatorname{Ext}^1(H_{k-1}(S^n),\Z/2\Z)\\
  &\cong
  \begin{cases}
    \Z/2\Z & k=0,n,\\
    0 & \text{otherwise},
  \end{cases}
\end{align*}
since $\Hom(\Z,\Z/2\Z)\cong\Z/2\Z$ and $\operatorname{Ext}^1(\Z,\Z/2\Z)=0$.
\end{example}

\begin{example}[UCT for the Klein bottle]
\index{Klein bottle!cohomology of}
The Klein bottle $K$ has $H_0(K)=\Z$, $H_1(K)=\Z\oplus\Z/2\Z$,
$H_k(K)=0$ for $k\geq 2$. The UCT with $G=\Z$ gives:
\begin{align*}
  H^0(K;\Z) &\cong \Hom(\Z,\Z) = \Z,\\
  H^1(K;\Z) &\cong \Hom(\Z\oplus\Z/2\Z,\Z)\oplus\operatorname{Ext}^1(\Z,\Z)
    = \Z\oplus 0 = \Z,\\
  H^2(K;\Z) &\cong \Hom(0,\Z)\oplus\operatorname{Ext}^1(\Z\oplus\Z/2\Z,\Z)
    = 0\oplus\Z/2\Z = \Z/2\Z.
\end{align*}
Comparing with homology: $H_1(K)\cong\Z\oplus\Z/2\Z$ but
$H^1(K;\Z)\cong\Z$. The torsion in $H_1$ ``migrates'' to $H^2$
via $\operatorname{Ext}^1$.
\end{example}

%% ──────────────────────────────────────────────────────────────────
\section{The cup product}
\label{sec:cup-product}
\index{cup product!definition}

The most important additional structure on cohomology is the cup product,
which makes $H^*(X;R)$ into a graded ring when $R$ is a commutative ring.

\begin{definition}[Cup product on cochains]
Let $R$ be a commutative ring with unity. For
$\varphi\in C^p(X;R)$ and $\psi\in C^q(X;R)$, the
\emph{cup product} $\varphi\smile\psi\in C^{p+q}(X;R)$ is defined by
\[
  (\varphi\smile\psi)(\sigma)
  = \varphi(\sigma\circ[v_0,\ldots,v_p])\cdot\psi(\sigma\circ[v_p,\ldots,v_{p+q}])
\]
for every singular $(p+q)$-simplex $\sigma\colon\Delta^{p+q}\to X$.
Here $[v_0,\ldots,v_p]$ denotes the front $p$-face and
$[v_p,\ldots,v_{p+q}]$ the back $q$-face of $\Delta^{p+q}$.
\end{definition}

\begin{lemma}[Leibniz rule]\label{lem:leibniz-cup}
\index{cup product!Leibniz rule}
For $\varphi\in C^p(X;R)$ and $\psi\in C^q(X;R)$,
\[
  \delta(\varphi\smile\psi)
  = (\delta\varphi)\smile\psi + (-1)^p\,\varphi\smile(\delta\psi).
\]
\end{lemma}

\begin{proof}
This is a direct (though somewhat lengthy) computation from the
definitions. One expands both sides using the formula for $\delta$ and
observes that the terms cancel in pairs except for those contributing
to the stated formula. The key identity is that the face maps satisfy
appropriate composition relations with the front and back face
inclusions.
\end{proof}

\begin{corollary}\label{cor:cup-cohomology}
The cup product on cochains descends to cohomology: if
$[\varphi]\in H^p(X;R)$ and $[\psi]\in H^q(X;R)$, then
$[\varphi]\smile[\psi] = [\varphi\smile\psi]\in H^{p+q}(X;R)$
is well defined.
\end{corollary}

\begin{proof}
If $\varphi$ and $\psi$ are cocycles ($\delta\varphi=0=\delta\psi$),
the Leibniz rule gives $\delta(\varphi\smile\psi)=0$, so
$\varphi\smile\psi$ is a cocycle. If $\varphi=\delta\alpha$ is a
coboundary, then $\varphi\smile\psi=\delta(\alpha\smile\psi)$ by
the Leibniz rule (since $\delta\psi=0$), so the product is a coboundary.
Similarly if $\psi$ is a coboundary.
\end{proof}

%% ──────────────────────────────────────────────────────────────────
\section{Properties of the cup product}
\label{sec:cup-properties}
\index{cup product!properties}

\begin{theorem}[Properties of the cup product]\label{thm:cup-properties}
\index{graded ring}
Let $R$ be a commutative ring with unity, and let $X$ be a topological
space. The cup product on $H^*(X;R)=\bigoplus_{n\geq 0}H^n(X;R)$
satisfies:
\begin{enumerate}[label=(\roman*)]
  \item\label{it:cup-assoc}
    \textbf{Associativity:}\index{cup product!associativity}
    $(\alpha\smile\beta)\smile\gamma = \alpha\smile(\beta\smile\gamma)$
    for all $\alpha\in H^p$, $\beta\in H^q$, $\gamma\in H^r$.

  \item\label{it:cup-unit}
    \textbf{Unit:}\index{cup product!unit}
    The element $1\in H^0(X;R)\cong R$ (the cocycle sending every
    $0$-simplex to $1_R$) satisfies
    $1\smile\alpha = \alpha = \alpha\smile 1$ for all $\alpha$.

  \item\label{it:cup-graded-comm}
    \textbf{Graded commutativity:}%
    \index{cup product!graded commutativity}
    $\alpha\smile\beta = (-1)^{pq}\,\beta\smile\alpha$
    for $\alpha\in H^p(X;R)$ and $\beta\in H^q(X;R)$.

  \item\label{it:cup-bilinear}
    \textbf{Bilinearity:}
    The cup product is $R$-bilinear:
    $(\alpha+\alpha')\smile\beta = \alpha\smile\beta+\alpha'\smile\beta$
    and $\alpha\smile(\beta+\beta')=\alpha\smile\beta+\alpha\smile\beta'$.

  \item\label{it:cup-naturality}
    \textbf{Naturality:}\index{cup product!naturality}
    If $f\colon X\to Y$ is continuous, then
    $f^*(\alpha\smile\beta) = f^*(\alpha)\smile f^*(\beta)$.
\end{enumerate}
\end{theorem}

\begin{proof}[Proof sketch]
\ref{it:cup-assoc} and \ref{it:cup-unit} follow directly from the
cochain-level definitions. \ref{it:cup-bilinear} is immediate from
linearity. \ref{it:cup-naturality} follows from the fact that
$f_\sharp$ preserves front and back faces.

The proof of \ref{it:cup-graded-comm} is the most substantial.
It does \emph{not} hold at the cochain level; rather, one constructs
an explicit chain homotopy showing that
$\varphi\smile\psi - (-1)^{pq}\psi\smile\varphi$ is a coboundary
whenever $\varphi$ and $\psi$ are cocycles. This uses the
\emph{Alexander--Whitney} and \emph{Eilenberg--Zilber} maps and the
acyclic models theorem.
\index{Alexander--Whitney map}
\index{Eilenberg--Zilber map}
\index{acyclic models theorem}
\end{proof}

\begin{example}[Cup product structure on a point]
For $X=\{*\}$, $H^0(\{*\};R)\cong R$ and $H^n(\{*\};R)=0$ for $n>0$.
The cohomology ring is simply $H^*(\{*\};R)\cong R$, concentrated in
degree zero.
\end{example}

\begin{example}[Cup product on $H^*(S^1;\Z)$]
We have $H^0(S^1;\Z)\cong\Z$ with generator $1$, and
$H^1(S^1;\Z)\cong\Z$ with generator $\alpha$. Since
$H^2(S^1;\Z)=0$, the product $\alpha\smile\alpha=0$.
The full ring is $H^*(S^1;\Z)\cong\Z[\alpha]/(\alpha^2)$ with
$\abs{\alpha}=1$.
\end{example}

%% ──────────────────────────────────────────────────────────────────
\section{The cohomology ring}
\label{sec:cohomology-ring}
\index{cohomology ring}

\begin{definition}[Cohomology ring]
Let $R$ be a commutative ring. The \emph{cohomology ring} of $X$ is
the graded $R$-algebra
\[
  H^*(X;R) = \bigoplus_{n\geq 0} H^n(X;R),
\]
with multiplication given by the cup product. By
Theorem~\ref{thm:cup-properties}, $H^*(X;R)$ is a graded-commutative,
associative, unital $R$-algebra.
\end{definition}

\begin{remark}
\index{cohomology ring!as invariant}
The cohomology ring is a strictly finer invariant than the cohomology
groups alone. For example, $\CP^2$ and $S^2\vee S^4$ have isomorphic
integral cohomology groups:
\[
  H^n(\CP^2;\Z)\cong H^n(S^2\vee S^4;\Z)
  \cong
  \begin{cases}
    \Z & n=0,2,4,\\
    0  & \text{otherwise.}
  \end{cases}
\]
However, their cohomology rings differ: in $H^*(\CP^2;\Z)$, the
square of the degree-$2$ generator is the degree-$4$ generator, whereas
in $H^*(S^2\vee S^4;\Z)$, all cup products in positive degrees vanish.
\end{remark}

%% ──────────────────────────────────────────────────────────────────
\section{Computations of cohomology rings}
\label{sec:cohomology-computations}

\subsection{Cohomology ring of $S^n$}
\label{subsec:coh-sphere}
\index{sphere!cohomology ring of}

\begin{proposition}[Cohomology ring of $S^n$]\label{prop:coh-ring-Sn}
For $n\geq 1$,
\[
  H^*(S^n;\Z) \cong \Z[\alpha]/(\alpha^2),
  \quad \abs{\alpha}=n,
\]
where $\abs{\alpha}$ denotes the degree of $\alpha$. That is,
$H^*(S^n;\Z)$ is the exterior algebra on one generator of degree $n$.
\end{proposition}

\begin{proof}
We have $H^k(S^n;\Z)\cong\Z$ for $k=0,n$ and vanishes otherwise.
Let $1\in H^0$ be the unit and $\alpha\in H^n$ a generator. The only
potentially nontrivial cup product is $\alpha\smile\alpha\in H^{2n}(S^n;\Z)$.
Since $H^{2n}(S^n;\Z)=0$ for $n\geq 1$, we must have $\alpha^2=0$.
\end{proof}

\subsection{Cohomology ring of the torus $T^2$}
\label{subsec:coh-torus}
\index{torus!cohomology ring of}

\begin{theorem}[Cohomology ring of $T^2$]\label{thm:coh-ring-T2}
\[
  H^*(T^2;\Z) \cong \Z[\alpha,\beta]/(\alpha^2,\,\beta^2,\,
  \alpha\beta+\beta\alpha),
  \quad \abs{\alpha}=\abs{\beta}=1.
\]
Equivalently, $H^*(T^2;\Z)\cong\bigwedge_\Z(\alpha,\beta)$,
the exterior algebra on two generators of degree $1$.
\end{theorem}

\begin{proof}
Since $T^2=S^1\times S^1$, we can use the K\"unneth formula for
cohomology (Theorem~\ref{thm:kunneth-cohomology} below). The cohomology
groups are:
\begin{align*}
  H^0(T^2;\Z)&\cong\Z, &
  H^1(T^2;\Z)&\cong\Z^2, &
  H^2(T^2;\Z)&\cong\Z.
\end{align*}
Let $\alpha,\beta\in H^1(T^2;\Z)$ be generators corresponding to the
two $S^1$ factors. By graded commutativity,
$\alpha\smile\beta = -\beta\smile\alpha$, and
$\alpha\smile\alpha = -\alpha\smile\alpha$ forces $\alpha^2=0$
(and similarly $\beta^2=0$).

To see that $\alpha\smile\beta$ generates $H^2(T^2;\Z)\cong\Z$:
by the K\"unneth formula for cohomology rings of products, the cross
product $\alpha\times\beta$ corresponds to $\alpha\smile\beta$ under
the identification $T^2=S^1\times S^1$, and this is a generator of
$H^2$.
\end{proof}

\subsection{Cohomology ring of $\RP^n$}
\label{subsec:coh-rpn}
\index{real projective space!cohomology ring}

\begin{theorem}[Cohomology ring of $\RP^n$ with $\Z/2\Z$ coefficients]
\label{thm:coh-ring-RPn}
\[
  H^*(\RP^n;\Z/2\Z) \cong (\Z/2\Z)[\alpha]/(\alpha^{n+1}),
  \quad \abs{\alpha}=1.
\]
That is, $H^*(\RP^n;\Z/2\Z)$ is the truncated polynomial algebra over
$\Z/2\Z$ on a single generator $\alpha$ of degree $1$.
\end{theorem}

\begin{proof}[Proof sketch]
The cohomology groups are $H^k(\RP^n;\Z/2\Z)\cong\Z/2\Z$ for
$0\leq k\leq n$ and vanish for $k>n$. Let $\alpha\in H^1(\RP^n;\Z/2\Z)$
be the unique nonzero element.

The key step is to show that $\alpha^k\neq 0$ in $H^k(\RP^n;\Z/2\Z)$
for all $1\leq k\leq n$. Since $H^k(\RP^n;\Z/2\Z)\cong\Z/2\Z$, any
nonzero element is the unique generator.

This is proved by induction on $n$ using the long exact sequence in
cohomology associated to the pair $(\RP^n,\RP^{n-1})$ and naturality
of the cup product with respect to the inclusion
$\RP^{n-1}\hookrightarrow\RP^n$. The crucial geometric input is that
$\RP^n/\RP^{n-1}\cong S^n$, which identifies the relevant maps.
\end{proof}

\begin{remark}
\index{Stiefel--Whitney class}
The generator $\alpha\in H^1(\RP^n;\Z/2\Z)$ is the first
Stiefel--Whitney class $w_1$ of the tautological line bundle over
$\RP^n$. The ring structure of $H^*(\RP^n;\Z/2\Z)$ plays a fundamental
role in the theory of characteristic classes and in the proof of the
Borsuk--Ulam theorem via cohomological methods.
\index{Borsuk--Ulam theorem}
\end{remark}

\begin{corollary}\label{cor:rpn-immersion}
\index{real projective space!immersion obstruction}
The cohomology ring of $\RP^n$ with $\Z/2\Z$ coefficients places
constraints on immersions and embeddings. For instance, one can use it
to show that $\RP^n$ cannot be immersed in $\R^m$ for certain small
values of $m$, via characteristic class arguments.
\end{corollary}

\subsection{Cohomology ring of $\CP^n$}
\label{subsec:coh-cpn}
\index{complex projective space!cohomology ring}

\begin{theorem}[Cohomology ring of $\CP^n$]\label{thm:coh-ring-CPn}
\[
  H^*(\CP^n;\Z) \cong \Z[\beta]/(\beta^{n+1}),
  \quad \abs{\beta}=2.
\]
That is, $H^*(\CP^n;\Z)$ is the truncated polynomial algebra over $\Z$
on a single generator $\beta$ of degree $2$.
\end{theorem}

\begin{proof}[Proof sketch]
The integral cohomology groups are $H^{2k}(\CP^n;\Z)\cong\Z$ for
$0\leq k\leq n$ and $H^{\text{odd}}(\CP^n;\Z)=0$. Let
$\beta\in H^2(\CP^n;\Z)$ be a generator.

The proof that $\beta^k$ generates $H^{2k}(\CP^n;\Z)\cong\Z$ for all
$1\leq k\leq n$ proceeds by induction on $n$. One uses the cofiber
sequence $\CP^{n-1}\hookrightarrow\CP^n\to\CP^n/\CP^{n-1}\cong S^{2n}$
and the associated long exact sequence in cohomology. The inductive
step shows that the restriction map
$H^{2k}(\CP^n;\Z)\to H^{2k}(\CP^{n-1};\Z)$ is an isomorphism for
$k<n$, and one verifies that $\beta^n$ generates $H^{2n}$ by analyzing
the Gysin sequence of the Hopf fibration $S^1\to S^{2n+1}\to\CP^n$.
\index{Gysin sequence}
\index{Hopf fibration}
\end{proof}

\begin{remark}
\index{Chern class}
The generator $\beta\in H^2(\CP^n;\Z)$ is the first Chern class $c_1$
of the tautological line bundle over $\CP^n$. Taking the limit
$n\to\infty$, one obtains $H^*(\CP^\infty;\Z)\cong\Z[\beta]$,
an honest polynomial ring. The space $\CP^\infty$ is the classifying
space $B\mathrm{U}(1)$ for complex line bundles.
\end{remark}

\begin{remark}[Comparison of $\RP^n$ and $\CP^n$]
\index{real projective space!vs.\ complex projective space}
The following table highlights the structural analogy:
\begin{center}
\renewcommand{\arraystretch}{1.2}
\begin{tabular}{lcc}
\toprule
& $\RP^n$ & $\CP^n$ \\
\midrule
Coefficient ring & $\Z/2\Z$ & $\Z$ \\
Generator degree & $1$ & $2$ \\
Truncation & $\alpha^{n+1}=0$ & $\beta^{n+1}=0$ \\
Ring & $(\Z/2\Z)[\alpha]/(\alpha^{n+1})$ & $\Z[\beta]/(\beta^{n+1})$ \\
Dimension & $n$ & $2n$ \\
Bundle interpretation & $w_1(\gamma^1)$ & $c_1(\gamma^1)$ \\
\bottomrule
\end{tabular}
\end{center}
Both are truncated polynomial algebras, but the parity and coefficient
choices differ in ways that reflect the real vs.\ complex geometry.
\end{remark}

%% ──────────────────────────────────────────────────────────────────
\section{The K\"unneth formula in cohomology}
\label{sec:kunneth-cohomology}
\index{K\"unneth formula!cohomology}

\begin{theorem}[K\"unneth formula for cohomology]
\label{thm:kunneth-cohomology}
Let $R$ be a principal ideal domain and let $X,Y$ be topological spaces
such that $H^*(Y;R)$ is a finitely generated free $R$-module in each
degree. Then there is a natural isomorphism of graded $R$-algebras
\[
  H^*(X\times Y;R) \cong H^*(X;R)\otimes_R H^*(Y;R),
\]
where the right-hand side has the product
\[
  (\alpha\otimes\beta)\cdot(\alpha'\otimes\beta')
  = (-1)^{\abs{\beta}\abs{\alpha'}}\,
    (\alpha\smile\alpha')\otimes(\beta\smile\beta').
\]
\end{theorem}

\begin{remark}
The sign $(-1)^{\abs{\beta}\abs{\alpha'}}$ is the Koszul sign rule%
\index{Koszul sign rule}: whenever two elements of degrees $p$ and $q$
are transposed, a sign $(-1)^{pq}$ is introduced.
\end{remark}

\begin{example}[Cohomology ring of $T^n$]
\index{torus!cohomology ring of}
Applying the K\"unneth formula iteratively to $T^n=(S^1)^n$:
\[
  H^*(T^n;R) \cong H^*(S^1;R)^{\otimes n}
  \cong \bigwedge\nolimits_R(\alpha_1,\ldots,\alpha_n),
  \quad \abs{\alpha_i}=1,
\]
the exterior algebra on $n$ generators of degree $1$.
In particular, $H^k(T^n;R)\cong R^{\binom{n}{k}}$.
\end{example}

\begin{example}[Cohomology ring of $S^m\times S^n$]
For $m,n\geq 1$,
\[
  H^*(S^m\times S^n;\Z) \cong \Z[\alpha,\beta]/(\alpha^2,\beta^2),
  \quad \abs{\alpha}=m,\;\abs{\beta}=n,
\]
with the relation $\alpha\smile\beta=(-1)^{mn}\beta\smile\alpha$.
\end{example}

%% ──────────────────────────────────────────────────────────────────
\section{Applications of the cohomology ring}
\label{sec:coh-applications}

\subsection{Distinguishing spaces with the same cohomology groups}

\begin{example}[$\CP^2$ vs.\ $S^2\vee S^4$]
\label{ex:cp2-vs-wedge}
\index{complex projective space!vs.\ wedge}
As mentioned in Section~\ref{sec:cohomology-ring}, both spaces have
\[
  H^k \cong
  \begin{cases}
    \Z & k=0,2,4,\\
    0  & \text{otherwise.}
  \end{cases}
\]
In $H^*(\CP^2;\Z)\cong\Z[\beta]/(\beta^3)$, the generator
$\beta\in H^2$ satisfies $\beta^2\neq 0$ (it generates $H^4$).

In $H^*(S^2\vee S^4;\Z)$, let $\gamma\in H^2$ and $\mu\in H^4$ be
generators. Since the inclusion $S^2\hookrightarrow S^2\vee S^4$
factors through the wedge, the cup product $\gamma^2$ is computed
in $H^4(S^2;\Z)=0$ and then pushed forward, giving $\gamma^2=0$.

Therefore $H^*(\CP^2;\Z)\not\cong H^*(S^2\vee S^4;\Z)$ as graded
rings, so $\CP^2\not\simeq S^2\vee S^4$.
\end{example}

\subsection{The Hopf invariant}
\label{subsec:hopf-invariant}
\index{Hopf invariant}

\begin{definition}[Hopf invariant]
Let $f\colon S^{2n-1}\to S^n$ be a continuous map, $n\geq 2$. Form the
mapping cone $C_f = S^n\cup_f D^{2n}$. Then
$H^k(C_f;\Z)\cong\Z$ for $k=0,n,2n$ and vanishes otherwise.
Let $\alpha\in H^n(C_f;\Z)$ and $\beta\in H^{2n}(C_f;\Z)$ be
generators. The \emph{Hopf invariant} of $f$ is the integer $h(f)$
defined by
\[
  \alpha\smile\alpha = h(f)\cdot\beta.
\]
\end{definition}

\begin{remark}
The Hopf invariant is a homotopy invariant of $f$. The classical Hopf
map $\eta\colon S^3\to S^2$ has Hopf invariant $h(\eta)=1$. The theorem
of Adams states that a map $S^{2n-1}\to S^n$ of Hopf invariant $\pm 1$
exists if and only if $n\in\{1,2,4,8\}$, corresponding to the real
numbers, complex numbers, quaternions, and octonions.
\index{Adams' theorem}
\end{remark}

%% ──────────────────────────────────────────────────────────────────
\section{Connection to de Rham cohomology}
\label{sec:derham}
\index{de Rham cohomology}

When $X=M$ is a smooth manifold, there is an alternative definition
of cohomology using differential forms.

\begin{definition}[de Rham cohomology]
Let $M$ be a smooth manifold. The \emph{de Rham cohomology} is
\[
  H^n_{\mathrm{dR}}(M) = \frac{\ker\bigl(d\colon\Omega^n(M)\to\Omega^{n+1}(M)\bigr)}
    {\im\bigl(d\colon\Omega^{n-1}(M)\to\Omega^n(M)\bigr)},
\]
where $\Omega^n(M)$ is the vector space of smooth $n$-forms on $M$.
\end{definition}

\begin{theorem}[de Rham's theorem]\label{thm:deRham}
\index{de Rham theorem}
For any smooth manifold $M$, there is a natural isomorphism of graded
$\R$-algebras
\[
  H^*_{\mathrm{dR}}(M) \cong H^*(M;\R).
\]
The isomorphism is induced by integration: a closed $n$-form $\omega$
corresponds to the singular cocycle
$\sigma\mapsto\int_\sigma\omega$ on smooth singular $n$-simplices.
Under this isomorphism, the wedge product of differential forms
corresponds to the cup product.
\end{theorem}

\begin{remark}
De Rham's theorem is a profound bridge between analysis (differential
forms, exterior derivative) and topology (singular cohomology, cup
product). It explains why topological invariants can be computed by
analytic methods, a theme central to the Hodge theory and the
Atiyah--Singer index theorem.
\index{Hodge theory}
\index{Atiyah--Singer index theorem}
\end{remark}

The following diagram summarizes the relationships:
\[
\begin{tikzcd}[column sep=huge, row sep=large]
  \Omega^*(M) \ar[r,"d"] \ar[d,"\int"]
    & \Omega^{*+1}(M) \ar[d,"\int"]\\
  C^*(M;\R) \ar[r,"\delta"]
    & C^{*+1}(M;\R)
\end{tikzcd}
\]
The commutativity of this diagram (Stokes' theorem:
$\int_\sigma d\omega = \int_{\partial\sigma}\omega$)
is what makes the de Rham map a cochain map.
\index{Stokes' theorem}

\begin{example}[de Rham cohomology of $S^1$]
The $1$-form $d\theta/(2\pi)$ on $S^1$ is closed (every $1$-form on
a $1$-manifold is closed) but not exact (its integral over $S^1$ is $1$).
It represents a generator of $H^1_{\mathrm{dR}}(S^1)\cong\R$.
\end{example}

\begin{example}[de Rham cohomology of $T^2$]
\index{torus!de Rham cohomology}
Let $\theta_1,\theta_2$ be angular coordinates on $T^2=S^1\times S^1$.
Then:
\begin{itemize}
  \item $H^0_{\mathrm{dR}}(T^2)\cong\R$, generated by the constant
    function $1$.
  \item $H^1_{\mathrm{dR}}(T^2)\cong\R^2$, generated by
    $[d\theta_1]$ and $[d\theta_2]$.
  \item $H^2_{\mathrm{dR}}(T^2)\cong\R$, generated by
    $[d\theta_1\wedge d\theta_2]$.
\end{itemize}
The wedge product $d\theta_1\wedge d\theta_2$ corresponds to the cup
product of the generators, confirming the ring structure
$H^*_{\mathrm{dR}}(T^2)\cong\bigwedge_\R(\alpha,\beta)$.
\end{example}

%% ──────────────────────────────────────────────────────────────────
\section{Relative cohomology and the long exact sequence}
\label{sec:relative-cohomology}
\index{relative cohomology}

\begin{definition}[Relative cohomology]
For a pair $(X,A)$ with $A\subseteq X$, define
\[
  C^n(X,A;G) = \ker\bigl(C^n(X;G)\xrightarrow{i^*} C^n(A;G)\bigr)
  \cong \Hom(C_n(X)/C_n(A),\,G) = \Hom(C_n(X,A),\,G).
\]
The relative cohomology groups are
$H^n(X,A;G) = H^n\bigl(C^\bullet(X,A;G)\bigr)$.
\end{definition}

\begin{theorem}[Long exact sequence in cohomology]
\label{thm:les-cohomology}
\index{long exact sequence!cohomology}
For a pair $(X,A)$ there is a long exact sequence
\[
\begin{tikzcd}[column sep=small]
  \cdots \ar[r]
  & H^n(X,A;G) \ar[r,"j^*"]
  & H^n(X;G) \ar[r,"i^*"]
  & H^n(A;G) \ar[r,"\delta^*"]
  & H^{n+1}(X,A;G) \ar[r]
  & \cdots
\end{tikzcd}
\]
where $\delta^*$ is the connecting homomorphism. Note that the maps
$i^*$ and $j^*$ go in the ``reverse'' direction compared to homology,
reflecting the contravariance of cohomology.
\end{theorem}

%% ──────────────────────────────────────────────────────────────────
\section{The Mayer--Vietoris sequence in cohomology}
\label{sec:mv-cohomology}
\index{Mayer--Vietoris theorem!cohomology version}

\begin{theorem}[Mayer--Vietoris in cohomology]\label{thm:mv-cohomology}
Let $X = A\cup B$ with $X=\mathring{A}\cup\mathring{B}$. There is a
long exact sequence
\[
\begin{tikzcd}[column sep=small]
  \cdots \ar[r]
  & H^n(X;G) \ar[r,"\Psi^*"]
  & H^n(A;G)\oplus H^n(B;G) \ar[r,"\Phi^*"]
  & H^n(A\cap B;G) \ar[r,"\delta^*"]
  & H^{n+1}(X;G) \ar[r]
  & \cdots
\end{tikzcd}
\]
where $\Psi^*(\alpha)=(j_A^*\alpha,\,j_B^*\alpha)$ and
$\Phi^*(\alpha,\beta) = i_A^*\alpha - i_B^*\beta$.
\end{theorem}

%% ──────────────────────────────────────────────────────────────────
\section{Exercises}
\label{sec:coh-exercises}

\begin{exercise}
Compute $H^*(S^n;G)$ for an arbitrary abelian group $G$ using the
universal coefficient theorem. Verify your answer for $G=\Z/m\Z$.
\end{exercise}

\begin{exercise}
Let $X=\RP^2$. Compute $H^*(X;\Z)$ and $H^*(X;\Z/2\Z)$ using the
universal coefficient theorem. Compare the two.
\end{exercise}

\begin{exercise}
Compute the cohomology ring $H^*(T^3;\Z)$, where $T^3=S^1\times S^1\times S^1$. Express your answer as an exterior algebra and write down all cup products of basis elements explicitly.
\end{exercise}

\begin{exercise}
Show that $S^2\times S^4$ and $\CP^3$ have the same integral cohomology
groups. Determine whether their cohomology rings are isomorphic.
\textit{Hint:} Consider the cube of a degree-$2$ generator.
\end{exercise}

\begin{exercise}
\index{suspension}
Let $\Sigma X$ denote the suspension of $X$. Prove that the cup product
of any two positive-degree cohomology classes in $H^*(\Sigma X;R)$
vanishes. \textit{Hint:} Use the Mayer--Vietoris sequence for the
decomposition of $\Sigma X$ into two cones.
\end{exercise}

\begin{exercise}
Use the cohomology ring structure to show that $\RP^2$ is not a
retract of $\RP^3$... wait, actually show that if $r\colon\RP^3\to\RP^2$
were a retraction, then $r^*\colon H^*(\RP^2;\Z/2\Z)\to H^*(\RP^3;\Z/2\Z)$
would lead to a contradiction. \textit{Hint:} Consider what happens
to $\alpha^3$ where $\alpha$ is the degree-$1$ generator.
\end{exercise}

\begin{exercise}
Compute $H^*(\Sigma_g;\Z)$ as a graded ring, where $\Sigma_g$ is the
closed orientable surface of genus $g$. Show that $H^*(\Sigma_g;\Z)
\cong\bigwedge_\Z(\alpha_1,\beta_1,\ldots,\alpha_g,\beta_g)/I$
where $I$ is a suitable ideal, and determine the cup product pairing
$H^1\otimes H^1\to H^2\cong\Z$.
\index{surface!cohomology ring}
\end{exercise}

\begin{exercise}
\index{K\"unneth formula!cohomology!with torsion}
Use the universal coefficient theorem to compute
$H^*(\RP^2\times\RP^2;\Z/2\Z)$ as a graded ring. Show that it is
isomorphic to
$(\Z/2\Z)[\alpha,\beta]/(\alpha^3,\beta^3)$
with $\abs{\alpha}=\abs{\beta}=1$.
\end{exercise}

\begin{exercise}
Let $f\colon S^3\to S^2$ be the Hopf map. Compute $H^*(C_f;\Z)$ as a
graded ring, where $C_f$ is the mapping cone of $f$, and verify that
the Hopf invariant of $f$ is $\pm 1$.
\end{exercise}

\begin{exercise}
\index{de Rham cohomology!computation}
Let $M=S^1\times S^2$.
\begin{enumerate}[label=(\alph*)]
  \item Compute the de Rham cohomology $H^*_{\mathrm{dR}}(M)$ using the
    K\"unneth formula.
  \item Find explicit closed differential forms representing generators of
    each cohomology group.
  \item Verify the ring structure by computing wedge products of your
    representative forms.
\end{enumerate}
\end{exercise}

\index{singular cohomology|)}
\index{cup product|)}

% === END OF CHUNK ch7-8 ===
% === START OF CHUNK ch9-11+end ===

%%%%%%%%%%%%%%%%%%%%%%%%%%%%%%%%%%%%%%%%%%%%%%%%%%%%%%%%%%%%%%%%%%%%
%  CHAPTER 9 — Poincaré Duality
%%%%%%%%%%%%%%%%%%%%%%%%%%%%%%%%%%%%%%%%%%%%%%%%%%%%%%%%%%%%%%%%%%%%
\chapter{Poincar\'e Duality}\label{ch:poincare-duality}
\index{Poincare duality@Poincar\'e duality}
\minitoc\bigskip

Poincar\'e duality is one of the most striking structural results in
algebraic topology: for a closed oriented $n$-manifold~$M$, the
cohomology group $H^p(M;R)$ is isomorphic to the homology group
$H_{n-p}(M;R)$.  This chapter develops the machinery needed to state
and prove this theorem, and explores its far-reaching consequences.

%% ─── 9.1 Orientability ─────────────────────────────────────────
\section{Orientability and the fundamental class}\label{sec:orientability}
\index{orientability}\index{manifold!orientable}

Throughout this chapter, $M$ denotes a connected, closed (compact,
without boundary) topological $n$-manifold\index{manifold!closed}.
We write $R$ for a commutative ring with unity (typically $\Z$, $\Q$,
or~$\R$).

\begin{definition}[Local orientation]
\index{orientation!local}
Let $x\in M$.  A \textbf{local orientation} of~$M$ at~$x$ is a choice
of generator $\mu_x$ of the local homology group
\[
  H_n(M, M\setminus\{x\};\Z) \;\cong\; \Z.
\]
\end{definition}

The local homology groups form a covering space over~$M$, which we now
describe precisely.

\begin{definition}[Orientation sheaf and orientability]
\index{orientation!sheaf}\index{orientation!double cover}
The \textbf{orientation double cover}\index{double cover!orientation}
$\widetilde{M}\to M$ is the two-sheeted covering whose fibre over~$x$
consists of the two generators of
$H_n(M,M\setminus\{x\};\Z)\cong\Z$.  The manifold~$M$ is
\textbf{orientable} if this covering is trivial, i.e., if
$\widetilde{M}$ has two connected components.  An \textbf{orientation}
of~$M$ is a choice of section of $\widetilde{M}\to M$, equivalently a
consistent choice of local orientations $\{\mu_x\}_{x\in M}$.
\end{definition}

\begin{example}[Orientable manifolds]
All spheres $S^n$, all complex projective spaces $\CP^n$, all tori
$T^n = (S^1)^n$, and all Lie groups are orientable.
\end{example}

\begin{example}[Non-orientable manifolds]
\index{manifold!non-orientable}
The real projective plane $\RP^2$\index{projective space!real}, the
M\"obius band\index{Mobius band@M\"obius band}, and the Klein
bottle\index{Klein bottle} are non-orientable.
\end{example}

\begin{theorem}[Fundamental class]\label{thm:fundamental-class}
\index{fundamental class}\index{$[M]$|see{fundamental class}}
Let $M$ be a closed, connected, oriented $n$-manifold.  Then there
exists a unique class $[M]\in H_n(M;\Z)$ such that for every
$x\in M$, the image of $[M]$ under the canonical map
$H_n(M;\Z)\to H_n(M,M\setminus\{x\};\Z)\cong\Z$ is the chosen local
orientation~$\mu_x$.  We call $[M]$ the \textbf{fundamental class}
of~$M$.

If $M$ is connected and closed but non-orientable, there still exists
a fundamental class $[M]\in H_n(M;\Z/2)$.
\end{theorem}

\begin{proof}[Proof sketch]
One proceeds by induction using a Mayer--Vietoris\index{Mayer--Vietoris}
argument.  Cover $M$ by open sets each homeomorphic to~$\R^n$.  For a
single chart, the statement is clear.  Gluing via Mayer--Vietoris and
a compactness argument (finitely many charts suffice) produces the
global class.  The key input is that the consistency of local
orientations forces the relevant connecting homomorphisms to be
isomorphisms.  For full details, see Hatcher~\cite{hatcher}
Theorem~3.26.
\end{proof}

\begin{remark}[Coefficients in $\Z/2$]
\index{Z/2 coefficients@$\Z/2$ coefficients}
Every closed connected $n$-manifold, orientable or not, admits a
$\Z/2$-fundamental class $[M]_2\in H_n(M;\Z/2)$.  Poincar\'e duality
with $\Z/2$ coefficients therefore holds universally for closed
manifolds.
\end{remark}

%% ─── 9.2 The cap product ───────────────────────────────────────
\section{The cap product}\label{sec:cap-product}
\index{cap product}

Poincar\'e duality is realized by capping with the fundamental class.
We first introduce the relevant algebraic operation.

\begin{definition}[Cap product]
\index{cap product!definition}
Let $X$ be a topological space and $R$ a commutative ring.  The
\textbf{cap product} is the bilinear pairing
\[
  \frown\;:\; H_n(X;R)\otimes H^p(X;R) \;\longrightarrow\;
  H_{n-p}(X;R),
\]
defined on the chain level as follows.  For an $n$-simplex
$\sigma\colon\Delta^n\to X$ and a cochain
$\varphi\in C^p(X;R)$, set
\[
  \sigma \frown \varphi
  \;=\;
  \varphi\bigl(\sigma|_{[v_0,\dots,v_p]}\bigr)
  \cdot
  \sigma|_{[v_p,\dots,v_n]}.
\]
\end{definition}

\begin{proposition}[Properties of the cap product]
\label{prop:cap-properties}
\index{cap product!properties}
The cap product satisfies:
\begin{enumerate}[label=(\roman*)]
  \item \textbf{Naturality.}  For a continuous map $f\colon X\to Y$,
    $f_*(\alpha\frown f^*\varphi) = f_*\alpha\frown\varphi$.
  \item \textbf{Associativity.}  $\alpha\frown(\varphi\smile\psi)
    = (\alpha\frown\varphi)\frown\psi$, where $\smile$ denotes the
    cup product\index{cup product}.
  \item \textbf{Unit.}  $\alpha\frown 1 = \alpha$, where
    $1\in H^0(X;R)$ is the unit of the cup product ring.
\end{enumerate}
\end{proposition}

\begin{proof}
These identities are verified at the cochain level using the
Alexander--Whitney diagonal approximation\index{Alexander--Whitney map}
and then pass to (co)homology.  See Hatcher~\cite{hatcher} \S\,3.3.
\end{proof}

The cap product relates cup product duality to homology in the
following commutative diagram:
\[
\begin{tikzcd}[column sep=large]
  H_n(X;R)\otimes H^p(X;R) \ar[r,"\frown"]
    \ar[d,"\mathrm{id}\otimes\Delta^*"']
  & H_{n-p}(X;R) \\
  H_n(X;R)\otimes H^p(X;R) \ar[ru,"\frown"']
\end{tikzcd}
\]

\begin{example}[Cap product on $S^n$]
\index{cap product!on spheres}
Let $\iota\in H^n(S^n;\Z)\cong\Z$ be the canonical generator and
$[S^n]\in H_n(S^n;\Z)\cong\Z$ the fundamental class.  Then
\[
  [S^n]\frown\iota \;=\; 1 \;\in\; H_0(S^n;\Z)\cong\Z.
\]
More generally, $[S^n]\frown 1 = [S^n]$ where
$1\in H^0(S^n;\Z)$ is the unit.
\end{example}

\begin{example}[Cap product on the torus]
\index{cap product!on the torus}
Let $T^2=S^1\times S^1$ with generators $\alpha,\beta\in H^1(T^2;\Z)$
and fundamental class $[T^2]\in H_2(T^2;\Z)$.  Then
$[T^2]\frown\alpha\in H_1(T^2;\Z)$ is the homology class of the
second $S^1$ factor (with appropriate sign), and
$[T^2]\frown(\alpha\smile\beta) = \langle\alpha\smile\beta,[T^2]\rangle
= \pm1\in H_0(T^2;\Z)\cong\Z$.  The cap product map $D_{T^2}$
provides the explicit Poincar\'e duality isomorphism
$H^1(T^2)\cong H_1(T^2)\cong\Z^2$.
\end{example}

\begin{proposition}[Cap product and evaluation]
\label{prop:cap-evaluation}
\index{cap product!and evaluation}
For a closed oriented $n$-manifold~$M$ and $\varphi\in H^n(M;R)$,
\[
  [M]\frown\varphi \;=\; \langle\varphi,[M]\rangle \;\in\;
  H_0(M;R)\cong R,
\]
where $\langle-,-\rangle\colon H^n(M;R)\otimes H_n(M;R)\to R$ is the
Kronecker pairing\index{Kronecker pairing}.
\end{proposition}

%% ─── 9.3 Poincaré Duality Theorem ──────────────────────────────
\section{The Poincar\'e Duality Theorem}\label{sec:pd-theorem}

\begin{theorem}[Poincar\'e Duality]\label{thm:poincare-duality}
\index{Poincare duality@Poincar\'e duality!theorem}
Let $M$ be a closed, oriented, connected $n$-manifold and let $R$ be a
commutative ring with unity.  Then the cap product with the fundamental
class\index{fundamental class}
\[
  D_M\;:\; H^p(M;R) \;\xrightarrow{\;\;[M]\frown -\;\;}
  \; H_{n-p}(M;R)
\]
is an isomorphism for every $0\le p\le n$.
\end{theorem}

Equivalently, the pairing
\[
  H^p(M;R)\otimes_R H^{n-p}(M;R) \;\xrightarrow{\;\smile\;}
  H^n(M;R) \;\xrightarrow{\;\langle -,[M]\rangle\;}
  R
\]
is non-degenerate when $R$ is a field\index{Poincare duality@Poincar\'e duality!over a field}.

\begin{proof}[Proof strategy]
\index{Poincare duality@Poincar\'e duality!proof}
The proof proceeds in several stages:

\smallskip\noindent
\textbf{Step 1} (Open subsets of $\R^n$).  For $U\subset\R^n$ open,
one establishes the isomorphism using the fact that $U$ is a union of
convex sets and applying Mayer--Vietoris.

\smallskip\noindent
\textbf{Step 2} (Finite good covers).  If $M = U\cup V$ and the
theorem holds for $U$, $V$, and $U\cap V$, then it holds for~$M$ by
the Five Lemma\index{Five Lemma} applied to the Mayer--Vietoris
sequences in homology and cohomology linked by cap product maps:
\[
\begin{tikzcd}[column sep=small, font=\small]
  \cdots \ar[r]
  & H^p(M) \ar[r]\ar[d,"D_M"]
  & H^p(U)\oplus H^p(V) \ar[r]\ar[d,"D_U\oplus D_V"]
  & H^p(U\cap V) \ar[r]\ar[d,"D_{U\cap V}"]
  & H^{p+1}(M) \ar[r]\ar[d,"D_M"]
  & \cdots \\
  \cdots \ar[r]
  & H_{n-p}(M) \ar[r]
  & H_{n-p}(U)\oplus H_{n-p}(V) \ar[r]
  & H_{n-p}(U\cap V) \ar[r]
  & H_{n-p-1}(M) \ar[r]
  & \cdots
\end{tikzcd}
\]

\smallskip\noindent
\textbf{Step 3} (Compactness).  Since $M$ is compact, it admits a
finite good cover\index{good cover}.  An induction on the number of
open sets in the cover, using Step~2, completes the proof.

The bookkeeping requires care with signs and orientations; see
Hatcher~\cite{hatcher} Theorem~3.30 or Bredon~\cite{bredon} VI.8.
\end{proof}

\begin{remark}[Non-orientable case]
\index{Poincare duality@Poincar\'e duality!non-orientable}
For a closed connected $n$-manifold that is \emph{not} orientable,
Poincar\'e duality still holds with $\Z/2$ coefficients:
$H^p(M;\Z/2)\cong H_{n-p}(M;\Z/2)$.  Alternatively, one can use
\emph{twisted coefficients}\index{twisted coefficients} via the
orientation sheaf.
\end{remark}

%% ─── 9.4 Applications ──────────────────────────────────────────
\section{Applications of Poincar\'e duality}\label{sec:pd-applications}

\subsection{Betti numbers and symmetry}
\index{Betti number!symmetry}

\begin{corollary}[Symmetry of Betti numbers]
\label{cor:betti-symmetry}
Let $M$ be a closed oriented $n$-manifold and $\F$ a field.  Then
$b_p(M;\F) = b_{n-p}(M;\F)$ for all $0\le p\le n$, where
$b_p = \dim_{\F} H_p(M;\F)$.
\end{corollary}

\begin{proof}
By the universal coefficient theorem over a field,
$H^p(M;\F)\cong H_p(M;\F)^*$, so
$\dim H^p(M;\F)=\dim H_p(M;\F)$.  Poincar\'e duality gives
$H^p(M;\F)\cong H_{n-p}(M;\F)$, and the result follows.
\end{proof}

\begin{corollary}[Euler characteristic of odd-dimensional manifolds]
\label{cor:euler-odd}
\index{Euler characteristic!odd-dimensional manifold}
If $M$ is a closed oriented manifold of odd dimension $n=2k+1$, then
$\chi(M)=0$.
\end{corollary}

\begin{proof}
We have $\chi(M)=\sum_{p=0}^n(-1)^p b_p$.  Pairing $b_p$ with
$b_{n-p}$ and using $b_p=b_{n-p}$, the terms cancel in pairs because
$(-1)^p+(-1)^{n-p}=(-1)^p(1+(-1)^n)=0$ when $n$ is odd.
\end{proof}

\subsection{Cohomology of spheres and tori}
\index{sphere!Poincare duality@Poincar\'e duality}

\begin{example}[Spheres]
For $S^n$ ($n\ge1$), Poincar\'e duality gives $H^p(S^n)\cong
H_{n-p}(S^n)$, confirming $H^0(S^n)\cong H^n(S^n)\cong\Z$ and
$H^p(S^n)=0$ for $0<p<n$.
\end{example}

\begin{example}[Tori]
\index{torus!cohomology}
For the $n$-torus $T^n=(S^1)^n$, the cohomology ring is the exterior
algebra $H^*(T^n;\Z)\cong\Lambda_\Z[\alpha_1,\dots,\alpha_n]$ with
each $\alpha_i\in H^1$.  Poincar\'e duality is reflected in the
isomorphism $\Lambda^p\cong\Lambda^{n-p}$, so
$b_p(T^n)=\binom{n}{p}$.
\end{example}

\subsection{Projective spaces}
\index{projective space!Poincare duality@Poincar\'e duality}

\begin{example}[Complex projective spaces]
For $\CP^n$, $H^{2k}(\CP^n;\Z)\cong\Z$ for $0\le k\le n$ and all
odd-degree cohomology vanishes.  Poincar\'e duality gives
$H^{2k}\cong H_{2n-2k}$, consistent with the ring structure
$H^*(\CP^n;\Z)\cong\Z[\alpha]/(\alpha^{n+1})$, $\abs{\alpha}=2$.
\end{example}

\begin{example}[Real projective spaces]
\index{projective space!real}
For $\RP^n$ with $\Z/2$ coefficients, $H^p(\RP^n;\Z/2)\cong\Z/2$
for all $0\le p\le n$, and Poincar\'e duality with $\Z/2$
coefficients yields the expected symmetry.  Over~$\Z$, one must use
twisted coefficients when $n$ is even (since $\RP^n$ is non-orientable
for even~$n$).
\end{example}

\subsection{The intersection form}\label{subsec:intersection-form}
\index{intersection form}

\begin{definition}[Intersection form]
Let $M$ be a closed oriented $4$-manifold.  The \textbf{intersection
form} of~$M$ is the symmetric bilinear form
\[
  Q_M\;:\; H^2(M;\Z)/\mathrm{tors}\;\times\;
  H^2(M;\Z)/\mathrm{tors} \;\longrightarrow\;\Z
\]
defined by $Q_M(\alpha,\beta) = \langle \alpha\smile\beta,\,[M]\rangle$.
Poincar\'e duality implies that $Q_M$ is
\textbf{unimodular}\index{unimodular form} (i.e., its matrix has
determinant $\pm1$).
\end{definition}

\begin{example}[Intersection forms of familiar 4-manifolds]
\index{intersection form!examples}
\begin{enumerate}[label=(\roman*)]
  \item $Q_{S^4}=0$ (the trivial form), since $H^2(S^4)=0$.
  \item $Q_{\CP^2}=(1)$, the $1\times1$ matrix with entry~$1$.
  \item $Q_{S^2\times S^2}=\bigl(\begin{smallmatrix}0&1\\1&0
    \end{smallmatrix}\bigr)$, the hyperbolic form.
  \item $Q_{\CP^2\#\overline{\CP^2}}=
    \bigl(\begin{smallmatrix}1&0\\0&-1\end{smallmatrix}\bigr)$.
\end{enumerate}
\end{example}

\begin{remark}[Donaldson and Freedman]
\index{Donaldson's theorem}\index{Freedman's theorem}
The intersection form is a key invariant in 4-manifold topology.
Freedman (1982) showed that the homeomorphism type of a simply
connected closed oriented topological 4-manifold is determined by its
intersection form.  Donaldson (1983) proved that if $Q_M$ is positive
definite and $M$ is smooth, then $Q_M$ is diagonalizable
over~$\Z$---a dramatic restriction that distinguishes the smooth and
topological categories in dimension four.
\end{remark}

\subsection{Poincar\'e--Lefschetz duality}
\index{Poincare--Lefschetz duality@Poincar\'e--Lefschetz duality}

\begin{theorem}[Poincar\'e--Lefschetz duality]
\label{thm:pl-duality}
Let $M$ be a compact oriented $n$-manifold with boundary
$\partial M$\index{manifold!with boundary}.  Then there is an
isomorphism
\[
  D\;:\; H^p(M,\partial M;R)
  \;\xrightarrow{\;\cong\;}\;
  H_{n-p}(M;R)
\]
for all $p$, induced by the cap product with the relative fundamental
class $[M,\partial M]\in H_n(M,\partial M;R)$.

Similarly, $H^p(M;R)\cong H_{n-p}(M,\partial M;R)$.
\end{theorem}

\begin{remark}[Hodge duality]
\index{Hodge duality}\index{de Rham cohomology}
On a closed oriented Riemannian $n$-manifold~$M$, the Hodge star
operator $*\colon\Omega^p(M)\to\Omega^{n-p}(M)$ induces an isomorphism
on harmonic forms $\mathcal{H}^p(M)\cong\mathcal{H}^{n-p}(M)$.  By
the Hodge theorem, $\mathcal{H}^p(M)\cong H^p_{\mathrm{dR}}(M;\R)$,
so this gives an analytic proof of Poincar\'e duality over~$\R$ for
smooth manifolds.  The relationship between the topological cap product
and the Hodge star is a deep theme in geometric analysis.
\end{remark}

%% ─── 9.5 Exercises ─────────────────────────────────────────────
\section{Exercises}

\begin{exercise}
\index{fundamental class!computation}
Compute the fundamental class $[T^2]\in H_2(T^2;\Z)$ by
triangulating $T^2$ as a quotient of a square and tracking the
orientation through the identifications.
\end{exercise}

\begin{exercise}
Show that a closed connected manifold $M$ satisfies
$H_n(M;\Z)\cong\Z$ if $M$ is orientable and $H_n(M;\Z)=0$ if $M$ is
non-orientable.
\end{exercise}

\begin{exercise}
Use Poincar\'e duality to show that if $M$ is a closed oriented
$3$-manifold, then $H_1(M;\Z)\cong H^2(M;\Z)$ and $b_1(M)=b_2(M)$.
\end{exercise}

\begin{exercise}
\index{cup product!Poincare duality@Poincar\'e duality}
Let $M$ be a closed oriented $4k$-manifold.  Show that the signature
$\sigma(M)$\index{signature} of the intersection form $Q_M$ on
$H^{2k}(M;\R)$ is a homotopy invariant.
\end{exercise}

\begin{exercise}
Verify Poincar\'e duality explicitly for $\CP^2$ by computing
$H^p(\CP^2;\Z)$ and $H_{4-p}(\CP^2;\Z)$ for all~$p$ and
exhibiting the cap product isomorphism.
\end{exercise}

\begin{exercise}
Let $M$ be a closed oriented $n$-manifold with $n\equiv 2\pmod{4}$.
Show that $\chi(M)$ is even.
\textit{Hint:} Use the non-degenerate skew-symmetric pairing on
$H^{n/2}(M;\R)$.
\end{exercise}

\begin{exercise}
\index{Klein bottle!duality}
The Klein bottle $K$ is a closed non-orientable $2$-manifold.  Compute
$H_*(K;\Z)$ and $H_*(K;\Z/2)$.  Verify Poincar\'e duality with
$\Z/2$ coefficients.
\end{exercise}

\begin{exercise}
\index{connected sum}
Let $M$ and $N$ be closed oriented $4$-manifolds.  Show that
$Q_{M\# N}\cong Q_M\oplus Q_N$.
\end{exercise}

\begin{exercise}
\index{Poincare--Lefschetz duality@Poincar\'e--Lefschetz duality!exercise}
Let $M$ be the closed unit disk $D^n\subset\R^n$, viewed as a compact
manifold with boundary $S^{n-1}$.  Verify Poincar\'e--Lefschetz
duality by computing all relative and absolute (co)homology groups.
\end{exercise}

\begin{exercise}
Prove that every closed oriented $3$-manifold is parallelizable.
\textit{Hint:} Use Poincar\'e duality and the fact that $w_2=0$
for oriented $3$-manifolds.\index{parallelizable}
\end{exercise}


%%%%%%%%%%%%%%%%%%%%%%%%%%%%%%%%%%%%%%%%%%%%%%%%%%%%%%%%%%%%%%%%%%%%
%  CHAPTER 10 — Higher Homotopy Groups: Introduction
%%%%%%%%%%%%%%%%%%%%%%%%%%%%%%%%%%%%%%%%%%%%%%%%%%%%%%%%%%%%%%%%%%%%
\chapter{Higher Homotopy Groups --- Introduction}\label{ch:higher-homotopy}
\index{homotopy group!higher}
\minitoc\bigskip

The fundamental group captures information about loops in a space.
Higher homotopy groups $\pi_n(X,x_0)$ generalize this to maps from
higher-dimensional spheres, providing a much finer (and harder to
compute) invariant.

%% ─── 10.1 Definition and first properties ──────────────────────
\section{Definition and group structure}\label{sec:pin-definition}

\begin{definition}[Higher homotopy groups]
\index{homotopy group!definition}
\index{$\pi_n$|see{homotopy group}}
Let $X$ be a topological space with basepoint $x_0\in X$ and let
$n\ge1$.  The \textbf{$n$-th homotopy group} of~$X$ is
\[
  \pi_n(X,x_0)
  \;=\;
  [S^n,X]_0
  \;=\;
  \set{f\colon (S^n,*)\to(X,x_0)} / \text{based homotopy},
\]
the set of based homotopy classes of based maps $S^n\to X$.

Equivalently, identifying $S^n$ with $I^n/\partial I^n$, an element
of $\pi_n(X,x_0)$ is represented by a map $f\colon I^n\to X$ sending
$\partial I^n$ to~$x_0$, modulo homotopies that also fix $\partial
I^n$ at~$x_0$.
\end{definition}

The group operation on $\pi_n$ is defined by ``stacking'' in the first
coordinate:

\begin{definition}[Group operation]
\index{homotopy group!group operation}
Given $[f],[g]\in\pi_n(X,x_0)$ (represented as maps $I^n\to X$ fixing
$\partial I^n$), define $[f]\cdot[g]$ to be the class of
\[
  (f\cdot g)(t_1,t_2,\dots,t_n)
  \;=\;
  \begin{cases}
    f(2t_1,t_2,\dots,t_n) & \text{if } 0\le t_1\le\tfrac12,\\
    g(2t_1-1,t_2,\dots,t_n) & \text{if } \tfrac12\le t_1\le 1.
  \end{cases}
\]
\end{definition}

\begin{proposition}[$\pi_n$ is a group]
\label{prop:pin-group}
For $n\ge1$, the operation above makes $\pi_n(X,x_0)$ into a group.
The identity is the class of the constant map, and the inverse of
$[f]$ is $[f^{-1}]$ where
$f^{-1}(t_1,t_2,\dots,t_n)=f(1-t_1,t_2,\dots,t_n)$.
\end{proposition}

\begin{theorem}[Eckmann--Hilton argument: commutativity for $n\ge2$]
\label{thm:eckmann-hilton}
\index{Eckmann--Hilton argument}\index{homotopy group!abelian}
For $n\ge2$, the group $\pi_n(X,x_0)$ is abelian.
\end{theorem}

\begin{proof}
For $n\ge2$ we can stack maps in either the first or the second
coordinate of~$I^n$.  Define two multiplications:
\begin{align*}
  (f*_1 g)(t_1,\dots,t_n)
  &=
  \begin{cases} f(2t_1,t_2,\dots) & t_1\le\tfrac12,\\
    g(2t_1-1,t_2,\dots) & t_1\ge\tfrac12,\end{cases} \\
  (f*_2 g)(t_1,\dots,t_n)
  &=
  \begin{cases} f(t_1,2t_2,\dots) & t_2\le\tfrac12,\\
    g(t_1,2t_2-1,\dots) & t_2\ge\tfrac12.\end{cases}
\end{align*}
Both are associative with the same unit (the constant map).  Moreover,
$*_1$ and $*_2$ satisfy the \emph{interchange law}:
$(f*_1 g)*_2(h*_1 k) = (f*_2 h)*_1(g*_2 k)$.  The Eckmann--Hilton
argument then gives $f*_1 g = f*_2 g = g*_1 f$; that is, $\pi_n$ is
abelian.
\end{proof}

%% ─── 10.2 Functoriality ────────────────────────────────────────
\section{Functoriality}\label{sec:pin-functoriality}
\index{homotopy group!functoriality}

\begin{proposition}[Functoriality of $\pi_n$]
A based map $f\colon (X,x_0)\to(Y,y_0)$ induces a group homomorphism
\[
  f_*\;:\;\pi_n(X,x_0)\to\pi_n(Y,y_0),
  \qquad
  [\alpha]\mapsto[f\circ\alpha].
\]
This assignment is functorial:
$(\mathrm{id}_X)_*=\mathrm{id}$ and
$(g\circ f)_*=g_*\circ f_*$.

In particular, if $f$ is a homotopy equivalence, then $f_*$ is an
isomorphism for all~$n$.
\end{proposition}

This means $\pi_n$ is a functor from the homotopy category of based
spaces to the category of groups (abelian groups for $n\ge2$).

\begin{example}[Functoriality and degree]
\index{degree of a map!and homotopy groups}
A map $f\colon S^n\to S^n$ of degree~$d$ induces the multiplication-by-$d$
map $f_*\colon\pi_n(S^n)\cong\Z\to\pi_n(S^n)\cong\Z$,
$f_*([{\mathrm{id}}])=d\cdot[{\mathrm{id}}]$.  In particular, the
antipodal map $a\colon S^n\to S^n$ has degree $(-1)^{n+1}$ and induces
$a_*=(-1)^{n+1}\cdot\mathrm{id}$ on~$\pi_n(S^n)$.
\end{example}

\begin{remark}[Dependence on basepoint]
\index{homotopy group!basepoint dependence}
If $X$ is path-connected, then $\pi_n(X,x_0)\cong\pi_n(X,x_1)$ for
any two basepoints $x_0,x_1\in X$, via conjugation by a path from
$x_0$ to~$x_1$.  This isomorphism depends on the choice of path, so
the basepoint cannot simply be ignored.  For $n\ge2$ the group is
abelian, so conjugation by different paths yields the same isomorphism
if and only if $\pi_1(X)$ acts trivially on~$\pi_n(X)$.  A space where
this action is trivial for all~$n$ is called \textbf{simple}%
\index{simple space}.
\end{remark}

%% ─── 10.3 Long exact sequence of a pair ────────────────────────
\section{Long exact sequence of a pair}\label{sec:les-pair}
\index{long exact sequence!of a pair (homotopy)}
\index{homotopy group!relative}

\begin{definition}[Relative homotopy group]
\index{homotopy group!relative}
For a pair $(X,A)$ with $x_0\in A$, the \textbf{relative homotopy
group} $\pi_n(X,A,x_0)$ for $n\ge1$ consists of homotopy classes
of maps $(I^n,\partial I^n,J^{n-1})\to(X,A,x_0)$, where
$J^{n-1}=\partial I^n\setminus I^{n-1}\times\{0\}$ is the
``upper'' part of the boundary.
\end{definition}

\begin{theorem}[Long exact sequence of a pair]
\label{thm:les-homotopy}
\index{long exact sequence!of a pair (homotopy)!theorem}
For a pair $(X,A)$ with $x_0\in A$, there is a long exact sequence
\[
\begin{tikzcd}[column sep=1.4em]
  \cdots \ar[r]
  & \pi_n(A) \ar[r,"i_*"]
  & \pi_n(X) \ar[r,"j_*"]
  & \pi_n(X,A) \ar[r,"\partial"]
  & \pi_{n-1}(A) \ar[r]
  & \cdots \ar[r]
  & \pi_0(A) \ar[r]
  & \pi_0(X)
\end{tikzcd}
\]
where $i\colon A\hookrightarrow X$ is the inclusion, $j$ comes from
$(X,x_0)\hookrightarrow(X,A)$, and $\partial$ is the boundary map.
\end{theorem}

\begin{proof}[Proof sketch]
Exactness at $\pi_n(X)$: $[f]\in\ker j_*$ means $f$ is homotopic
(as a map of pairs) to a map into~$A$, which is exactly $\im i_*$.
The other exactness arguments are analogous; see
Hatcher~\cite{hatcher} Theorem~4.3.
\end{proof}

%% ─── 10.4 Fibrations ───────────────────────────────────────────
\section{Fibrations and the homotopy exact sequence}
\label{sec:fibrations}
\index{fibration}\index{fibration!homotopy lifting property}

\begin{definition}[Fibration]
A map $p\colon E\to B$ is a \textbf{(Hurewicz) fibration} if it has
the \textbf{homotopy lifting property} (HLP) with respect to all
spaces: for every space~$Y$, every map $\tilde f_0\colon Y\to E$,
and every homotopy $F\colon Y\times I\to B$ with
$F(-,0)=p\circ\tilde f_0$, there exists a homotopy
$\tilde F\colon Y\times I\to E$ with
$p\circ\tilde F=F$ and $\tilde F(-,0)=\tilde f_0$.
\end{definition}

\begin{definition}[Serre fibration]
\index{fibration!Serre}
A map $p\colon E\to B$ is a \textbf{Serre fibration} if it has the HLP
with respect to all disks $D^n$ (equivalently, all CW complexes).
\end{definition}

\begin{theorem}[Long exact sequence of a fibration]
\label{thm:les-fibration}
\index{long exact sequence!of a fibration}
If $F\hookrightarrow E\xrightarrow{p}B$ is a fibration sequence with
fibre $F=p^{-1}(b_0)$, there is a long exact sequence of homotopy
groups:
\[
\begin{tikzcd}[column sep=1.4em]
  \cdots \ar[r]
  & \pi_n(F) \ar[r,"i_*"]
  & \pi_n(E) \ar[r,"p_*"]
  & \pi_n(B) \ar[r,"\partial"]
  & \pi_{n-1}(F) \ar[r]
  & \cdots \ar[r]
  & \pi_0(E) \ar[r]
  & \pi_0(B)
\end{tikzcd}
\]
\end{theorem}

\begin{proof}
This is a consequence of Theorem~\ref{thm:les-homotopy} combined with
the identification $\pi_n(E,F)\cong\pi_n(B)$ coming from the HLP.
See Hatcher~\cite{hatcher} Theorem~4.41.
\end{proof}

%% ─── 10.5 Computations ─────────────────────────────────────────
\section{Computations of homotopy groups}\label{sec:pin-computations}

\subsection{Cellular approximation and vanishing results}
\index{cellular approximation}\index{homotopy group!vanishing}

\begin{theorem}[Cellular approximation]
\label{thm:cellular-approx}
Every map between CW complexes is homotopic to a cellular map.
Consequently, if $X$ has no cells in dimensions $1,2,\dots,n-1$
(beyond the basepoint), then $\pi_k(X)=0$ for $k<n$.
\end{theorem}

\begin{corollary}
\label{cor:pin-sphere-vanishing}
$\pi_n(S^m)=0$ for all $n<m$.\index{sphere!homotopy groups}
\end{corollary}

\begin{proof}
$S^m$ has a CW structure with one $0$-cell and one $m$-cell.  By
cellular approximation, any map $S^n\to S^m$ with $n<m$ is homotopic
to a map into the $(m-1)$-skeleton, which is a point.
\end{proof}

\subsection{The Hurewicz theorem}
\index{Hurewicz theorem}

\begin{definition}[Hurewicz homomorphism]
\index{Hurewicz homomorphism}
The \textbf{Hurewicz map}\index{Hurewicz map|see{Hurewicz homomorphism}}
$h\colon\pi_n(X,x_0)\to H_n(X;\Z)$ sends $[f\colon S^n\to X]$ to
$f_*(\iota_n)$, where $\iota_n\in H_n(S^n;\Z)\cong\Z$ is the
canonical generator.
\end{definition}

\begin{theorem}[Hurewicz theorem]\label{thm:hurewicz}
\index{Hurewicz theorem!statement}
Let $X$ be a path-connected space with $x_0\in X$.
\begin{enumerate}[label=(\roman*)]
  \item If $n=1$: the Hurewicz map $h\colon\pi_1(X,x_0)\to H_1(X;\Z)$
    is the abelianization\index{abelianization} map
    $\pi_1\twoheadrightarrow\pi_1^{\mathrm{ab}}\cong H_1$.
  \item If $n\ge2$ and $\pi_k(X)=0$ for all $1\le k<n$ (i.e., $X$ is
    $(n{-}1)$-connected)\index{connected!$(n-1)$-connected}: then
    $\widetilde{H}_k(X;\Z)=0$ for $k<n$ and
    \[
      h\;:\;\pi_n(X,x_0)\;\xrightarrow{\;\cong\;}\;H_n(X;\Z)
    \]
    is an isomorphism.
\end{enumerate}
\end{theorem}

\begin{proof}[Proof for $n=1$]
\index{Hurewicz theorem!proof for $n=1$}
For $n=1$, $\pi_1(X,x_0)$ consists of homotopy classes of loops
$\gamma\colon S^1\to X$.  Viewing $\gamma$ as a singular $1$-cycle, we
obtain a well-defined map $h\colon\pi_1(X)\to H_1(X;\Z)$.  This is
a homomorphism since the concatenation of loops corresponds to the sum
of cycles (up to a boundary).

\textbf{Surjectivity.} Every $1$-cycle in a path-connected space is
homologous to a sum of loops (using paths to the basepoint), so $h$
is surjective.

\textbf{Kernel.} If $h([\gamma])=0$, then $\gamma$ (as a $1$-chain)
is a boundary $\gamma=\partial c$ for some $2$-chain~$c$.  Each
singular $2$-simplex in~$c$ gives a relation among loops, and these
relations are precisely the commutator relations.  Hence
$\ker h=[\pi_1,\pi_1]$.
\end{proof}

\begin{corollary}\label{cor:pin-sn}
\index{sphere!$\pi_n(S^n)\cong\Z$}
$\pi_n(S^n)\cong\Z$ for all $n\ge1$.
\end{corollary}

\begin{proof}
For $n=1$, $\pi_1(S^1)\cong\Z$ is known from covering space theory.
For $n\ge2$, $S^n$ is $(n{-}1)$-connected (by
Corollary~\ref{cor:pin-sphere-vanishing}), so the Hurewicz theorem
gives $\pi_n(S^n)\cong H_n(S^n;\Z)\cong\Z$.
\end{proof}

\subsection{The Hopf fibration}
\index{Hopf fibration}

\begin{theorem}[Hopf fibration and $\pi_3(S^2)$]
\label{thm:hopf}
There is a fibration $S^1\hookrightarrow S^3\xrightarrow{\eta}S^2$
(the \textbf{Hopf fibration}), defined by
\[
  \eta(z_1,z_2)=[z_1:z_2]\in\CP^1\cong S^2,
  \qquad (z_1,z_2)\in S^3\subset\C^2.
\]
Moreover, $\pi_3(S^2)\cong\Z$, generated by $[\eta]$.
\end{theorem}

\begin{proof}
That $\eta$ is a fibration with fibre $S^1$ follows from the
structure of the $\mathrm{U}(1)$-action on~$S^3$.  The long exact
sequence of the fibration gives:
\[
\begin{tikzcd}[column sep=1.6em]
  \pi_3(S^1) \ar[r]
  & \pi_3(S^3) \ar[r,"\eta_*"]
  & \pi_3(S^2) \ar[r,"\partial"]
  & \pi_2(S^1) \ar[r]
  & \pi_2(S^3)
\end{tikzcd}
\]
Since $\pi_3(S^1)=0$, $\pi_2(S^1)=0$ (using the universal cover
$\R\to S^1$), and $\pi_3(S^3)\cong\Z$, exactness yields
$\pi_3(S^2)\cong\pi_3(S^3)\cong\Z$.
\end{proof}

\begin{remark}[Other Hopf fibrations]
\index{Hopf fibration!quaternionic}\index{Hopf fibration!octonionic}
There are analogous Hopf fibrations:
\[
  S^3\hookrightarrow S^7\to S^4,
  \qquad
  S^7\hookrightarrow S^{15}\to S^8,
\]
arising from quaternionic and octonionic multiplication.  These show
$\pi_7(S^4)\cong\Z\oplus\Z_{12}$ (the $\Z$ factor) and
$\pi_{15}(S^8)\cong\Z$.
\end{remark}

\begin{remark}[Homotopy groups of spheres]
\index{sphere!homotopy groups!table}
The computation of $\pi_n(S^m)$ in general is one of the deepest
problems in algebraic topology.  Some known values:
\begin{center}
\renewcommand{\arraystretch}{1.2}
\begin{tabular}{c|cccccc}
  $\pi_i(S^n)$ & $i=1$ & $2$ & $3$ & $4$ & $5$ & $6$ \\
  \hline
  $S^1$ & $\Z$ & $0$ & $0$ & $0$ & $0$ & $0$ \\
  $S^2$ & $0$ & $\Z$ & $\Z$ & $\Z_2$ & $\Z_2$ & $\Z_{12}$ \\
  $S^3$ & $0$ & $0$ & $\Z$ & $\Z_2$ & $\Z_2$ & $\Z_{12}$ \\
  $S^4$ & $0$ & $0$ & $0$ & $\Z$ & $\Z_2$ & $\Z_2$ \\
\end{tabular}
\end{center}
The appearance of torsion and the lack of a general closed formula
distinguish homotopy groups from homology groups.
\end{remark}

%% ─── 10.6 The Whitehead theorem ─────────────────────────────────
\section{The Whitehead theorem}\label{sec:whitehead}
\index{Whitehead theorem}

\begin{theorem}[Whitehead]\label{thm:whitehead}
Let $f\colon X\to Y$ be a map between connected CW
complexes\index{CW complex}.  If $f_*\colon\pi_n(X)\to\pi_n(Y)$ is an
isomorphism for all $n\ge1$, then $f$ is a homotopy equivalence.
\end{theorem}

\begin{proof}[Proof sketch]
By the homotopy extension property for CW pairs, one constructs the
homotopy inverse cell by cell.  The hypothesis that $f$ induces
isomorphisms on all $\pi_n$ allows one to extend over each skeleton.
The passage from a weak homotopy equivalence to a homotopy equivalence
uses the fact that CW complexes have the homotopy type of
CW complexes.  See Hatcher~\cite{hatcher} Theorem~4.5.
\end{proof}

\begin{remark}[CW hypothesis is essential]
\index{Whitehead theorem!counterexample}
The CW hypothesis cannot be dropped.  There exist spaces that have the
same homotopy groups as a point but are not contractible (e.g., the
Warsaw circle\index{Warsaw circle}).  The issue is that non-CW spaces
may have pathological local behavior that prevents the construction of
a homotopy inverse.
\end{remark}

\begin{example}[Application of the Whitehead theorem]
\index{Whitehead theorem!application}
Let $X$ be a simply connected CW complex with
$H_n(X;\Z)\cong H_n(S^k;\Z)$ for all~$n$.  Then $X\simeq S^k$.

Indeed, since $X$ is $(k{-}1)$-connected (by the Hurewicz theorem),
there exists a map $f\colon S^k\to X$ representing a generator of
$\pi_k(X)\cong H_k(X)\cong\Z$.  By the Hurewicz theorem and a
comparison of homology, $f$ induces isomorphisms on all homology
groups.  By Corollary~\ref{cor:whitehead-homology}, $f$ is a homotopy
equivalence.
\end{example}

\begin{corollary}\label{cor:whitehead-homology}
If $f\colon X\to Y$ is a map between simply connected CW complexes
that induces isomorphisms on all homology groups, then $f$ is a
homotopy equivalence.
\end{corollary}

\begin{proof}
By the Hurewicz theorem, the hypothesis on homology combined with
simple connectivity implies isomorphisms on homotopy groups.  Then
apply Theorem~\ref{thm:whitehead}.
\end{proof}

%% ─── 10.7 Exercises ────────────────────────────────────────────
\section{Exercises}

\begin{exercise}
\index{Eckmann--Hilton argument!exercise}
Carry out the Eckmann--Hilton argument in full detail: given a set~$S$
with two binary operations $*_1,*_2$, both associative, both sharing
the same unit~$e$, and satisfying the interchange law, prove that
$*_1=*_2$ and both are commutative.
\end{exercise}

\begin{exercise}
\index{homotopy group!product space}
Show that $\pi_n(X\times Y)\cong\pi_n(X)\times\pi_n(Y)$ for
all $n\ge1$.  Deduce $\pi_n(T^k)$ for all $n,k$.
\end{exercise}

\begin{exercise}
Use the long exact sequence of the covering $\R\to S^1$ to prove that
$\pi_n(S^1)=0$ for all $n\ge2$.
\end{exercise}

\begin{exercise}
\index{Hopf fibration!exercise}
From the Hopf fibration $S^1\hookrightarrow S^3\to S^2$, determine
$\pi_2(S^2)$ and $\pi_4(S^3)$ (the latter requires knowing
$\pi_3(S^1)$ and $\pi_4(S^2)\cong\Z_2$).
\end{exercise}

\begin{exercise}
\index{loop space}
Let $\Omega X$ denote the loop space of~$X$.  Show that
$\pi_n(X)\cong\pi_{n-1}(\Omega X)$ for all $n\ge1$.
\end{exercise}

\begin{exercise}
\index{suspension}
Let $\Sigma X$ denote the (reduced) suspension of~$X$.  Show that
there is a natural homomorphism
$E\colon\pi_n(X)\to\pi_{n+1}(\Sigma X)$ (the \textbf{suspension
homomorphism}\index{Freudenthal suspension theorem}).  State (without
proof) the Freudenthal suspension theorem.
\end{exercise}

\begin{exercise}
\index{Hurewicz theorem!exercise}
Let $X$ be a simply connected CW complex with $H_2(X;\Z)\cong\Z^3$
and $H_k(X;\Z)=0$ for $k=1$.  What can you conclude about $\pi_2(X)$?
\end{exercise}

\begin{exercise}
\index{Whitehead theorem!exercise}
Construct an example of a map $f\colon X\to Y$ between
path-connected spaces (not CW complexes) such that $f$ induces
isomorphisms on all homotopy groups but $f$ is not a homotopy
equivalence.
\end{exercise}

\begin{exercise}
\index{fibration!exercise}
Consider the fibration $\mathrm{SO}(n)\hookrightarrow
\mathrm{SO}(n+1)\to S^n$.  Use the long exact sequence to show that
$\pi_1(\mathrm{SO}(3))\cong\Z/2$.
\end{exercise}

\begin{exercise}
Show that $\pi_4(S^2)\cong\Z_2$.
\textit{Hint:} Use the long exact sequence of the Hopf fibration and
the known value $\pi_4(S^3)\cong\Z_2$.
\end{exercise}


%%%%%%%%%%%%%%%%%%%%%%%%%%%%%%%%%%%%%%%%%%%%%%%%%%%%%%%%%%%%%%%%%%%%
%  CHAPTER 11 — Obstruction Theory and Applications
%%%%%%%%%%%%%%%%%%%%%%%%%%%%%%%%%%%%%%%%%%%%%%%%%%%%%%%%%%%%%%%%%%%%
\chapter{Obstruction Theory and Applications}\label{ch:obstruction-theory}
\index{obstruction theory}
\minitoc\bigskip

Obstruction theory provides a systematic framework for answering a
fundamental question in topology: \emph{when can a map defined on a
subspace be extended to the whole space?}  The obstructions to
extension lie in cohomology groups with coefficients in homotopy
groups of the target, revealing a profound interaction between
cohomology and homotopy.

%% ─── 11.1 The extension problem ────────────────────────────────
\section{The extension problem}\label{sec:extension-problem}
\index{extension problem}

\begin{definition}[Extension problem]
Given a CW pair $(X,A)$, a space~$Y$, and a map $f\colon A\to Y$, the
\textbf{extension problem} asks: does there exist a continuous map
$\tilde f\colon X\to Y$ with $\tilde f|_A=f$?
\[
\begin{tikzcd}
  A \ar[r,"f"] \ar[d,hookrightarrow] & Y \\
  X \ar[ru,dashed,"\tilde f"']
\end{tikzcd}
\]
\end{definition}

The key difficulty is that, unlike homological questions (which can
often be resolved by algebraic machinery), the extension problem
requires understanding the interplay between the cell structure of~$X$
and the homotopy type of~$Y$.  Obstruction theory reduces this
geometric problem to a sequence of algebraic computations in
cohomology.

\begin{remark}[Lifting vs.\ extension]
\index{lifting problem}
A closely related question is the \textbf{lifting problem}: given a
fibration $p\colon E\to B$ and a map $f\colon X\to B$, does there
exist a lift $\tilde f\colon X\to E$ with $p\circ\tilde f=f$?
\[
\begin{tikzcd}
  & E \ar[d,"p"] \\
  X \ar[r,"f"] \ar[ru,dashed,"\tilde f"]
  & B
\end{tikzcd}
\]
Obstruction theory applies to lifting problems as well, with
obstructions lying in $H^{n+1}(X;\pi_n(F))$, where $F$ is the fibre.
\end{remark}

The strategy of obstruction theory is to extend $f$ one skeleton at a
time.

\begin{proposition}[Skeletal extension]
\label{prop:skeletal-extension}
\index{extension problem!skeletal}
Let $(X,A)$ be a relative CW complex and $f\colon A\to Y$ a continuous
map.
\begin{enumerate}[label=(\roman*)]
  \item $f$ always extends to $A\cup X^{(0)}$ (by choosing images for
    the new $0$-cells arbitrarily, in the same path component of~$Y$).
  \item If $Y$ is path-connected, $f$ extends to $A\cup X^{(1)}$.
  \item If $Y$ is additionally simply connected, $f$ extends to
    $A\cup X^{(2)}$.
  \item In general, the obstruction to extending from the $n$-skeleton
    to the $(n{+}1)$-skeleton lives in a cohomology group.
\end{enumerate}
\end{proposition}

%% ─── 11.2 The obstruction cocycle ──────────────────────────────
\section{The obstruction cocycle}\label{sec:obstruction-cocycle}
\index{obstruction cocycle}

Suppose $f$ has been extended to a map $f_n\colon A\cup X^{(n)}\to Y$.
For each $(n{+}1)$-cell $e^{n+1}_\alpha$ of $X$ relative to~$A$, the
attaching map
$\varphi_\alpha\colon S^n\to X^{(n)}$ composed with $f_n$ gives
an element
\[
  [f_n\circ\varphi_\alpha]\;\in\;\pi_n(Y).
\]
The map $f_n$ extends over $e^{n+1}_\alpha$ if and only if this class
is zero (because extending is equivalent to filling in the sphere).

\begin{definition}[Obstruction cocycle]
\index{obstruction cocycle!definition}
The \textbf{obstruction cocycle} is the cellular cochain
\[
  \mathfrak{o}^{n+1}(f_n)
  \;\in\;
  C^{n+1}(X,A;\pi_n(Y))
\]
defined by
$\mathfrak{o}^{n+1}(f_n)(e^{n+1}_\alpha)
= [f_n\circ\varphi_\alpha]\in\pi_n(Y)$.
\end{definition}

\begin{theorem}[The obstruction cocycle is a cocycle]
\label{thm:obstruction-cocycle}
\index{obstruction cocycle!is a cocycle}
$\mathfrak{o}^{n+1}(f_n)$ is a cocycle, i.e.,
$\delta\mathfrak{o}^{n+1}(f_n)=0$ in
$C^{n+2}(X,A;\pi_n(Y))$.  Moreover:
\begin{enumerate}[label=(\roman*)]
  \item $f_n$ extends to $A\cup X^{(n+1)}$ if and only if
    $\mathfrak{o}^{n+1}(f_n)=0$ as a cochain.
  \item The cohomology class
    $[\mathfrak{o}^{n+1}(f_n)]\in H^{n+1}(X,A;\pi_n(Y))$ depends only
    on the homotopy class of $f_n|_{A\cup X^{(n-1)}}$.
  \item $f_{n-1}=f_n|_{A\cup X^{(n-1)}}$ can be extended to some
    $f'_n\colon A\cup X^{(n)}\to Y$ that further extends to
    $A\cup X^{(n+1)}$ if and only if $[\mathfrak{o}^{n+1}(f_n)]=0$ in
    $H^{n+1}(X,A;\pi_n(Y))$.
\end{enumerate}
\end{theorem}

\begin{proof}[Proof sketch]
That $\mathfrak{o}^{n+1}(f_n)$ is a cocycle follows from the fact
that the boundary of a boundary is zero: the attaching map of an
$(n{+}2)$-cell, restricted to any $(n{+}1)$-cell in its boundary, and
then mapped by~$f_n$, yields a sum of homotopy classes that must
vanish because the total attaching sphere is null-homotopic in
$X^{(n)}$.

For (iii), modifying $f_n$ on individual $n$-cells changes
$\mathfrak{o}^{n+1}$ by a coboundary.  See Whitehead~\cite{spanier}
or Davis--Kirk for the detailed argument.
\end{proof}

%% ─── 11.3 The main obstruction theorem ─────────────────────────
\section{The main obstruction theorem}\label{sec:main-obstruction}
\index{obstruction theory!main theorem}

\begin{theorem}[Obstruction theorem]\label{thm:obstruction-main}
Let $(X,A)$ be a relative CW complex, $Y$ a simple space
(i.e., $\pi_1(Y)$ acts trivially on $\pi_n(Y)$ for all~$n$), and
$f\colon A\to Y$ a continuous map.  Then:
\begin{enumerate}[label=(\roman*)]
  \item The first (and possibly only) obstruction to extending $f$ to
    all of~$X$ lies in $H^{n+1}(X,A;\pi_n(Y))$ for the smallest~$n$
    such that $\pi_n(Y)\ne0$.
  \item If all obstruction classes vanish successively, $f$ extends to
    all of~$X$.
  \item If $X$ is finite-dimensional, the process terminates in
    finitely many steps.
\end{enumerate}
\end{theorem}

\begin{remark}[Primary obstruction]
\index{obstruction!primary}
When $Y$ is $(n{-}1)$-connected (i.e., $\pi_k(Y)=0$ for $k<n$), the
first potential obstruction lies in $H^{n+1}(X,A;\pi_n(Y))$.  This is
called the \textbf{primary obstruction}.  In favorable cases (e.g.,
$\dim(X\setminus A)\le n+1$), the primary obstruction is the only one.
\end{remark}

\begin{example}[Sections of sphere bundles]
\index{sphere bundle!section}
\index{Euler class}
Let $\xi\colon E\to M$ be an oriented $S^{n-1}$-bundle over a CW
complex~$M$ (i.e., the unit sphere bundle of an oriented rank-$n$
vector bundle).  A global section $s\colon M\to E$ exists if and only
if the successive obstructions vanish.  The primary obstruction is
a class
\[
  e(\xi)\;\in\; H^n(M;\pi_{n-1}(S^{n-1}))\;=\;H^n(M;\Z),
\]
which is precisely the \textbf{Euler class}\index{Euler class} of the
underlying vector bundle.  In particular, $e(\xi)=0$ is necessary
(and, for $\dim M\le n$, sufficient) for the existence of a nowhere-zero
section of the vector bundle.
\end{example}

\begin{example}[Obstruction to a map $\RP^2\to S^2$]
\index{projective space!maps to spheres}
Consider the problem of finding a map $f\colon\RP^2\to S^2$ inducing
a nontrivial map on $H^2(-;\Z/2)$.  The CW structure of $\RP^2$ has
cells in dimensions $0,1,2$.  The map on the $0$-skeleton is trivial.
Extending to the $1$-skeleton requires $\pi_1(S^2)=0$, which is
satisfied, so there is no obstruction.  Extending to the $2$-skeleton
requires the obstruction in $H^2(\RP^2;\pi_1(S^2))=0$ to vanish,
which it does.  So any map $\RP^2\to S^2$ must be null-homotopic,
since $\pi_2(S^2)\cong\Z$ and $H^2(\RP^2;\Z)\cong\Z/2$ yields
$[\RP^2,S^2]=H^2(\RP^2;\pi_2(S^2))$ by obstruction theory (noting
$S^2$ is $1$-connected).
\end{example}

%% ─── 11.4 Applications ─────────────────────────────────────────
\section{Applications}\label{sec:obstruction-applications}

\subsection{The Borsuk--Ulam theorem revisited}
\index{Borsuk--Ulam theorem}

\begin{theorem}[Borsuk--Ulam]\label{thm:borsuk-ulam}
For every continuous map $f\colon S^n\to\R^n$, there exists a point
$x\in S^n$ with $f(x)=f(-x)$.
\end{theorem}

\begin{proof}[Proof via obstruction theory for $n=2$]
Suppose no such $x$ exists.  Then
$g(x)=\frac{f(x)-f(-x)}{\|f(x)-f(-x)\|}\colon S^2\to S^1$ is a
well-defined odd map (i.e., $g(-x)=-g(x)$).  This induces a map
$\bar g\colon\RP^2\to\RP^1\cong S^1$ on quotients.

On fundamental groups, $\bar g_*$ must send the generator of
$\pi_1(\RP^2)\cong\Z/2$ to the generator of $\pi_1(S^1)\cong\Z$,
which is impossible since there is no nonzero homomorphism
$\Z/2\to\Z$.
\end{proof}

\begin{remark}
The obstruction-theoretic proof in general dimensions uses the fact
that an equivariant map $S^n\to S^{n-1}$ (with respect to the
antipodal action) would yield a section of a certain sphere bundle
over $\RP^n$, whose obstruction class is nonzero.
\end{remark}

\subsection{The Lusternik--Schnirelmann theorem}
\index{Lusternik--Schnirelmann theorem}

\begin{theorem}[Lusternik--Schnirelmann]
\label{thm:lusternik-schnirelmann}
If $S^n$ is covered by $n+1$ closed sets $A_1,\dots,A_{n+1}$, each
of which contains no pair of antipodal points, then
$\bigcap_{i=1}^{n+1}A_i\ne\emptyset$... but more importantly, $S^n$
cannot be covered by $n+1$ closed sets, each containing no antipodal
pair.  Equivalently, any cover of $\RP^n$ by $n+1$ open sets must have
at least one set that is not contractible within $\RP^n$.
\end{theorem}

\begin{proof}[Proof sketch]
This follows from the Borsuk--Ulam theorem.  If
$S^n = A_1\cup\cdots\cup A_{n+1}$ with no $A_i$ containing an
antipodal pair, define $f\colon S^n\to\R^n$ by
$f_i(x)=d(x,A_i)$ for $i=1,\dots,n$.  By Borsuk--Ulam, there exists
$x$ with $f(x)=f(-x)$, which forces $x$ and $-x$ to lie in
$A_{n+1}$, a contradiction.
\end{proof}

\subsection{Vector bundle classification}
\index{vector bundle!classification}
\index{classifying space}

\begin{remark}[Classification of vector bundles --- preview]
Obstruction theory underlies the classification of vector
bundles\index{vector bundle}.  An oriented rank-$k$ real vector bundle
$E\to X$ is classified (up to isomorphism) by a homotopy class
of maps $X\to\mathrm{BSO}(k)$\index{BSO(k)@$\mathrm{BSO}(k)$},
where $\mathrm{BSO}(k)$ is the classifying space.  The obstructions
to finding a global nonvanishing section, a global frame, or a
spin structure are cohomology classes known as
\textbf{characteristic classes}\index{characteristic class}
(Stiefel--Whitney\index{Stiefel--Whitney class},
Chern\index{Chern class},
Pontryagin\index{Pontryagin class} classes).  This is the subject of
Milnor--Stasheff~\cite{milnor-stasheff}.
\end{remark}

%% ─── 11.5 Eilenberg-MacLane spaces ─────────────────────────────
\section{Eilenberg--MacLane spaces}\label{sec:eilenberg-maclane}
\index{Eilenberg--MacLane space}

\begin{definition}[Eilenberg--MacLane space]
\index{K(G,n)@$K(G,n)$|see{Eilenberg--MacLane space}}
Let $G$ be a group (abelian if $n\ge2$) and $n\ge1$.  An
\textbf{Eilenberg--MacLane space} $K(G,n)$ is a connected CW
complex\index{CW complex} with
\[
  \pi_k(K(G,n))
  \;=\;
  \begin{cases}
    G & \text{if } k=n,\\
    0 & \text{if } k\ne n.
  \end{cases}
\]
\end{definition}

\begin{theorem}[Existence and uniqueness]
\label{thm:em-existence}
\index{Eilenberg--MacLane space!existence}
For every group $G$ (abelian if $n\ge2$) and every $n\ge1$, a
$K(G,n)$ exists and is unique up to homotopy equivalence.
\end{theorem}

\begin{proof}[Proof sketch]
\textbf{Existence.}  Start with a presentation of~$G$ and build a
$2$-complex with the correct $\pi_1$ (or an $(n{+}1)$-complex with
the correct $\pi_n$ via attaching cells to~$S^n$).  Then
systematically attach cells in dimensions $\ge n+2$ to kill all higher
homotopy groups.

\textbf{Uniqueness.}  Any two $K(G,n)$'s have isomorphic homotopy
groups, and a map inducing these isomorphisms can be constructed
by obstruction theory.  By the Whitehead theorem
(Theorem~\ref{thm:whitehead}), this map is a homotopy equivalence.
\end{proof}

\begin{example}[Familiar Eilenberg--MacLane spaces]
\index{Eilenberg--MacLane space!examples}
\begin{enumerate}[label=(\roman*)]
  \item $K(\Z,1)\simeq S^1$, since $\pi_1(S^1)\cong\Z$ and
    $\pi_n(S^1)=0$ for $n\ge2$.
  \item $K(\Z,2)\simeq\CP^\infty$, the infinite complex projective
    space.
  \item $K(\Z/2,1)\simeq\RP^\infty$, the infinite real projective space.
  \item $K(\Z,n)$ has no simple closed-form model for $n\ge3$; it
    must be constructed by cell attachment.
\end{enumerate}
\end{example}

%% ─── 11.6 Brown representability ───────────────────────────────
\section{Representability of cohomology}\label{sec:brown-representability}
\index{Brown representability}

The deepest connection between Eilenberg--MacLane spaces and
cohomology is:

\begin{theorem}[Brown representability for ordinary cohomology]
\label{thm:brown-representability}
\index{Brown representability!theorem}
For any abelian group~$G$, $n\ge0$, and CW complex~$X$, there is a
natural bijection
\[
  [X,\,K(G,n)]
  \;\cong\;
  H^n(X;G),
\]
where $[X,K(G,n)]$ denotes the set of (unbased) homotopy classes of
maps from $X$ to $K(G,n)$.
\end{theorem}

\begin{proof}[Proof sketch]
The map $[X,K(G,n)]\to H^n(X;G)$ sends $[f]$ to $f^*(\iota_n)$,
where $\iota_n\in H^n(K(G,n);G)\cong\Hom(G,G)$ is the
\textbf{fundamental class}\index{fundamental class!of $K(G,n)$}
(corresponding to $\mathrm{id}_G$).

That this is a bijection is proved by verifying the
Eilenberg--Steenrod axioms: both sides satisfy the Mayer--Vietoris
property and have the same values on points, so they agree on all
CW complexes by induction on the number of cells.  This is an instance
of the \textbf{Brown representability theorem}\index{Brown
representability!general}, which states that any contravariant functor
from the homotopy category of connected CW complexes to sets that
satisfies the wedge axiom and the Mayer--Vietoris axiom is
representable.
\end{proof}

The following diagram summarizes the representability:
\[
\begin{tikzcd}[column sep=large]
  X \ar[r,"f"] \ar[rd,"f^*(\iota_n)"']
  & K(G,n) \\
  & H^n(X;G) \ar[u,phantom,"\scriptstyle\cong"
    description,sloped]
\end{tikzcd}
\qquad\longleftrightarrow\qquad
[f]\in[X,K(G,n)]\;\longleftrightarrow\;
f^*(\iota_n)\in H^n(X;G).
\]

\begin{remark}[Postnikov towers]
\index{Postnikov tower}
Every connected CW complex $X$ can be approximated by a tower of
fibrations
\[
\begin{tikzcd}
  \vdots \ar[d] \\
  X_3 \ar[d,"p_3"] \\
  X_2 \ar[d,"p_2"] \\
  X_1
\end{tikzcd}
\]
where $X_n$ has the same $\pi_k$ as $X$ for $k\le n$ and
$\pi_k(X_n)=0$ for $k>n$, and the fibre of $p_n$ is $K(\pi_n(X),n)$.
The ``$k$-invariants'' that glue successive stages are cohomology
classes, bringing obstruction theory full circle.
\end{remark}

%% ─── 11.7 Stable homotopy and spectra ──────────────────────────
\section{Stable homotopy --- a preview}\label{sec:stable-homotopy}
\index{stable homotopy}\index{spectrum}

\begin{definition}[Stable homotopy groups]
\index{stable homotopy group}
The \textbf{$k$-th stable homotopy group of spheres} is
\[
  \pi_k^s
  \;=\;
  \lim_{n\to\infty}\pi_{n+k}(S^n),
\]
where the maps in the direct system are given by the suspension
homomorphism $E\colon\pi_m(S^n)\to\pi_{m+1}(S^{n+1})$.
The Freudenthal suspension theorem\index{Freudenthal suspension theorem}
guarantees that this limit stabilizes (i.e., $E$ is an isomorphism for
$n>m+1$).
\end{definition}

\begin{example}[Low-dimensional stable homotopy groups]
\index{stable homotopy group!table}
\begin{center}
\renewcommand{\arraystretch}{1.2}
\begin{tabular}{c|cccccccc}
  $k$ & $0$ & $1$ & $2$ & $3$ & $4$ & $5$ & $6$ & $7$ \\
  \hline
  $\pi_k^s$ & $\Z$ & $\Z_2$ & $\Z_2$ & $\Z_{24}$ & $0$ & $0$
  & $\Z_2$ & $\Z_{240}$
\end{tabular}
\end{center}
\end{example}

\begin{remark}[Spectra and generalized cohomology]
\index{spectrum!definition}\index{generalized cohomology}
To systematically study stable phenomena, one introduces the notion
of a \textbf{spectrum} $\mathbf{E}=\{E_n,\sigma_n\}$, a sequence of
spaces with structure maps $\sigma_n\colon\Sigma E_n\to E_{n+1}$.
Every spectrum defines a \textbf{generalized cohomology theory}
$h^n(X)=[X,E_n]$, and the Brown representability theorem guarantees
that every generalized cohomology theory arises this way.  Important
examples include:
\begin{itemize}
  \item The \textbf{Eilenberg--MacLane spectrum}
    $H\Z=\{K(\Z,n)\}$\index{Eilenberg--MacLane spectrum}, representing
    ordinary cohomology.
  \item \textbf{$K$-theory}\index{K-theory} spectra $KU$ and $KO$,
    representing complex and real topological $K$-theory.
  \item \textbf{Cobordism} spectra $MU$, $MSO$,
    $\dots$\index{cobordism}.
\end{itemize}
The stable homotopy category is the natural setting for duality
phenomena (Spanier--Whitehead duality\index{Spanier--Whitehead duality},
Atiyah duality\index{Atiyah duality}) and is foundational to modern
homotopy theory.
\end{remark}

%% ─── 11.8 Connections to mathematical physics ──────────────────
\section{Connections to mathematical physics}
\label{sec:physics-connections}
\index{mathematical physics}

\begin{remark}[Topology in physics]
Algebraic topology, and obstruction theory in particular, appears
naturally throughout mathematical physics:
\begin{enumerate}[label=(\roman*)]
  \item \textbf{Gauge theory.}\index{gauge theory}  Principal
    $G$-bundles over a manifold~$M$ are classified by $[M,BG]$.
    Characteristic classes provide topological invariants of gauge
    fields.  The Chern--Weil homomorphism\index{Chern--Weil
    homomorphism} realizes characteristic classes via curvature forms.
  \item \textbf{Topological defects.}\index{topological defect}
    In condensed matter physics and cosmology, topological defects
    (vortices, monopoles, domain walls) are classified by homotopy
    groups of the vacuum manifold~$M$:
    $\pi_0(M)$ classifies domain walls, $\pi_1(M)$ classifies strings,
    and $\pi_2(M)$ classifies monopoles.
  \item \textbf{Anomalies.}\index{anomaly}  Topological obstructions
    (related to index theory\index{index theory} and $\eta$-invariants)
    detect quantum anomalies in gauge theories.  The famous chiral
    anomaly\index{chiral anomaly} is related to the index of the Dirac
    operator via the Atiyah--Singer index
    theorem\index{Atiyah--Singer index theorem}.
  \item \textbf{Topological quantum field theory.}
    \index{topological quantum field theory}
    A TQFT in the sense of Atiyah assigns to each closed
    $(n{-}1)$-manifold a vector space and to each cobordism a linear
    map.  The algebraic structure mirrors the composition of
    cobordisms, and homotopy-theoretic methods (spectra, $\infty$-categories)
    provide the modern framework.
\end{enumerate}
\end{remark}

%% ─── 11.9 Exercises ────────────────────────────────────────────
\section{Exercises}

\begin{exercise}
\index{extension problem!exercise}
Let $X=S^2$ and $A=\{*\}$.  Show that every map $f\colon A\to S^1$
extends to $X$.  What is the obstruction group?
\end{exercise}

\begin{exercise}
\index{obstruction cocycle!exercise}
Let $X$ be a CW complex of dimension $\le n$ and $Y$ an
$(n{-}1)$-connected space.  Show that the primary obstruction to
extending any map $f\colon X^{(n)}\to Y$ to $X^{(n+1)}$ is the only
obstruction, and lives in $H^{n+1}(X;\pi_n(Y))$.
\end{exercise}

\begin{exercise}
\index{Eilenberg--MacLane space!exercise}
Prove that $\Omega K(G,n)\simeq K(G,n-1)$ for $n\ge2$.
\textit{Hint:} Use the long exact sequence of the path-space fibration
$\Omega K(G,n)\hookrightarrow PK(G,n)\to K(G,n)$.
\end{exercise}

\begin{exercise}
\index{Brown representability!exercise}
Use the Brown representability theorem to show that
$[X,\CP^\infty]\cong H^2(X;\Z)$ for any CW complex~$X$.  Interpret
this in terms of complex line bundles.
\end{exercise}

\begin{exercise}
\index{Borsuk--Ulam theorem!exercise}
Use the Borsuk--Ulam theorem to prove that $\R^n$ and $\R^m$ are not
homeomorphic for $n\ne m$.
\end{exercise}

\begin{exercise}
\index{Postnikov tower!exercise}
Describe the first two stages of the Postnikov tower for $S^2$.
What are the relevant $k$-invariants?
\end{exercise}

\begin{exercise}
\index{Lusternik--Schnirelmann!exercise}
Use the Lusternik--Schnirelmann theorem to show that any continuous map
$f\colon S^2\to\R^2$ maps some pair of antipodal points to the same
value.
\end{exercise}

\begin{exercise}
\index{stable homotopy!exercise}
Using the Freudenthal suspension theorem, determine $\pi_k^s$ for
$k=0$ and $k=1$.
\end{exercise}

\begin{exercise}
Show that $K(\Z\times\Z,1)\simeq T^2$, the $2$-torus.  More
generally, show that $K(\Z^n,1)\simeq T^n$.
\end{exercise}

\begin{exercise}
\index{characteristic class!exercise}
Let $E\to S^2$ be a complex line bundle.  The first Chern class
$c_1(E)\in H^2(S^2;\Z)\cong\Z$ classifies $E$ up to isomorphism.
Identify this with the clutching construction\index{clutching
construction}: line bundles over $S^2$ correspond to maps
$S^1\to\mathrm{GL}_1(\C)=\C^*$, i.e., elements of
$\pi_1(\C^*)\cong\Z$.
\end{exercise}


%%%%%%%%%%%%%%%%%%%%%%%%%%%%%%%%%%%%%%%%%%%%%%%%%%%%%%%%%%%%%%%%%%%%
%  APPENDIX
%%%%%%%%%%%%%%%%%%%%%%%%%%%%%%%%%%%%%%%%%%%%%%%%%%%%%%%%%%%%%%%%%%%%
\appendix

\chapter{Homological Algebra Review}\label{app:homological-algebra}
\index{homological algebra}
\minitoc\bigskip

This appendix collects the main tools from homological algebra used
throughout the text.

%% ─── A.1 Chain complexes ────────────────────────────────────────
\section{Chain complexes}\label{sec:chain-complexes}
\index{chain complex}

\begin{definition}[Chain complex]
A \textbf{chain complex} $(C_\bullet,d)$ over a ring~$R$ is a sequence
of $R$-modules and $R$-module homomorphisms
\[
  \cdots \;\xrightarrow{d_{n+2}}\;
  C_{n+1} \;\xrightarrow{d_{n+1}}\;
  C_n \;\xrightarrow{d_n}\;
  C_{n-1} \;\xrightarrow{d_{n-1}}\;
  \cdots
\]
satisfying $d_n\circ d_{n+1}=0$ for all~$n$.

The \textbf{$n$-th homology} of~$C_\bullet$ is
\[
  H_n(C_\bullet)
  \;=\;
  \ker d_n\,/\,\im d_{n+1}
  \;=\;
  Z_n/B_n,
\]
where $Z_n=\ker d_n$ is the module of \textbf{$n$-cycles}\index{cycle}
and $B_n=\im d_{n+1}$ the module of
\textbf{$n$-boundaries}\index{boundary}.
\end{definition}

\begin{definition}[Cochain complex]
\index{cochain complex}
A \textbf{cochain complex} $(C^\bullet,\delta)$ is a sequence
\[
  \cdots \;\xrightarrow{\delta^{n-2}}\;
  C^{n-1} \;\xrightarrow{\delta^{n-1}}\;
  C^n \;\xrightarrow{\delta^n}\;
  C^{n+1} \;\xrightarrow{\delta^{n+1}}\;
  \cdots
\]
with $\delta^{n+1}\circ\delta^n=0$.  The \textbf{$n$-th cohomology}
is $H^n(C^\bullet)=\ker\delta^n/\im\delta^{n-1}$.
\end{definition}

\begin{definition}[Chain map]
\index{chain map}
A \textbf{chain map} $f\colon C_\bullet\to D_\bullet$ is a family of
homomorphisms $f_n\colon C_n\to D_n$ commuting with the
differentials: $d_n^D\circ f_n = f_{n-1}\circ d_n^C$.  Equivalently:
\[
\begin{tikzcd}
  \cdots \ar[r] & C_{n+1} \ar[r,"d_{n+1}^C"]\ar[d,"f_{n+1}"]
  & C_n \ar[r,"d_n^C"]\ar[d,"f_n"]
  & C_{n-1} \ar[r]\ar[d,"f_{n-1}"] & \cdots \\
  \cdots \ar[r] & D_{n+1} \ar[r,"d_{n+1}^D"]
  & D_n \ar[r,"d_n^D"]
  & D_{n-1} \ar[r] & \cdots
\end{tikzcd}
\]
A chain map induces homomorphisms
$f_*\colon H_n(C_\bullet)\to H_n(D_\bullet)$.
\end{definition}

\begin{definition}[Chain homotopy]
\index{chain homotopy}
Two chain maps $f,g\colon C_\bullet\to D_\bullet$ are \textbf{chain
homotopic}, written $f\simeq g$, if there exist homomorphisms
$s_n\colon C_n\to D_{n+1}$ satisfying
$f_n-g_n = d_{n+1}^D\circ s_n + s_{n-1}\circ d_n^C$ for all~$n$.
\end{definition}

\begin{proposition}
Chain homotopic maps induce the same map on homology: if
$f\simeq g$, then $f_*=g_*\colon H_n(C)\to H_n(D)$.
\end{proposition}

%% ─── A.2 Exact sequences ───────────────────────────────────────
\section{Exact sequences}\label{sec:exact-sequences}
\index{exact sequence}

\begin{definition}[Exact sequence]
A sequence of $R$-module homomorphisms
$\cdots\to A\xrightarrow{f}B\xrightarrow{g}C\to\cdots$
is \textbf{exact} at~$B$ if $\im f=\ker g$.  The sequence is
\textbf{exact} if it is exact at every interior term.
\end{definition}

\begin{definition}[Short exact sequence]
\index{exact sequence!short}
A \textbf{short exact sequence} (SES) is an exact sequence
\[
  0\to A\xrightarrow{f}B\xrightarrow{g}C\to 0,
\]
meaning $f$ is injective, $g$ is surjective, and $\im f=\ker g$.
\end{definition}

\begin{theorem}[Long exact sequence in homology]
\label{thm:les-homology}
\index{long exact sequence!in homology}
A short exact sequence of chain complexes
$0\to A_\bullet\to B_\bullet\to C_\bullet\to 0$ induces a long exact
sequence in homology:
\[
\begin{tikzcd}[column sep=1.4em]
  \cdots \ar[r]
  & H_n(A) \ar[r]
  & H_n(B) \ar[r]
  & H_n(C) \ar[r,"\partial"]
  & H_{n-1}(A) \ar[r]
  & \cdots
\end{tikzcd}
\]
The connecting homomorphism $\partial$ is constructed via diagram
chasing (the ``snake lemma''\index{snake lemma}).
\end{theorem}

\begin{lemma}[Five Lemma]\label{lem:five-lemma}
\index{Five Lemma}
Given a commutative diagram of $R$-modules with exact rows:
\[
\begin{tikzcd}
  A_1\ar[r]\ar[d,"f_1"] & A_2\ar[r]\ar[d,"f_2"]
  & A_3\ar[r]\ar[d,"f_3"] & A_4\ar[r]\ar[d,"f_4"]
  & A_5\ar[d,"f_5"] \\
  B_1\ar[r] & B_2\ar[r] & B_3\ar[r] & B_4\ar[r] & B_5
\end{tikzcd}
\]
If $f_1,f_2,f_4,f_5$ are isomorphisms, then $f_3$ is an isomorphism.
\end{lemma}

\begin{lemma}[Snake Lemma]\label{lem:snake-lemma}
\index{snake lemma}
Given a commutative diagram with exact rows:
\[
\begin{tikzcd}
  & A \ar[r,"f"]\ar[d,"\alpha"]
  & B \ar[r,"g"]\ar[d,"\beta"]
  & C \ar[r]\ar[d,"\gamma"]
  & 0 \\
  0 \ar[r]
  & A' \ar[r,"f'"]
  & B' \ar[r,"g'"]
  & C'
\end{tikzcd}
\]
there is an exact sequence
$\ker\alpha\to\ker\beta\to\ker\gamma
\xrightarrow{\partial}\coker\alpha\to\coker\beta\to\coker\gamma$.
\end{lemma}

%% ─── A.3 Universal coefficients and Künneth ────────────────────
\section{Universal coefficient theorems}\label{sec:uct}
\index{universal coefficient theorem}

\begin{theorem}[Universal Coefficient Theorem for cohomology]
\label{thm:uct}
\index{universal coefficient theorem!for cohomology}
For a chain complex $C_\bullet$ of free abelian groups and an abelian
group~$G$, there is a natural short exact sequence
\[
  0 \;\to\;
  \mathrm{Ext}^1_\Z(H_{n-1}(C),G) \;\to\;
  H^n(C;G) \;\to\;
  \Hom(H_n(C),G) \;\to\; 0.
\]
This sequence splits (non-naturally).
\end{theorem}

\begin{theorem}[K\"unneth formula]
\label{thm:kunneth}
\index{Kunneth formula@K\"unneth formula}
For spaces $X,Y$ with $H_*(X;\Z)$ free, there is a natural short exact
sequence
\[
  0 \;\to\;
  \bigoplus_{i+j=n}H_i(X)\otimes H_j(Y) \;\to\;
  H_n(X\times Y) \;\to\;
  \bigoplus_{i+j=n-1}\mathrm{Tor}_1(H_i(X),H_j(Y)) \;\to\; 0.
\]
Over a field, the Tor term vanishes and homology of a product is the
tensor product of homologies.
\end{theorem}

%% ─── A.4 Derived functors ──────────────────────────────────────
\section{Derived functors}\label{sec:derived-functors}
\index{derived functor}

\begin{definition}[Projective and injective resolutions]
\index{projective resolution}\index{injective resolution}
A \textbf{projective resolution} of an $R$-module $M$ is an exact
sequence
\[
  \cdots\to P_2\to P_1\to P_0\to M\to 0,
\]
where each $P_i$ is projective\index{projective module}.
Dually, an \textbf{injective resolution} is an exact sequence
$0\to M\to I^0\to I^1\to\cdots$ with each $I^j$
injective\index{injective module}.
\end{definition}

\begin{definition}[Tor and Ext]
\index{Tor}\index{Ext}
Let $M,N$ be $R$-modules.  Choose a projective resolution
$P_\bullet\to M\to 0$.  Then:
\begin{itemize}
  \item $\mathrm{Tor}_n^R(M,N) = H_n(P_\bullet\otimes_R N)$.
  \item $\mathrm{Ext}^n_R(M,N) = H^n(\Hom_R(P_\bullet,N))$.
\end{itemize}
These are independent of the choice of resolution (up to canonical
isomorphism).  The fact that different resolutions yield the same
derived functors follows from the \textbf{comparison theorem}%
\index{comparison theorem}: any two projective resolutions of the same
module are chain homotopy equivalent.
\end{definition}

\begin{example}[Computing Ext and Tor over $\Z$]
\index{Ext!computation}\index{Tor!computation}
To compute $\mathrm{Ext}^1_\Z(\Z/m,\Z/n)$, use the projective
resolution $0\to\Z\xrightarrow{\cdot m}\Z\to\Z/m\to 0$.  Applying
$\Hom_\Z(-,\Z/n)$:
\[
  0\to\Hom(\Z,\Z/n)\xrightarrow{\cdot m}\Hom(\Z,\Z/n)\to 0
  \quad\Longrightarrow\quad
  0\to\Z/n\xrightarrow{\cdot m}\Z/n\to 0.
\]
Hence $\mathrm{Ext}^1_\Z(\Z/m,\Z/n)\cong\Z/n\,/\,m(\Z/n)\cong
\Z/\gcd(m,n)$.

Similarly, $\mathrm{Tor}_1^\Z(\Z/m,\Z/n)\cong\ker(\Z/n\xrightarrow
{\cdot m}\Z/n)\cong\Z/\gcd(m,n)$.
\end{example}

\begin{proposition}[Basic properties of Tor and Ext]
\index{Tor!properties}\index{Ext!properties}
\begin{enumerate}[label=(\roman*)]
  \item $\mathrm{Tor}_0^R(M,N)\cong M\otimes_R N$ and
    $\mathrm{Ext}^0_R(M,N)\cong\Hom_R(M,N)$.
  \item A short exact sequence $0\to A\to B\to C\to0$ induces long
    exact sequences in both Tor and Ext.
  \item If $R=\Z$ and $M$ is free, then $\mathrm{Tor}_n(M,N)=0$ and
    $\mathrm{Ext}^n(M,N)=0$ for all $n\ge1$.
  \item $\mathrm{Ext}^1_\Z(\Z/m,G)\cong G/mG$ and
    $\mathrm{Tor}_1^\Z(\Z/m,G)\cong\set{g\in G\mid mg=0}$.
\end{enumerate}
\end{proposition}


%%%%%%%%%%%%%%%%%%%%%%%%%%%%%%%%%%%%%%%%%%%%%%%%%%%%%%%%%%%%%%%%%%%%
%  BIBLIOGRAPHY
%%%%%%%%%%%%%%%%%%%%%%%%%%%%%%%%%%%%%%%%%%%%%%%%%%%%%%%%%%%%%%%%%%%%
\begin{thebibliography}{99}

\bibitem{hatcher}
\index{Hatcher, A.}
A.~Hatcher,
\textit{Algebraic Topology},
Cambridge University Press, 2002.
Available at \url{https://pi.math.cornell.edu/~hatcher/AT/ATpage.html}.

\bibitem{spanier}
\index{Spanier, E.H.}
E.\,H.~Spanier,
\textit{Algebraic Topology},
Springer-Verlag, 1966 (corrected reprint 1981).

\bibitem{bredon}
\index{Bredon, G.E.}
G.\,E.~Bredon,
\textit{Topology and Geometry},
Graduate Texts in Mathematics \textbf{139}, Springer, 1993.

\bibitem{may}
\index{May, J.P.}
J.\,P.~May,
\textit{A Concise Course in Algebraic Topology},
Chicago Lectures in Mathematics, University of Chicago Press, 1999.

\bibitem{tomdieck}
\index{tom Dieck, T.}
T.~tom Dieck,
\textit{Algebraic Topology},
EMS Textbooks in Mathematics, European Mathematical Society, 2008.

\bibitem{eilenberg-steenrod}
\index{Eilenberg, S.}\index{Steenrod, N.E.}
S.~Eilenberg and N.\,E.~Steenrod,
\textit{Foundations of Algebraic Topology},
Princeton University Press, 1952.

\bibitem{bott-tu}
\index{Bott, R.}\index{Tu, L.W.}
R.~Bott and L.\,W.~Tu,
\textit{Differential Forms in Algebraic Topology},
Graduate Texts in Mathematics \textbf{82}, Springer, 1982.

\bibitem{milnor-stasheff}
\index{Milnor, J.W.}\index{Stasheff, J.D.}
J.\,W.~Milnor and J.\,D.~Stasheff,
\textit{Characteristic Classes},
Annals of Mathematics Studies \textbf{76}, Princeton University Press,
1974.

\bibitem{davis-kirk}
\index{Davis, J.F.}\index{Kirk, P.}
J.\,F.~Davis and P.~Kirk,
\textit{Lecture Notes in Algebraic Topology},
Graduate Studies in Mathematics \textbf{35}, AMS, 2001.

\bibitem{switzer}
\index{Switzer, R.M.}
R.\,M.~Switzer,
\textit{Algebraic Topology --- Homotopy and Homology},
Classics in Mathematics, Springer, 2002.

\bibitem{whitehead}
\index{Whitehead, G.W.}
G.\,W.~Whitehead,
\textit{Elements of Homotopy Theory},
Graduate Texts in Mathematics \textbf{61}, Springer, 1978.

\end{thebibliography}

%%%%%%%%%%%%%%%%%%%%%%%%%%%%%%%%%%%%%%%%%%%%%%%%%%%%%%%%%%%%%%%%%%%%
%  INDEX
%%%%%%%%%%%%%%%%%%%%%%%%%%%%%%%%%%%%%%%%%%%%%%%%%%%%%%%%%%%%%%%%%%%%
\printindex

\end{document}
% === END OF CHUNK ch9-11+end ===
